% !TeX root=../main.tex
\chapter{روش تحقیق}
\label{cmp:ch:method}
% Comparison-only compatibility label: chapter 4 is temporarily replaced here.
\label{chap:results}
\begingroup
% Chapter-local equation and definition numbering.
\renewcommand{\theequation}{\thechapter-\arabic{equation}}
\renewcommand{\thedefinition}{\thechapter-\arabic{definition}}

هدف این پژوهش، ارائهٔ روشی برای مهندسی ویژگی و استخراج بازنمایی‌های جدید از نمونه‌های یک مجموعه‌داده بر پایهٔ مقایسهٔ مقادیر ویژگی‌های آن‌هاست. برخلاف رویکردهایی که شباهت دو نمونه را مستقیماً از ترکیب فاصله‌های ویژگی‌های مشترک محاسبه می‌کنند، در روش پیشنهادی ابتدا ساختار ارتباطی نمونه‌ها به‌شکل یک گراف نمایش داده می‌شود؛ سپس روابط مستقیم و غیرمستقیم آن‌ها از طریق گشت روی این گراف استخراج می‌شود.

ایده اصلی آن است که هر نمونه را به مجموعه‌ای از زوج‌های نمونه--ویژگی تجزیه کنیم. روابط مقایسه‌ای مقادیر هر ویژگی در نمونه‌های مختلف و ارتباط میان ویژگی‌های متعلق به یک نمونه، دو نوع اتصال در این گراف ایجاد می‌کنند. ترکیب این اتصال‌ها امکان حرکت میان نمونه‌ها از مسیرهای متفاوت را فراهم می‌آورد؛ به‌گونه‌ای که حتی دو نمونه فاقد ویژگی مشترک نیز ممکن است از طریق نمونه‌های واسط با یکدیگر مرتبط شوند.

همچنین می‌توان اهمیت نسبی ویژگی‌ها در ایجاد این ارتباطات را به‌جای تعیین پیش‌فرض، با سیاستی قابل یادگیری مشخص کرد. حاصل این ساختار، بازنمایی احتمالاتی همسایگی نمونه‌هاست که می‌توان آن را در مهندسی ویژگی و مسائل یادگیری ماشین به کار گرفت.

مسیر این فصل از \emph{تعریف مدل} تا \emph{شیوهٔ اجرای آن}
به‌ترتیب سازمان یافته است. ابتدا گراف نمونه--ویژگی،
گشت تصادفیِ بدون بازگشت به مبدأ و خانواده‌های
سیاست انتخاب ویژگی را معرفی می‌کنیم؛ سپس
محاسبهٔ امید دقیق گشت، بازنمایی احتمالاتی نمونه‌ها
و سازوکار پیش‌بینی با مراکز توزیعی طبقات
را صورت‌بندی می‌کنیم. در گام بعد،
معیارهای کیفیت بازنمایی، حاشیهٔ تصمیم
و توابع زیان، مبنای آموزش سیاست را
فراهم می‌کنند. در پایان مسیر اجرایی،
الگوریتم‌های آموزش و استنتاج برای
نمونهٔ جدید با مرجع آموزشی ثابت
مشخص می‌شوند.

در کنار این مسیر اصلی، بعضی بخش‌ها به
\emph{خواص و کران‌های نظری تکمیلی}
اختصاص دارند؛ مطالعهٔ آن‌ها برای
فهم تعریف‌ها و اجرای روش اصلی
ضروری نیست و در مقدمهٔ همان بخش‌ها
مسیر ادامهٔ مطالعه مشخص شده است.
پس از آموزش و استنتاج، نگاشت سیاست
به مراکز را از جهت دستیابی‌پذیری و
چندگانگی بررسی می‌کنیم و سپس اثر
ویژگی‌های نویزی، زائد و وابسته را
تحلیل می‌کنیم. سرانجام ظرفیت
بازنمایی‌های آموخته‌شده برای
تحلیل پساآموزش و توسعه‌های مقاوم
صورت‌بندی می‌شود؛ پیشنهادهای
این بخش پایانی، نتیجهٔ تجربی
تأییدشدهٔ پژوهش حاضر نیستند.

\section{زنجیره‌های مارکوف درهم‌تنیده: از گراف نمونه--ویژگی به گراف نمونه‌ها}
\label{cmp:sec:ch3-interwoven}
در این بخش، پس از بیان تعاریف و فرض‌های اولیه، دو خانواده زیرگراف را معرفی می‌کنیم و با درهم‌تنیدن آن‌ها، گراف نمونه--ویژگی را می‌سازیم. سپس حرکت تصادفی روی این گراف را تعریف کرده و نشان می‌دهیم چگونه هر گشت یک بردار احتمال و میانگین گشت‌ها یک توزیع همسایگی ایجاد می‌کند. در انتها، ماتریس توزیع همسایگی و بازنمایی فشرده مبتنی بر برچسب را معرفی می‌کنیم.

\subsection{تعاریف اولیه و فرض‌های مسئله}
\label{cmp:subsec:ch3-definitions}

\begin{definition}[فضای ویژگی نمونه]
فرض کنیم مجموعه‌داده $D=\{X_1,X_2,\ldots,X_n\}$ شامل $n$ نمونه باشد. به ازای هر نمونه $X_i\in D$، مجموعه ویژگی‌هایی را که در آن نمونه دارای مقدار مشخص هستند، با $A_{X_i}$ نشان می‌دهیم. فضای ویژگی کل مجموعه‌داده، اجتماع فضای ویژگی تمام نمونه‌هاست:
\begin{equation}\label{cmp:eq:ch3-feature-space}
A_D=\bigcup_{X_i\in D}A_{X_i}.
\end{equation}
\end{definition}
در این تعریف، الزامی وجود ندارد که تمام نمونه‌ها مجموعه ویژگی یکسانی داشته باشند. بنابراین ساختار پیشنهادی امکان نمایش نمونه‌های دارای ویژگی‌های ناموجود یا مقادیر گمشده را فراهم می‌کند. ناموجود بودن یک ویژگی با صفر بودن مقدار آن متفاوت است؛ ویژگی دارای مقدار صفر همچنان عضو $A_{X_i}$ است.

\begin{definition}[مجموعه نمونه‌های دارای یک ویژگی]
برای هر ویژگی $a\in A_D$، مجموعه نمونه‌هایی که مقدار ویژگی $a$ در آن‌ها موجود است، با $D_a$ نمایش داده می‌شود:
\begin{equation}\label{cmp:eq:ch3-feature-samples}
D_a=\bigl\{X_i\in D\mid a\in A_{X_i}\bigr\}.
\end{equation}
\end{definition}
تعاریف فوق دو دیدگاه مکمل فراهم می‌کنند: ویژگی‌های موجود در یک نمونه و نمونه‌های دارای یک ویژگی مشخص. بر اساس این تعاریف، رابطه زیر برقرار است:
\begin{equation}\label{cmp:eq:ch3-incidence}
a\in A_{X_i}\quad\Longleftrightarrow\quad X_i\in D_a.
\end{equation}

\begin{definition}[نمونه--ویژگی]
هر زوج مرتب $(X_i,a)$، به شرط $X_i\in D$ و $a\in A_{X_i}$، یک \textbf{نمونه--ویژگی} نامیده می‌شود. مقدار آن با نماد $a(X_i)$ مشخص می‌شود.
\end{definition}
برای نمونه، اگر فردی به نام علی دارای ویژگی‌های قد، وزن و رنگ چشم باشد، سه نمونه--ویژگی متناظر با این سه ویژگی برای او تعریف می‌شود. در این نمایش، مجموعه‌داده دیگر صرفاً مجموعه‌ای از بردارهای ویژگی نیست؛ بلکه به مجموعه‌ای از واحدهای قابل ارتباط و مقایسه تبدیل می‌شود.

\noindent\textbf{مثال ۳-۱ (فضای ویژگی ناهمسان).}
فرض کنیم $X_1$ دارای ویژگی‌های قد، وزن و فشار خون باشد و $X_2$ دارای ویژگی‌های قد، وزن و چربی خون. در این صورت، $A_{X_1}$ و $A_{X_2}$ هرکدام سه عضو دارند، اما $A_D$ شامل چهار ویژگی قد، وزن، فشار خون و چربی خون است. همچنین $D_{\text{قد}}=D_{\text{وزن}}=\{X_1,X_2\}$، درحالی‌که $D_{\text{فشار خون}}=\{X_1\}$ و $D_{\text{چربی خون}}=\{X_2\}$ است. اگر مقدار فشار خون نمونه اول $12.6$ باشد، با قراردادن $a=\text{فشار خون}$، می‌نویسیم $a(X_1)=12.6$. این مثال نشان می‌دهد که وجود فضای ویژگی یکسان برای همه نمونه‌ها ضروری نیست. مثال حاضر صرفاً برای توضیح نمادگذاری است و قرار نیست فرض حداقل تعداد نمونه‌های دارای هر ویژگی را برای اجرای گشت برآورده کند.

\medskip
\noindent\textbf{فرض اول (ثابت بودن فضای ویژگی‌ها).}
فضای ویژگی $A_D$ در محدوده مسئله ثابت و از پیش مشخص فرض می‌شود. با ورود نمونه‌های جدید، ویژگی تازه‌ای خارج از این فضا به مدل اضافه نمی‌شود. این فرض نه به معنای یکسان بودن ویژگی‌های موجود در همه نمونه‌هاست و نه به معنای ثابت بودن مقادیر آن‌ها.

\medskip
\noindent\textbf{فرض دوم (وجود متریک برای مقایسه مقادیر هر ویژگی).}
برای هر ویژگی $a\in A_D$، یک تابع متریک $\mu_a$ بر فضای مقادیر آن وجود دارد. در نتیجه، فاصله میان $a(X_i)$ و $a(X_j)$ برای $X_i,X_j\in D_a$ قابل محاسبه است. این متریک لزوماً برای ویژگی‌های متفاوت یکسان نیست. مقادیر ویژگی‌ها نیز لازم نیست اسکالر باشند: برای متن، تصویر یا صوت می‌توان از بازنمایی مناسب و متریک تعریف‌شده روی آن استفاده کرد. وجود متریک، امکان تعریف وزن شباهت و سپس احتمال انتقال در زیرگراف ویژگی را فراهم می‌کند؛ البته خود متریک به‌تنهایی احتمال انتقال یکتایی تعیین نمی‌کند.

\medskip
\noindent\textbf{فرض سوم (اتصال‌پذیری مجموعه‌داده).}
فرض می‌کنیم نتوان $D$ را به دو زیرمجموعه ناتهی و مجزای $D_1$ و $D_2$ افراز کرد که فضای ویژگی آن‌ها هیچ اشتراکی نداشته باشد. فضای ویژگی هر زیرمجموعه از داده‌ها نیز، مشابه تعریف فضای ویژگی کل مجموعه‌داده، برابر اجتماع فضای ویژگی نمونه‌های عضو آن زیرمجموعه در نظر گرفته می‌شود. یعنی برای هر افراز غیر بدیهی $D=D_1\mathbin{\dot\cup}D_2$، داریم:
\begin{equation}\label{cmp:eq:ch3-connectivity}
A_{D_1}\cap A_{D_2}\neq\varnothing.
\end{equation}
در نتیجه، میان هر دو نمونه می‌توان زنجیره‌ای از نمونه‌های واسط یافت که هر دو نمونه متوالی آن حداقل یک ویژگی مشترک دارند. این شرط اتصال‌پذیری \emph{ساختاری} را بیان می‌کند و به‌تنهایی تضمین نمی‌کند پس از اعمال ممنوعیت بازگشت به مبدأ، در هر وضعیت مقصد مجاز وجود داشته باشد.

\medskip
\noindent\textbf{فرض چهارم (طبقه‌بندی تک‌برچسبی).}
در بخش‌های نظارت‌شده این پژوهش، هر نمونه دقیقاً یک برچسب $y_i\in\{1,\ldots,K\}$ دارد، به‌طوری‌که $K\ge2$ و هر طبقه در مجموعه مرجع دارای حداقل یک نمونه است. بنابراین ماتریس برچسب‌های یک‌داغ و مراکز طبقه‌ها برای مسئله تک‌برچسبی تعریف می‌شوند و تعمیم آن‌ها به طبقه‌بندی چندبرچسبی در این فصل ادعا نمی‌شود.

\medskip
\noindent\textbf{فرض پنجم (حداقل پوشش ویژگی‌ها).}
برای هر ویژگی $a\in A_D$ فرض می‌کنیم $|D_a|\ge3$. با توجه به مثبت بودن وزن ارتباط، این فرض تضمین می‌کند که پس از کنارگذاشتن نمونه جاری و نمونه آغازین، برای هر ویژگی موجود دست‌کم یک مقصد مجاز باقی بماند.

\subsection{ساخت گراف نمونه--ویژگی}
\label{cmp:subsec:ch3-graph}
گراف مورد نظر از درهم‌تنیدن دو خانواده زیرگراف تشکیل می‌شود: زیرگراف‌های عمودی که مقادیر یک ویژگی را در نمونه‌های مختلف مرتبط می‌کنند و زیرگراف‌های افقی که ارتباط میان ویژگی‌های یک نمونه را نمایش می‌دهند.

\paragraph{زیرگراف‌های عمودی.}
برای هر $a\in A_D$، مجموعه نمونه--ویژگی‌های $\{(X_i,a):X_i\in D_a\}$ را رأس‌های یک گراف کامل $G_a$ در نظر می‌گیریم. فاصله میان دو رأس از متریک $\mu_a$ به‌دست می‌آید. برای تبدیل آن به وزن ارتباط، تابع مثبت و نزولی $\kappa_a$ را به کار می‌بریم:
\begin{equation}\label{cmp:eq:ch3-edge-weight}
w_a(X_i,X_j)=\kappa_a\!\left(\mu_a\bigl(a(X_i),a(X_j)\bigr)\right),
\qquad X_i,X_j\in D_a.
\end{equation}
بدین ترتیب، نمونه‌های نزدیک‌تر از دید ویژگی $a$ وزن ارتباط بیشتری دارند. متریک متقارن و تابع تبدیل یکسان، وزن خام متقارن ایجاد می‌کنند؛ ازاین‌رو این زیرگراف‌ها را می‌توان بدون جهت نمایش داد. مجموعه زیرگراف‌های عمودی را با $\mathcal G_A=\{G_a:a\in A_D\}$ نشان می‌دهیم.

\paragraph{زیرگراف‌های افقی.}
به ازای هر نمونه $X_i$، گراف جهت‌دار و کامل $G_{X_i}$ را روی رأس‌های $\{(X_i,a):a\in A_{X_i}\}$ در نظر می‌گیریم. در این گراف، یال‌ها امکان جابه‌جایی از یک ویژگی به ویژگی دیگر را نشان می‌دهند. چون مقادیر دو ویژگی متفاوت الزاماً قابل مقایسه مستقیم نیستند، وزن این یال‌ها از پیش معلوم نیست و به وسیله \textbf{سیاست انتخاب ویژگی} تعیین می‌شود. جهت‌دار بودن یال‌ها اجازه می‌دهد احتمال انتخاب $b$ پس از $a$ با احتمال انتخاب $a$ پس از $b$ متفاوت باشد؛ طوقه‌ها نیز امکان انتخاب مجدد همان ویژگی را فراهم می‌کنند. مجموعه زیرگراف‌های افقی را با $\mathcal G_D=\{G_{X_i}:X_i\in D\}$ نمایش می‌دهیم.

\paragraph{درهم‌تنیدن زیرگراف‌ها.}
اکنون هر زوج نمونه--ویژگی موجود را یک رأس گراف اصلی در نظر می‌گیریم. هر $G_a$ رأس‌های متعلق به ویژگی $a$ را در نمونه‌های متفاوت و هر $G_{X_i}$ رأس‌های متعلق به نمونه $X_i$ را به یکدیگر متصل می‌کند. این دو خانواده مانند تار و پود در یک گراف واحد درهم‌تنیده می‌شوند. بدین ترتیب، حرکت می‌تواند به‌طور متناوب با انتخاب ویژگی در یک نمونه و سپس انتقال به نمونه‌ای دیگر از طریق همان ویژگی انجام شود.

\noindent\textbf{مثال ۳-۲ (گراف درهم‌تنیده).}
شکل~\ref{cmp:fig:ch3-sample-feature-graph} گراف مجموعه‌ای با سه نمونه و چهار ویژگی را نشان می‌دهد. هر نمونه در ناحیه‌ای مجزا نمایش داده شده و دایره‌های درون آن نمونه--ویژگی‌ها هستند. یال‌های پیوسته رنگی، ارتباطات عمودی و یال‌های خط‌چین، ارتباطات افقی را مشخص می‌کنند.

در شکل~\ref{cmp:fig:ch3-missing-features} برخی ویژگی‌ها ناموجودند؛ دایره‌های توخالی به این وضعیت اشاره دارند. اتصال عمودی فقط بین نمونه--ویژگی‌های موجود ساخته می‌شود و سیاست افقی نیز هنگام انتخاب ویژگی، گزینه‌های ناموجود را حذف می‌کند. یال‌های قرمز خط‌چین در شکل نمایانگر انتقال‌های غیرمجاز هستند. این دو شکل برای نمایش \emph{ساختار گراف} هستند؛ به‌ویژه شکل دوم با تنها سه نمونه، لزوماً فرض حداقل سه نمونه برای هر ویژگی را برآورده نمی‌کند و نمونه عددیِ اجرای گشت تحت همه فرض‌های بعدی نیست.

% Two captions/labels remain separate; one float keeps both figures on a single page.
\begin{figure}[!htbp]
\centering
\includegraphics[width=0.79\linewidth,height=0.32\textheight,keepaspectratio]{chap3/sample-feature-graph.png}
\caption{نمونه‌ای از درهم‌تنیدن زیرگراف‌های عمودی و افقی برای سه نمونه و چهار ویژگی.}
\label{cmp:fig:ch3-sample-feature-graph}

\medskip
\includegraphics[width=0.79\linewidth,height=0.32\textheight,keepaspectratio]{chap3/missing-features.png}
\caption{گراف نمونه--ویژگی با برخی ویژگی‌های ناموجود و حذف انتقال‌های غیرمجاز.}
\label{cmp:fig:ch3-missing-features}
\end{figure}

\subsection{گام‌برداری تصادفی و پرهیز از بازگشت به نمونه آغازین}
\label{cmp:subsec:ch3-walk}
هر گشت تصادفی از نمونه آغازین $X_s$ شروع می‌شود. ابتدا یک ویژگی از میان ویژگی‌های مجاز این نمونه، مطابق توزیع احتمال اولیه انتخاب می‌شود. پس از تعیین رأس نمونه--ویژگی آغازین، حرکت با تکرار گام‌های دومرحله‌ای ادامه می‌یابد. توزیع اولیه نیز می‌تواند به سیاست انتخاب ویژگی وابسته باشد.

یک گام کامل حرکت تصادفی از دو مرحله تشکیل می‌شود. \textbf{در مرحله اول}، گام‌بردار از رأس $(X_i,a)$ با استفاده از سیاست انتخاب ویژگی به $(X_i,b)$ می‌رود؛ در این مرحله $b\in A_{X_i}$ است و طوقه $b=a$ نیز می‌تواند انتخاب شود. \textbf{در مرحله دوم}، با استفاده از زیرگراف عمودی $G_b$، گام‌بردار از $(X_i,b)$ به نمونه--ویژگی $(X_j,b)$ در نمونه‌ای دیگر منتقل می‌شود. در نتیجه احتمال جابه‌جایی هم به سیاست انتخاب ویژگی و هم به شباهت نمونه‌ها در ویژگی انتخاب‌شده وابسته است.

فرض کنیم گشت از نمونه $X_s$ آغاز شده باشد. در تمام مراحل انتقال، مشاهده مجدد $X_s$ ممنوع است؛ اما بازدید دوباره از سایر نمونه‌ها مجاز است. برای یک ویژگی $b$، مجموعه مقصدهای مجاز از $X_i$ را $D_b\setminus\{X_i,X_s\}$ در نظر می‌گیریم. اگر این مجموعه ناتهی باشد، احتمال انتقال عمودی با نرمال‌سازی وزن‌های \eqref{cmp:eq:ch3-edge-weight} تعریف می‌شود:
\begin{equation}\label{cmp:eq:ch3-vertical-transition}
q_b^{(s)}(j\mid i)=
\begin{cases}
\displaystyle\frac{w_b(X_i,X_j)}
{\sum\limits_{X_k\in D_b\setminus\{X_i,X_s\}}w_b(X_i,X_k)},
& X_j\in D_b\setminus\{X_i,X_s\},\\[6pt]
0,&\text{در غیر این صورت}.
\end{cases}
\end{equation}
در خانواده‌های سیاستِ شرطی به طبقه، \emph{برچسب نمونهٔ آغازین}
$X_s$ در تعیین سیاست انتخاب ویژگی نقش دارد؛ این وابستگی
به برچسب نمونهٔ جاری $X_i$ یا نمونهٔ مقصد $X_j$ نیست.
برای صورت‌بندی یکسان آموزش و آزمون، $c_s\in\{1,\ldots,K\}$
را \emph{طبقهٔ منتسب به نمونهٔ آغازین} می‌نامیم.
در ساخت مرکز طبقه با داده‌های آموزشیِ برچسب‌دار،
$c_s=y_s$ است و از سیاست متناظر با برچسب واقعی مبدأ استفاده می‌شود.
اما هنگام پیش‌بینی، برچسب واقعی نمونهٔ آغازین ناشناخته است؛
در این حالت یا هر مقدار نامزد $c_s$ را به‌عنوان فرضیه بررسی می‌کنیم،
یا از سیاست میانگینِ مستقل از برچسب مجهول بهره می‌بریم.
این دو شیوه و تفاوت کاربرد سیاست در آموزش و آزمون در
زیربخش~\ref{cmp:subsec:ch3-parametric-feature-policy} توضیح داده می‌شوند.

با ثابت‌گرفتن مبدأ $X_s$ و فرضیهٔ طبقهٔ آن $c_s$ در طول گشت،
احتمال یک گام کامل از حاصل‌ضرب احتمال انتخاب ویژگی و
احتمال انتقال عمودی به‌دست می‌آید:
\begin{equation}\label{cmp:eq:ch3-full-transition}
\Pr\bigl((X_j,b)\mid(X_i,a),X_s,c_s\bigr)
=\underbrace{\pi_s^{(c_s)}(b\mid i,a)}_{\text{سیاست مشروط به طبقهٔ مبدأ}}\,
q_b^{(s)}(j\mid i).
\end{equation}
وابستگی به $c_s$ در عامل افقی $\pi_s^{(c_s)}$ وارد می‌شود؛
عامل عمودی $q_b^{(s)}$ با ثابت‌بودن مبدأ، به فرضیهٔ طبقه وابسته نیست.
در خانواده‌های مشترک و یکنواخت، سیاست افقی نیز مستقل از $c_s$
است و شرط‌گذاری بالا در آن‌ها اثری بر احتمال انتخاب ویژگی ندارد.
در این رابطه، سیاست انتخاب ویژگی فقط روی ویژگی‌های موجود با مقصد عمودی مجاز نرمال می‌شود. بر اساس فرض پنجم و مثبت بودن وزن‌ها، این مقصدها برای ویژگی‌های موجود در وضعیت‌های قابل دسترس وجود دارند و نیازی به تعریف قاعده توقف زودهنگام نیست. شرط اتصال‌پذیری ساختاری به‌تنهایی برای این نتیجه کافی نبود.

از آنجا که ممنوعیت بازگشت به مبدأ به نمونه $X_s$ وابسته است و سیاست انتخاب ویژگی ممکن است به ویژگی قبلی وابسته باشد، حرکت به‌طور کلی یک زنجیره مارکوف همگن روی \emph{فضای نمونه‌ها به‌تنهایی} نیست. با ثابت نگه داشتن $X_s$ و در نظر گرفتن زوج $(X_i,a)$ به‌عنوان حالت، می‌توان آن را روی فضای حالت گسترش‌یافته به‌صورت مارکوفی نمایش داد.

\paragraph{طول گشت و وزن‌دهی مشاهدات.}
طول گشت را با $L$ نشان می‌دهیم و برای آن توزیع هندسی مستقل از انتخاب مسیر در نظر می‌گیریم. اگر $\beta\in(0,1]$ احتمال توقف پس از هر گام باشد، آنگاه:
\begin{equation}\label{cmp:eq:ch3-geometric-length}
\Pr(L=\ell)=\beta(1-\beta)^{\ell-1},\qquad
\mathbb E[L]=\frac{1}{\beta},\qquad \ell=1,2,\ldots.
\end{equation}
به‌عنوان نمونه، $\beta=1/3$ متوسط طول گشت را برابر سه قرار می‌دهد. مستقل از این انتخاب، پارامتر $0<\lambda\le1$ وزن مشاهده‌های دورتر را کنترل می‌کند: وزن مشاهده گام $t$ام برابر $\lambda^{t-1}$ است. مجموع وزن‌های گشتی با طول $\ell$ به‌صورت زیر تعریف می‌شود:
\begin{equation}\label{cmp:eq:ch3-walk-weight-normalizer}
\mathcal N_\ell(\lambda)=\sum_{t=1}^{\ell}\lambda^{t-1}
=1+\lambda+\cdots+\lambda^{\ell-1}.
\end{equation}
برای محاسبه این مجموع، اگر $0<\lambda<1$ باشد، با ضرب رابطه بالا در $\lambda$ و تفریق دو تساوی، جمله‌های میانی حذف می‌شوند و داریم:
\begin{equation}\label{cmp:eq:ch3-walk-weight-normalizer-telescoping}
(1-\lambda)\mathcal N_\ell(\lambda)=1-\lambda^\ell.
\end{equation}
در نتیجه، عامل نرمال‌سازی دقیق هر گشت برابر است با:
\begin{equation}\label{cmp:eq:ch3-walk-weight-normalizer-closed}
\mathcal N_\ell(\lambda)=
\begin{cases}
\displaystyle\frac{1-\lambda^\ell}{1-\lambda},&0<\lambda<1,\\[5pt]
\ell,&\lambda=1.
\end{cases}
\end{equation}
پس $\beta$ احتمال وقوع طول‌های مختلف گشت را تعیین می‌کند، اما $\lambda$ اهمیت نسبی مشاهدات در طول \emph{همان گشت} را مشخص می‌سازد. این دو نقش متمایزند.

\paragraph{بردار همسایگی حاصل از یک گشت.}
به ازای هر نمونه $X_j$، بردار سطری یکه $e_j\in\mathbb R^{1\times n}$ را تعریف می‌کنیم که مؤلفه $j$ام آن یک و سایر مؤلفه‌ها صفر است. فرض کنیم گشت $W$ از $X_s$ آغاز شود و در $L$ گام خود، به‌ترتیب نمونه‌های $X_{J_1},\ldots,X_{J_L}$ را مشاهده کند. بردار حاصل از این گشت را چنین تعریف می‌کنیم:
\begin{equation}\label{cmp:eq:ch3-walk-vector}
R(W)=\frac{1}{\mathcal N_L(\lambda)}
\sum_{t=1}^{L}\lambda^{t-1}e_{J_t}\in\mathbb R^{1\times n}.
\end{equation}
این بردار نامنفی است و مجموع مؤلفه‌های آن یک است. در یک گشت تک‌گامی، $R(W)$ دقیقاً بردار یکه نمونه مشاهده‌شده خواهد بود؛ در گشت‌های طولانی‌تر، وزن بین نمونه‌های مشاهده‌شده تقسیم می‌شود. اگر نمونه‌ای چند بار دیده شود، سهم آن مجموع وزن تمام بازدیدهایش است. هر گشت، صرف‌نظر از طولش، \emph{فقط یک بردار احتمال نرمال‌شده} تولید می‌کند. چون نمونه آغازین در گشت دوباره مشاهده نمی‌شود، مؤلفه $s$ام این بردار همواره صفر است.

\subsection{ماتریس‌های توزیع همسایگی نمونه‌ها و برچسب‌ها}
\label{cmp:subsec:ch3-representations}
\paragraph{میانگین‌گیری از گشت‌ها.}
یک گشت تنها یکی از مسیرهای ممکن را مشاهده می‌کند. اگر $N$ گشت مستقل $W_1,\ldots,W_N$ از نمونه $X_s$ اجرا کنیم، توزیع همسایگی با میانگین‌گیری بردارهای حاصل تخمین زده می‌شود:
\begin{equation}\label{cmp:eq:ch3-h-estimate}
\widehat h_s=\frac{1}{N}\sum_{r=1}^{N}R(W_r),
\qquad h_s=\mathbb E[R(W)\mid X_s].
\end{equation}
به دلیل نرمال‌سازی جداگانه هر گشت، گشت‌های طولانی‌تر صرفاً به علت تعداد بیشتر مشاهدات، وزن کلی بیشتری نمی‌گیرند. $h_s$ یک بردار سطری احتمالی با $n$ مؤلفه است؛ مؤلفه $s$ام آن صفر باقی می‌ماند.

\paragraph{توزیع انتقال مرتبه $t$ام.}
برای مشخص کردن نقش عمق گشت، بردار $P_{X_s,t}$ را توزیع نمونه مشاهده‌شده در گام $t$ام، به شرط شروع حرکت از $X_s$ و ادامه یافتن فرایند تا آن گام، تعریف می‌کنیم:
\begin{equation}\label{cmp:eq:ch3-step-observation-distribution}
P_{X_s,t}(j)=\Pr(J_t=j\mid J_0=s),
\qquad j=1,\ldots,n.
\end{equation}
در این تعریف، متغیر توقف مستقل از مسیر است و $J_t$ را می‌توان در فرایند ادامه‌یافته حرکت در نظر گرفت. هر $P_{X_s,t}$ یک بردار احتمال سطری است و مؤلفه مربوط به نمونه آغازین در آن صفر است. این نماد به معنای توان $t$ام یک ماتریس گذار روی نمونه‌ها نیست. در زنجیره مارکوف همگن معمولی روی نمونه‌ها، توزیع‌های چندگامی از توان ماتریس گذار به‌دست می‌آیند؛ اما در اینجا ممنوعیت بازگشت به مبدأ و وابستگی سیاست به ویژگی قبلی، مانع استفاده مستقیم از یک ماتریس ثابت روی فضای نمونه‌ها می‌شود. برای مثال، حتی اگر قطر یک ماتریس گذار معمولی $M$ صفر باشد، ممکن است $(M^2)_{ss}=\sum_jM_{sj}M_{js}>0$ شود؛ درحالی‌که این بازگشت در روش ما ممنوع است.

با تعریف \eqref{cmp:eq:ch3-walk-vector} و توزیع هندسی طول گشت، بردار همسایگی را می‌توان به‌صورت ترکیب محدب توزیع‌های مشاهده در عمق‌های مختلف نوشت:
\begin{equation}\label{cmp:eq:ch3-h-alpha}
\boxed{\displaystyle
h_s=\sum_{t=1}^{\infty}\alpha_t(\beta,\lambda)\,P_{X_s,t},
\qquad
\alpha_t(\beta,\lambda)=\lambda^{t-1}
\sum_{\ell=t}^{\infty}
\frac{\beta(1-\beta)^{\ell-1}}{\mathcal N_\ell(\lambda)}.}
\end{equation}
ضرایب $\alpha_t$ نامنفی‌اند و مجموع آن‌ها برابر یک است. اثبات رابطه \eqref{cmp:eq:ch3-h-alpha} و خاصیت نرمال‌بودن این ضرایب در پیوست~\ref{app:ch3-weighted-walk-proof} آمده است. باید توجه داشت که به‌سبب نرمال‌سازی جداگانه هر گشت با $\mathcal N_L(\lambda)$، این ضرایب الزاماً یک دنباله هندسی ساده نیستند؛ ازاین‌رو رابطه را نمی‌توان بدون فرض اضافی به سری توان‌های یک ماتریس انتقال معمولی فروکاست.

\paragraph{ماتریس توزیع همسایگی.}
با قراردادن بردار همسایگی هر نمونه در سطر متناظر، ماتریس توزیع همسایگی روی نمونه‌ها را تعریف می‌کنیم:
\begin{equation}\label{cmp:eq:ch3-H}
\mathscr H=
\begin{bmatrix}h_1\\h_2\\\vdots\\h_n\end{bmatrix}
=[h_{ij}]_{n\times n},
\qquad h_{ij}\ge0,\quad
\sum_{j=1}^{n}h_{ij}=1,\quad h_{ii}=0.
\end{equation}
درایه $h_{ij}$ سهم نمونه $X_j$ در همسایگی نمونه $X_i$ را نشان می‌دهد. ماتریس $\mathscr H$ لزوماً متقارن نیست؛ زیرا احتمال مشاهده یک نمونه با شروع از نمونه دیگر به جهت حرکت و سیاست انتخاب ویژگی بستگی دارد. هر بردار تک‌گشت معمولاً تنک است، اما ماتریس امید ریاضی $\mathscr H$ الزاماً تنک باقی نمی‌ماند و نگهداری کامل آن برای مجموعه‌داده‌های بزرگ می‌تواند پرهزینه باشد.

\paragraph{ماتریس توزیع برچسب همسایگی.}
در کاربردهای دارای برچسب، گاهی به‌جای احتمال مشاهده تک‌تک نمونه‌ها، توزیع برچسب نمونه‌های مشاهده‌شده مورد نیاز است. مطابق فرض تک‌برچسبی، ماتریس برچسب‌های یک‌داغ $Y\in\{0,1\}^{n\times K}$ را تعریف می‌کنیم که در هر سطر آن دقیقاً یک مؤلفه برابر یک است و این مؤلفه برچسب نمونه متناظر را مشخص می‌کند. \textbf{ماتریس توزیع برچسب همسایگی} را با نماد $\mathscr Y$ نمایش می‌دهیم. این ماتریس از تجمیع احتمال‌های مشاهده نمونه‌ها بر اساس برچسب آن‌ها به‌دست می‌آید و می‌توان آن را با ضرب ماتریس توزیع همسایگی در ماتریس برچسب‌های یک‌داغ محاسبه کرد:
\begin{equation}\label{cmp:eq:ch3-Y-script}
\boxed{\mathscr Y=\mathscr H Y},
\qquad \mathscr Y\in\mathbb R^{n\times K}.
\end{equation}
درایهٔ $(i,k)$ این ماتریس، مجموع احتمال مشاهدهٔ نمونه‌های دارای برچسب $k$ است:
\begin{equation}\label{cmp:eq:ch3-Y-script-entry}
\mathscr Y_{ik}=\sum_{\substack{j=1\\ y_j=k}}^{n}h_{ij}
=\sum_{\substack{j=1\\ j\ne i,\;y_j=k}}^{n}h_{ij}.
\end{equation}
تساوی دوم حاصل صفر بودن قطر $\mathscr H$ است؛ بنابراین \textbf{برچسب خود نمونه $X_i$ مستقیماً در بازنمایی همسایگی آن مشارکت ندارد}. البته برچسب نمونه‌های دیگر، حتی اگر با برچسب $X_i$ یکسان باشند، می‌تواند سهم داشته باشد. در نتیجه، از توزیع $n$بعدی روی نمونه‌ها به توزیعی $K$بعدی روی برچسب‌ها می‌رسیم. در کاربرد آزمون، برچسب نمونه جدید در ماتریس مرجع $Y$ وارد نمی‌شود؛ $Y$ و مراکز طبقه‌ها تنها با برچسب‌های مجموعه مرجع آموزشی ساخته می‌شوند.

علاوه بر این، برای تخمین $\mathscr Y$ نیازی به ساخت کامل $\mathscr H$ نیست: در هر گشت می‌توان به‌جای جمع‌کردن بردارهای یکه نمونه‌ها، سطر متناظر آن‌ها از ماتریس $Y$ را با همان وزن‌ها جمع کرد و پس از نرمال‌سازی، میانگین گشت‌ها را محاسبه کرد. این روش دو خروجی روشن برای مهندسی ویژگی فراهم می‌کند: بازنمایی عمومی همسایگی نمونه‌ها در $\mathscr H$ و بازنمایی فشرده وابسته به برچسب در $\mathscr Y$. بررسی کیفیت سیاست انتخاب ویژگی و کاربردهای این بازنمایی‌ها موضوع بخش‌های بعدی خواهد بود.

\section{سیاست انتخاب ویژگی}
\label{cmp:sec:ch3-feature-selection-policy}
در بخش~\ref{cmp:subsec:ch3-walk}، احتمال انتخاب ویژگی با
$\pi_s^{(c_s)}(b\mid i,a)$ معرفی شد، اما چگونگی تعیین آن هنوز مشخص نشده است.
در این بخش، نخست سیاست را با \emph{ماتریسی از توابع پارامتری}
صورت‌بندی می‌کنیم؛ سپس دو خانوادهٔ گره‌محورِ ناپارامتری،
سیاست‌های ثابت ماتریسی و سیاست یکنواخت را بررسی می‌کنیم.
در تمام این حالت‌ها، نرمال‌سازی انتخاب ویژگی روی مجموعهٔ مجاز نمونهٔ جاری انجام می‌شود.
نقش سیاست در محاسبهٔ دقیق امید گشت و بازنمایی همسایگی در
بخش~\ref{cmp:sec:ch3-exact-walk-calculation} و کاربرد این بازنمایی
برای پیش‌بینی برچسب و تعریف مراکز طبقات در
بخش~\ref{cmp:sec:ch3-neighborhood-to-label-prediction} بررسی می‌شود.

\subsection{ماتریس توابع پارامتری}
\label{cmp:subsec:ch3-parametric-feature-policy}
برای محاسبهٔ احتمال انتخاب ویژگی در رابطه~\eqref{cmp:eq:ch3-full-transition}،
امتیاز هر انتخاب را از توابعی به‌دست می‌آوریم که به
\emph{فرضیهٔ طبقهٔ مبدأ} و مقادیر ویژگی‌های \emph{نمونهٔ جاری} وابسته‌اند.
در هر گشت، نمونهٔ مبدأ $X_s$ و فرضیهٔ $c_s$ ثابت می‌مانند،
اما با ورود به نمونهٔ $X_i$ امتیازهای افقی از مقادیر ویژگی‌های
همان نمونه محاسبه می‌شوند؛ بنابراین فرضیهٔ مبدأ با برچسب
نمونهٔ جاری $y_i$ جایگزین نمی‌شود.

فرض کنیم مجموعه ویژگی‌های کل داده $A_D$ دارای $d$ عضو باشد.
برای هر مقدار ممکن طبقه فرضی مبدأ $c_s$، توابع میانی $F_{c_s,a,b}(\cdot,\cdot;\Theta)$ برای
$a,b\in A_D$ و توابع آغاز $F_{c_s,b}(\cdot;\Theta)$ برای $b\in A_D$
را تعریف می‌کنیم. با قرارداد $F_{c_s,\bot,b}(\cdot;\Theta):=F_{c_s,b}(\cdot;\Theta)$،
می‌توان این خانواده را مانند یک ماتریس از توابع نمایش داد:
\begin{equation}\label{cmp:eq:ch3-parametric-score-matrix}
\boxed{
\mathcal F_{c_s}(\Theta)=
\bigl[F_{c_s,a,b}\bigr]_{a\in A_D\cup\{\bot\},\,b\in A_D},
\qquad
\operatorname{size}\bigl(\mathcal F_{c_s}(\Theta)\bigr)=(d+1)\times d.
}
\end{equation}
این یک \emph{ماتریس توابع} است، نه ماتریسی از احتمال‌ها یا امتیازهای
ثابت. پارامترهای $\Theta$ می‌توانند میان توابع مختلف مشترک یا مستقل باشند.
هر سطر میانی به ویژگی انتخاب‌شده در گام قبلی و هر ستون به ویژگی پیشنهادی
برای گام بعدی اختصاص دارد. سطر اضافی $\bot$ تنها در آغاز گشت به کار می‌رود،
چون در آن لحظه هنوز ویژگی قبلی وجود ندارد.

در نمونه جاری $X_i$، وقتی ویژگی قبلی $a$ و ویژگی پیشنهادی $b$
هر دو موجود باشند، امتیاز عددی انتقال از ارزیابی تابع میانی حاصل می‌شود:
\begin{equation}\label{cmp:eq:ch3-parametric-feature-score}
\boxed{
z_{c_s,a,b}(X_i;\Theta)=
F_{c_s,a,b}\bigl(a(X_i),b(X_i);\Theta\bigr),
\qquad a,b\in A_{X_i}.
}
\end{equation}
اما در \emph{نخستین} انتخاب از نمونه مبدأ $X_s$، تنها مقدار ویژگی مقصد
لازم است:
\begin{equation}\label{cmp:eq:ch3-parametric-start-score}
\boxed{
z_{c_s,\bot,b}(X_s;\Theta)=F_{c_s,b}\bigl(b(X_s);\Theta\bigr),
\qquad b\in A_{X_s}.
}
\end{equation}
امتیازهای به‌دست‌آمده را می‌توان در یک ماتریس عددی موضعی
$Z_{c_s}(X_i;\Theta)$ قرار داد؛ تنها درایه‌های مربوط به ویژگی‌های موجود
و انتقال‌های مجاز ارزیابی می‌شوند. مهم آن است که این ماتریس عددی
با تغییر نمونه جاری تغییر می‌کند: $Z_{c_s}(X_i;\Theta)$ لزوماً
با $Z_{c_s}(X_j;\Theta)$ برابر نیست.

مطابق محدودیت‌های بخش~\ref{cmp:subsec:ch3-walk}،
اگر در وضعیت جاری $X_i$ قرار داشته باشیم،
ویژگی‌ای مجاز است که پس از حذف نمونه جاری و مبدأ
هنوز دست‌کم یک مقصد عمودی داشته باشد:
\begin{equation}\label{cmp:eq:ch3-feature-eligibility}
B_s(X_i)=\bigl\{b\in A_{X_i}\mid
D_b\setminus\{X_i,X_s\}\ne\varnothing\bigr\}
=A_{X_i}.
\end{equation}
برابری اخیر دقیقاً از \emph{فرض پنجم}
($|D_b|\ge3$ برای هر $b$) نتیجه می‌شود.
چون $b\in A_{X_i}$، نمونه جاری $X_i$ در $D_b$ قرار دارد؛
حذف $X_i$ و $X_s$ حداکثر \emph{دو نمونه متمایز}
از $D_b$ را کنار می‌گذارد، بنابراین
$\bigl|D_b\setminus\{X_i,X_s\}\bigr|
\ge |D_b|-2\ge1$.
در گام نخست که $i=s$ است فقط یک نمونه حذف می‌شود،
پس نتیجه حتی قوی‌تر است. این استدلال با فرض مثبت بودن
وزن‌های عمودی، نرمال‌سازی انتقال را نیز ممکن می‌سازد.
پس تحت فرض پنجم، مجموعهٔ مجاز انتخاب ویژگی در هر وضعیت
همان $A_{X_i}$ است؛ بدون این فرض، باید از $B_s(X_i)$ استفاده کرد.

در آغاز گشت، سطر $\bot$ از امتیازهای نمونه $X_s$ و ماسک
$A_{X_s}$ استفاده می‌کند. در گام‌های بعدی، اگر گشت به وضعیت
$(X_i,a)$ رسیده باشد، $a\in A_{X_i}$ معلوم است و تنها
امتیازهای $z_{c_s,a,b}(X_i;\Theta)$ برای $b\in A_{X_i}$
ارزیابی و نرمال می‌شوند. در هر دو حالت، نرمال‌سازی با بیشینهٔ هموار روی
\emph{مجموعه مجاز همان نمونه جاری} صورت می‌گیرد:
\begin{equation}\label{cmp:eq:ch3-fixed-policy-softmax}
\boxed{
\pi_s^{(c_s)}(b\mid i,a)
=[\Pi_i^{(s,c_s)}]_{ab}
=\frac{\exp\bigl(z_{c_s,a,b}(X_i;\Theta)\bigr)}
{\sum\limits_{u\in A_{X_i}}
\exp\bigl(z_{c_s,a,u}(X_i;\Theta)\bigr)},
\qquad b\in A_{X_i}.
}
\end{equation}
ماتریس احتمال سیاست $\Pi_i^{(s,c_s)}$ برای نمونهٔ جاری
$X_i$، دارای $d+1$ سطر (ویژگی‌های قبلی به‌علاوهٔ سطر آغاز
$\bot$) و $d$ ستون (ویژگی‌های پیشنهادی) است؛ بنابراین
$\Pi_i^{(s,c_s)}\in\mathbb R^{(d+1)\times d}$.
درایه‌های مربوط به ویژگی‌های ناموجود در $X_i$ صفر
در نظر گرفته می‌شوند و فقط سطرهای قابل استفاده،
یعنی سطرهای ویژگی‌های موجود و در گام نخست سطر $\bot$،
به‌صورت توزیع احتمال نرمال می‌شوند.

برای یک سطر سیاست با امتیازهای ثابت، اگر ماسک مجاز از $A$ به $A'$
تغییر کند، نسبت احتمال انتخاب دو ویژگیِ مشترک $b,u\in A\cap A'$
ثابت می‌ماند:
\begin{equation}\label{cmp:eq:ch3-masked-softmax-log-ratio}
\log\frac{p_A(b)}{p_A(u)}
=\log\frac{p_{A'}(b)}{p_{A'}(u)}
=z_b-z_u.
\end{equation}
در اینجا $p_A$ و $p_{A'}$ حاصل نرمال‌سازی \emph{همان امتیازها}
روی دو مجموعهٔ مجازند؛ بنابراین احتمال‌های مطلق ممکن است تغییر کنند،
ولی نسبت و ترتیب دو گزینهٔ باقی‌مانده تغییر نمی‌کند.
این نتیجه در صورت تغییر خود امتیازها لزوماً برقرار نیست.
اثبات و شرایط دقیق آن در پیوست~\ref{app:ch3-softmax-properties}،
رابطه~\eqref{eq:app-ch3-softmax-mask-ratio} آمده است.

در رابطه~\eqref{cmp:eq:ch3-fixed-policy-softmax}، برای $i=s$ و $a=\bot$ از امتیاز شروع
رابطه~\eqref{cmp:eq:ch3-parametric-start-score} و در سایر وضعیت‌ها
از امتیاز میانی رابطه~\eqref{cmp:eq:ch3-parametric-feature-score}
استفاده می‌شود. وابستگی به $\Theta$ در نماد $\Pi_i^{(s,c_s)}$
ضمنی است. هر سطر قابل استفاده یک توزیع احتمال معتبر است:
\begin{equation}\label{cmp:eq:ch3-policy-stochastic}
[\Pi_i^{(s,c_s)}]_{ab}\ge0,\qquad
\sum_{b\in A_{X_i}}[\Pi_i^{(s,c_s)}]_{ab}=1.
\end{equation}
ناموجود بودن ویژگی $b$ در \emph{نمونه آغازین} $X_s$، مانع
انتخاب آن در یک \emph{نمونه میانی} $X_i$ که $b\in A_{X_i}$
باشد نیست؛ ماسک مبدأ در گام‌های میانی دوباره اعمال نمی‌شود.
همچنین برای ویژگی ناموجود در خود نمونه جاری، تابع امتیاز نیازی
به ارزیابی یا جایگذاری مقدار مصنوعی ندارد. احتمال انتقال کامل
به نمونه مقصد نیز مطابق رابطه~\eqref{cmp:eq:ch3-full-transition}
با ضرب احتمال انتخاب $b$ در توزیع عمودی $q_b^{(s)}(j\mid i)$
به‌دست می‌آید.

\paragraph{تفاوت استفاده از سیاست در آموزش و آزمون.}
در آموزش، برچسب واقعی نمونهٔ آغازین برای ساخت مرکز طبقهٔ متناظر
و محاسبهٔ زیان نظارت‌شده معلوم است. در ساخت \emph{مراکز شرطی}،
برای هر نمونهٔ آموزشی $c_s=y_s$ قرار می‌گیرد؛ اما اگر قاعدهٔ تصمیم
\emph{استنتاج مبتنی بر فرضیهٔ طبقه} باشد، برای محاسبهٔ حاشیه و
پیش‌بینی \emph{هر نمونه، حتی نمونهٔ آموزشی} باید تمام فرضیه‌های
$c_s\in\{1,\ldots,K\}$ ارزیابی شوند. در آزمون نیز برچسب نمونهٔ
جدید ناشناخته است و نمی‌توان آن را برای انتخاب سیاست به کار برد.
دو شیوهٔ استنتاج را در نظر می‌گیریم.

راه نخست، تشکیل یک \emph{سیاست میانگین} از سیاست‌های شرطی طبقه‌هاست. اگر $\pi_s^{(c_s)}(b\mid i,a)$ احتمال انتخاب ویژگی تحت فرض طبقه مبدأ $c_s$ و $\rho_{c_s}$ وزن‌های نامنفی با مجموع یک باشند، می‌توان نوشت:
\begin{equation}\label{cmp:eq:ch3-mean-policy-inference}
\overline\pi_s(b\mid i,a)=
\sum_{c_s=1}^{K}\rho_{c_s}\,\pi_s^{(c_s)}(b\mid i,a),
\qquad \sum_{c_s=1}^{K}\rho_{c_s}=1.
\end{equation}
این وزن‌ها می‌توانند برای مثال از فراوانی طبقه‌های مجموعه آموزش به‌دست آیند. در این روش، سیاست حاصل به برچسب ناشناخته آزمون وابسته نیست و مراکز طبقه‌ها نیز باید با همین سیاست میانگین محاسبه شوند. میانگین‌گیری در اینجا روی \emph{احتمال‌های انتخاب ویژگی} صورت می‌گیرد و الزاماً معادل میانگین گرفتن از امتیازهای خام و سپس اعمال نرمال‌سازی با بیشینهٔ هموار نیست.

همچنین اجرای گشت با سیاست میانگین $\overline\pi_s$ عموماً با
میانگین‌گیری از بازنمایی‌های حاصل از گشت‌های جداگانه،
$\sum_{c=1}^{K}\rho_c\,\mathscr Y_s^{(c)}$، برابر نیست.
در حالت دوم، یک فرضیهٔ طبقه در سراسر هر گشت ثابت است و
بازنمایی‌ها پس از محاسبه میانگین گرفته می‌شوند؛
اما سیاست میانگین در \emph{هر گام} احتمال انتخاب ویژگی را
از ترکیب سیاست‌های شرطی به‌دست می‌آورد و در گشت چندگامی
احتمال مسیر به حاصل‌ضرب احتمال‌های انتقال وابسته است.
سیاست میانگین را نباید با \emph{سیاست مشترک} نیز یکسان دانست:
سیاست مشترک پارامترهای مستقل خود را دارد و تنها از طبقهٔ
مبدأ مستقل است. این تمایزها هم برای سیاست‌های مبتنی بر
توابع پارامتری و هم برای سیاست‌های ثابت ماتریسی برقرارند.


راه دوم، \emph{استنتاج مبتنی بر فرضیه طبقه} است: برای هر مقدار نامزد $c_s$، گشت از نمونه آزمون را \emph{با فرض آزمایشی طبقه مبدأ $c_s$} اجرا کرده و بازنمایی متناظر $\mathscr Y_x^{(c_s)}$ را می‌سازیم. سپس این بازنمایی را با مرکز $\mathcal V_{c_s}^{(c_s)}$ که از نمونه‌های آموزشی دارای برچسب $c_s$ و با همان سیاست شرطی تشکیل شده است مقایسه می‌کنیم. با انتخاب معیار فاصله مشترک $d$، قاعده تصمیم چنین خواهد بود:
\begin{equation}\label{cmp:eq:ch3-class-hypothesis-inference}
\widehat y(x)=\arg\min_{c_s\in\{1,\ldots,K\}}
d\bigl(\mathscr Y_x^{(c_s)},\mathcal V_{c_s}^{(c_s)}\bigr).
\end{equation}
در روش دوم، مقدار $c_s$ فقط یک \emph{فرضیه} برای اجرای گشت است و برچسب واقعی نمونه آزمون استفاده نمی‌شود. این روش وابستگی سیاست به طبقه آغازین را حفظ می‌کند، اما مقایسه امتیازهای حاصل از سیاست‌های متفاوت نیازمند ارزیابی تجربی است و برتری آن نسبت به سیاست میانگین از تعریف‌ها به‌تنهایی نتیجه نمی‌شود. قواعد تصمیم مبتنی بر یک بازنمایی واحد در ادامه فصل برای سیاست مشترک یا میانگین به همان شکل نوشته می‌شوند.

\subsection{خانواده‌های سیاست ناپارامتری}
\label{cmp:subsec:ch3-node-policy-families}
در ماتریس توابع پارامتریِ رابطه~\eqref{cmp:eq:ch3-parametric-score-matrix}،
امتیازها از مقادیر ویژگی‌های نمونهٔ جاری محاسبه می‌شوند.
در مقابل، در دو خانوادهٔ گره‌محور \emph{ناپارامتری}،
برای هر گرهٔ مرجع، امتیازهای عددیِ آزاد و قابل‌آموزش تعریف می‌کنیم.
منظور از «ناپارامتری» در اینجا آن است که تعداد این امتیازهای آزاد
با شمار نمونه‌های مرجع رشد می‌کند، نه اینکه مدل فاقد پارامتر باشد.
اگر $n$ شمار نمونه‌های مرجع، $d$ شمار ویژگی‌ها و $K$ شمار طبقات باشد،
آرایه‌های امتیاز دو خانوادهٔ ناپارامتری ابعاد زیر را دارند:
\begin{equation}\label{cmp:eq:ch3-policy-family-shapes}
\begin{aligned}
Z^{\mathrm{node}}&\in\mathbb R^{n\times(d+1)\times d},\\
Z^{\mathrm{node,class}}&\in\mathbb R^{K\times n\times(d+1)\times d}.
\end{aligned}
\end{equation}

در خانوادهٔ \emph{گره‌محور} ($Z^{\mathrm{node}}$)، اندیس گره
به نمونهٔ جاری اشاره دارد و سیاست به طبقهٔ مبدأ وابسته نیست.
در خانوادهٔ \emph{گره--طبقه} ($Z^{\mathrm{node,class}}$)،
علاوه بر گرهٔ جاری، اندیس $c_s$ نیز سیاست را تعیین می‌کند؛
این اندیس به طبقهٔ \emph{نمونهٔ آغازین} تعلق دارد و در طول
گشت ثابت می‌ماند، نه اینکه از برچسب نمونه‌های میانی گرفته شود.
در ساخت مراکز طبقات از نمونه‌های آموزشیِ برچسب‌دار،
$c_s=y_s$ است؛ اما برای پیش‌بینی نمونهٔ جدید، $c_s$ برچسبی
معلوم نیست و باید مطابق یکی از شیوه‌های استنتاج
زیربخش~\ref{cmp:subsec:ch3-parametric-feature-policy}
با آن برخورد کرد.
سطر اضافی $\bot$ برای انتخاب نخستین ویژگی از نمونهٔ آغازین
و سطرهای دیگر برای انتقال‌های میانی‌اند. در همهٔ حالت‌ها،
امتیازها فقط روی ویژگی‌های موجود و مجاز نمونهٔ جاری نرمال می‌شوند.

برای سیاست‌های آزاد $Z^{\mathrm{node}}$ و
$Z^{\mathrm{node,class}}$، هنگام ورود نمونهٔ جدید $X_\star$
ماتریس مختص آن نمونه در میان پارامترهای آموخته‌شده وجود ندارد.
انتقال‌های میانی پس از ورود به گره‌های مرجع همچنان تعریف‌شده‌اند،
اما توزیع انتخاب ویژگی \emph{در آغاز گشت} به یک قاعدهٔ
برون‌گراف نیاز دارد. این مسئله حتی برای خانوادهٔ
$Z^{\mathrm{node}}$، که به برچسب مبدأ وابسته نیست، وجود دارد.
در خانوادهٔ $Z^{\mathrm{node,class}}$ مسئلهٔ دیگری نیز مطرح است:
برچسب واقعی $X_\star$ مجهول است و نمی‌توان آن را برای تعیین
$c_s$ به کار برد. پس علاوه بر تعیین سیاست آغاز برون‌گراف،
باید فرضیه‌های طبقه را جداگانه ارزیابی کرد یا از سیاست میانگین
با مراکز سازگار با آن استفاده کرد.
بدون تعیین قاعدهٔ آغاز برون‌گراف، دقت آزمون این
خانواده‌ها \emph{تعریف‌نشده} است، نه اینکه صفر باشد.
این تمایز، خانوادهٔ آزاد را از سیاست تابعی و قابل‌تعمیم جدا می‌کند.

یک خانواده مقید میانی را نیز با \emph{آغاز مشترکِ شرطی طبقه}
و \emph{انتقال‌های میانیِ اختصاصی گره} تعریف می‌کنیم:
\begin{equation}\label{cmp:eq:ch3-class-tied-start}
\boxed{
Z_{c_s,j,\bot b}=A_{c_s,b},\qquad
Z_{c_s,j,ab}=B_{c_s,j,a,b},\quad a,b\in A_D.
}
\end{equation}
در اینجا $c_s$ اندیس فرضیه آغازین است و $A\in\mathbb R^{K\times d}$ برای همه نمونه‌های مبدأ
مشترک است، اما $B\in\mathbb R^{K\times n\times d\times d}$
به گره جاری بستگی دارد. تعداد پارامترهای خام مستقل برابر
$Kd+Knd^2$ است؛ در مقایسه، خانواده گره--طبقه آزاد
$Kn(d+1)d$ امتیاز خام دارد. در هر دو مدل، ماسک ویژگی‌ها
پس از امتیازگذاری اعمال می‌شود. قید بالا باید در \emph{فرایند
آموزش} برقرار باشد؛ میانگین‌گرفتن سطر آغاز از سیاست آزادِ
ازپیش‌آموخته‌شده، عموماً راه‌حل همان مسئله بهینه‌سازی مقید
نیست. در این خانوادهٔ مقید، سطر آغاز $A_{c_s,:}$ به گرهٔ
آغازین وابسته نیست و ازاین‌رو برای نمونهٔ جدید نیز تعریف شده است؛
اما \emph{انتخاب اندیس طبقهٔ آن} همچنان باید مشخص شود.
هنگام ساخت مراکز آموزشی، $c_s=y_s$ است. هنگام پیش‌بینی
نمونهٔ جدید، یا برای هر فرضیهٔ $c_s$ سطر آغاز
$A_{c_s,:}$ و انتقال‌های میانی متناظر به کار می‌روند،
یا احتمال‌های انتخاب ویژگیِ سیاست‌های شرطی، مطابق
زیربخش~\ref{cmp:subsec:ch3-parametric-feature-policy}،
میانگین گرفته می‌شوند و مراکز نیز با همان سیاست میانگین
ساخته می‌شوند. در حالت دوم، میانگین‌گیری روی احتمال‌ها
\emph{پس از اعمال ماسک و نرمال‌سازی با بیشینهٔ هموار} انجام می‌شود، نه روی
سطرهای امتیاز خام $A_{c_s,:}$.
بنابراین قید آغاز مشترک مشکل نبودن سطر آغاز برای گرهٔ جدید
را رفع می‌کند، نه مسئلهٔ مجهول‌بودن طبقهٔ آن را؛
سیاست میانی نیز همچنان به گره‌های مرجع آموزش وابسته است.

\subsection{سیاست ثابت ماتریسی}
\label{cmp:subsec:ch3-shared-matrix-policy}
در سیاست ثابت ماتریسی، امتیازهای انتخاب ویژگی به
مقادیر ویژگی‌های نمونهٔ جاری یا اندیس گرهٔ جاری وابسته نیستند.
سطرهای آرایهٔ امتیاز با ویژگی قبلی $a$ و ستون‌ها با
ویژگی پیشنهادی $b$ مشخص می‌شوند؛ سطر اضافی $\bot$
برای انتخاب نخستین ویژگی است. دو خانوادهٔ ثابت را
به‌صورت زیر تعریف می‌کنیم:
\begin{equation}\label{cmp:eq:ch3-fixed-policy-matrix}
\begin{aligned}
Z^{\mathrm{shared}}=Z
&\in\mathbb R^{(d+1)\times d},\\
Z^{\mathrm{class}}=(Z^{(c)})_{c=1}^{K}
&\in\mathbb R^{K\times(d+1)\times d}.
\end{aligned}
\end{equation}
در هر دو خانواده، اندیس سطر $a\in A_D\cup\{\bot\}$ و
اندیس ستون $b\in A_D$ است. امتیاز سطر $a$ و ستون $b$
در سیاست مشترک با $z_{ab}=[Z]_{ab}$ و در سیاست
شرطی طبقهٔ $c$ با $z^{(c)}_{ab}=[Z^{(c)}]_{ab}$
نشان داده می‌شود.

در خانوادهٔ \emph{مشترک} ($Z^{\mathrm{shared}}$)، یک ماتریس
امتیاز برای همهٔ نمونه‌ها و همهٔ فرضیه‌های طبقهٔ مبدأ به کار
می‌رود. در خانوادهٔ \emph{شرطی طبقه} ($Z^{\mathrm{class}}$)،
برای هر فرضیهٔ $c_s$ ماتریس $Z^{(c_s)}$ انتخاب می‌شود؛
$c_s$ به طبقهٔ نمونهٔ آغازین تعلق دارد و در طول گشت ثابت است.

در وضعیت $(X_i,a)$، امتیاز انتخاب ویژگی مجاز $b\in A_{X_i}$
در خانوادهٔ مشترک برابر $z_{ab}$ و در خانوادهٔ شرطی طبقه
برابر $z^{(c_s)}_{ab}$ است. در نخستین انتخاب، $a=\bot$
قرار می‌گیرد. احتمال‌های انتخاب با نرمال‌سازی با بیشینهٔ هموار
روی ویژگی‌های مجاز \emph{نمونهٔ جاری}، مطابق
رابطه~\eqref{cmp:eq:ch3-fixed-policy-softmax}، به دست می‌آیند.
بنابراین، با وجود ثابت‌بودن امتیازهای خام در طول گشت،
ماتریس احتمال $\Pi_i^{(s,c_s)}$ می‌تواند به‌دلیل تفاوت
مجموعهٔ ویژگی‌های موجود از نمونه‌ای به نمونهٔ دیگر تغییر کند.

\subsection{سیاست یکنواخت}
\label{cmp:subsec:ch3-uniform-policy}
حالت خاصی را در نظر می‌گیریم که در آن تمام ویژگی‌های قابل استفاده با احتمال یکسان انتخاب شوند. این سیاست با انتخاب $Z=0$ و امتیازهای یکسان برای ویژگی‌های شروع، یا به‌طور کلی برابر گرفتن امتیازهای هر سطر روی ویژگی‌های مجاز، از رابطه~\eqref{cmp:eq:ch3-fixed-policy-softmax} به‌دست می‌آید. در این حالت، احتمال انتخاب ویژگی $b$ به ویژگی قبلی $a$ وابسته نیست و برای هر وضعیت دارای مجموعه مجاز ناتهی، برابر است با:
\begin{equation}\label{cmp:eq:ch3-uniform-policy}
\pi_s^{U}(b\mid i,a)=\frac{1}{|A_{X_i}|},
\qquad b\in A_{X_i}.
\end{equation}
استقلال انتخاب ویژگی از ویژگی قبلی اجازه می‌دهد دو مرحله انتخاب ویژگی و انتقال عمودی را در یک انتقال روی \emph{فضای نمونه‌ها} ادغام کنیم. با استفاده از توزیع انتقال عمودی $q_b^{(s)}$ در رابطه~\eqref{cmp:eq:ch3-vertical-transition}، ماتریس انتقال متناظر با نمونه آغازین $X_s$ را به‌صورت زیر تعریف می‌کنیم:
\begin{equation}\label{cmp:eq:ch3-uniform-T}
P_s^{U}(i,j)=\frac{1}{|A_{X_i}|}
\sum_{b\in A_{X_i}}q_b^{(s)}(j\mid i).
\end{equation}
ماتریس $P_s^{U}$ یک ماتریس گذار سطری است و درایه‌های آن نامنفی‌اند و مجموع هر سطر آن برابر یک است. در نتیجه، هرچند انتخاب ویژگی در سیاست یکنواخت انجام می‌شود، توزیع انتخاب \emph{نمونه مقصد} لزوماً یکنواخت نیست؛ زیرا توزیع‌های $q_b^{(s)}$ به وزن شباهت مقادیر ویژگی‌ها وابسته‌اند. به بیان دقیق‌تر، پس از انتخاب هر ویژگی، نمونه‌های نزدیک‌تر از دید همان ویژگی وزن انتقال بیشتری دارند و $P_s^{U}$ میانگین توزیع‌های انتقال ویژگی‌های قابل استفاده را نشان می‌دهد.

چون انتخاب ویژگی بعدی به ویژگی قبلی وابسته نیست، برای نمونه آغازین ثابت، حرکت روی نمونه‌ها با ماتریس $P_s^{U}$ مارکوفی است. بنابراین توزیع مشاهده در گام $t$ام، با نماد رابطه~\eqref{cmp:eq:ch3-step-observation-distribution}، برابر است با:
\begin{equation}\label{cmp:eq:ch3-uniform-step-distribution}
P_{X_s,t}=e_s(P_s^{U})^t,\qquad t\ge1.
\end{equation}
این برابری ویژه سیاستی است که انتخاب ویژگی در آن به ویژگی قبلی وابسته نیست؛ اثبات با شرط‌کردن روی نمونه جاری و جمع‌زدن ویژگی‌های انتخابی در پیوست~\ref{app:ch3-matrix-walk-proofs} آمده است. در حالت عمومی چنین کاهشی به یک ماتریس گذار روی نمونه‌ها الزاماً ممکن نیست. همچنین $P_s^{U}$ به نمونه آغازین $X_s$ وابسته است؛ زیرا بازگشت به این نمونه در تمام مراحل ممنوع شده و بنابراین یک ماتریس گذار مشترک برای همه نمونه‌های آغازین نداریم.

با استفاده از رابطه~\eqref{cmp:eq:ch3-h-alpha}، توزیع همسایگی حاصل از سیاست یکنواخت را می‌توان به‌صورت زیر نوشت:
\begin{equation}\label{cmp:eq:ch3-uniform-h}
h_s^{U}=\sum_{t=1}^{\infty}
\alpha_t(\beta,\lambda)\,e_s(P_s^{U})^t.
\end{equation}
این رابطه نشان می‌دهد بازنمایی همسایگی حتی بدون یادگیری سیاست انتخاب ویژگی، بر اساس ترکیبی از شباهت‌های محلی و ارتباطات چندگامی ساخته می‌شود. در گام اول نمونه‌های مشابه از دید ویژگی انتخاب‌شده احتمال مشاهده بیشتری دارند و در گام‌های بعدی نمونه‌های واسط نیز در ساخت بازنمایی مشارکت می‌کنند. بنابراین سیاست یکنواخت یک مبنای طبیعی برای روش پیشنهادی است، اما لزوماً اهمیت متفاوت ویژگی‌ها را برای هدف یادگیری مورد نظر منعکس نمی‌کند.




\section{محاسبهٔ دقیق امید گشت تصادفی در گراف نمونه--ویژگی}
\label{cmp:sec:ch3-exact-walk-calculation}
در بخش‌های پیشین، گشت تصادفی روی گراف نمونه--ویژگی،
قانون انتخاب ویژگی و بازنمایی همسایگی حاصل از مشاهدهٔ
نمونه‌های مقصد تعریف شدند. محاسبهٔ این بازنمایی با
نمونه‌گیری مسیرهای تصادفی امکان‌پذیر است، اما خطای
مونت‌کارلو دارد و برای کاهش آن به تکرار گشت‌ها
نیاز پیدا می‌کند. در این بخش نشان می‌دهیم که
برای گراف مرجع، مبدأ و سیاست ثابت، می‌توان امید
بازنمایی را مستقیماً و بدون نمونه‌گیری مسیر
محاسبه کرد.

برای این منظور، گشت را \emph{پس از نخستین انتقال}
به‌صورت یک زنجیرهٔ مارکوف روی حالت‌های
نمونه--ویژگی بیان می‌کنیم؛ زیرا انتخاب ویژگی
بعدی ممکن است به ویژگی مرحلهٔ پیشین وابسته باشد.
ابتدا فضای حالت، عامل‌های افقی و عمودی ماتریس
انتقال، توزیع نخستین گام و عملگرهای تجمیع را
معرفی می‌کنیم. سپس بازگشت پیش‌رونده و
پس‌روندهٔ محاسبهٔ امید، تقریب کنترل‌شده با
عمق محدود، کران‌های محاسباتی و امکان تنک‌سازی
انتقال‌های عمودی را بررسی می‌کنیم. در پایان،
شبه‌کدهای محاسبهٔ دقیق امید برای گراف کامل یا
تنک‌شده ارائه می‌شوند؛ مقصود از «دقیق» در هر
حالت، حذف خطای نمونه‌گیری نسبت به \emph{قانون
انتقال همان گراف} است.

\subsection{ماتریس بلوکی انتقال روی فضای نمونه--ویژگی}
\label{cmp:subsec:ch3-block-matrix-transition}
برای یک نمونه آغازین ثابت $X_s$، ممنوعیت بازگشت به مبدأ را می‌توان مستقیماً در فضای حالت اعمال کرد. مجموعه حالت‌های پس از نخستین انتقال را چنین تعریف می‌کنیم:
\begin{equation}\label{cmp:eq:ch3-state-space}
\mathcal S_s=\{(i,a)\mid i\ne s,\ a\in A_{X_i}\},
\qquad m_s=|\mathcal S_s|.
\end{equation}
در این فضا، نمونه مبدأ حضور ندارد، اما بازدید مجدد از سایر نمونه‌ها مجاز است.
فضای حالت تهی نیست: طبق فرض پنجم، هر ویژگی موجود
در مجموعهٔ داده دست‌کم در سه نمونه مشاهده شده است؛
همچنین اتصال‌پذیری فرض سوم اقتضا می‌کند که هر
نمونه دست‌کم یک ویژگی موجود داشته باشد، وگرنه
از دیگر نمونه‌ها جدا می‌ماند. بنابراین پس از
کنارگذاشتن نمونهٔ مبدأ نیز دست‌کم یک حالت
نمونه--ویژگی باقی می‌ماند.
در نتیجه، برای $s$ و $c_s$ ثابت و سیاستی که در طول
گشت تغییر نمی‌کند، یک ماتریس انتقال همگن روی
$\mathcal S_s$ وجود دارد. منظور از ثابت بودن،
ثابت‌ماندن قانون سیاست و فرض طبقه \emph{مبدأ}
است، نه یکسان‌بودن سطرهای سیاست برای تمام
نمونه‌های جاری یا ویژگی‌های ورودی.

در نمایش بلوکی با ترتیب نمونه‌ها، برای $i,j\ne s$ ماتریس قطری زیر را، با صفرگذاشتن درایه‌های مربوط به ویژگی‌های ناموجود، تعریف می‌کنیم:
\begin{equation}\label{cmp:eq:ch3-block-diagonal-D}
D_{ij}^{(s)}
=\operatorname{diag}\!\left(q_b^{(s)}(j\mid i)\right)_{b\in A_D}
\in\mathbb R^{d\times d}.
\end{equation}
در این زیر‌بخش، \emph{نمونه مبدأ $s$ و طبقه فرضی آن $c_s$ ثابت‌اند}.
برای جلوگیری از ازدحام اندیس‌ها در توان‌های ماتریسی،
از نماد کوتاه $\Pi_i^{(s)}$ برای
\emph{زیرماتریس سطرهای انتقال میانی} استفاده می‌کنیم،
یعنی سطر آغاز $\bot$ را که فقط در مبدأ کاربرد دارد،
از $\Pi_i^{(s,c_s)}$ کنار می‌گذاریم:
\begin{equation}\label{cmp:eq:ch3-intermediate-policy-matrix}
\Pi_i^{(s)}
:=\bigl([\Pi_i^{(s,c_s)}]_{ab}\bigr)_{a,b\in A_D}
\in\mathbb R^{d\times d}.
\end{equation}
اگر $d_i=|A_{X_i}|$ باشد، زیرماتریس مؤثرِ انتقال‌های
میانی روی ویژگی‌های موجود نمونهٔ جاری، اندازهٔ
$d_i\times d_i$ دارد. برای ساخت بلوک‌های انتقال،
این زیرماتریس را با افزودن سطرها و ستون‌های صفر
برای ویژگی‌های ناموجود به فضای مشترک $A_D$ گسترش
می‌دهیم؛ ازاین‌رو اندازهٔ $\Pi_i^{(s)}$ در نمایش
بلوکی $d\times d$ است، نه لزوماً $d_i\times d_i$.
درایه‌های مربوط به $a,b\in A_{X_i}$ برابر
$[\Pi_i^{(s)}]_{ab}=\pi_s^{(c_s)}(b\mid i,a)$ هستند.
باید توجه داشت که $\Pi_i^{(s)}$ و
$\Pi_i^{(s,c_s)}$ \emph{به‌صورت ماتریسی برابر نیستند}:
اولی فقط سطرهای میانی را دارد و دومی سطر
شروع $\bot$ را نیز شامل می‌شود.
\emph{این حذف نماد اختیاری نیست}: تمام ماتریس‌های
$\Pi_i^{(s)}$، $\mathsf P_s$ و بردار $\rho_s$
به همان فرض \emph{ثابت} $c_s$ وابسته‌اند و اگر
فرض طبقه مبدأ تغییر کند، باید آن‌ها را دوباره محاسبه کرد.
در این نمایش، تنها سطرهای $a\in A_{X_i}$
در فضای حالت استفاده خواهند شد.
رابطه~\eqref{cmp:eq:ch3-full-transition}
اکنون به‌صورت بلوکی درمی‌آید:
\begin{equation}\label{cmp:eq:ch3-block-Tij}
\boxed{\mathsf P_{s,ij}=\Pi_i^{(s)}D_{ij}^{(s)},\qquad
[\mathsf P_{s,ij}]_{ab}
=\pi_s^{(c_s)}(b\mid i,a)q_b^{(s)}(j\mid i).}
\end{equation}
در این نمایش، ابتدا بلوک‌ها را روی فضای مشترک ویژگی‌ها
$A_D$ می‌سازیم و سپس سطرها و ستون‌ها را به حالت‌های
واقعاً موجود در $\mathcal S_s$ محدود می‌کنیم.
این محدودسازی با ضرب عامل‌ها سازگار است؛ زیرا در
تمام درایه‌های متناظر با ویژگی‌های ناموجود، انتقال
صفر تعریف شده است.
اگر $F_s=\operatorname{blockdiag}_{i\ne s}(\Pi_i^{(s)})$
عامل انتخاب افقی ویژگی و
$S_s=[D_{ij}^{(s)}]_{i,j\ne s}$
عامل انتقال عمودی باشد (هر دو پس از محدودسازی
به $\mathcal S_s$)، آنگاه:
\begin{equation}\label{cmp:eq:ch3-block-factorization}
\boxed{\mathsf P_s=F_sS_s,\qquad
\mathsf P_s\in\mathbb R^{m_s\times m_s}.}
\end{equation}
بنابراین $F_s$ در ترتیب نمونه‌ای بلوکی‌قطری است و بلوک‌های $S_s$ قطری‌اند؛ با مرتب‌کردن حالت‌ها برحسب ویژگی، خود $S_s$ بلوکی‌قطری می‌شود. هر سطر قابل دسترس $\mathsf P_s$ مجموع یک دارد؛
زیرا برای وضعیت $(i,a)$:
\[
\sum_{(j,b)\in\mathcal S_s}[\mathsf P_s]_{(i,a),(j,b)}
=\sum_{b\in A_{X_i}}\pi_s^{(c_s)}(b\mid i,a)
\underbrace{\sum_{j\ne i,s}q_b^{(s)}(j\mid i)}_{=1}
=1.
\]
این برابری هم به فرض پوشش ویژگی‌ها و هم به
نرمال‌بودن توزیع حاصل از نرمال‌سازی با بیشینهٔ هموار وابسته است. صورت بلوکی
و نتیجه فوق در پیوست~\ref{app:ch3-matrix-walk-proofs}
نیز با جزئیات اثبات می‌شوند.

\paragraph{تمایز گام نخست.}
از نمونه مبدأ هنوز ویژگی قبلی نداریم و باید سطر آغاز $\bot$ را به کار ببریم. توزیع حالت‌های بلافاصله پس از نخستین انتقال را با بردار سطری $\rho_s\in\mathbb R^{1\times m_s}$ تعریف می‌کنیم:
\begin{equation}\label{cmp:eq:ch3-block-first-step}
\rho_s(j,b)=\pi_s^{(c_s)}(b\mid s,\bot)\,q_b^{(s)}(j\mid s),
\qquad (j,b)\in\mathcal S_s.
\end{equation}
این بردار واقعاً یک توزیع احتمال است؛ چون طبق
فرض پنجم برای هر ویژگی موجود در مبدأ،
$\sum_{j\ne s}q_b^{(s)}(j\mid s)=1$ و از نرمال‌سازی
سطر $\bot$ نیز داریم
\[
\sum_{(j,b)\in\mathcal S_s}\rho_s(j,b)
=\sum_{b\in A_{X_s}}\pi_s^{(c_s)}(b\mid s,\bot)
\underbrace{\sum_{j\ne s}q_b^{(s)}(j\mid s)}_{=1}
=1.
\]
برای تبدیل توزیع حالت‌های $(j,b)$ به توزیع نمونه‌های
مشاهده‌شده و سپس به توزیع برچسب‌ها، دو
\emph{ماتریس تجمیع حالت} تعریف می‌کنیم:
\begin{equation}\label{cmp:eq:ch3-state-aggregation-matrices}
\boxed{
\begin{aligned}
\mathsf M_s^{(X)}
&\in\{0,1\}^{m_s\times n},
&
[\mathsf M_s^{(X)}]_{(j,b),u}
&=\mathbf1\{j=u\},\\
\mathsf M_s^{(Y)}
&=\mathsf M_s^{(X)}Y
\in\{0,1\}^{m_s\times K},
&
[\mathsf M_s^{(Y)}]_{(j,b),c}
&=\mathbf1\{y_j=c\}.
\end{aligned}}
\end{equation}
ماتریس $\mathsf M_s^{(X)}$ مؤلفه‌های مربوط به
ویژگی‌های مختلف \emph{یک نمونه مقصد} را با هم جمع می‌کند،
و $\mathsf M_s^{(Y)}$ همان جرم احتمال را بر اساس
\emph{برچسب نمونه مقصد} تجمیع می‌کند.
اندیس $c$ در درایه $[\mathsf M_s^{(Y)}]_{(j,b),c}$
\emph{برچسب نمونه مقصد} است و با $c_s$،
یعنی \emph{فرض طبقه مبدأ}، تفاوت دارد.
این دو ماتریس صرفاً عملگر تجمیع‌اند،
نه ماتریس سیاست یا بردار خطا؛
هر سطرشان دقیقاً یک درایه یک دارد.
از آنجا که تمام حالت‌های $\mathcal S_s$ دارای $j\ne s$
هستند، حتی اگر سطر برچسب مبدأ در ماتریس آموزشی $Y$
موجود باشد، در $\mathsf M_s^{(Y)}$ و خروجی
گشت آغازشده از $X_s$ سهمی نخواهد داشت.
اکنون توزیع مشاهده در گام $t$ام را بدون
نمونه‌گیری تصادفی می‌توان محاسبه کرد:
\begin{equation}\label{cmp:eq:ch3-block-t-step}
\boxed{P_{X_s,t}
=\rho_s\mathsf P_s^{\,t-1}\mathsf M_s^{(X)},\qquad
P_{X_s,t}Y
=\rho_s\mathsf P_s^{\,t-1}\mathsf M_s^{(Y)}.}
\end{equation}
توان $t-1$ از آن روست که $\rho_s$ \emph{پس از} نخستین گام تعریف شده است؛ استقرا روی تعداد گام‌ها و نقش دو ماتریس تجمیع در پیوست~\ref{app:ch3-matrix-walk-proofs} آمده است. اگر ویژگی بعدی به ویژگی ورودی وابسته نباشد، فضای حالت را می‌توان مانند حالت یکنواختِ رابطه~\eqref{cmp:eq:ch3-uniform-T} به نمونه‌ها فروکاست؛ در حالت عمومی، جمع‌زدن بلوک‌ها به‌تنهایی برای محاسبه صحیح توان‌های چندگامی کافی نیست. همچنین حذف مبدأ باید پیش از نرمال‌سازی مقصدها \emph{در هر ویژگی} انجام شود؛ نرمال‌کردن مستقیم یک ماتریس انتقال تجمیع‌شده عموماً معادل آن نیست.


\paragraph{تجمیع برچسب یک‌گامی هر ویژگی.}
برای تمایز توزیع مقصد \emph{نمونه‌ای} $P_{X_s,t}$ از توزیع برچسب مقصد بعد از انتخاب ویژگی $b$ در وضعیت $i$، دومی را با $\mathsf Q_{i,b}^{(s)}$ نشان می‌دهیم:
\begin{equation}\label{cmp:eq:ch3-feature-one-step-Q}
\boxed{\mathsf Q_{i,b}^{(s)}
=\sum_{j\ne i,s}q_b^{(s)}(j\mid i)Y_{j,:}
\in\mathbb R^{1\times K}.}
\end{equation}
در حالت $(i,a)$، توزیع برچسب مقصد یک گام کامل برابر
$\sum_{b\in A_{X_i}}\pi_s^{(c_s)}(b\mid i,a)\mathsf Q_{i,b}^{(s)}$
است. این رابطه می‌تواند برای ساخت تقریب کم‌هزینه و مقداردهی آغازین سیاست استفاده شود؛ ولی چون شناسه نمونه مقصد و ادامه مسیر را حذف می‌کند، به‌تنهایی بازنمایی دقیق چندگامی یا مراکز حاصل از آن را تعیین نمی‌کند.

\subsection{امید دقیق ماتریسی و تقریب با عمق محدود}
\label{cmp:subsec:ch3-exact-matrix-walk}
از رابطه~\eqref{cmp:eq:ch3-h-alpha} و ماتریس حالت رابطه~\eqref{cmp:eq:ch3-block-factorization} نتیجه می‌شود:
\begin{equation}\label{cmp:eq:ch3-exact-matrix-neighborhood}
\boxed{\begin{aligned}
h_s&=\rho_s\left(\sum_{t\ge1}
\alpha_t(\beta,\lambda)\mathsf P_s^{\,t-1}\right)\mathsf M_s^{(X)},\\
\mathscr Y_s&=\rho_s\left(\sum_{t\ge1}
\alpha_t(\beta,\lambda)\mathsf P_s^{\,t-1}\right)\mathsf M_s^{(Y)}.
\end{aligned}}
\end{equation}
این رابطه، \emph{امید دقیق} بازنمایی تعریف‌شده با گشت‌های تصادفی است؛ بنابراین برای محاسبه آن به تولید گشت‌های مونت‌کارلو نیاز نداریم. مقصود از دقیق، حذف خطای نمونه‌گیری گشت برای گراف مرجع و سیاست ثابت است، نه حذف خطای ناشی از محدود بودن داده‌ها یا تقریب وزن‌های ارتباط.

در پیاده‌سازی با حداکثر عمق $H$، طول واقعی $L_H=\min(L,H)$ است. با قرار دادن
$p_\ell^{(H)}=\beta(1-\beta)^{\ell-1}$ برای $1\le\ell<H$
و $p_H^{(H)}=(1-\beta)^{H-1}$، ضرایب سازگار با نرمال‌سازی \emph{درون هر گشت} برابرند با:
\begin{equation}\label{cmp:eq:ch3-capped-alpha}
\boxed{\alpha_t^{(H)}(\beta,\lambda)=
\lambda^{t-1}\sum_{\ell=t}^{H}
\frac{p_\ell^{(H)}}{\mathcal N_\ell(\lambda)},\quad
1\le t\le H,\qquad
\sum_{t=1}^{H}\alpha_t^{(H)}=1.}
\end{equation}
باید میان سقف‌گذاری طول گشت و صرفاً قطع کردن یک سری نامتناهی تمایز گذاشت؛ در حالت اول، احتمال تمام طول‌های بیشتر از $H$ در $p_H^{(H)}$ جمع می‌شود. شرط‌کردن روی طول سقف‌دار و اثبات اینکه مجموع ضرایب $\alpha_t^{(H)}$ برابر یک است در پیوست~\ref{app:ch3-matrix-walk-proofs} آمده است.

برای پرهیز از محاسبه صریح توان‌های ماتریس، کافی است از $\nu_{s,1}=\rho_s$ شروع و رابطه
$\nu_{s,t+1}=\nu_{s,t}\mathsf P_s$ را تکرار کنیم. در این صورت:
\begin{equation}\label{cmp:eq:ch3-matrix-forward-recursion}
\boxed{\mathscr Y_s
=\sum_{t=1}^{H}\alpha_t^{(H)}\,\nu_{s,t}\mathsf M_s^{(Y)}.}
\end{equation}
رابطه بازگشتی معادل برای تجمیع از انتهای افق نیز وجود دارد:
\begin{equation}\label{cmp:eq:ch3-matrix-backward-recursion}
G_{s,H+1}=0,\qquad
G_{s,t}=\alpha_t^{(H)}\mathsf M_s^{(Y)}+\mathsf P_sG_{s,t+1},
\qquad \mathscr Y_s=\rho_sG_{s,1}.
\end{equation}
این دو روش، در حساب دقیق، یک پاسخ دارند؛
زیرا برای $t\le H$ از استقرا روی بازگشت پس‌رونده داریم
$G_{s,t}=\sum_{u=t}^{H}\alpha_u^{(H)}
\mathsf P_s^{u-t}\mathsf M_s^{(Y)}$.
در نتیجه $\rho_sG_{s,1}$ همان جمع پیش‌رونده است.
این برابری و اعتبار عامل‌گذاری ماتریسی در
پیوست~\ref{app:ch3-matrix-walk-proofs} بررسی می‌شود.
محاسبه ضرایب $\alpha_t^{(H)}$ نیز با یک جمع پس‌رونده
روی $p_\ell^{(H)}/\mathcal N_\ell$ در $O(H)$ عملیات انجام می‌شود.

\paragraph{تقریب کنترل‌شده.}
اگر فقط $r<H$ گام اول را محاسبه کنیم و
$\varepsilon_r=\sum_{t=r+1}^{H}\alpha_t^{(H)}$
را جرم وزنِ حذف‌شده بنامیم، توزیع نرمال‌شده تقریبی را چنین می‌سازیم:
\begin{equation}\label{cmp:eq:ch3-matrix-truncated}
\widehat{\mathscr Y}_{s,r}
=\frac{\sum_{t=1}^{r}\alpha_t^{(H)}
\nu_{s,t}\mathsf M_s^{(Y)}}{1-\varepsilon_r}.
\end{equation}
برای فاصله تغییر کل $\operatorname{TV}(P,Q)=\|P-Q\|_1/2$، مستقل از سیاست و ساختار گراف داریم:
\begin{equation}\label{cmp:eq:ch3-matrix-truncation-bound}
\boxed{\operatorname{TV}\!\left(
\widehat{\mathscr Y}_{s,r},\mathscr Y_s\right)
\le\varepsilon_r
\le\Pr(L_H>r)=(1-\beta)^r,\qquad r<H.}
\end{equation}
همین کران برای توزیع همسایگی نمونه‌ای و، با میانگین‌گیری از نمونه‌های طبقه، برای مراکز طبقه برقرار است. اثبات آن در پیوست~\ref{app:ch3-matrix-walk-proofs} آمده است. در مقابل، استفاده از $Q$ یک‌گامی برای تقریب کل گشت چنین کران کوچکی را بدون لحاظ وزن گام‌های بعدی تضمین نمی‌کند.

\paragraph{عمق مؤثر و محدودیت ناشی از پارامترها.}
پارامتر $\beta$ طول گشت و $\lambda$ وزن نسبی گام‌ها را تعیین می‌کند؛ هیچ‌کدام به‌تنهایی بیانگر عمق بازنمایی نیستند. دو شاخص مناسب عبارت‌اند از:
\begin{equation}\label{cmp:eq:ch3-effective-depth}
d_{\mathrm{eff}}=\sum_{t=1}^{H}t\,\alpha_t^{(H)},
\qquad
r_\varepsilon=\min\left\{r:\sum_{t=1}^{r}
\alpha_t^{(H)}\ge1-\varepsilon\right\}.
\end{equation}
انتخاب عمق دو، سه یا چهار گام باید با جرم تجمعی $\alpha_t^{(H)}$ و نقش نمونه‌های واسط سنجیده شود، نه صرفاً با مقدار $\mathbb E[L]$. افزایش عمق الزاماً کیفیت بازنمایی را بهبود نمی‌دهد؛ زیرا ممکن است گام‌های بیشتر تفاوت توزیع‌های همسایگی را کاهش دهند. تأثیر تجربی این انتخاب بر مراکز در فصل~\ref{chap:results} بررسی خواهد شد.

\paragraph{مرتبه محاسباتی.}
برای شمارش گره‌ها و انتقال‌های بالقوه، تعریف می‌کنیم:
\begin{equation}\label{cmp:eq:ch3-graph-size}
\begin{aligned}
d_i&=|A_{X_i}|,\qquad m=\sum_{i=1}^{n}d_i,\\
e_v&=\sum_{b\in A_D}|D_b|(|D_b|-1),\qquad
e_h=\sum_{i=1}^{n}d_i^2.
\end{aligned}
\end{equation}
در این رابطه، $m$ تعداد گره‌های موجود نمونه--ویژگی، $e_v$ تعداد انتقال‌های عمودی جهت‌دار بالقوه و $e_h$ تعداد انتقال‌های افقی بالقوه است. در نمایش تنکِ دو مرحله‌ای، ابتدا برای هر نمونه $X_i$
حداکثر $d_i^2$ انتخاب افقی انجام می‌شود و سپس
برای هر ویژگی $b$، احتمال‌ها روی حداکثر
$|D_b|(|D_b|-1)$ انتقال عمودی پخش می‌شوند.
ممنوعیت بازگشت به مبدأ فقط برخی یال‌ها را حذف می‌کند.
بنابراین یک ضرب بردار حالت در $\mathsf P_s=F_sS_s$
حداکثر $O(e_h+e_v)$ عملیات نیاز دارد؛ در نتیجه
هزینه محاسبه دقیق بازنمایی یک مبدأ تا عمق $H$ برابر
$O(H(e_h+e_v))$ و تقریب $r$گامی برابر $O(r(e_h+e_v))$ است، پس از ساخت وزن‌ها و ساختار گراف. برای تمام $n$ مبدأ و \emph{یک قرارداد سیاست ثابت}
(مانند سیاست مشترک یا یک فرض طبقه مشخص)،
کران بالای مستقیم $O(nH(e_h+e_v))$ به‌دست می‌آید.
اگر برای هر مبدأ لازم باشد تمام $K$ فرض طبقه $c_s$
جداگانه امتیازگذاری شوند، کران مستقیم
$O(KnH(e_h+e_v))$ خواهد بود.
این کران بدون هزینه تولید امتیازهای تابعی
و مشتق‌گیری پارامترهاست.
در گراف متراکم با $d$ ویژگی موجود برای همه نمونه‌ها، این کران به
$O(H(n^3d+n^2d^2))$ تبدیل می‌شود؛ هزینه روش بر حسب یال‌ها برای گراف‌های تنک مناسب‌تر است. نگهداری وزن‌ها و عامل‌های انتقالِ یک مبدأ در هر لحظه به $O(e_v+e_h+m)$ حافظه ساختاری نیاز دارد، نه ذخیره هم‌زمان $n$ ماتریس انتقال بزرگ. شمارش دقیق عامل‌های تنک و کران هر ضرب در پیوست~\ref{app:ch3-matrix-walk-proofs} توضیح داده شده است. هزینه و حافظه لازم برای مشتق‌گیری سیاست در بخش~\ref{cmp:sec:ch3-inverse-center-policy} بررسی می‌شود.

\paragraph{تنک‌سازی انتقال‌های عمودی برای داده‌های بزرگ.}
در تعریف اصلی، هر ویژگی $b$ می‌تواند نمونهٔ جاری را به
تمام نمونه‌های مجاز در $D_b$ متصل کند؛ اما برای محاسبهٔ
امید گشت لازم نیست این زیرگراف‌های عمودی حتماً کامل باشند.
یک گونهٔ محاسباتیِ تنک را می‌توان با تعیین سقف
$k_{\mathrm{nbr}}\ge1$ برای شمار مقصدهای هر جفت
نمونه--ویژگی ساخت. برای مبدأ ثابت $s$، مجموعهٔ ناتهی
$\mathcal N_{i,b}^{(s)}\subseteq
D_b\setminus\{X_i,X_s\}$ را با حداکثر
$k_{\mathrm{nbr}}$ عضو انتخاب می‌کنیم؛ برای مثال،
نمونه‌هایی را برمی‌گزینیم که بیشترین وزن
$w_b(X_i,X_j)$ را دارند. سپس به‌جای توزیع عمودی
رابطه~\eqref{cmp:eq:ch3-vertical-transition} می‌نویسیم:
\begin{equation}\label{cmp:eq:ch3-sparse-vertical-transition}
\widetilde q_b^{(s)}(j\mid i)=
\begin{cases}
\displaystyle
\frac{w_b(X_i,X_j)}
{\sum_{X_u\in\mathcal N_{i,b}^{(s)}}w_b(X_i,X_u)},
&X_j\in\mathcal N_{i,b}^{(s)},\\[6pt]
0,&\text{در غیر این صورت}.
\end{cases}
\end{equation}
به‌سبب فرض پنجم و مثبت‌بودن وزن‌ها، برای هر ویژگی
موجود می‌توان چنین مجموعهٔ ناتهی‌ای انتخاب کرد.
\emph{صرف صفرکردن یال‌های حذف‌شده کافی نیست}؛
احتمال‌های مقصدهای باقی‌مانده باید دوباره نرمال شوند
تا مجموع هر سطر انتقال برابر یک بماند. ممنوعیت بازگشت
به مبدأ نیز در انتخاب $\mathcal N_{i,b}^{(s)}$
رعایت می‌شود. برای آماده‌سازی همسایگان مستقل از مبدأ،
می‌توان برای هر $(i,b)$ تا
$k_{\mathrm{nbr}}+1$ مقصد برترِ متفاوت از $i$
را ذخیره کرد و هنگام تعیین مبدأ $s$، آن را حذف
و حداکثر $k_{\mathrm{nbr}}$ مقصد باقی‌مانده را نگه داشت؛
وجود دست‌کم سه نمونه برای هر ویژگی تضمین می‌کند
این فهرست پس از حذف مبدأ تهی نشود.

اگر $\widetilde e_v^{(s)}$ تعداد انتقال‌های عمودی
نگهداری‌شده برای مبدأ $s$ باشد، داریم
$\widetilde e_v^{(s)}\le k_{\mathrm{nbr}}m$.
بنابراین پس از آماده‌سازی فهرست همسایگان،
هزینهٔ بازگشت پیش‌رونده تا عمق $H$ حداکثر
$O\bigl(H(e_h+k_{\mathrm{nbr}}m)\bigr)$
و حافظهٔ ساختاری عامل‌ها و بردار حالت
$O(e_h+k_{\mathrm{nbr}}m+m)$ است.
هزینهٔ یافتن همسایگان و محاسبهٔ وزن‌های آن‌ها
در این کران‌ها منظور نشده و به روش ساخت فهرست‌ها
بستگی دارد. با کوچک نگه‌داشتن
$k_{\mathrm{nbr}}$ نسبت به $|D_b|$،
این روش امکان کاهش چشمگیر حافظه و هزینهٔ ضرب‌های
ماتریسی را فراهم می‌کند؛ بااین‌حال
اتصال‌پذیری و کیفیت بازنمایی ممکن است تغییر کند.
محاسبهٔ ماتریسی همچنان \emph{امید دقیق گشت روی
گراف تنک‌شده} را می‌دهد، نه لزوماً امید گشت روی
گراف کامل اولیه. کران خطای قطع عمق در
رابطه~\eqref{cmp:eq:ch3-matrix-truncation-bound}
برای گراف تنک‌شده نیز برقرار است، اما خطای ناشی از
\emph{خود تنک‌سازی} را محدود نمی‌کند.

همین تنک‌سازی برای \emph{نمونه‌گیری مستقیم گشت تصادفی}
نیز کاربرد دارد: پس از انتخاب ویژگی $b$، مقصد عمودی
به‌جای تمام $D_b$ از فهرست کوتاه
$\mathcal N_{i,b}^{(s)}$ نمونه‌گیری می‌شود.
پس از آماده‌سازی فهرست‌ها و احتمال‌های نرمال‌شده،
هزینهٔ انتخاب مقصد با جست‌وجوی تجمعی حداکثر
$O(k_{\mathrm{nbr}})$ است و با ساخت جدول نمونه‌گیری
مناسب برای هر فهرست می‌تواند به $O(1)$ در هر انتخاب
کاهش یابد؛ هزینهٔ ساخت این جدول‌ها جداگانه است.
در وزن‌دهی مبتنی بر فاصله، احتمال مقصدهای دور
ممکن است بسیار کوچک باشد، اما با تابع وزن مثبت
این احتمال‌ها از نظر ریاضی الزاماً صفر نیستند.
ملاک مناسب برای حذف آن‌ها، \emph{جرم احتمال
تجمعیِ حذف‌شده} است، نه کوچک‌بودن تک‌تک احتمال‌ها.
اگر برای یک سطر انتقال، این جرم را
$\delta_{i,b}^{(s)}
=\sum_{X_j\notin\mathcal N_{i,b}^{(s)}}
q_b^{(s)}(j\mid i)$
بنامیم، با نرمال‌سازی مجدد روی مجموعهٔ نگهداری‌شده
داریم
$\operatorname{TV}(q_b^{(s)}(\cdot\mid i),
\widetilde q_b^{(s)}(\cdot\mid i))
=\delta_{i,b}^{(s)}$.
ازاین‌رو در صورت امکانِ محاسبه یا برآورد این جرم،
می‌توان شمار همسایگان را به‌صورت تطبیقی و با
آستانه‌ای برای $\delta_{i,b}^{(s)}$ انتخاب کرد.
تنک‌سازی، چه در گشت نمونه‌گیری‌شده و چه در محاسبهٔ
ماتریسی، قانون انتقال را تغییر می‌دهد و صرفاً
یک بهینه‌سازیِ بی‌اثر بر توزیع گشت نیست.

\paragraph{شبه‌کد محاسبهٔ امید گشت با افق محدود.}
برای مشخص‌کردن تمام مراحل محاسبه، سه الگوریتم زیر
را تعریف می‌کنیم: محاسبهٔ ضرایب وزن‌دار گشت،
پیش‌برد یک‌گام توزیع حالت، و روال اصلی تجمیع امید.
در هر سه الگوریتم، مبدأ $s$ و فرض طبقهٔ $c_s$
در طول یک اجرا ثابت‌اند. برای یکسان‌کردن حالت
کامل و تنک، مجموعهٔ مقصدهای مجاز را با
$\mathcal J_{i,b}^{(s)}$ و احتمال انتقال
عمودی متناظر را با $\bar q_b^{(s)}(j\mid i)$
نشان می‌دهیم. در گراف کامل،
$\mathcal J_{i,b}^{(s)}
=D_b\setminus\{X_i,X_s\}$ و
$\bar q_b^{(s)}=q_b^{(s)}$ است؛
در گراف تنک،
$\mathcal J_{i,b}^{(s)}=\mathcal N_{i,b}^{(s)}$
و $\bar q_b^{(s)}=\widetilde q_b^{(s)}$
خواهد بود. در هر دو حالت، مجموعهٔ مقصدها
ناتهی است و $\bar q_b^{(s)}$ روی همان
مجموعه نرمال شده است.

الگوریتم~\ref{cmp:alg:ch3-capped-walk-weights}
ضرایب رابطه~\eqref{cmp:eq:ch3-capped-alpha}
را با جمع پس‌رونده و بدون محاسبهٔ
جداگانهٔ مجموع‌های تودرتو می‌سازد.
متغیر $r$ در این الگوریتم جرم طول گشت
باقی‌مانده پیش از اعمال توقف است؛
به همین دلیل حالت $H=1$ نیز مستقیماً
پوشش داده می‌شود.
\begin{algorithm}[!htbp]
\caption{زیرروال $\operatorname{CappedWeights}$ برای محاسبهٔ ضرایب گشت}
\label{cmp:alg:ch3-capped-walk-weights}
\begin{algorithmic}[1]
\REQUIRE $H\ge1$، $0<\beta\le1$ و $0<\lambda\le1$
\ENSURE $\alpha_1^{(H)},\ldots,\alpha_H^{(H)}$
\STATE $r\gets1$؛ $\mathcal N_0\gets0$
\FOR{$\ell=1,\ldots,H$}
    \STATE $\mathcal N_\ell\gets1+\lambda\mathcal N_{\ell-1}$
    \IF{$\ell<H$}
        \STATE $p_\ell^{(H)}\gets\beta r$
        \STATE $r\gets(1-\beta)r$
    \ELSE
        \STATE $p_H^{(H)}\gets r$
    \ENDIF
\ENDFOR
\STATE $z\gets0$
\FOR{$t=H,H-1,\ldots,1$}
    \STATE $z\gets z+p_t^{(H)}/\mathcal N_t$
    \STATE $\alpha_t^{(H)}\gets\lambda^{t-1}z$
\ENDFOR
\STATE \textbf{بازگردان} $(\alpha_1^{(H)},\ldots,\alpha_H^{(H)})$
\end{algorithmic}
\end{algorithm}

الگوریتم~\ref{cmp:alg:ch3-advance-walk-state}
پیش‌برد یک‌گام توزیع حالت را به‌صورت دو مرحلهٔ
افقی و عمودی اجرا می‌کند، بدون آنکه
ماتریس انتقال یا حتی عامل‌های $F_s,S_s$
به‌صورت ماتریس ساخته شوند. در این روال،
$r_b$ جرم حاصل از انتخاب ویژگی $b$
در نمونهٔ جاری است و سپس روی مقصدهای
مجاز همان ویژگی پخش می‌شود.
\begin{algorithm}[!htbp]
\caption{پیش‌برد یک گام توزیع حالت $\operatorname{Advance}$}
\label{cmp:alg:ch3-advance-walk-state}
\begin{algorithmic}[1]
\REQUIRE $\nu$ روی $\mathcal S_s$؛ $s,c_s$؛ سیاست $\pi_s^{(c_s)}$؛ $\bar q,\mathcal J$
\ENSURE $\nu^+$، توزیع حالت پس از یک گام با قانون $\bar q$
\STATE $\nu^+(j,b)\gets0$ برای هر $(j,b)\in\mathcal S_s$
\FORALL{$i\in\{1,\ldots,n\}\setminus\{s\}$}
    \FORALL{$b\in A_{X_i}$}
        \STATE $r_b\gets0$
        \FORALL{$a\in A_{X_i}$}
            \STATE $r_b\gets r_b+\nu(i,a)\pi_s^{(c_s)}(b\mid i,a)$
        \ENDFOR
        \FORALL{$X_j\in\mathcal J_{i,b}^{(s)}$}
            \STATE $\nu^+(j,b)\gets\nu^+(j,b)+r_b\bar q_b^{(s)}(j\mid i)$
        \ENDFOR
    \ENDFOR
\ENDFOR
\STATE \textbf{بازگردان} $\nu^+$
\end{algorithmic}
\end{algorithm}

الگوریتم~\ref{cmp:alg:ch3-exact-walk-expectation}
از دو زیرروال بالا استفاده می‌کند.
توزیع اولیهٔ $\nu=\rho_s$ با سطر آغاز
$\bot$ ساخته می‌شود؛ در گام‌های بعدی
سطرهای انتقال میانی سیاست به کار می‌روند.
تجمیع مستقیم در مؤلفهٔ نمونهٔ $j$ و
برچسب $y_j$ دقیقاً معادل ضرب در
$\mathsf M_s^{(X)}$ و $\mathsf M_s^{(Y)}$
است، ولی نیازی به ذخیرهٔ این دو ماتریس ندارد.
\begin{algorithm}[!htbp]
\caption{محاسبهٔ دقیق امید گشت روی گراف انتخاب‌شده}
\label{cmp:alg:ch3-exact-walk-expectation}
\begin{algorithmic}[1]
\REQUIRE $s,c_s,H,\beta,\lambda$؛ سیاست $\pi_s^{(c_s)}$؛ $\bar q,\mathcal J$؛ برچسب‌های $Y$
\ENSURE بازنمایی نمونه‌ای $h_s$ و بازنمایی برچسبی $\mathscr Y_s$
\STATE $\alpha^{(H)}\gets\operatorname{CappedWeights}(H,\beta,\lambda)$
\STATE $\mathcal S_s\gets\{(j,b)\mid j\ne s,\ b\in A_{X_j}\}$
\STATE $\nu(j,b)\gets0$ برای هر $(j,b)\in\mathcal S_s$
\FORALL{$b\in A_{X_s}$}
    \FORALL{$X_j\in\mathcal J_{s,b}^{(s)}$}
        \STATE $\nu(j,b)\gets\nu(j,b)+\pi_s^{(c_s)}(b\mid s,\bot)\bar q_b^{(s)}(j\mid s)$
    \ENDFOR
\ENDFOR
\STATE $h_s\gets0_{1\times n}$؛ $\mathscr Y_s\gets0_{1\times K}$
\FOR{$t=1,\ldots,H$}
    \FORALL{$(j,b)\in\mathcal S_s$}
        \STATE $h_s[j]\gets h_s[j]+\alpha_t^{(H)}\nu(j,b)$
        \STATE $\mathscr Y_s[y_j]\gets\mathscr Y_s[y_j]+\alpha_t^{(H)}\nu(j,b)$
    \ENDFOR
    \IF{$t<H$}
        \STATE $\nu\gets\operatorname{Advance}(\nu,s,c_s,\pi_s^{(c_s)},\bar q,\mathcal J)$
    \ENDIF
\ENDFOR
\STATE \textbf{بازگردان} $(h_s,\mathscr Y_s)$
\end{algorithmic}
\end{algorithm}

این سه الگوریتم برای افق سقف‌دار $H$ و
قانون انتقال انتخاب‌شده، بدون خطای نمونه‌گیری
مونت‌کارلو امید بازنمایی را محاسبه می‌کنند.
در حالت تنک، دقت محاسبه نسبت به \emph{گراف
تنک‌شده} است؛ خطای تغییر گراف از خطای
محاسبهٔ امید متمایز می‌ماند.
هزینهٔ محاسبهٔ ضرایب $O(H)$ است و
هر فراخوانی $\operatorname{Advance}$،
پس از آماده‌سازی وزن‌ها و همسایگان،
در حالت کامل حداکثر $O(e_h+e_v)$ و
در حالت تنک حداکثر
$O(e_h+k_{\mathrm{nbr}}m)$ عملیات دارد.
برای مبدأهای مختلف، فضای حالت و بردارهای
وابسته به هر مبدأ را می‌توان پس از اتمام
همان اجرا آزاد کرد؛ برای سیاست شرطی به طبقه
نیز هر فرض $c_s$ باید با سیاست متناظر
ارزیابی شود.


\section{پیش‌بینی برچسب نمونه}
\label{cmp:sec:ch3-neighborhood-to-label-prediction}
در بخش‌های پیشین، گشت روی گراف به بازنمایی‌های
$\mathscr H$ و $\mathscr Y=\mathscr H Y$ منتهی شد.
مؤلفه‌های $\mathscr Y_i$ برچسب \emph{نمونه‌های مشاهده‌شده در همسایگی}
را توصیف می‌کنند، نه احتمال پسین برچسب خود نمونه $X_i$.
برای رسیدن از این بازنمایی به پیش‌بینی برچسب نمونه مبدأ،
ابتدا تعبیر احتمالاتی آن را دقیق می‌کنیم،
سپس نشان می‌دهیم چگونه الگوهای همسایگی هر طبقه
در یک مرکز توزیعی خلاصه و برای مقایسه نمونه‌ها به کار گرفته می‌شوند.

\subsection{تعبیر احتمالاتی توزیع برچسب همسایگی}
\label{cmp:subsec:ch3-label-neighborhood-interpretation}

در بخش~\ref{cmp:subsec:ch3-representations}، ماتریس توزیع برچسب همسایگی با رابطه~\eqref{cmp:eq:ch3-Y-script} تعریف شد. در اینجا بردار سطری متناظر با نمونه آغازین $X_i$ را مطابق نمادگذاری پیشین با $\mathscr Y_i$ نشان می‌دهیم؛ یعنی $\mathscr Y_i$ سطر $i$ام ماتریس $\mathscr Y=\mathscr H Y$ است. مؤلفه مربوط به طبقه $c$، مجموع وزن مشاهده نمونه‌های دارای برچسب $c$ در گشت‌های آغازشده از $X_i$ را بیان می‌کند.

برای بیان دقیق تعبیر احتمالاتی این مؤلفه، فرض کنیم پس از تولید یک گشت با طول $L$، اندیس یکی از گام‌های آن را با متغیر تصادفی $\xi$ انتخاب کنیم، به گونه‌ای که:
\begin{equation}\label{cmp:eq:ch3-weighted-observation-position}
\mathrm P(\xi=t\mid L=\ell)
=\frac{\lambda^{t-1}}{\mathcal N_\ell(\lambda)},
\qquad t=1,\ldots,\ell.
\end{equation}
در این صورت، با نماد $J_\xi$ برای نمونه مشاهده‌شده در گام انتخاب‌شده و $y_j$ برای برچسب نمونه $X_j$، داریم:
\begin{equation}\label{cmp:eq:ch3-label-neighborhood-probability}
\boxed{\mathscr Y_i(c)
=\mathrm P\bigl(y_{J_\xi}=c\mid J_0=i\bigr)
=\sum_{j=1}^{n}h_{ij}Y_{jc}.}
\end{equation}
در این رابطه، $\mathrm P$ عملگر \textbf{احتمال} است و به معنای \textbf{احتمال مشاهده برچسب $c$ در همسایگی نمونه آغازین $X_i$} تحت فرایند گام‌برداری و وزن‌دهی تعریف‌شده است. این مقدار را نباید با احتمال آنکه \emph{برچسب خود نمونه $X_i$} برابر $c$ باشد، یکسان دانست. به بیان دیگر، $\mathscr Y_i$ توزیع برچسب \emph{همسایگی} نمونه است، نه توزیع پسین برچسب واقعی خود آن نمونه.

همچنین چون بازگشت به نمونه آغازین ممنوع است و $h_{ii}=0$، برچسب خود نمونه $X_i$ در تشکیل $\mathscr Y_i$ مشارکت ندارد. در مقابل، برچسب سایر نمونه‌های مشاهده‌شده، حتی اگر با برچسب واقعی نمونه آغازین برابر باشد، در این توزیع سهم خواهد داشت.

یکی از راه‌های ساده استفاده از این بازنمایی برای تشخیص برچسب، انتخاب طبقه‌ای است که بیشترین سهم را در توزیع همسایگی دارد:
\begin{equation}\label{cmp:eq:ch3-simple-label-decision}
\widehat y_i=\arg\max_{c\in\{1,\ldots,K\}}\mathscr Y_i(c).
\end{equation}
اما درستی این تصمیم مستلزم آن است که همسایگی نمونه عمدتاً از نمونه‌های هم‌برچسب آن تشکیل شده باشد؛ چنین شرطی از تعریف گراف و سیاست یکنواخت نتیجه نمی‌شود. برای نمونه، ممکن است توزیع‌های همسایگی دو طبقه به‌ترتیب به‌صورت $(0.1,0.9)$ و $(0.9,0.1)$ باشند. در این حالت، انتخاب بزرگ‌ترین مؤلفه می‌تواند برچسب هر دو طبقه را به اشتباه پیش‌بینی کند، درحالی‌که دو بردار به‌خوبی از یکدیگر قابل تفکیک‌اند.

بنابراین، در کاربرد مهندسی ویژگی، ارزش بازنمایی $\mathscr Y_i$ لزوماً در برابر بودن بیشترین مؤلفه آن با برچسب واقعی نمونه نیست؛ بلکه در این است که توزیع‌های همسایگی نمونه‌های متعلق به طبقه‌های متفاوت، الگوهای قابل تشخیصی داشته باشند. برای استفاده از این الگوها، در ادامه توزیع نماینده هر طبقه را معرفی می‌کنیم.


\subsection{چارچوب عمومی پیش‌بینی برچسب نمونه}
\label{cmp:subsec:ch3-class-centers}
از آنجا که $\mathscr Y_i$ الگوی برچسب‌های مشاهده‌شده در همسایگی
را نمایش می‌دهد، می‌توان برای هر طبقه یک نمایندهٔ توزیعی ساخت،
بازنمایی نمونه را با نماینده‌های طبقات مقایسه کرد
و امتیازهای حاصل را به احتمال تخمینی و برچسب نهایی تبدیل کرد.
در اینجا ابتدا صورت عمومیِ این طبقه‌بند را، بدون انتخاب
واگرایی یا متریک مشخص، معرفی می‌کنیم.

\noindent\textbf{ساخت مدل و استنتاج نمونهٔ جدید.}
در مرحلهٔ آموزش، گراف مرجع، برچسب‌های
نمونه‌های آموزشی و مراکز طبقات
فقط از دادهٔ آموزش ساخته می‌شوند.
هنگام پیش‌بینی برای نمونهٔ جدید
$x$، پارامترهای سیاست و مراکز
آموزش‌دیده ثابت می‌مانند؛ نمونهٔ
پرس‌وجو، با ویژگی‌های مشاهده‌شدهٔ خود،
گام نخست را به سمت مرجع آموزشی
برمی‌دارد و برچسب ناشناختهٔ آن
هرگز در بازنمایی یا انتخاب سیاست
به کار نمی‌رود. برای سیاست مستقل
از فرضیهٔ طبقه، یک بازنمایی مشترک
با مراکز مقایسه می‌شود؛ اما در
سیاست شرطی به طبقه، برای هر نامزد
یک گشت فرضیه‌محور اجرا و بازنمایی
آن با مرکز سازگار همان طبقه
مقایسه می‌شود. در سیاست‌های گره‌محور
آزاد نیز قاعدهٔ شروع برای نمونهٔ
خارج از مرجع باید جداگانه مشخص شود.
در این زیربخش فقط قاعدهٔ ریاضی
امتیازدهی و تصمیم را معرفی می‌کنیم؛
قرارداد دقیق گام نخست، وضعیت‌های
نامعتبر و الگوریتم پیش‌بینی یک نمونه
در زیربخش~\ref{cmp:subsec:ch3-heldout-validation-algorithm}
آمده است.

\paragraph{نمایندهٔ احتمالی هر طبقه.}
برای هر طبقهٔ $c\in\{1,\ldots,K\}$ یک
\emph{توزیع نماینده} با نماد $\mathcal V_c$ در نظر می‌گیریم.
چون بازنمایی‌های همسایگی $\mathscr Y_i$ بردار توزیع احتمال
روی $K$ برچسب‌اند، نمایندهٔ هر طبقه نیز باید از همین جنس باشد:
\[
\mathcal V_c\in\Delta_K,\qquad
\Delta_K:=\bigl\{P\in\mathbb R_{\ge0}^{K}:
\textstyle\sum_{u=1}^{K}P(u)=1\bigr\}.
\]
مجموعهٔ این نماینده‌ها را در سطرهای ماتریس
$\mathcal V\in\mathbb R^{K\times K}$ قرار می‌دهیم.
در این مرحله \emph{شیوهٔ ساخت نماینده} را تعیین نمی‌کنیم؛
فقط فرض می‌کنیم نماینده‌های طبقات از داده‌های مرجع
به دست می‌آیند. روش مشخص ساخت آن‌ها را
در زیربخش بعد انتخاب و توجیه می‌کنیم.

\paragraph{ناهمانندی نمونه و مرکز.}
برای مقایسهٔ دو توزیع احتمال از سنجهٔ ناهمانندی
$d(P,Q)\in[0,+\infty]$ استفاده می‌کنیم که مقدار کوچک‌تر
آن بیانگر سازگاری بیشتر است. لازم نیست این سنجه
متقارن باشد یا نامساوی مثلثی را ارضا کند.
در حالت بازنمایی مشترک، ناهمانندی نمونهٔ $X_i$
با مرکز طبقهٔ نامزد $c$ را تعریف می‌کنیم:
\begin{equation}\label{cmp:eq:ch3-general-center-dissimilarity}
d_{i,c}:=d(\mathscr Y_i,\mathcal V_c).
\end{equation}
در تمام نامزدها، همین سنجه و همین قرارداد محاسبه به کار می‌رود.

\paragraph{نگاشت ناهمانندی به احتمال تخمینی و برچسب.}
امتیاز طبقهٔ $c$ را برابر $-d_{i,c}$ قرار می‌دهیم
و با نرمال‌سازی با بیشینهٔ هموار (\lr{Softmax}) تعریف می‌کنیم:
\begin{equation}\label{cmp:eq:ch3-NLL-softmax}
\boxed{
\widehat p_{i,c}
:=\widehat{\mathrm P}(y_i=c\mid\mathscr Y_i)
=\frac{\exp(-d_{i,c})}
{\sum_{r=1}^{K}\exp(-d_{i,r})}.
}
\end{equation}
در این رابطه، فرض می‌کنیم دست‌کم یک ناهمانندی متناهی است؛
برای مقدار نامتناهی، قرارداد $\exp(-\infty)=0$
به کار می‌رود. خروجی یک توزیع احتمالِ \emph{تخمینی مدل} است،
نه اثبات برابری با احتمال پسین واقعی برچسب نمونه.
با حفظ قرارداد رفع تساوی، برچسب پیش‌بینی‌شده عبارت است از:
\begin{equation}\label{cmp:eq:ch3-class-center-decision}
\boxed{
\widehat y_i
=\underset{c}{\arg\max}\;\widehat p_{i,c}
=\underset{c}{\arg\min}\;d(\mathscr Y_i,\mathcal V_c).
}
\end{equation}
این قاعده به غالب‌بودن برچسب واقعی در همسایگی نمونه
وابسته نیست. شیوهٔ ساخت نمایندهٔ هر طبقه و سنجهٔ ناهمانندی
دو تصمیم جداگانه‌اند. در ادامه نخست نحوهٔ ساخت
نماینده‌ها و سپس مفهوم ناهمانندی و شباهت را بررسی می‌کنیم.

\subsection{توجیه مرکز میانگین و هزینهٔ تجمیع}
\label{cmp:subsec:ch3-mean-center-justification}
\paragraph{انتخاب نمایندهٔ میانگین.}
مجموعهٔ اندیس نمونه‌های طبقهٔ $c$ را
$C_c=\{i\mid y_i=c\}$ می‌نامیم و فرض می‌کنیم
$|C_c|>0$. در این پژوهش نمایندهٔ عمومیِ معرفی‌شده در
بخش~\ref{cmp:subsec:ch3-class-centers} را برابر
میانگین بازنمایی‌های اعضای طبقه انتخاب می‌کنیم:
\begin{equation}\label{cmp:eq:ch3-class-center}
\boxed{\mathcal V_c=\frac{1}{|C_c|}\sum_{i\in C_c}\mathscr Y_i.}
\end{equation}
با استفاده از ماتریس برچسب یک‌داغ $Y$، صورت ماتریسی
این انتخاب برابر است با:
\begin{equation}\label{cmp:eq:ch3-class-centers-matrix}
\boxed{\mathcal V=(Y^\top Y)^{-1}Y^\top\mathscr Y
=(Y^\top Y)^{-1}Y^\top\mathscr H Y.}
\end{equation}
زیرا $Y^\top Y=\operatorname{diag}(|C_1|,\ldots,|C_K|)$
و ضرب در وارون آن، جمع بازنمایی‌های هر طبقه را
به میانگین حسابی تبدیل می‌کند.

\paragraph{حفظ ساختار احتمالاتی مرکز.}
نخست، میانگین حسابیِ توزیع‌های احتمال با وزن‌های نامنفی
و مجموع یک، خود یک توزیع احتمال معتبر است.
در نتیجه برای مرکز هر طبقه داریم:
\[
\mathcal V_c(u)\ge0,\qquad
\sum_{u=1}^{K}\mathcal V_c(u)=1.
\]
این ویژگی برای مقایسهٔ مراکز و بازنمایی‌ها در یک فضای
احتمالی مشترک ضروری است؛ هرچند به‌تنهایی اثبات نمی‌کند
میانگین برای هر تابع ناهمانندی بهینه است.
مؤلفهٔ $u$ام مرکز، میانگین احتمال مشاهدهٔ برچسب $u$
در همسایگی اعضای طبقهٔ $c$ است. اگر نمونهٔ آغازین
به‌طور یکنواخت از اعضای $C_c$ انتخاب شود،
این تعبیر را می‌توان چنین نوشت:
\begin{equation}\label{cmp:eq:ch3-class-center-prob}
\boxed{\mathcal V_c(u)=
\mathrm P\bigl(y_{J_\xi}=u\mid y_{J_0}=c\bigr)
=\frac{1}{|C_c|}\sum_{i\in C_c}\mathscr Y_i(u).}
\end{equation}
شرط داخل احتمال دربارهٔ برچسب \emph{نمونهٔ آغازین} است، درحالی‌که متغیر سمت چپ رویدادِ مشاهده برچسب در \emph{همسایگی} را توصیف می‌کند. همچنین انتخاب یکنواخت نمونه آغازین از یک طبقه را نباید با سیاست یکنواخت انتخاب ویژگی در طول گشت یکسان دانست.


\paragraph{بهینگی نمایندهٔ میانگین از دید درست‌نمایی.}
برای توجیه انتخاب میانگین، مسئلهٔ بیشینه‌سازی درست‌نمایی
در مدل‌کردن چند توزیع با یک نمایندهٔ مشترک را بررسی می‌کنیم.
فرض کنیم مشاهدات از توزیع‌های $P_1,\ldots,P_g$ با
سهم‌های $\eta_r\ge0$ و $\sum_r\eta_r=1$ به‌دست آیند،
ولی بخواهیم همهٔ آن‌ها را تنها با یک توزیع مشترک $Q$
مدل کنیم. میانگین این توزیع‌ها را با
$\overline P=\sum_r\eta_rP_r$ نشان می‌دهیم.
\emph{میانگین لگاریتم درست‌نمایی مورد انتظار}
مشاهده‌ها تحت مدل مشترک چنین بازنویسی می‌شود:
\begin{equation}\label{cmp:eq:ch3-mixture-log-likelihood}
\begin{aligned}
\mathcal L(Q)
&=\sum_{r=1}^{g}\eta_r\sum_{u=1}^{K}P_r(u)\log Q(u)\\
&=\sum_{u=1}^{K}\overline P(u)\log Q(u)\\
&=\sum_{u=1}^{K}\overline P(u)\log\overline P(u)\\
&\quad-\sum_{u=1}^{K}\overline P(u)
 \log\frac{\overline P(u)}{Q(u)}\\
&=-H(\overline P)-D_{\mathrm{KL}}(\overline P\|Q).
\end{aligned}
\end{equation}
در این محاسبه، نخست سهم توزیع‌های مختلف در
$\overline P$ تجمیع می‌شود؛ سپس با افزودن و کاستن
جملهٔ $\sum_u\overline P(u)\log\overline P(u)$،
دو کمیت آشنا از خودِ بازنویسی درست‌نمایی ظاهر می‌شوند:
آنتروپی شانون $H$ و واگرایی کولبک--لایبلر
$D_{\mathrm{KL}}$. برای دو توزیع احتمال $P$ و $Q$
روی $K$ برچسب، این کمیت‌ها به‌ترتیب عبارت‌اند از:
\begin{equation}\label{cmp:eq:ch3-KL-definition}
\begin{aligned}
H(P)&=-\sum_{u=1}^{K}P(u)\log P(u),\\
D_{\mathrm{KL}}(P\|Q)
&=\sum_{u=1}^{K}P(u)\log\frac{P(u)}{Q(u)}.
\end{aligned}
\end{equation}
در این روابط، جملهٔ دارای ضریب صفر را با حد پیوسته
صفر در نظر می‌گیریم؛ اگر $P(u)>0$ و $Q(u)=0$ باشد،
واگرایی KL بی‌نهایت و لگاریتم درست‌نمایی متناظر
برابر منفی بی‌نهایت است. با قراردادهای فوق،
$H(\overline P)$ به $Q$ وابسته نیست و
$D_{\mathrm{KL}}(\overline P\|Q)\ge0$ است؛
ازاین‌رو بیشینه‌سازی درست‌نمایی، نمایندهٔ میانگین
$\overline P$ را نتیجه می‌دهد، نه انتخاب پیشینیِ
یک سنجهٔ فاصله برای تصمیم طبقه‌بندی.
بنابراین:
\begin{equation}\label{cmp:eq:ch3-mixture-likelihood-optimum}
\boxed{\underset{Q}{\arg\max}\;\mathcal L(Q)=\overline P
=\sum_{r=1}^{g}\eta_rP_r.}
\end{equation}
این استدلال، بهینگی \emph{درست‌نماییِ مدل مشترک} را نشان می‌دهد، نه بهینگی میانگین برای هر تابع فاصله دلخواه. اثبات دقیق در حالت نمونه‌های تصادفی و شرایط مربوط به مؤلفه‌های صفر در بخش~\ref{app:ch3-mixture-likelihood} از پیوست~\ref{app:ch3-proofs} آمده است.

\paragraph{هزینه مدل‌کردن یک طبقه با توزیع مشترک.}
حال برای طبقه $c$، توزیع‌های $P_r$ را همان $\mathscr Y_i$ها و وزن‌ها را $1/|C_c|$ می‌گیریم؛ در نتیجه $\overline P=\mathcal V_c$. با قراردادن توزیع بهینه در رابطه درست‌نمایی، \emph{کمترین میانگین منفی لگاریتم درست‌نمایی مورد انتظار} برای مدل مشترک به‌دست می‌آید:
\begin{equation}\label{cmp:eq:ch3-class-shared-NLL}
\boxed{\operatorname{NLL}^{\mathrm{shared}}_{c,\min}
=-\mathcal L_c(\mathcal V_c)=H(\mathcal V_c).}
\end{equation}
برای مقایسه، اگر توزیع مولد اختصاصی هر مشاهده معلوم باشد و مشاهدات متناظر با $\mathscr Y_i$ را با خود همان توزیع مدل کنیم، کمترین هزینه مورد انتظار برابر است با:
\begin{equation}\label{cmp:eq:ch3-class-individual-NLL}
\boxed{\operatorname{NLL}^{\mathrm{individual}}_{c,\min}
=\frac{1}{|C_c|}\sum_{i\in C_c}H(\mathscr Y_i).}
\end{equation}
اختلاف این دو مقدار، هزینه اضافی \emph{تجمیع} توزیع‌های نمونه‌های طبقه در یک نماینده مشترک را می‌سنجد. این اختلاف را با $R_c$ نشان می‌دهیم:
\begin{equation}\label{cmp:eq:ch3-R-entropy}
\boxed{
\begin{aligned}
R_c
&=\operatorname{NLL}^{\mathrm{shared}}_{c,\min}
-\operatorname{NLL}^{\mathrm{individual}}_{c,\min}\\
&=H(\mathcal V_c)
-\frac1{|C_c|}\sum_{i\in C_c}H(\mathscr Y_i)\\
&=\frac1{|C_c|}\sum_{i\in C_c}
D_{\mathrm{KL}}(\mathscr Y_i\|\mathcal V_c).
\end{aligned}}
\end{equation}
برابری آخر از تعریف آنتروپی و میانگین‌بودن $\mathcal V_c$ حاصل می‌شود. بنابراین $R_c\ge0$؛ اگر توزیع همسایگی همه نمونه‌های طبقه یکسان باشد، این هزینه صفر است. برای مثال، دو توزیع $(1,0)$ و $(0,1)$ و نیز دو توزیع یکسان $(1/2,1/2)$ هر دو مرکز $(1/2,1/2)$ دارند؛ اما هزینه تجمیع در حالت اول $\log2$ و در حالت دوم صفر است. پس آنتروپی مرکز به‌تنهایی پراکندگی توزیع‌های اعضای طبقه را مشخص نمی‌کند.

\paragraph{اثر انتخاب نماینده‌ای غیر از میانگین.}
هویت زیر نشان می‌دهد در صورت انتخاب یک توزیع مشترک دیگر مانند $U$، علاوه بر هزینه اجتناب‌ناپذیر تجمیع، هزینه‌ای اضافی نیز پرداخت می‌شود:
\begin{theorem}[تجزیه هزینه مدل‌سازی مشترک]\label{cmp:thm:ch3-KL-decomposition}
برای هر طبقه $c$ و هر توزیع احتمال $U$ که واگرایی‌های مربوط به آن متناهی باشند، داریم:
\begin{equation}\label{cmp:eq:ch3-KL-decomposition}
\boxed{
\frac{1}{|C_c|}\sum_{i\in C_c}
D_{\mathrm{KL}}(\mathscr Y_i\|U)
=R_c+D_{\mathrm{KL}}(\mathcal V_c\|U).}
\end{equation}
\end{theorem}
در نتیجه، جمله $R_c$ هزینه ناشی از خلاصه‌کردن توزیع‌های متفاوت با \emph{یک} نماینده است و جمله $D_{\mathrm{KL}}(\mathcal V_c\|U)$ زیان اضافیِ برگزیدن نماینده‌ای غیر از میانگین بهینه را نشان می‌دهد. به‌ویژه:
\begin{equation}\label{cmp:eq:ch3-KL-opt-center}
\boxed{\mathcal V_c=
\underset{U}{\arg\min}\;
\frac{1}{|C_c|}\sum_{i\in C_c}
D_{\mathrm{KL}}(\mathscr Y_i\|U).}
\end{equation}
تجزیه عمومی این رابطه و اثبات کامل صورت درون‌طبقه‌ای آن در بخش~\ref{app:ch3-mixture-likelihood} از پیوست آمده‌اند.

این تفسیر روشن می‌کند چرا $R_c$ را در ارزیابی بازنمایی اندازه می‌گیریم: این معیار هزینه فشرده‌سازی الگوهای همسایگی هر طبقه در یک توزیع مشترک است. البته این هزینه، \emph{منفی لگاریتم درست‌نماییِ برچسب واقعی نمونه} یا خطای تصمیم طبقه‌بند نیست. ممکن است $R_c$ برای همه طبقه‌ها صفر باشد ولی مراکز آن‌ها نیز با یکدیگر برابر باشند و امکان تفکیک وجود نداشته باشد. بنابراین، بهینگی نمایندگی درون هر طبقه را باید از کیفیت \emph{تمایز میان طبقه‌ها} جدا کرد؛ معیارهای این دو جنبه در بخش~\ref{cmp:sec:ch3-quality-criteria} بررسی می‌شوند.



\subsection{ناهمانندی در فضای بردارهای توزیع احتمال}
\label{cmp:subsec:ch3-probability-distance-similarity}
\label{cmp:subsec:ch3-probability-dissimilarity-families}
توزیع‌های برچسب همسایگی و نماینده‌های طبقات،
اعضای فضای $\Delta_K$ هستند که در
بخش~\ref{cmp:subsec:ch3-class-centers}
معرفی شد. این فضا زیرمجموعه‌ای از
$\mathbb R^K$ است، اما مؤلفه‌های بردارهای آن
نامنفی و دارای مجموع یک‌اند و ممکن است صفر باشند. برای سنجش نزدیکی
دو توزیع و مقایسه نمونه با مرکز،
سنجه‌ای لازم است که با این تعبیر احتمالاتی
سازگار باشد. البته فاصله اقلیدسی نیز در این فضا
قابل تعریف است؛ ما سنجه‌های اطلاعاتی را
به دلیل معنای احتمالاتی آن‌ها بررسی می‌کنیم.
\emph{واگرایی} را نباید با \emph{متریک فاصله}
یا \emph{تابع شباهت} یکسان دانست:
یک متریک علاوه بر نامنفی‌بودن و صفرشدن
فقط برای دو توزیع برابر، تقارن و
نامساوی مثلثی را نیز ارضا می‌کند.
تابع شباهت، برخلاف سنجهٔ ناهمانندی،
مقدار بزرگ‌تر را نشانهٔ سازگاری بیشتر می‌گیرد.
هر سنجه باید با توجه به جهت، دامنهٔ مقدارها،
رفتار مؤلفه‌های صفر و نقش آن در تصمیم انتخاب شود.
در ادامه، نخست محدودیت‌های KL را می‌گوییم،
سپس با یک نمایش هیلبرتی به واگرایی جنسن--شنون،
متریک جذر آن و شباهت متناظر می‌رسیم.
انتخاب سنجه، تعریف مرکز میانگین را تغییر نمی‌دهد.

\paragraph{واگرایی KL و محدودیت‌های آن.}
واگرایی KL نخست از تجزیهٔ لگاریتم درست‌نمایی
در بخش~\ref{cmp:subsec:ch3-mean-center-justification}
ظاهر شد و صورت ریاضی آن در
رابطه~\eqref{cmp:eq:ch3-KL-definition} آمده است.
در اینجا آن را به‌عنوان یکی از انتخاب‌های ممکن
برای \emph{مقایسهٔ بازنمایی نمونه و مرکز} بررسی می‌کنیم؛
نقش آن در توجیه مرکز میانگین با انتخاب فاصلهٔ تصمیم
یکسان نیست.
این سنجه نامنفی است و فقط برای $P=Q$ صفر می‌شود،
اما جهت‌دار و عموماً نامتقارن است،
نامساوی مثلثی را ندارد و با وجود
$p_j>0$ و $q_j=0$ می‌تواند بی‌نهایت شود.
پس KL با وجود تعبیر اطلاعاتی مهم آن،
متریک فاصله محسوب نمی‌شود.

\paragraph{انگیزهٔ نمایش هیلبرتی.}
محدودیت‌های KL ما را به بررسی ناهمانندی‌ای
متقارن، کراندار و قابل‌تعریف حتی در حضور
مؤلفه‌های صفر هدایت می‌کنند. بااین‌حال،
صرفِ متقارن‌بودن یک واگرایی برای داشتن
نامساوی مثلثی کافی نیست. راه مناسب آن است
که توزیع‌ها را به فضایی هیلبرتی ببریم و
فاصلهٔ اقلیدسیِ تعمیم‌یافتهٔ تصاویر را
به‌عنوان متریک در نظر بگیریم.
برای حفظ تعبیر شباهت، مطلوب است تصاویر
\emph{نرم ثابت} نیز داشته باشند؛ در این صورت
ضرب داخلی، پس از نرمال‌سازی، شباهتی سازگار
با فاصلهٔ تصاویر به دست می‌دهد.

نگاشت زیر انتخابی دلخواه از میان نگاشت‌های
هیلبرتی نیست: شکل نماییِ مؤلفه‌های آن با
بازنمایی انتگرالی عبارت‌های لگاریتمی سازگار
است و عامل مقیاس به‌گونه‌ای برگزیده می‌شود
که نرم مربع هر مؤلفه با احتمال همان مؤلفه
متناسب باشد. بنابراین قید مجموع یکِ
توزیع‌ها، نرم کل تصویر را ثابت می‌کند.
پس از ساخت نگاشت، نشان می‌دهیم مربع
فاصلهٔ تصاویر دقیقاً برابر JSD است؛
در نتیجه جذر JSD متریک و شباهت نرمال‌شده
دارای تفسیر ضرب داخلی خواهد بود.
جزئیات هویت انتگرالی و شرایط مرزی
در پیوست~\ref{app:ch3-js-error-proof} آمده‌اند.

\paragraph{نگاشت توزیع‌ها به فضای هیلبرت.}
فضای هیلبرت $\mathcal H=L^2((0,\infty);\mathbb R^K)$ را در نظر گرفته و هر توزیع $P=(p_1,\ldots,p_K)$ را به تابع برداری زیر نگاشت می‌کنیم:
\begin{equation}\label{cmp:eq:ch3-js-Hilbert-map}
\boxed{
\Phi(P)(t)=
\bigl(\phi_{p_1}(t),\ldots,\phi_{p_K}(t)\bigr),
\qquad
\phi_a(t)=\frac{1-e^{-at}}{2t},\quad t>0.
}
\end{equation}
این نگاشت تعریف‌شده و یک‌به‌یک است و خواص زیر را دارد: برای هر $t>0$، مؤلفه $\phi_a(t)$ نسبت به احتمال $a$ اکیداً صعودی و مقعر است؛ برای هر $a>0$، نسبت به $t$ نزولی بوده و بیشینه حدی آن در $t\to0^+$ برابر $a/2$ است. بنابراین نگاشت، ترتیب مؤلفه‌های احتمال را حفظ می‌کند و دو توزیع متفاوت را به دو تصویر متفاوت می‌برد. افزون بر این، نرم تصویر همه توزیع‌های احتمال برابر است:
\[
\|\Phi(P)\|_{\mathcal H}^2=\frac{\log2}{2}.
\]
پس تصاویر توزیع‌ها روی کره‌ای با شعاع ثابت
$r=\sqrt{\log2/2}$ در فضای هیلبرت قرار می‌گیرند.
بااین‌حال، \emph{این کره تمام فضای قابل‌دستیابی نیست}.
زیرا برای هر $a\ge0$ و $t>0$،
$\phi_a(t)\ge0$ است و همه مؤلفه‌های
$\Phi(P)(t)$ نامنفی‌اند. افزون بر آن،
مؤلفه‌ها باید هم‌زمان از یک بردار احتمال
با مجموع یک حاصل شده باشند. بنابراین:
\begin{equation}\label{cmp:eq:ch3-js-positive-spherical-image}
\boxed{
\begin{aligned}
\Phi(\Delta_K)\subseteq\bigl\{f\in\mathcal H:\;&
\|f\|_{\mathcal H}^{2}=\tfrac{\log2}{2},\\[-1pt]
& f_j(t)\ge0\
\text{برای تقریباً همه }t>0\bigr\}.
\end{aligned}
}
\end{equation}
این مجموعه، بخشی محدود از ناحیه نامنفی
کره هیلبرتی است، نه تمام آن؛
همچنین حتی همه نقاط نامنفی کره نیز
لزوماً تصویر توزیع احتمال نیستند.
تعبیر تصویریِ «بخشی شبیه یک هشتم کره»
تنها یک تشبیه برای محدودیت علامت مؤلفه‌هاست:
فضای هیلبرت حاضر بی‌نهایت‌بعدی است و
توزیع‌های روی $K$ برچسب فقط
$K-1$ درجه آزادی دارند.
از نامنفی‌بودن تصاویر،
$\langle\Phi(P),\Phi(Q)\rangle_{\mathcal H}\ge0$
و درنتیجه زاویه میان تصاویر
بیش از $\pi/2$ نخواهد بود.

\paragraph{از فاصلهٔ تصاویر به JSD و جذر آن.}
برای دو توزیع $P,Q\in\Delta_K$، مخلوط مشترک
$M=(P+Q)/2$ را در نظر می‌گیریم. واگرایی
جنسن--شنون بر حسب همین مخلوط تعریف می‌شود:
\begin{equation}\label{cmp:eq:ch3-JSD-definition}
\operatorname{JSD}(P,Q)
=\frac12D_{\mathrm{KL}}\!\left(P\middle\|M\right)
+\frac12D_{\mathrm{KL}}\!\left(Q\middle\|M\right),
\qquad M=\frac{P+Q}{2}.
\end{equation}
محاسبهٔ فاصلهٔ دو تصویر هیلبرتی با استفاده
از ضرب داخلی و نرم ثابت آن‌ها نشان می‌دهد:
\begin{equation}\label{cmp:eq:ch3-js-Hilbert-metric}
\boxed{
\|\Phi(P)-\Phi(Q)\|_{\mathcal H}^2
=\frac12D_{\mathrm{KL}}(P\|M)
+\frac12D_{\mathrm{KL}}(Q\|M)
=\operatorname{JSD}(P,Q),
\quad M=\frac{P+Q}{2}.
}
\end{equation}
بنابراین فاصله موردنظر به‌صورت زیر تعریف می‌شود:
\begin{equation}\label{cmp:eq:ch3-js-distance}
\boxed{
d_{\mathrm{JS}}(P,Q)
=\|\Phi(P)-\Phi(Q)\|_{\mathcal H}
=\sqrt{\operatorname{JSD}(P,Q)}.
}
\end{equation}
این فاصله متقارن، یک‌به‌یک در تشخیص توزیع‌های متفاوت، کراندار در بازه $[0,\sqrt{\log2}]$ و واجد نامساوی مثلثی است؛ بنابراین یک متریک واقعی است. ساختار متریک جذر JSD و نمایش هیلبرتی آن در پژوهش‌های پیشین نیز بررسی شده‌اند \cite{EndresSchindelin2003,FugledeTopsoe2004}. این تعریف دوتوزیعی با وزن‌های برابر است؛
در تعمیم چندتوزیعی باید وزن‌ها و مخلوط مربوط
به‌صورت مستقل مشخص شوند.

\paragraph{خواص احتمالاتی و کران‌های JSD.}
از تعریف واگرایی جنسن--شنون می‌توان با
جایگذاری رابطهٔ KL، نمایش آنتروپی آن را نتیجه گرفت:
\begin{equation}\label{cmp:eq:ch3-JSD-entropy}
\operatorname{JSD}(P,Q)
=H\left(\frac{P+Q}{2}\right)
-\frac{H(P)+H(Q)}{2}.
\end{equation}
ازاین‌رو، JSD عدم‌تجانس دو توزیع
نسبت به مخلوط آن‌ها را اندازه می‌گیرد؛
با این تفاوت که برخلاف KL مستقیم،
به دلیل واردشدن همان مخلوط مشترک،
در صورت وجود مؤلفه‌های صفر نیز متناهی می‌ماند.
زیرا $M(j)\ge P(j)/2$ و $M(j)\ge Q(j)/2$،
هر یک از واگرایی‌های نسبت به مخلوط
حداکثر $\log2$ است. با نامنفی‌بودن KL
و شرط تساوی آن نتیجه می‌گیریم:
\[
0\le\operatorname{JSD}(P,Q)\le\log2,
\qquad
\operatorname{JSD}(P,Q)=0
\ \Longleftrightarrow\ P=Q.
\]
تقارن مستقیماً از تعریف به‌دست می‌آید؛
بیشینه $\log2$ در حالت تکیه‌گاه‌های
جدا به‌دست می‌آید. مثلاً برای
$P=(1,0)$ و $Q=(0,1)$،
KL مستقیم نامتناهی ولی JSD برابر $\log2$ است.
خود JSD عموماً نامساوی مثلثی را ارضا نمی‌کند،
اما جذر آن، مطابق نمایش هیلبرتی بالا، متریک است.


\paragraph{شباهت تصاویر و شباهت نرمال‌شده.}
از ثابت‌بودن نرم تصاویر و بسط مربع فاصله، رابطه زیر حاصل می‌شود:
\begin{equation}\label{cmp:eq:ch3-js-Hilbert-cosine}
\operatorname{JSD}(P,Q)
=\log2-2\langle\Phi(P),\Phi(Q)\rangle_{\mathcal H}.
\end{equation}
در نتیجه شباهت نرمال‌شده زیر دقیقاً برابر شباهت کسینوسی تصاویر است:
\begin{equation}\label{cmp:eq:ch3-js-similarity}
\boxed{
s(P,Q)=1-\frac{\operatorname{JSD}(P,Q)}{\log2}
=\frac{\langle\Phi(P),\Phi(Q)\rangle_{\mathcal H}}
{\|\Phi(P)\|_{\mathcal H}\|\Phi(Q)\|_{\mathcal H}},
\quad 0\le s(P,Q)\le1.
}
\end{equation}
این شباهت متقارن است؛ برای توزیع‌های یکسان مقدار یک
و برای تکیه‌گاه‌های جدا مقدار صفر دارد.
از آنجا که $s$ ضرب داخلی تصاویر هیلبرتی
نرمال‌شده است، یک \emph{کرنل مثبت نیمه‌معین}
محسوب می‌شود. به‌علاوه، JSD نسبت به هر توزیع
با ثابت‌بودن دیگری محدب است؛ بنابراین
$s=1-\operatorname{JSD}/\log2$
نسبت به هر ورودی به‌تنهایی مقعر خواهد بود.
برای مقایسه با بردارهای یک‌داغ نیز
شباهت $s(e_c,P)$ تابعی اکیداً صعودی
از مؤلفه $p_c$ است و ترتیب مؤلفه‌های
توزیع را حفظ می‌کند؛ البته این حکم خاص
مقایسه با بردارهای یک‌داغ است و درباره
مراکز عمومی به همان صورت ادعا نمی‌شود.
اثبات ویژگی‌های نگاشت، فرمول‌های انتگرالی،
کرنل مثبت نیمه‌معین، تقعر و خاصیت یک‌داغ
در پیوست~\ref{app:ch3-js-error-proof} آمده است.
پس متریک $d_{\mathrm{JS}}$ و شباهت $s$
هر دو از \emph{همان تصاویر هیلبرتی با نرم ثابت}
به‌دست می‌آیند، اما نقش متفاوتی دارند:
اولی فاصله و دومی شباهت مبتنی بر ضرب داخلی است؛
$s$ خود متریک فاصله نیست.

\paragraph{دورترین توزیع از یک توزیع دلخواه.}
یکی از پیامدهای مهم این هندسه آن است
که برای هر توزیع ثابت $P\in\Delta_K$،
دست‌کم یک توزیع \emph{یک‌داغ}
بیشترین فاصله متریک جنسن--شنون را
از $P$ دارد. اگر $e_j$ رأس متناظر با
برچسب $j$ باشد، آنگاه
\begin{equation}\label{cmp:eq:ch3-js-farthest-vertex}
\boxed{
\begin{aligned}
\max_{Q\in\Delta_K}d_{\mathrm{JS}}(P,Q)
&=\max_{1\le j\le K}d_{\mathrm{JS}}(P,e_j)\\
&=d_{\mathrm{JS}}(P,e_{j_*}),
\qquad j_*\in\arg\min_j p_j.
\end{aligned}}
\end{equation}
دلیل بخش نخست آن است که
$Q=\sum_jq_je_j$ و واگرایی
$\operatorname{JSD}(P,Q)$ نسبت به $Q$
\emph{محدب} است؛ پس
\[
\operatorname{JSD}(P,Q)
\le\sum_jq_j\operatorname{JSD}(P,e_j)
\le\max_j\operatorname{JSD}(P,e_j).
\]
برای تشخیص اینکه کدام رأس دورتر است،
از رابطه صعودی‌بودن
$s(P,e_j)$ نسبت به $p_j$ استفاده می‌کنیم:
\[
s(P,e_j)=
\frac{(1+p_j)\log(1+p_j)-p_j\log p_j}
{2\log2}.
\]
بنابراین کمترین مؤلفه $P$،
کمترین شباهت و بیشترین فاصله را
تولید می‌کند. در صورت برابر بودن
چند مؤلفه کمینه، چند رأس دورترین‌اند.
این گزاره \emph{وجود یک رأس بیشینه‌ساز}
را تضمین می‌کند، نه اینکه همه
بیشینه‌سازها حتماً رأس باشند:
مثلاً اگر $P=(1,0,0)$ باشد،
توزیع $Q=(0,\tfrac12,\tfrac12)$
نیز همان فاصله بیشینه $\sqrt{\log2}$
را دارد. اثبات شرایط محدب‌بودن و
مرز سیمپلکس در
پیوست~\ref{app:ch3-js-error-proof} آمده است.



\section{معیارهای تراکم بازنمایی‌ها و جدایی مراکز}
\label{cmp:sec:ch3-quality-criteria}

در بخش پیش، مراکز توزیعی طبقات را تعریف کردیم
و $R_c$ را به‌عنوان هزینهٔ تجمیع بازنمایی‌های
هم‌طبقه حول مرکز مشترک تفسیر کردیم.
اکنون کیفیت این بازنمایی‌ها و مراکز حاصل از
سیاست انتخاب ویژگی را می‌سنجیم. معیارها را بر اساس \emph{آنچه
اندازه می‌گیرند} از یکدیگر جدا می‌کنیم:
(۱) \emph{تراکم درون‌طبقه‌ای}، اختلاف
توزیع‌های برچسب همسایگی نمونه‌های هم‌طبقه
با مرکز همان طبقه را می‌سنجد؛
(۲) \emph{جدایی مراکز}، تفاوت توزیع‌های
نماینده طبقات مختلف را ارزیابی می‌کند؛
سنجه‌های اختلاف و شباهت میان توزیع‌های احتمال
در بخش پیش معرفی شدند. اکنون تنها معیارهای
تراکم و جدایی را صورت‌بندی می‌کنیم.
برتری امتیاز طبقهٔ واقعی بر رقیبان،
در بخش مستقل حاشیهٔ تصمیم بررسی می‌شود؛
پیوندهای کمی این معیارها با NLL نیز پس از
تعریف تابع زیان در بخش توابع هدف خواهد آمد.
تحلیل‌های اثباتی و روابط تکمیلی خود معیارها
در بخش~\ref{cmp:sec:ch3-quality-criteria-theory}
آمده‌اند؛ برای دنبال‌کردن سازوکار اصلی مدل
مطالعهٔ آن بخش الزامی نیست.

\paragraph{دو رویکرد ارزیابی کیفیت بازنمایی.}
برای سنجش کیفیت بازنمایی حاصل از گراف، دو رویکرد اصلی وجود دارد. تفاوت این رویکردها در آن است که آیا ابتدا یک افراز جدید از نمونه‌ها استخراج می‌کنیم یا افراز متناظر با برچسب‌های معلوم را ثابت در نظر می‌گیریم.

\noindent\textbf{رویکرد نخست: خوشه‌بندی و ارزیابی خارجی.}
در این رویکرد، برچسب نمونه‌ها در مرحله خوشه‌بندی استفاده نمی‌شود. هر نمونه با بردار توزیع همسایگی خود، یعنی سطر متناظر در ماتریس $\mathscr H$، نمایش داده می‌شود و یک الگوریتم خوشه‌بندی، افرازی از نمونه‌ها را بر اساس این بردارها به‌دست می‌آورد. سپس با استفاده از برچسب‌های واقعی که فقط برای ارزیابی نگه داشته شده‌اند، میزان انطباق خوشه‌های حاصل با طبقات توسط یک \textbf{معیار ارزیابی خارجی} سنجیده می‌شود. در این رویکرد نمی‌توان $\mathscr Y$ یا $\mathcal V$ را به‌عنوان ورودی خوشه‌بندی بدون برچسب به کار برد، زیرا این دو از ماتریس برچسب $Y$ ساخته شده‌اند. همچنین برای تولید $\mathscr H$ بدون برچسب نمونه آغازین، باید سیاستی مستقل از برچسب به کار رود؛ برای مثال می‌توان از سیاست میانگین رابطه~\eqref{cmp:eq:ch3-mean-policy-inference} استفاده کرد. علاوه بر این، نتیجه ارزیابی به انتخاب الگوریتم خوشه‌بندی، تعداد خوشه‌ها و پارامترهای آن نیز وابسته خواهد بود.

\noindent\textbf{رویکرد دوم: افراز ثابت طبقات و ارزیابی داخلی.}
در این رویکرد، چون برچسب نمونه‌های مجموعه آموزش مشخص است، افراز مجموعه‌داده را از ابتدا مطابق همین برچسب‌ها در نظر می‌گیریم و نیازی به اجرای الگوریتم خوشه‌بندی برای یافتن افراز جدید نداریم. سپس کیفیت ساختار همسایگیِ این افراز ثابت را با \textbf{معیارهای ارزیابی داخلی} می‌سنجیم؛ یعنی معیارهایی که بر اساس توزیع‌ها و روابط میان نمونه‌های متعلق به یک طبقه و طبقات متفاوت محاسبه می‌شوند. در اینجا برچسب‌ها برای مشخص کردن عضویت طبقات به کار می‌روند، اما کیفیت از روی ساختار بازنمایی حاصل از سیاست انتخاب ویژگی سنجیده می‌شود. به دلیل امکان تجمیع برچسب‌ها، این معیارها را می‌توان در بسیاری از موارد مستقیماً با ماتریس توزیع برچسب همسایگی $\mathscr Y$ و توزیع‌های نماینده طبقات در $\mathcal V$ محاسبه کرد.

در این پژوهش، \textbf{رویکرد دوم} را انتخاب می‌کنیم؛ زیرا برچسب نمونه‌های آموزش در دسترس است و استفاده از $\mathscr Y$ و $\mathcal V$، ارزیابی کیفیت را بدون وابستگی به انتخاب و اجرای یک الگوریتم خوشه‌بندی جداگانه امکان‌پذیر می‌کند. بنابراین معیارهایی که در ادامه این بخش معرفی می‌شوند، همگی از نوع معیارهای داخلی برای افراز ثابت مبتنی بر برچسب هستند و هدف آن‌ها فراهم کردن مبنایی برای بهبود سیاست انتخاب ویژگی است.



\subsection{معیارهای تراکم درون‌طبقه‌ای}
\label{cmp:subsec:ch3-within-class-compactness}
پس از معرفی سنجه‌های اختلاف توزیع‌ها،
تراکم را بر اساس فاصله توزیع برچسب همسایگی
هر نمونه از مرکز طبقه خودش می‌سنجیم.
در تمام معیارهای زیر مراکز $\mathcal V_c$
همان میانگین‌های حسابی پیشین‌اند؛
فقط سنجه اختلاف عوض می‌شود.

\paragraph{معیار مبتنی بر KL.}

در بخش~\ref{cmp:subsec:ch3-mean-center-justification}،
$R_c$ به‌عنوان هزینه اضافیِ مدل‌کردن
توزیع‌های برچسب همسایگی اعضای یک طبقه
با یک توزیع مشترک معرفی شد؛
تجزیه KL و آنتروپی در همان بخش نشان داد
که این هزینه برابر میانگین واگرایی
اعضا از مرکز حسابی طبقه است.
اکنون همان رابطه را \emph{با تعبیر تراکم}
دوباره می‌نویسیم:
\begin{equation}\label{cmp:eq:ch3-R-class}
R_c=\frac{1}{|C_c|}\sum_{i\in C_c}
D_{\mathrm{KL}}\bigl(\mathscr Y_i\|\mathcal V_c\bigr).
\end{equation}
در این رابطه $R_c$ یک معیار \emph{در سطح طبقه} است؛ برای ارزیابی کل مجموعه‌داده، میانگین وزن‌دار پراکندگی طبقات را به کار می‌بریم. اگر بردار وزن طبقات را با $\Omega=(\omega_1,\ldots,\omega_K)$ نشان دهیم، به‌طوری که $\omega_c\ge0$ و $\sum_c\omega_c=1$، معیار میانگین پراکندگی برابر است با:
\begin{equation}\label{cmp:eq:ch3-R-average}
\boxed{R=\sum_{c=1}^{K}\omega_c R_c
=\sum_{c=1}^{K}\frac{\omega_c}{|C_c|}
\sum_{i\in C_c}D_{\mathrm{KL}}\bigl(\mathscr Y_i\|\mathcal V_c\bigr).}
\end{equation}
برای وزن‌دهی یکسان به طبقات،
$\omega_c=1/K$ انتخاب می‌شود؛
نمایش وزن‌دار امکان اولویت‌دهی
به طبقات مختلف را نیز فراهم می‌کند.
از نامنفی بودن KL، $R\ge0$ نتیجه می‌شود.
کاهش این مقدار، نزدیک‌تر شدن توزیع‌های
برچسب همسایگی اعضای هر طبقه به
مرکز آن را \emph{در معنای KL} بیان می‌کند.
پس در اینجا تعریف مستقلی ساخته نشده است،
بلکه همان هزینه تجمیع پیشین را
به‌عنوان معیار تراکم تعبیر می‌کنیم:
$R_{c,\mathrm{KL}}:=R_c$ و
$R_{\mathrm{KL}}:=R$.


\paragraph{قرارداد تکیه‌گاه KL در برابر هموارسازی عددی.}
\phantomsection\label{cmp:par:ch3-KL-legacy-clip}
در تعریف دقیق از قرارداد $0\log(0/q)=0$ برای $q\ge0$
استفاده می‌شود و هرگاه $p_u>0$ ولی $q_u=0$ باشد،
$D_{\mathrm{KL}}(P\|Q)=+\infty$ است.
برای $R_{\mathrm{KL}}$ که مرکز همان میانگین حسابی
نمونه‌های هر طبقه است، چنین ناسازگاریِ تکیه‌گاه
در جهت $\mathscr Y_i\|\mathcal V_c$ رخ نمی‌دهد
و مقدار معیار متناهی است؛ اما در فاصله نمونه به
\emph{مرکز طبقه رقیب} یا هدف خارجی این تضمین وجود ندارد.
جایگزینی صفرها با کران پایین مثبت در محاسبهٔ
\lr{clip} یک قرارداد عددیِ متفاوت می‌سازد:
$R_{\mathrm{KL}}^{\mathrm{legacy,clip}}$
را نباید بدون باز‌محاسبه با تعریف دقیق
$R_{\mathrm{KL}}$ یکی گرفت. این تفاوت در
تفسیر خروجی‌های ذخیره‌شدهٔ فصل~\ref{chap:results}
اهمیت دارد، نه در درستی هویت‌های دقیقِ KL.

\paragraph{معیارهای مبتنی بر جنسن--شنون.}

به‌جای سنجیدن اختلاف با KL،
می‌توان آن را با واگرایی جنسن--شنون
و با خود متریک جذر آن اندازه گرفت.
با حفظ \emph{همان مراکز حسابی} و وزن طبقات،
دو معیار دیگر را تعریف می‌کنیم:
\begin{equation}\label{cmp:eq:ch3-R-JS-forms}
\boxed{
\begin{aligned}
R_{\mathrm{JS}}
&=\sum_{c=1}^{K}\frac{\omega_c}{|C_c|}
\sum_{i\in C_c}\operatorname{JSD}(\mathscr Y_i,\mathcal V_c),\\
R_{\sqrt{\mathrm{JS}}}
&=\sum_{c=1}^{K}\frac{\omega_c}{|C_c|}
\sum_{i\in C_c}\sqrt{\operatorname{JSD}(\mathscr Y_i,\mathcal V_c)}.
\end{aligned}
}
\end{equation}
در تحلیل طبقه $c$ نیز نمادهای $R_{c,\mathrm{JS}}$
و $R_{c,\sqrt{\mathrm{JS}}}$ را برای میانگین‌های
درون همان طبقه به کار می‌بریم. در این سه تعریف،
فقط تابع سنجش فاصله متفاوت است؛ تعریف
$\mathcal V_c$ و وزن‌های $\Omega$ تغییر نمی‌کنند.
با لگاریتم طبیعی، همواره
$0\le R_{\mathrm{JS}}\le\log2$ و
$0\le R_{\sqrt{\mathrm{JS}}}\le\sqrt{\log2}$ است.

به دلیل تقعر تابع ریشه دوم،
این دو روش تجمیع به‌طور کلی برابر نیستند:
\[
R_{\sqrt{\mathrm{JS}}}
=\mathbb E_{\Omega,C}\!\left[\sqrt{\operatorname{JSD}
(\mathscr Y_i,\mathcal V_c)}\right]
\le\sqrt{\mathbb E_{\Omega,C}\!\left[
\operatorname{JSD}(\mathscr Y_i,\mathcal V_c)\right]}
=\sqrt{R_{\mathrm{JS}}}.
\]
امید در اینجا دقیقاً به معنای وزن‌دهی طبقات با
$\omega_c$ و نمونه‌های هر طبقه به‌طور یکسان است.
تساوی تنها در حالت ثابت‌بودن تقریباً همهٔ
مقادیر JSD دارای وزن مثبت برقرار می‌شود؛
بنابراین ریشه‌گرفتن از یک ستون میانگین
جایگزین محاسبهٔ فاصلهٔ تک‌نمونه‌ای نیست.

\paragraph{تعبیر هندسی و شباهتی دو معیار.}
از برابری
$d_{\mathrm{JS}}(P,Q)^2
=\operatorname{JSD}(P,Q)
=\log2\,[1-s(P,Q)]$
می‌توان دو معیار را، بدون تغییر تعریف
مراکز، مستقیماً در قالب فاصله و شباهت نوشت.
با تعریف میانگین وزن‌دار شباهت به مرکز
\[
\overline s_{\mathrm{in}}
:=\sum_{c=1}^{K}\frac{\omega_c}{|C_c|}
\sum_{i\in C_c}s(\mathscr Y_i,\mathcal V_c),
\]
داریم:
\begin{equation}\label{cmp:eq:ch3-R-JS-similarity-interpretation}
\boxed{
\begin{aligned}
R_{\mathrm{JS}}
&=\sum_c\frac{\omega_c}{|C_c|}
\sum_{i\in C_c}
d_{\mathrm{JS}}(\mathscr Y_i,\mathcal V_c)^2
=\log2\,(1-\overline s_{\mathrm{in}}),\\
R_{\sqrt{\mathrm{JS}}}
&=\sum_c\frac{\omega_c}{|C_c|}
\sum_{i\in C_c}
d_{\mathrm{JS}}(\mathscr Y_i,\mathcal V_c)\\
&=\sqrt{\log2}\,
\sum_c\frac{\omega_c}{|C_c|}
\sum_{i\in C_c}
\sqrt{1-s(\mathscr Y_i,\mathcal V_c)}.
\end{aligned}}
\end{equation}
پس کمینه‌کردن $R_{\mathrm{JS}}$
\emph{دقیقاً معادل} بیشینه‌کردن
میانگین شباهت نرمال‌شده نمونه‌ها
به مراکز طبقات است.
اما $R_{\sqrt{\mathrm{JS}}}$
میانگین \emph{جذرِ یک منهای شباهت}
را می‌سنجد، نه «یک منهای
میانگین جذر شباهت»؛
جابه‌جایی عمل جذر و میانگین
در حالت کلی مجاز نیست.
در هندسه هیلبرتی،
هر دو به فاصله از تصویر
$\Phi(\mathcal V_c)$ مربوط‌اند.
بااین‌حال، چون نگاشت $\Phi$
خطی نیست، این تصویر لزوماً
\emph{میانگین هیلبرتی}
$\Phi(\mathscr Y_i)$های همان طبقه
نیست؛ ازاین‌رو نباید
$R_{\mathrm{JS}}$ را بدون توضیح
«واریانس حول میانگین هیلبرتی»
تعبیر کرد.


\noindent
در جمع‌بندی تعریف‌ها، $R_{\mathrm{KL}}$ میانگین
هزینهٔ KL نمونه‌ها نسبت به مرکز حسابی طبقه است؛
$R_{\mathrm{JS}}$ میانگین \emph{مربع فاصلهٔ متریک}
جنسن--شنون و $R_{\sqrt{\mathrm{JS}}}$ میانگین
\emph{خود این فاصله} تا همان مراکز را می‌سنجد.
این سه مقدار مقیاس و تعبیر یکسانی ندارند.
هویت‌های آنتروپی، کران‌ها و روابط دقیق میان
این معیارها در بخش~\ref{cmp:sec:ch3-quality-criteria-theory}
بررسی می‌شوند.

\paragraph{چرا کمینه‌سازی $R$ به‌تنهایی کافی نیست؟}
هر سه معیار تراکم بالا فقط نزدیک‌شدن توزیع‌ها \emph{درون هر طبقه} را می‌سنجند و هیچ شرط مستقیمی بر تمایز مراکز طبقات مختلف تحمیل نمی‌کند. برای مثال، فرض کنیم سیاستی موجب شود همه نمونه‌ها، بدون توجه به برچسب آغازین، توزیع همسایگی یکسانی مانند $q$ داشته باشند:
\begin{equation}\label{cmp:eq:ch3-R-collapse}
\mathscr Y_i=q\quad(\forall i)
\quad\Longrightarrow\quad
\mathcal V_1=\cdots=\mathcal V_K=q
\quad\Longrightarrow\quad
R_{\mathrm{KL}}=R_{\mathrm{JS}}=R_{\sqrt{\mathrm{JS}}}=0.
\end{equation}
در این وضعیت، هر سه معیار تراکم به کمترین مقدار ممکن رسیده‌اند، اما هیچ اطلاعاتی برای تشخیص طبقه نمونه آغازین از روی توزیع همسایگی آن باقی نمی‌ماند. بنابراین \emph{فشردگی درون‌طبقه‌ای} باید در کنار \emph{جدایی مراکز طبقات} بررسی شود. در ادامه، ابتدا معیارهای ساده مبتنی بر انتساب برچسب‌های مشاهده‌شده به طبقات آغازین را معرفی می‌کنیم و سپس معیاری را بررسی خواهیم کرد که جدایی توزیع‌های نماینده را به‌طور مستقیم اندازه می‌گیرد.


\subsection{معیارهای جدایی مراکز}
\label{cmp:subsec:ch3-center-separation-criteria}
معیارهای تراکم درون‌طبقه‌ای، اختلاف توزیع برچسب همسایگی
نمونه‌های یک طبقه با مرکز همان طبقه را کاهش می‌دهند؛
اما حتی با صفرشدن همه آن‌ها، ممکن است تمام مراکز
$\mathcal V_c$ یکسان باشند و تشخیص طبقات ممکن نشود.
ازاین‌رو در کنار تراکم، به معیارهایی برای
\emph{تمایز الگوهای برچسب همسایگی طبقات مختلف}
نیاز داریم. این معیارها روی ماتریس مراکز
$\mathcal V\in\mathbb R^{K\times K}$ تعریف می‌شوند:
سطر $c$ام توزیع احتمال مشاهده برچسب‌های همسایگی
برای مبدأیی از طبقه $c$ است و درایه
$\mathcal V_{cd}$ احتمال مشاهده برچسب مقصد $d$
را نشان می‌دهد. در این مرحله \emph{خود مراکز}
را می‌سنجیم؛ پراکندگی نمونه‌ها حول هر مرکز
باید جداگانه بررسی شود.

برای ساخت معیارهای جدایی، از ساده‌ترین
ایده آغاز می‌کنیم: \emph{افزایش احتمال
مؤلفه‌های برگزیده هر مرکز}.
معیار خطی $T$ مؤلفه‌های قطر اصلی و
$BT$ مؤلفه‌های متناظر با بهترین
جایگشت ستون‌ها را تقویت می‌کند.
سپس با تبدیل همین احتمال‌ها به
هزینه منفی لگاریتمی، به $TNL$ و $BTNL$
می‌رسیم؛ این دو علاوه بر تعبیر
درست‌نماییِ الگوی مرجع، معادل
واگرایی KL بردارهای یک‌داغ هدف
از مراکز هستند و در برابر
احتمال‌های بسیار کوچک حساسیت بیشتری دارند.
در پایان، الزام انتخاب یک مؤلفه
از هر مرکز را کنار می‌گذاریم
و از طریق معیار اطلاعات متقابل $I$،
تمایز \emph{تمام توزیع‌های مرکزی}
را در یک تعبیر اطلاعاتی بررسی می‌کنیم.
در هر مرحله تعریف، خواص و محدودیت معیار
را پیش از رفتن به مرحله بعد بیان خواهیم کرد.

\subsubsection{معیار خطی هم‌برچسبی: \texorpdfstring{$T$}{T}}
\label{cmp:subsubsec:ch3-T}
ساده‌ترین معیار، میانگین وزن‌دار احتمال مشاهده برچسبی برابر با برچسب طبقه آغازین است:
\begin{equation}\label{cmp:eq:ch3-T-on-V}
\boxed{T(\mathcal V)=\sum_{c=1}^{K}\omega_c\mathcal V_{cc}
=\operatorname{tr}\!\bigl(\operatorname{diag}(\Omega)\mathcal V\bigr).}
\end{equation}
از آنجا که هر درایه $\mathcal V_{cc}$ در بازه $[0,1]$ قرار دارد، $0\le T(\mathcal V)\le1$. بیشینه‌سازی $T$، سهم مشاهده نمونه‌های هم‌برچسب با طبقه آغازین را افزایش می‌دهد. تعبیر این معیار روی $\mathscr Y$ و سپس روی $\mathscr H$، با استفاده از تعریف مراکز طبقات، چنین است:
\begin{equation}\label{cmp:eq:ch3-T-on-Y-H}
T(\mathcal V)=
\sum_{c=1}^{K}\frac{\omega_c}{|C_c|}
\sum_{i\in C_c}\mathscr Y_i(c)
=
\sum_{c=1}^{K}\frac{\omega_c}{|C_c|}
\sum_{i\in C_c}\sum_{j\in C_c}h_{ij}.
\end{equation}
بنابراین همین معیار را می‌توان بدون نگهداری $\mathscr H$ و فقط از روی $\mathscr Y$ به‌دست آورد. بااین‌حال، $T$ به‌طور ضمنی فرض می‌کند که بهترین الگوی همسایگی هر طبقه باید مشاهده غالب \emph{همان} برچسب باشد. چنین فرضی برای همه مجموعه‌داده‌ها ضروری نیست؛ همان‌طور که در بخش~\ref{cmp:subsec:ch3-label-neighborhood-interpretation} توضیح داده شد، ممکن است توزیع‌های همسایگی طبقات به‌خوبی از هم قابل تشخیص باشند، ولی مؤلفه قطری بزرگ‌ترین سهم آن‌ها نباشد.

\subsubsection{معیار خطی جایگشتی: \texorpdfstring{$BT$}{BT}}
\label{cmp:subsubsec:ch3-BT}
برای حذف الزام انطباق مستقیم برچسب مشاهده‌شده با برچسب طبقه آغازین، به هر طبقه آغازین یک برچسب مقصد اختصاص می‌دهیم؛ به شرط آنکه این انتساب برای تمام طبقات یک‌به‌یک باشد. اگر $S_K$ مجموعه همه جایگشت‌های $K$ برچسب و $\sigma(c)$ برچسب مقصد منتسب به طبقه $c$ باشد، معیار بهترین اثر جایگشتی را به شکل زیر تعریف می‌کنیم:
\begin{equation}\label{cmp:eq:ch3-BT-on-V}
\boxed{BT(\mathcal V)=
\max_{\sigma\in S_K}
\sum_{c=1}^{K}\omega_c\mathcal V_{c,\sigma(c)}.}
\end{equation}
در این تعریف، معیار $T$ حالت خاصی است که جایگشت همانی $\sigma(c)=c$ به کار رود؛ بنابراین $BT(\mathcal V)\ge T(\mathcal V)$. تعبیر این معیار روی دو بازنمایی پیشین برابر است با:
\begin{equation}\label{cmp:eq:ch3-BT-on-Y-H}
\begin{aligned}
BT(\mathcal V)
&=\max_{\sigma\in S_K}
\sum_{c=1}^{K}\frac{\omega_c}{|C_c|}
\sum_{i\in C_c}\mathscr Y_i\bigl(\sigma(c)\bigr)\\
&=\max_{\sigma\in S_K}
\sum_{c=1}^{K}\frac{\omega_c}{|C_c|}
\sum_{i\in C_c}\sum_{j\in C_{\sigma(c)}}h_{ij}.
\end{aligned}
\end{equation}
در نتیجه، به‌جای اجبار به مشاهده هم‌برچسب، این معیار می‌پذیرد که هر طبقه الگویی از همسایگی داشته باشد که برچسب مقصد غالب در آن ممکن است با برچسب آغازین متفاوت باشد. قید جایگشتی، از اختصاص یک برچسب مقصد واحد به چند طبقه آغازین جلوگیری می‌کند و می‌تواند به تفکیک مراکز کمک کند. مسئله یافتن جایگشت بهینه یک مسئله انتساب بیشینه‌وزن است و لازم نیست همه $K!$ جایگشت به‌صورت مستقیم بررسی شوند.

برای تعیین کران پایین، کافی است توجه کنیم که میانگین امتیاز جایگشت‌ها، حتی با وزن‌های نامساوی نرمال‌شده، برابر $1/K$ است؛ در نتیجه بهترین جایگشت نمی‌تواند امتیازی کمتر از این میانگین داشته باشد. اثبات شمارشی کامل این نتیجه، از جمله حالت وزن یکسان $\omega_c=1/K$، در بخش~\ref{app:ch3-BT-bound-proof} از پیوست~\ref{app:ch3-proofs} آمده است. بنابراین:
\begin{equation}\label{cmp:eq:ch3-BT-bounds}
\boxed{\frac{1}{K}\le BT(\mathcal V)\le1.}
\end{equation}
چنانچه تمام وزن‌های طبقات مثبت باشند، $BT(\mathcal V)=1$ دقیقاً زمانی برقرار است که $\mathcal V$ یک ماتریس جایگشتی باشد؛ یعنی در هر سطر و ستون دقیقاً یک درایه برابر یک و سایر درایه‌ها صفر باشند. اثبات این شرط نیز در بخش~\ref{app:ch3-BT-bound-proof} از پیوست~\ref{app:ch3-proofs} آمده است. در این حالت، برچسب همسایگی هر طبقه، حتی اگر با نام خود آن طبقه یکسان نباشد، می‌تواند آن طبقه را به‌طور یکتا مشخص کند.


\subsubsection{معیار لگاریتمی هم‌برچسبی: \texorpdfstring{$TNL$}{TNL}}
\label{cmp:subsubsec:ch3-TNL}
معیار $T$ میانگین حسابی احتمال مشاهده برچسبی برابر با طبقه آغازین را اندازه می‌گیرد. برای سنجش میزان انطباق مراکز با یک \emph{الگوی ایدئال هم‌برچسبی}، فرض کنیم در یک آزمایش ارزیابی، برچسب مطلوبِ مشاهده‌شده در همسایگی نمونه‌ای از طبقه $c$ به‌طور قطعی همان $c$ باشد؛ توزیع این برچسب مطلوب را با بردار یک‌داغ $e_c$ نشان می‌دهیم. این فرض، \emph{معیار مرجع} را تعریف می‌کند و به معنای آن نیست که توزیع واقعی برچسب‌های همسایگی یا احتمال پسین برچسب نمونه الزاماً $e_c$ است.

از سوی دیگر، مرکز $\mathcal V_c$ توزیع احتمالی برچسب‌های مشاهده‌شده در همسایگی نمونه‌های طبقه $c$ است. بنابراین در آزمایش فرضی بالا، اگر تنها برچسب $c$ مشاهده شود، احتمال اختصاص‌یافته به آن از سوی مدل مرکزی برابر $\mathcal V_c(c)=\mathcal V_{cc}$ خواهد بود. میانگین \emph{منفی لگاریتم درست‌نمایی مورد انتظار} این مشاهده ایدئال تحت مدل $\mathcal V_c$، برابر $-\log\mathcal V_{cc}$ است. با وزن‌دهی به طبقات توسط $\Omega=(\omega_1,\ldots,\omega_K)$، معیار زیر را تعریف می‌کنیم:
\begin{equation}\label{cmp:eq:ch3-TNL-on-V}
\boxed{
\begin{aligned}
TNL(\mathcal V)
&=-\sum_{c=1}^{K}\omega_c
\sum_{d=1}^{K}e_c(d)\log\mathcal V_c(d)\\
&=-\sum_{c=1}^{K}\omega_c\log\mathcal V_{cc}\\
&=\sum_{c=1}^{K}\omega_c
D_{\mathrm{KL}}(e_c\|\mathcal V_c).
\end{aligned}}
\end{equation}
در این تعریف وزن‌ها نامنفی و دارای مجموع یک‌اند؛
برای بیان شرایط تساوی و مقایسه معیارهای لگاریتمی
در ادامه، فرض می‌کنیم $\omega_c>0$ برای همه طبقات.
قرارداد $-\log0=+\infty$ را نیز به کار می‌بریم.
این معیار نخست از درست‌نماییِ \emph{مشاهدات فرضی
مطابق الگوی ایدئال} به‌دست آمد؛
برابری آخر در رابطه~\eqref{cmp:eq:ch3-TNL-on-V}
\emph{تعبیر معادل} آن را برحسب KL روشن می‌کند:
چون $H(e_c)=0$ و
$D_{\mathrm{KL}}(e_c\|\mathcal V_c)
=-\log\mathcal V_{cc}$،
هر طبقه به اندازه واگرایی
\emph{بردار یک‌داغ مرجع از مرکز آن طبقه}
جریمه می‌شود. جهت این واگرایی مهم است:
$D_{\mathrm{KL}}(\mathcal V_c\|e_c)$
در حالت کلی نامتناهی است و نمی‌تواند
جایگزین رابطه بالا شود.

تفاوت این تعبیر با بخش~\ref{cmp:subsec:ch3-mean-center-justification} مهم است: در آنجا کمترین منفی لگاریتم درست‌نمایی برای مدل‌کردن \emph{مشاهدات واقعی همسایگی} با یک نماینده مشترک برابر $H(\mathcal V_c)$ بود و $R_c$ افزایش هزینه ناشی از تجمیع توزیع‌های نمونه‌ها را می‌سنجید؛ اما $TNL$ می‌سنجد که مراکز فعلی تا چه اندازه مشاهده \emph{هدفِ قطعیِ هم‌برچسبی} را محتمل می‌دانند. بنابراین $TNL$ نه هزینه تجمیع $R$ است و نه منفی لگاریتم درست‌نمایی برچسب واقعی در یک طبقه‌بند.

از آنجا که $0\le\mathcal V_{cc}\le1$، معیار $TNL$ نامنفی است و باید کمینه شود. با فرض مثبت بودن همه وزن‌ها، $TNL(\mathcal V)=0$ دقیقاً زمانی رخ می‌دهد که $\mathcal V$ ماتریس همانی باشد. همچنین:
\begin{equation}\label{cmp:eq:ch3-TNL-geometric-mean}
\boxed{\exp\!\bigl(-TNL(\mathcal V)\bigr)
=\prod_{c=1}^{K}\mathcal V_{cc}^{\omega_c}.}
\end{equation}
پس کمینه‌کردن $TNL$ معادل بیشینه‌کردن میانگین هندسی وزن‌دارِ احتمال‌های هم‌برچسبی است؛ در مقابل، $T$ میانگین حسابی وزن‌دار همان احتمال‌ها را بیشینه می‌کند.

برای محاسبه معیار از بازنمایی‌های همسایگی، درایه‌های قطری مراکز را با میانگین‌های متناظر جایگزین می‌کنیم:
\begin{equation}\label{cmp:eq:ch3-TNL-on-Y-H}
\begin{aligned}
TNL(\mathcal V)
&=-\sum_{c=1}^{K}\omega_c\log
\left(\frac{1}{|C_c|}
\sum_{i\in C_c}\mathscr Y_i(c)\right)\\
&=-\sum_{c=1}^{K}\omega_c\log
\left(\frac{1}{|C_c|}
\sum_{i\in C_c}\sum_{j\in C_c}h_{ij}\right).
\end{aligned}
\end{equation}
توجه می‌کنیم که لگاریتم \emph{پس از میانگین‌گیری درون هر طبقه} محاسبه می‌شود؛ بنابراین این معیار با میانگین منفی لگاریتم احتمال هم‌برچسبی تک‌تک نمونه‌ها یکسان نیست.

برای مثال، با وزن برابر برای سه طبقه، مقادیر قطری $(1,1,0.4)$ و $(0.8,0.8,0.8)$ هر دو $T=0.8$ دارند؛ اما $TNL$ آن‌ها به‌ترتیب $-\log(0.4)/3$ و $-\log(0.8)$ است. بنابراین معیار لگاریتمی، افت شدید احتمال در یک طبقه را بیشتر جریمه می‌کند. از تقعر لگاریتم و نامساوی ینسن نیز داریم:
\begin{equation}\label{cmp:eq:ch3-TNL-Jensen}
\boxed{TNL(\mathcal V)\ge-\log T(\mathcal V).}
\end{equation}
این نامساوی کرانی میان دو معیار است، نه تضمینی برای یکسان بودن ترتیب آن‌ها در همه ماتریس‌های مراکز. افزون بر این، الگوی هم‌برچسبی قطعی فقط یک هدف ارزیابی است؛ ممکن است مراکز متمایزی داشته باشیم که مؤلفه غالب آن‌ها با برچسب طبقه آغازین یکسان نباشد. ازاین‌رو $TNL$ نیز به‌تنهایی معیار کامل کیفیت طبقه‌بندی نیست.

\subsubsection{معیار لگاریتمی جایگشتی: \texorpdfstring{$BTNL$}{BTNL}}
\label{cmp:subsubsec:ch3-BTNL}
مانند تعمیم $T$ به $BT$، می‌توان الزام هم‌برچسب‌بودن را از معیار لگاریتمی برداشت. برای هر جایگشت $\sigma\in S_K$، الگوی مطلوب فرضی طبقه $c$ را توزیع قطعی $e_{\sigma(c)}$ در نظر می‌گیریم. در این صورت مدل مرکزی $\mathcal V_c$ به برچسب مطلوب $\sigma(c)$ احتمال $\mathcal V_{c,\sigma(c)}$ می‌دهد و منفی لگاریتم درست‌نمایی وزن‌دار آن برابر $-\sum_c\omega_c\log\mathcal V_{c,\sigma(c)}$ است. انتخاب جایگشت با کمترین هزینه، معیار زیر را به‌دست می‌دهد:
\begin{equation}\label{cmp:eq:ch3-BTNL-on-V}
\boxed{
\begin{aligned}
BTNL(\mathcal V)
&=\min_{\sigma\in S_K}
\left[-\sum_{c=1}^{K}\omega_c
\log\mathcal V_{c,\sigma(c)}\right]\\
&=\min_{\sigma\in S_K}
\sum_{c=1}^{K}\omega_c
D_{\mathrm{KL}}(e_{\sigma(c)}\|\mathcal V_c).
\end{aligned}}
\end{equation}
این تعریف نیز از درست‌نمایی \emph{الگوی هدف جایگشتی}
به‌دست می‌آید، نه از فرض قطعی‌بودن
توزیع واقعی همسایگی طبقات.
از آنجا که $H(e_{\sigma(c)})=0$،
هر هزینه لگاریتمی برابر
$D_{\mathrm{KL}}(e_{\sigma(c)}\|\mathcal V_c)$
است؛ بنابراین $BTNL$ از میان
تمام انتساب‌های یک‌به‌یکِ بردارهای
یک‌داغ مرجع به مراکز، انتسابی را
برمی‌گزیند که \emph{کمترین مجموع
واگرایی جهت‌دار از الگوهای هدف}
را داشته باشد. در $TNL$ این انتساب
از پیش ثابت و همانی است.
پس با انتخاب جایگشت همانی
به $TNL$ بازمی‌گردیم و
$BTNL(\mathcal V)\le TNL(\mathcal V)$.
اگر احتمال انتخاب‌شده در یک جایگشت
برای طبقه‌ای با وزن مثبت صفر باشد،
هزینه آن جایگشت بی‌نهایت است؛
اما وجود جایگشتی با تمام احتمال‌های
انتخاب‌شده مثبت، برای متناهی‌بودن
$BTNL$ کافی است.

نمایش معادل این معیار روی $\mathscr Y$ و $\mathscr H$ چنین است:
\begin{equation}\label{cmp:eq:ch3-BTNL-on-Y-H}
\begin{aligned}
BTNL(\mathcal V)
&=\min_{\sigma\in S_K}
\left[-\sum_{c=1}^{K}\omega_c\log
\left(\frac{1}{|C_c|}
\sum_{i\in C_c}\mathscr Y_i(\sigma(c))\right)\right]\\
&=\min_{\sigma\in S_K}
\left[-\sum_{c=1}^{K}\omega_c\log
\left(\frac{1}{|C_c|}
\sum_{i\in C_c}\sum_{j\in C_{\sigma(c)}}h_{ij}\right)\right].
\end{aligned}
\end{equation}
پس این معیار نیز از ماتریس کوچک‌تر $\mathcal V$ محاسبه می‌شود و مسئله یافتن جایگشت کمینه‌کننده، یک مسئله انتساب با هزینه‌های $-\omega_c\log\mathcal V_{cd}$ است.

\paragraph{ویژگی‌های معیار جایگشتی لگاریتمی.}
از نامنفی بودن منفی لگاریتم احتمال، داریم $BTNL(\mathcal V)\ge0$. اگر همه وزن‌ها مثبت باشند، این معیار دقیقاً زمانی به صفر می‌رسد که $\mathcal V$ یک ماتریس جایگشتی باشد. بنابراین $BT$ و $BTNL$ در فضای همه ماتریس‌های احتمالی سطری، به‌ترتیب بیشینه و کمینه جهانی خود را روی مجموعه یکسانی از ماتریس‌های جایگشتی به‌دست می‌آورند. این اشتراک در جواب ایدئال، به معنای یکسان بودن رفتار آن‌ها در وضعیت‌های غیربهینه نیست؛ $BT$ بر میانگین حسابی احتمال‌ها و $BTNL$ بر میانگین لگاریتم آن‌ها تکیه دارد.

همچنین، با اعمال نامساوی ینسن برای هر جایگشت و سپس گرفتن کمینه روی جایگشت‌ها، داریم:
\begin{equation}\label{cmp:eq:ch3-BTNL-Jensen}
\boxed{
BTNL(\mathcal V)\ge-\log BT(\mathcal V).
}
\end{equation}
این نامساوی نشان می‌دهد که هزینه لگاریتمی جایگشتی، از پایین به وسیله بهترین میانگین حسابی جایگشتی کنترل می‌شود؛ اما تضمین نمی‌کند جایگشتی که $BT$ را بیشینه می‌کند، لزوماً همان جایگشتی باشد که $BTNL$ را کمینه می‌کند.

این چهار معیار، با وجود تفاوت در شکل و حساسیت، همگی بر انتخاب \emph{یک مؤلفه از هر توزیع مرکزی} با انتساب ثابت یا جایگشتی تکیه می‌کنند. در ادامه، برای استفاده از اطلاعات کامل توزیع‌های مرکزی و سنجش تمایز آن‌ها بدون انتخاب مؤلفه اختصاصی، معیار اطلاعات متقابل را معرفی می‌کنیم.


\subsubsection{معیار جدایی اطلاعاتی: \texorpdfstring{$I$}{I}}
\label{cmp:subsubsec:ch3-I-criterion}
چهار معیار پیشین، یک مؤلفه معین
از هر توزیع مرکزی را با الگوی هدف
ثابت یا جایگشتی مقایسه می‌کنند.
اما ممکن است تفاوت طبقات
در \emph{تمام الگوی توزیع برچسب‌های همسایگی}
نهفته باشد، بدون آنکه برچسب مقصد
خاصی برای هر طبقه مطلوب باشد.
ازاین‌رو در گام پایانی،
\textbf{معیار اطلاعات متقابل} $I$
را برای سنجش جدایی اطلاعاتی
\emph{تمام مراکز} به کار می‌بریم.
اگر مراکز یکسان شوند، مشاهده
برچسب همسایگی تمایزی میان
طبقات آغازین ایجاد نمی‌کند؛
اما متفاوت‌بودن مراکز می‌تواند
عدم‌قطعیت درباره طبقه آغازین
را کاهش دهد. ابتدا $I$ را
از میانگین واگرایی مراکز تا
توزیع مشترک تعریف می‌کنیم،
سپس تعبیر دقیق آن به‌عنوان
اطلاعات متقابل را نشان می‌دهیم.

اگر $\Omega=(\omega_1,\ldots,\omega_K)$ بردار وزن طبقات باشد، میانگین وزن‌دار مراکز را با $\overline{\mathcal V}$ نشان می‌دهیم:
\begin{equation}\label{cmp:eq:ch3-global-center}
\overline{\mathcal V}
=\sum_{c=1}^{K}\omega_c\mathcal V_c
=\Omega\mathcal V.
\end{equation}
اکنون میانگین واگرایی مراکز طبقات از این توزیع مشترک را به‌صورت زیر تعریف می‌کنیم:
\begin{equation}\label{cmp:eq:ch3-I-on-V}
\boxed{I(\mathcal V)=
\sum_{c=1}^{K}\omega_c
D_{\mathrm{KL}}\bigl(\mathcal V_c\|\overline{\mathcal V}\bigr).}
\end{equation}
این معیار همان واگرایی جنسن--شنون وزن‌دار میان توزیع‌های نماینده طبقات است. برخلاف $T$ و نسخه‌های جایگشتی و لگاریتمی آن، برای محاسبه $I$ همه مؤلفه‌های مراکز طبقات در نظر گرفته می‌شوند و نیازی به انتخاب برچسب مقصد اختصاصی برای هر طبقه نیست.

\paragraph{تعبیر اطلاعاتی: اطلاعات متقابل طبقات.}
فرض کنیم $C_0$ برچسب \emph{نمونه آغازین} و $C_\xi$ برچسب \emph{نمونه مشاهده‌شده در همسایگی} باشد. بر اساس تعریف $\mathcal V$، طبقه آغازین با احتمال $\omega_c$ انتخاب می‌شود و احتمال مشاهده برچسب $d$ به شرط شروع از طبقه $c$ برابر $\mathcal V_{cd}$ است. بنابراین توزیع مشترک و توزیع حاشیه‌ای برچسب مشاهده‌شده دقیقاً عبارت‌اند از:
\begin{equation}\label{cmp:eq:ch3-I-joint-distribution}
\mathrm P(C_0=c,C_\xi=d)=\omega_c\mathcal V_{cd},
\qquad
\mathrm P(C_\xi=d)=\sum_{c=1}^{K}\omega_c\mathcal V_{cd}
=\overline{\mathcal V}(d).
\end{equation}
اگر دو برچسب آغازین و مشاهده‌شده مستقل باشند، توزیع مشترک آن‌ها برابر حاصل‌ضرب توزیع‌های حاشیه‌ای خواهد بود؛ یعنی احتمال زوج $(c,d)$ برابر $\omega_c\overline{\mathcal V}(d)$ می‌شود. اطلاعات متقابل، واگرایی KL توزیع مشترک واقعی از توزیع حاصل از این فرض استقلال است. بنابراین:
\begin{equation}\label{cmp:eq:ch3-I-mutual-info}
\begin{aligned}
\mathcal I(C_0;C_\xi)
&=D_{\mathrm{KL}}\!\left(
\mathrm P(C_0,C_\xi)
\,\middle\|\,
\mathrm P(C_0)\mathrm P(C_\xi)\right)\\
&=\sum_{c=1}^{K}\sum_{d=1}^{K}
\omega_c\mathcal V_{cd}
\log\frac{\omega_c\mathcal V_{cd}}
{\omega_c\overline{\mathcal V}(d)}\\
&=\sum_{c=1}^{K}\omega_c
D_{\mathrm{KL}}(\mathcal V_c\|\overline{\mathcal V})
=\boxed{I(\mathcal V)}.
\end{aligned}
\end{equation}
در عبارت نخست، $\mathrm P(C_0,C_\xi)$ توزیع مشترک دو متغیر و $\mathrm P(C_0)\mathrm P(C_\xi)$ حاصل‌ضرب توزیع‌های حاشیه‌ای آن‌ها را نشان می‌دهد.
این محاسبه دلیل نام‌گذاری \textbf{معیار اطلاعات متقابل} را روشن می‌کند: $I$ بیان می‌کند دانستن برچسب \emph{نمونه مشاهده‌شده در همسایگی} تا چه اندازه درباره طبقه \emph{نمونه آغازین} اطلاعات فراهم می‌کند؛ این تعریف هنوز به معنای مشاهده مستقیم برچسب واقعی نمونه آغازین یا یک احتمال پیش‌بینی‌شده از کل بردار $\mathscr Y_i$ نیست.

\subsection{چرا تراکم درون‌طبقه‌ای و جدایی مراکز هر دو لازم‌اند؟}
\label{cmp:subsec:ch3-compactness-separation}
اکنون می‌توان دو ویژگی مطلوب بازنمایی را با \emph{یک مقیاس فاصله} کنار هم قرار داد. ویژگی نخست، \textbf{فشردگی درون‌طبقه‌ای} است؛ یعنی توزیع‌های همسایگی نمونه‌های یک طبقه به مرکز همان طبقه نزدیک باشند. با استفاده از مربع فاصله جنسن--شنون، پراکندگی هر طبقه را به‌صورت زیر تعریف می‌کنیم:
\begin{equation}\label{cmp:eq:ch3-js-within-dispersion}
\boxed{
W_c^{\mathrm{JS}}=R_{c,\mathrm{JS}}
=\frac1{|C_c|}\sum_{i\in C_c}
d_{\mathrm{JS}}(\mathscr Y_i,\mathcal V_c)^2.
}
\end{equation}
مقدار کوچک‌تر $W_c^{\mathrm{JS}}$ نشان می‌دهد بازنمایی‌های طبقه $c$ در اطراف مرکز تعریف‌شده آن فشرده‌ترند. توجه کنیم $\mathcal V_c$ از اصل درست‌نمایی و میانگین حسابی به‌دست آمده است؛ در اینجا ادعا نمی‌کنیم که لزوماً کمینه‌کننده مجموع فاصله‌های مربعی جنسن--شنون نیز هست.

ویژگی دوم، \textbf{جدایی مراکز} است. در طبقه‌بندی نزدیک‌ترین مرکز، فاصله از نزدیک‌ترین رقیب اهمیت دارد؛ ازاین‌رو برای هر طبقه تعریف می‌کنیم:
\begin{equation}\label{cmp:eq:ch3-js-center-separation}
\boxed{\delta_c
=\min_{r\ne c}d_{\mathrm{JS}}(\mathcal V_c,\mathcal V_r),
\qquad K\ge2.}
\end{equation}
بزرگ بودن $\delta_c$ به معنی دور بودن مرکز طبقه $c$ از همه مراکز رقیب است. درحالی‌که معیار اطلاعات متقابل $I$ میزان تمایز کلی مراکز را می‌سنجد، به‌تنهایی کوچک نبودن همه $\delta_c$ها را تضمین نمی‌کند؛ برای نمونه، با مثبت بودن وزن همه طبقات، ممکن است $\mathcal V_1=\mathcal V_2\ne\mathcal V_3$ و در نتیجه $I>0$ باشد، ولی $\delta_1=\delta_2=0$ بماند.

این دو کمیت جنبه‌های متفاوت کیفیت را می‌سنجند و باید \emph{هم‌زمان} بررسی شوند: کوچک بودن $W_c^{\mathrm{JS}}$ بدون جدایی مراکز برای تصمیم کافی نیست، و بزرگ بودن $\delta_c$ نیز اگر نمونه‌های طبقه از مرکز خود دور باشند، تضمین‌کننده تشخیص درست نیست. در واقع از نامساوی مثلثی نتیجه می‌شود که هر نمونه از طبقه $c$ با شرط
$d_{\mathrm{JS}}(\mathscr Y_i,\mathcal V_c)<\delta_c/2$
به مرکز واقعی خود از همه مراکز رقیب نزدیک‌تر است. بنابراین اندازه پراکندگی در مقایسه با \emph{مجذور فاصله نزدیک‌ترین رقیب}، نه هر کدام به‌تنهایی، برای کنترل خطای تصمیم اهمیت دارد. صورت کمی کران تجربی خطا، که در آزمایش‌های این مرحله به کار نمی‌رود، در بخش~\ref{app:ch3-js-error-proof} از پیوست، ذیل اثبات کران هندسی خطای تصمیم، آمده است.

این چارچوب همچنین جایگاه معیارهای پیشین را روشن می‌کند: $R$ هزینه تجمیع توزیع‌های نمونه‌ها را با زیان لگاریتمی می‌سنجد؛ $T$، $BT$، $TNL$ و $BTNL$ انطباق مراکز با الگوهای هدف مشخص را ارزیابی می‌کنند؛ و $I$ جدایی اطلاعاتی مراکز را بدون بررسی مستقیم پراکندگی تک‌تک نمونه‌ها اندازه می‌گیرد. هیچ‌یک از این کمیت‌ها به‌تنهایی هر دو شرط فشردگی و جدایی را در قاعده تصمیم نزدیک‌ترین مرکز کنترل نمی‌کنند.

در این مرحله معیارهای تراکم و جدایی را معرفی کردیم.
اکنون باید ارتباط آن‌ها با درست‌بودن تصمیم
هر نمونه بررسی شود. این پرسش در
بخش~\ref{cmp:sec:ch3-normalized-margin}
با مفهوم حاشیهٔ تصمیم دنبال می‌شود.
بخش بعد صرفاً به روابط نظری تکمیلی
معیارهای معرفی‌شده اختصاص دارد.



\section{خواص نظری و روابط تکمیلی معیارهای تراکم و جدایی}
\label{cmp:sec:ch3-quality-criteria-theory}
این بخش هویت‌های آنتروپی، کران‌ها، روابط
مقایسه‌ای و شرایط برابری معیارهای تعریف‌شده
در بخش~\ref{cmp:sec:ch3-quality-criteria}
را بررسی می‌کند. این نتایج برای
\emph{تعریف و اجرای سازوکار اصلی مدل}
ضروری نیستند؛ خواننده‌ای که صرفاً روش
طبقه‌بندی و یادگیری سیاست را دنبال می‌کند،
می‌تواند مستقیماً به
بخش~\ref{cmp:sec:ch3-normalized-margin}
برود. برای مطالعهٔ تحلیلی معیارها و
تفسیر دقیق حدود آن‌ها، روابط زیر مکمل
تعاریف پیشین‌اند.

\subsection{هویت‌ها و کران‌های معیارهای تراکم}
\paragraph{تعبیر آنتروپی هزینه تجمیع.}
طبق رابطه~\eqref{cmp:eq:ch3-R-entropy}، هزینه هر طبقه از اختلاف آنتروپی مرکز و میانگین آنتروپی نمونه‌های همان طبقه به‌دست می‌آید. با وزن‌دهی به طبقات بر اساس $\Omega$، رابطه مشابه زیر برای $R$ حاصل می‌شود:
\begin{equation}\label{cmp:eq:ch3-R-average-entropy}
\boxed{R=\sum_{c=1}^{K}\omega_c H(\mathcal V_c)
-\sum_{c=1}^{K}\frac{\omega_c}{|C_c|}
\sum_{i\in C_c}H(\mathscr Y_i).}
\end{equation}
این رابطه نشان می‌دهد پراکندگی درون‌طبقه‌ای را نمی‌توان صرفاً از آنتروپی مرکز تعیین کرد؛ تفاوت آن با میانگین آنتروپی نمونه‌ها اهمیت دارد. همچنین چون $\mathcal V_c$ میانگین توزیع‌های طبقه است، تقعر آنتروپی نامنفی بودن $R_c$ را نیز نتیجه می‌دهد. اثبات هویت آنتروپی مخلوط که مبنای
این رابطه است، در بخش~\ref{app:ch3-entropy-mixture}
از پیوست~\ref{app:ch3-proofs} آمده است.

\paragraph{کران‌های تراکم KL و شرایط برابری.}
از رابطه~\eqref{cmp:eq:ch3-R-class} و نامنفی بودن
واگرایی KL، بی‌درنگ $R_c\ge0$ نتیجه می‌شود.
همچنین از تعبیر آنتروپی
رابطه~\eqref{cmp:eq:ch3-R-entropy}، با توجه به
$0\le H(P)\le\log K$ برای هر $P\in\Delta_K$، داریم:
\begin{equation}\label{cmp:eq:ch3-R-KL-bounds}
\boxed{
0\le R_c
=H(\mathcal V_c)
-\frac1{|C_c|}\sum_{i\in C_c}H(\mathscr Y_i)
\le H(\mathcal V_c)\le\log K.
}
\end{equation}
در این پایان‌نامه لگاریتم طبیعی به کار می‌رود؛
ازاین‌رو کران بالای دقیقِ حاصل از آنتروپی
$\log K$ است، نه $\sqrt K$.
کران پایین زمانی به‌دست می‌آید که همه
توزیع‌های برچسب همسایگی اعضای طبقه یکسان باشند.
برای دستیابی به کران بالای $\log K$،
لازم است همه این توزیع‌ها یک‌داغ باشند
(آنتروپی صفر) و در عین حال مرکز آن‌ها
یکنواخت روی $K$ برچسب باشد
(آنتروپی بیشینه).
به عبارت دیگر، برچسب‌های یک‌داغ باید
به نسبت برابر در میان اعضای آن طبقه دیده شوند؛
با وزن یکسان نمونه‌ها، این تساوی دقیق تنها
وقتی تعداد نمونه‌ها امکان چنین تقسیم برابری را
فراهم کند، قابل‌دستیابی است.
از آنجا که $R=\sum_c\omega_cR_c$
با وزن‌های نامنفی و مجموع یک است،
\begin{equation}\label{cmp:eq:ch3-R-weighted-bounds}
\boxed{0\le R\le\log K.}
\end{equation}
در صورت مثبت‌بودن همه وزن‌های طبقات،
تساوی در کران بالای $R$ مستلزم آن است
که هر طبقه دارای همان ساختار یک‌داغ
و مرکز یکنواخت باشد. این نتیجه از
ویژگی کلی میانگین وزن‌دار ناشی می‌شود
و وابستگی تازه‌ای به سیاست گشت ندارد.
به‌خصوص هرچند KL برای دو توزیع دلخواه
ممکن است نامتناهی باشد، در تعریف $R_c$،
مرکز $\mathcal V_c$ میانگین خود
$\mathscr Y_i$هاست؛ بنابراین
$D_{\mathrm{KL}}(\mathscr Y_i\|\mathcal V_c)$
برای هر عضو طبقه متناهی خواهد بود.
استدلال شرایط برابری و این نکته در
پیوست~\ref{app:ch3-compactness-entropy-relations} آمده است.


\paragraph{رابطه آنتروپی تراکم جنسن--شنون با هزینه تجمیع.}
برای نشان‌دادن ارتباط دقیق این معیار با
$R_c$، در هر طبقه $c$ از توزیع‌های
برچسب همسایگی و مرکز همان طبقه،
توزیع‌های \emph{میانگین جفتی} می‌سازیم:
\begin{equation}\label{cmp:eq:ch3-R-JS-midpoint-distributions}
\boxed{
\mathscr B_{c,i}
=\frac{\mathscr Y_i+\mathcal V_c}{2},
\quad i\in C_c,
\qquad
\overline{\mathscr B}_c
=\frac1{|C_c|}\sum_{i\in C_c}\mathscr B_{c,i}
=\mathcal V_c.
}
\end{equation}
در اینجا $\mathscr B_{c,i}$ یک توزیع احتمال
مصنوعی برای \emph{تحلیل همین معیار} است؛
منظور افزودن نمونه به داده‌ها یا اجرای
گشت جدید نیست. هزینه تجمیع KL این
مجموعه جدید را دقیقاً مشابه $R_c$ تعریف می‌کنیم:
\begin{equation}\label{cmp:eq:ch3-R-midpoint-KL}
\boxed{
\begin{aligned}
R_{B,c}
&:=\frac1{|C_c|}\sum_{i\in C_c}
D_{\mathrm{KL}}(\mathscr B_{c,i}\|\mathcal V_c)\\
&=H(\mathcal V_c)
-\frac1{|C_c|}\sum_{i\in C_c}H(\mathscr B_{c,i}).
\end{aligned}}
\end{equation}
اکنون از هویت آنتروپی جنسن--شنون
در رابطه~\eqref{cmp:eq:ch3-JSD-entropy}
و میانگین‌بودن $\mathcal V_c$ نتیجه می‌شود:
\begin{equation}\label{cmp:eq:ch3-R-JS-KL-midpoint-identity}
\boxed{
R_{c,\mathrm{JS}}
=\frac12R_c-R_{B,c},
\qquad
R_{\mathrm{JS}}
=\frac12R-\sum_{c=1}^{K}\omega_cR_{B,c}.
}
\end{equation}
ضریب $\frac12$ ضروری است:
در تعریف JSD، آنتروپی هر یک از
دو توزیع با وزن نصف وارد می‌شود.
بنابراین برای \emph{تعریف نمونه--مرکز}
رابطه $R_{c,\mathrm{JS}}=R_c-R_{B,c}$
برقرار نیست؛ رابطه بدون ضریب نصف
به \emph{تعریف زوجی نمونه--نمونه} در ادامه مربوط است.
از نامنفی‌بودن $R_{B,c}$
و $R_{c,\mathrm{JS}}$، به‌ویژه داریم:
\[
0\le R_{c,\mathrm{JS}}\le\tfrac12R_c,
\qquad
0\le R_{\mathrm{JS}}\le\tfrac12R.
\]
پس JS بخشی از هزینه پراکندگی KL را می‌سنجد
که پس از تشکیل میانگین جفتی با مرکز،
از اختلاف آنتروپی‌ها باقی می‌ماند.
این رابطه نشان‌دهنده ارتباط ریاضی
دو معیار است، اما یکسانی مقدار
یا جهت گرادیان‌هایشان را تضمین نمی‌کند.
محاسبات سطر‌به‌سطر و قرارداد مؤلفه‌های صفر
در پیوست~\ref{app:ch3-compactness-entropy-relations}
اثبات می‌شوند.

\paragraph{تعبیر زوجی تراکم جنسن--شنون.}
تا اینجا در تعریف $R_{c,\mathrm{JS}}$
\emph{هر نمونه با مرکز همان طبقه} مقایسه شد.
تعریف دیگری نیز وجود دارد که
\emph{هر دو نمونه از یک طبقه} را با یکدیگر
مقایسه می‌کند. برای شفاف‌بودن تفاوت،
این معیار را با بالانویس $\mathrm{pair}$
مشخص می‌کنیم و تمام زوج‌های مرتب
$i,j\in C_c$، شامل $i=j$، را با
وزن برابر در نظر می‌گیریم:
\begin{equation}\label{cmp:eq:ch3-R-JS-pairwise-distributions}
\boxed{
\mathscr B_{c,ij}^{\mathrm{pair}}
=\frac{\mathscr Y_i+\mathscr Y_j}{2},
\qquad
\frac1{|C_c|^2}\sum_{i,j\in C_c}
\mathscr B_{c,ij}^{\mathrm{pair}}
=\mathcal V_c.
}
\end{equation}
هزینه تجمیع KL مجموعه جدید
و میانگین واگرایی جنسن--شنون تمام زوج‌ها را
به‌ترتیب چنین تعریف می‌کنیم:
\begin{equation}\label{cmp:eq:ch3-R-JS-pairwise-definitions}
\boxed{
\begin{aligned}
R_{B,c}^{\mathrm{pair}}
&=\frac1{|C_c|^2}\sum_{i,j\in C_c}
D_{\mathrm{KL}}\!
\left(\mathscr B_{c,ij}^{\mathrm{pair}}\middle\|
\mathcal V_c\right),\\
R_{c,\mathrm{JS}}^{\mathrm{pair}}
&=\frac1{|C_c|^2}\sum_{i,j\in C_c}
\operatorname{JSD}(\mathscr Y_i,\mathscr Y_j).
\end{aligned}}
\end{equation}
چون مرکز مجموعه زوجی نیز
همان $\mathcal V_c$ است،
با استفاده از هویت آنتروپی مخلوط
و تعریف JSD، رابطه \emph{دقیق}
زیر حاصل می‌شود:
\begin{equation}\label{cmp:eq:ch3-R-JS-pairwise-entropy-identity}
\boxed{
R_{c,\mathrm{JS}}^{\mathrm{pair}}
=R_c-R_{B,c}^{\mathrm{pair}}.
}
\end{equation}
در این رابطه، برخلاف تعریف نمونه--مرکز،
\emph{ضریب $\frac12$ در برابر $R_c$
وجود ندارد}، زیرا در میانگین روی
همه زوج‌های مرتب، هر نمونه یک‌بار
در نقش مؤلفه نخست و یک‌بار در نقش
مؤلفه دوم ظاهر می‌شود و مجموع دو
وزن $\frac12$ به یک می‌رسد.
برابری بالا نشان می‌دهد پراکندگی
زوجی JS همان بخشی از هزینه KL
اولیه است که پس از تجمیع میانگین‌های
جفتی باقی نمی‌ماند.
در نتیجه
$0\le R_{c,\mathrm{JS}}^{\mathrm{pair}}\le R_c$.
این معیار را نباید با
$R_{c,\mathrm{JS}}$ که فاصله از مرکز
را می‌سنجد، یکسان دانست؛
هر دو مفهوم معتبرند، اما
ممکن است مقادیر و گرادیان‌های
متفاوتی تولید کنند.
اثبات کامل و توضیح درباره حذف
زوج‌های خودی $i=j$ در
پیوست~\ref{app:ch3-compactness-entropy-relations}
آمده است.

\paragraph{تعبیر و مقایسه معیارهای تراکم.}
$R_{\mathrm{KL}}$ هزینه اضافی مدل مشترک
برای توزیع‌های یک طبقه را می‌سنجد؛
هرچند KL به‌طور کلی جهت‌دار و بالقوه نامتناهی است،
$R_c$ در این مسئله به‌دلیل انتخاب مرکز حسابی
همواره متناهی و طبق
رابطه~\eqref{cmp:eq:ch3-R-KL-bounds}
کراندار است. تعبیر آن از هویت
آنتروپی مخلوط حاصل می‌شود؛
رابطه~\eqref{cmp:eq:ch3-R-JS-KL-midpoint-identity}
نیز نشان می‌دهد معیار JS تعبیر
آنتروپی مرتبط ولی متفاوتی دارد.
در مقابل، $R_{\mathrm{JS}}$ میانگین
\emph{مربع فاصله متریک} و
$R_{\sqrt{\mathrm{JS}}}$ میانگین
\emph{خود فاصله متریک} نمونه‌ها
از مراکز طبقات است. بنابراین اولی
انحراف‌های دورتر را به‌صورت مربعی
و دومی فاصله‌ها را به‌صورت خطی
در میانگین وارد می‌کند.
با درنظرگرفتن توزیع وزنی نمونه‌ها
که وزن هر عضو طبقه $c$ در آن
$\omega_c/|C_c|$ است،
از ینسن و کرانداری
$d_{\mathrm{JS}}\le\sqrt{\log2}$
به‌دست می‌آید:
\[
\bigl(R_{\sqrt{\mathrm{JS}}}\bigr)^2
\le R_{\mathrm{JS}}
\le \sqrt{\log2}\,R_{\sqrt{\mathrm{JS}}}.
\]
نامساوی نخست از
$(\mathbb E[d])^2\le\mathbb E[d^2]$
و دومی از $d^2\le\sqrt{\log2}\,d$
نتیجه می‌شود.
این پیوند به معنای
یکسان‌بودن معیارها یا اثرشان
در بهینه‌سازی نیست.
همچنین میانگین حسابی طبقه،
اگرچه برای $R_{\mathrm{KL}}$ بهینه است،
لزوماً کمینه‌کننده
$R_{\mathrm{JS}}$ یا
$R_{\sqrt{\mathrm{JS}}}$ نیست.
پس مقدارهای این سه معیار
هم‌واحد نیستند و مقایسه عددی‌شان
بدون توجه به تعبیر و مقیاس
موجه نخواهد بود.


\subsection{خواص نظری و روابط معیارهای جدایی}
\paragraph{مقایسه چهار معیار.}
اگر هدف معیارهای خطی را بیشینه و هدف معیارهای لگاریتمی را کمینه در نظر بگیریم، روابط زیر همواره برقرارند:
\begin{equation}\label{cmp:eq:ch3-four-criteria-comparison}
\boxed{
\begin{gathered}
T(\mathcal V)\le BT(\mathcal V),\qquad
BTNL(\mathcal V)\le TNL(\mathcal V),\\
TNL(\mathcal V)\ge-\log T(\mathcal V),\qquad
BTNL(\mathcal V)\ge-\log BT(\mathcal V).
\end{gathered}}
\end{equation}
دو نامساوی نخست حاصل مجاز بودن جایگشت همانی در مسئله انتساب و دو نامساوی بعدی حاصل تقعر لگاریتم هستند. شرط مثبت بودن وزن طبقات تضمین می‌کند که $T=1$ و $TNL=0$ فقط در ماتریس همانی، و $BT=1$ و $BTNL=0$ فقط در ماتریس‌های جایگشتی رخ دهند. بااین‌حال، یکسان بودن جواب‌های ایده‌آل به معنای یکسان بودن رتبه‌بندی حالت‌های میانی نیست: معیار لگاریتمی در برابر یک احتمال کوچک حساس‌تر است و بهترین جایگشت در $BT$ لزوماً با بهترین جایگشت در $BTNL$ یکسان نیست. اثبات دقیق نامساوی‌ها و شرایط تساوی در بخش~\ref{app:ch3-four-criteria-comparison} از پیوست~\ref{app:ch3-proofs} آمده است.

\paragraph{تجزیه بر حسب آنتروپی مراکز.}
اگر $H(U)=-\sum_{d=1}^{K}U(d)\log U(d)$ آنتروپی توزیع احتمال $U$ باشد، هویت آنتروپی مخلوط که پیش‌تر برای تعبیر $R$ به کار رفت، اکنون با انتخاب توزیع‌های $\mathcal V_c$ و وزن‌های $\omega_c$ نتیجه می‌دهد:
\begin{equation}\label{cmp:eq:ch3-I-entropy}
\boxed{I(\mathcal V)=H(\overline{\mathcal V})
-\sum_{c=1}^{K}\omega_c H(\mathcal V_c).}
\end{equation}
اثبات این هویت مشترک در بخش~\ref{app:ch3-entropy-mixture} از پیوست~\ref{app:ch3-proofs} آمده است.
در این رابطه، افزایش آنتروپی میانگین مراکز و کاهش میانگین آنتروپی خود مراکز هر دو می‌توانند مقدار $I$ را افزایش دهند. بااین‌حال، این دو جمله مستقل از یکدیگر نیستند؛ زیرا میانگین مراکز از خود آن‌ها ساخته شده است.

\paragraph{حد پایین و بالای $I$.}
از تعریف اطلاعات متقابل، می‌توان عدم‌قطعیت باقی‌مانده درباره طبقه آغازین پس از مشاهده برچسب همسایگی را به‌صورت زیر بیان کرد:
\begin{equation}\label{cmp:eq:ch3-I-conditional}
\boxed{I(\mathcal V)=H(C_0)-H(C_0\mid C_\xi)
=H(\Omega)-H(C_0\mid C_\xi).}
\end{equation}
برای روشن شدن منشأ این هویت، توجه می‌کنیم که طبق رابطه~\eqref{cmp:eq:ch3-I-joint-distribution}، در هر برچسب مشاهده‌شده $d$ با احتمال مثبت، احتمال شرطی طبقه آغازین برابر
$\mathrm P(C_0=c\mid C_\xi=d)=\omega_c\mathcal V_{cd}/\overline{\mathcal V}(d)$ است. با جایگذاری این احتمال در تعریف آنتروپی شرطی و باز کردن لگاریتم کسر، یک جمله برابر $H(\Omega)$ و جمله دیگر برابر منفی اطلاعات متقابل در رابطه~\eqref{cmp:eq:ch3-I-mutual-info} می‌شود. محاسبه جبری این تساوی در بخش~\ref{app:ch3-I-bounds-proof} از پیوست~\ref{app:ch3-proofs} نیز آمده است.

از نامنفی بودن اطلاعات متقابل و آنتروپی شرطی در رابطه~\eqref{cmp:eq:ch3-I-conditional}، و نیز از بیشینه بودن آنتروپی توزیع یکنواخت بر روی $K$ طبقه، نتیجه می‌شود:
\begin{equation}\label{cmp:eq:ch3-I-bounds}
\boxed{0\le I(\mathcal V)\le H(\Omega)\le\log K.}
\end{equation}
حد پایین زمانی به‌دست می‌آید که همه مراکز $\mathcal V_c$ برای طبقات دارای وزن مثبت یکسان باشند؛ زیرا در این صورت برچسب مشاهده‌شده هیچ اطلاعاتی درباره طبقه آغازین فراهم نمی‌کند.

اگر همه وزن‌ها مثبت باشند، رسیدن به کران بالای $H(\Omega)$ به این معناست که با دانستن برچسب مشاهده‌شده، طبقه آغازین بدون ابهام مشخص شود. چون تعداد برچسب‌های آغازین و مقصد هر دو $K$ است، این وضعیت دقیقاً برای ماتریس‌های جایگشتی رخ می‌دهد. بنابراین در فضای تمام ماتریس‌های احتمالی سطری، $I$ و $BT$ نقاط بهینه جهانی یکسانی دارند؛ هرچند شکل تابع هدف و رفتار آن‌ها در وضعیت‌های غیرایده‌آل متفاوت است. اثبات دقیق کران‌ها و شرایط تساوی در بخش~\ref{app:ch3-I-bounds-proof} از پیوست~\ref{app:ch3-proofs} آمده است. البته ممکن است همه این ماتریس‌های بهینه با سیاست انتخاب ویژگیِ قابل یادگیری روی گراف موجود دست‌یافتنی نباشند.

\paragraph{تمایز از معیار پراکندگی $R$.}
قضیه تجزیه KL در بخش~\ref{cmp:subsec:ch3-mean-center-justification} نشان می‌دهد که میانگین وزن‌دار واگرایی توزیع نمونه‌ها از مرکز مشترک $\overline{\mathcal V}$ به دو جزء تجزیه می‌شود:
\begin{equation}\label{cmp:eq:ch3-R-I-decomposition}
\boxed{
\sum_{c=1}^{K}\frac{\omega_c}{|C_c|}
\sum_{i\in C_c}
D_{\mathrm{KL}}\bigl(\mathscr Y_i\|\overline{\mathcal V}\bigr)
=R+I(\mathcal V).
}
\end{equation}
این رابطه تمایز میان دو مفهوم را روشن می‌کند: $R$ میزان پراکندگی توزیع‌های نمونه‌ها حول \emph{مرکز طبقه خودشان} و $I$ میزان جدایی \emph{مراکز طبقات} از میانگین مشترک را اندازه می‌گیرد. کمینه شدن $R$ به‌تنهایی مانع یکسان شدن همه مراکز نمی‌شود؛ در مقابل، مثبت بودن $I$ مستلزم وجود تفاوت میان دست‌کم دو مرکز دارای وزن مثبت است.

\paragraph{ارتباط معیار اطلاعات متقابل با $T$ و $BT$.}
معیار اطلاعات متقابل $I$، تفاوت میان توزیع مشترک طبقه آغازین و برچسب مشاهده‌شده، یعنی $\mathrm P(C_0,C_\xi)$، و توزیع حاصل از فرض استقلال آن‌ها را اندازه می‌گیرد. برای مقایسه این معیار با $T$، لازم نیست تمام زوج‌های ممکن برچسب‌ها را جداگانه بررسی کنیم. کافی است تنها رویداد \emph{برابر بودن دو برچسب} را در نظر بگیریم و سایر حالت‌ها را در یک گروه قرار دهیم. نامساوی پردازش داده‌ها می‌گوید که چنین خلاصه‌سازی‌ای نمی‌تواند واگرایی KL دو توزیع را افزایش دهد؛ بنابراین واگرایی حاصل ازاین‌رویداد ساده، کران پایینی برای $I$ خواهد بود.

متغیر دودویی $B_{\mathrm{eq}}$ را برابر یک می‌گیریم اگر $C_0=C_\xi$ باشد و در غیر این صورت برابر صفر. طبق تعریف~\eqref{cmp:eq:ch3-T-on-V}، احتمال $B_{\mathrm{eq}}=1$ در توزیع مشترک واقعی برابر $T(\mathcal V)$ است. در توزیع فرض استقلال، همین احتمال برابر $q=\sum_{c=1}^{K}\omega_c\overline{\mathcal V}(c)$ خواهد بود. با تبدیل هر دو توزیع به توزیع‌های برنولی متناظر، نتیجه می‌شود:
\begin{equation}\label{cmp:eq:ch3-I-vs-T}
\boxed{
I(\mathcal V)\ge
D_{\mathrm{KL}}\!\left(
\operatorname{Bern}\bigl(T(\mathcal V)\bigr)
\middle\|
\operatorname{Bern}(q)
\right),\qquad
q=\sum_{c=1}^{K}\omega_c\overline{\mathcal V}(c).
}
\end{equation}
در اینجا $\operatorname{Bern}(p)$ توزیع برنولی با احتمال موفقیت $p$ است و $q$، احتمال هم‌برچسب بودن \emph{در فرض استقلال} را نشان می‌دهد؛ نه احتمال واقعی آن را. به‌ویژه، اگر وزن طبقات یکنواخت باشد، از مجموع یک بودن مؤلفه‌های $\overline{\mathcal V}$ داریم $q=1/K$.

معیار $T$ تنها برابر بودن برچسب‌ها را بررسی می‌کند؛ اما در معیار جایگشتی $BT$، برچسب مشاهده‌شده لازم نیست با برچسب طبقه آغازین یکسان باشد. برای هر جایگشت $\sigma$، به‌جای رویداد قبلی، رویداد $B_\sigma=1$ را معادل $C_\xi=\sigma(C_0)$ تعریف می‌کنیم. احتمال این رویداد در توزیع واقعی و در فرض استقلال به‌ترتیب برابر است با:
\begin{equation}\label{cmp:eq:ch3-I-BT-event-probs}
t_\sigma=\sum_{c=1}^{K}\omega_c\mathcal V_{c,\sigma(c)},
\qquad
q_\sigma=\sum_{c=1}^{K}\omega_c\overline{\mathcal V}\bigl(\sigma(c)\bigr).
\end{equation}
همان استدلال پردازش داده‌ها برای هر جایگشت برقرار است و $I(\mathcal V)\ge D_{\mathrm{KL}}(\operatorname{Bern}(t_\sigma)\|\operatorname{Bern}(q_\sigma))$ را نتیجه می‌دهد. در حالت وزن‌های یکنواخت $\omega_c=1/K$، چون $\sigma$ یک جایگشت است، جمع احتمال‌های مقصد در $q_\sigma$ برابر یک می‌شود و \emph{برای هر جایگشت} داریم $q_\sigma=1/K$. اگر $\sigma^\star$ جایگشتی باشد که امتیاز $BT$ را بیشینه می‌کند، آنگاه $t_{\sigma^\star}=BT(\mathcal V)$ و در نتیجه:
\begin{equation}\label{cmp:eq:ch3-I-vs-BT}
\boxed{
I(\mathcal V)\ge
D_{\mathrm{KL}}\!\left(
\operatorname{Bern}\bigl(BT(\mathcal V)\bigr)
\middle\|
\operatorname{Bern}(1/K)
\right),\qquad \omega_c=\frac{1}{K}.
}
\end{equation}
این کران ارتباط مشخصی میان دو معیار فراهم می‌کند. در بخش~\ref{cmp:subsubsec:ch3-BT} ثابت کردیم $1/K\le BT(\mathcal V)\le1$. از سوی دیگر، تابع $D_{\mathrm{KL}}(\operatorname{Bern}(x)\|\operatorname{Bern}(1/K))$ روی بازه $[1/K,1]$ اکیداً صعودی است. بنابراین، \textbf{با افزایش $BT$، کران پایین تضمین‌شده برای $I$ نیز افزایش می‌یابد}؛ هرچند از این نتیجه نمی‌توان نتیجه گرفت که مقدار واقعی $I$ در هر مرحله بهینه‌سازی الزاماً افزایش پیدا می‌کند.

در نقطه ایده‌آل $BT(\mathcal V)=1$، سمت راست رابطه~\eqref{cmp:eq:ch3-I-vs-BT} به $\log K$ می‌رسد. چون برای وزن‌های یکنواخت $I(\mathcal V)\le\log K$ است، در این حالت الزاماً $I(\mathcal V)=\log K$ خواهد بود. بدین‌ترتیب، هر دو معیار بر روی مجموعه ماتریس‌های احتمالی سطری، در ماتریس‌های جایگشتی به پاسخ ایده‌آل یکسانی می‌رسند؛ ولی شکل تابع هدف و مسیر بهینه‌سازی آن‌ها می‌تواند متفاوت باشد. اثبات تفصیلی کران‌ها و صعودی بودن آن‌ها در بخش~\ref{app:ch3-I-bounds-proof} از پیوست~\ref{app:ch3-proofs} آمده است. برای وزن‌های غیریکنواخت نیز کران برنولی با پارامتر $q_\sigma$ برقرار می‌ماند، اما این پارامتر دیگر لزوماً برابر $1/K$ نیست.

\paragraph{ارتباط اطلاعات متقابل با خطای تشخیص.}
یکی از مزایای تعبیر اطلاعاتی معیار $I$ آن است که می‌توان از نتایج شناخته‌شده نظریه اطلاعات، از جمله نامساوی فانو، استفاده کرد. اگر هدف، تشخیص طبقه آغازین $C_0$ تنها با مشاهده برچسب همسایگی $C_\xi$ باشد و $P_e=\mathrm P(\widehat C_0\ne C_0)$ احتمال خطای یک تصمیم‌گیرنده مبتنی بر $C_\xi$ باشد، از هویت~\eqref{cmp:eq:ch3-I-conditional} و نامساوی فانو داریم \cite{CoverThomas2006}:
\begin{equation}\label{cmp:eq:ch3-I-Fano}
\boxed{H(\Omega)-I(\mathcal V)
=H(C_0\mid C_\xi)
\le h_2(P_e)+P_e\log(K-1),}
\end{equation}
که در آن $K\ge2$ و $h_2(p)=-p\log p-(1-p)\log(1-p)$ آنتروپی دودویی است. این نامساوی محدودیتی اطلاعاتی بر خطای تشخیص قرار می‌دهد؛ به‌ویژه، اگر اطلاعات متقابل کوچک باشد، خطای تشخیص نمی‌تواند به‌دلخواه کوچک شود. بااین‌حال، \textbf{فانو کران پایین خطا را فراهم می‌کند، نه کران بالای خطای یک طبقه‌بند معین}. افزون بر این، مشاهده تصادفی $C_\xi$ با در اختیار داشتن کل بردار $\mathscr Y_i$ یکسان نیست. ازاین‌رو این نامساوی را نباید بدون فرض‌های اضافی با کران بالای درست‌نمایی لگاریتمی منفی برای مدل پیش‌بینی مبتنی بر $\mathscr Y_i$ یکی دانست.

در نتیجه، $I$ علاوه بر بیان تمایز میان مراکز، با میزان اطلاعات
قابل استفاده برای تشخیص طبقه نیز ارتباط نظری دارد؛
اما این رابطه به‌تنهایی کیفیت تصمیم‌گیری از روی
تمام بردار بازنمایی و حاشیه تک‌تک نمونه‌ها را تعیین نمی‌کند.

در مجموع، معیار اطلاعات متقابل $I$ بدون تحمیل هم‌برچسب بودن یا انتخاب یک جایگشت خاص، تمایز توزیع‌های نماینده طبقات را اندازه می‌گیرد و مانند چهار معیار پیشین، صرفاً بر اساس $\mathcal V$ قابل محاسبه است. تعبیر اطلاعات متقابل و ارتباط آن با عدم‌قطعیت طبقه آغازین، انگیزه نظری مناسبی برای استفاده از $I$ در ارزیابی و بهبود سیاست انتخاب ویژگی فراهم می‌کند.



\section{حاشیهٔ تصمیم و کیفیت طبقه‌بندی}
\label{cmp:sec:ch3-normalized-margin}
\phantomsection\label{cmp:subsec:ch3-label-estimation-to-margin}
در چارچوب نزدیک‌ترین نماینده، برای سنجش استحکام
و درستی تصمیم باید اختلاف امتیاز طبقهٔ واقعی با
قوی‌ترین رقیب بررسی شود. این اختلاف،
\emph{حاشیهٔ تصمیم} نام دارد.
در این بخش ابتدا حاشیه و خطای متوازن تجربی،
سپس هندسهٔ حاشیه و ارتباط آن با پراکندگی و جدایی،
کران‌های آماری و پایداری تصمیم را تحلیل می‌کنیم.
ساخت تابع زیان از حاشیه و مقایسهٔ آن با NLL
در بخش توابع هدف انجام می‌شود.

\subsection{حاشیهٔ تصمیم و خطای متوازن تجربی}
\label{cmp:subsec:ch3-decision-margin-empirical}
برای یک نمونه با برچسب واقعی $y_i$،
امتیاز نامزد $c$ را با $s_{i,c}$
نشان می‌دهیم؛ این امتیاز می‌تواند
شباهت $s(\mathscr Y_i,\mathcal V_c)$
باشد یا در استنتاج شرطی، از
بازنمایی مخصوص نامزد $c$ محاسبه شود.
حاشیه به اختلاف امتیاز طبقه واقعی
با قوی‌ترین رقیب اطلاق می‌شود؛
در صورت یکتابودن برنده، علامت
این اختلاف دقیقاً وضعیت درست یا
نادرست بودن تصمیم را مشخص می‌کند.
ابتدا \emph{حاشیهٔ عمومی} را مستقل از انتخاب
سنجهٔ فاصله یا شباهت تعریف می‌کنیم:
\begin{equation}\label{cmp:eq:ch3-general-decision-margin}
\boxed{
m_i:=s_{i,y_i}-\max_{r\ne y_i}s_{i,r}
=\min_{r\ne y_i}(s_{i,y_i}-s_{i,r}).
}
\end{equation}
در حالت امتیازهای منفی فاصله، این حاشیه
برابر $\min_{r\ne y_i}(d_{i,r}-d_{i,y_i})$
است. مقدار مثبت حاکی از برتری اکید طبقهٔ درست،
مقدار منفی بیانگر وجود رقیب برتر، و مقدار صفر
نشان‌دهندهٔ تساوی در بالاترین امتیاز است.
در ادامه، برای تحلیل هندسی این فصل،
شباهت نرمال‌شدهٔ جنسن--شنون
تعریف‌شده در بخش پیشین را
مبنای امتیاز قرار می‌دهیم؛
$M_i^{\mathrm{JS}}$ حالت خاص همین حاشیه است.

\paragraph{دقت دقیق، خطای صفر--یک و علامت حاشیه تصمیم.}
معیار نهایی طبقه‌بندی، \emph{درست یا نادرست بودن پیش‌بینی هر نمونه} است،
نه صرفاً کوچکی فاصله از مرکز طبقه درست یا دوربودن دو مرکز.
اگر $\widehat y_i$ خروجی قاعده تصمیم نزدیک‌ترین مرکز
با یک قرارداد مشخص برای رفع تساوی باشد،
\emph{دقت متوازن تجربی دقیق} و خطای متناظر آن برابرند با
\begin{equation}\label{cmp:eq:ch3-js-exact-balanced-accuracy}
\boxed{
\begin{aligned}
\widehat{\mathrm{BA}}
&=\frac1K\sum_{c=1}^{K}\frac1{|C_c|}
\sum_{i\in C_c}\mathbf1\{\widehat y_i=c\},\\
\widehat E_{\mathrm{bal}}
&=1-\widehat{\mathrm{BA}}
=\frac1K\sum_{c=1}^{K}\frac1{|C_c|}
\sum_{i\in C_c}\mathbf1\{\widehat y_i\ne c\}.
\end{aligned}}
\end{equation}
این تعریف سهم هر طبقه را مستقل از تعداد نمونه‌هایش یکسان می‌گیرد
و \emph{معیار دقیق ارزیابی} ماست؛ برخلاف تابع زیان هموار،
هنگام تغییر پارامترهای سیاست فقط با جابه‌جایی مرز تصمیم
و تغییر برچسب پیش‌بینی‌شده عوض می‌شود.

اکنون حاشیهٔ عمومی
رابطه~\eqref{cmp:eq:ch3-general-decision-margin}
را برای امتیازهای مبتنی بر شباهت جنسن--شنون
به‌شکل زیر می‌نویسیم:
\begin{equation}\label{cmp:eq:ch3-js-margin-distance-difference}
\boxed{
\begin{aligned}
M_i^{\mathrm{JS}}
&=s(\mathscr Y_i,\mathcal V_{y_i})
-\max_{r\ne y_i}s(\mathscr Y_i,\mathcal V_r)\\
&=\frac{\displaystyle
\min_{r\ne y_i}\operatorname{JSD}(\mathscr Y_i,\mathcal V_r)
-\operatorname{JSD}(\mathscr Y_i,\mathcal V_{y_i})}
{\log 2}.
\end{aligned}}
\end{equation}
این حاشیه، \emph{اختلاف واقعی امتیاز تصمیم}
بین طبقه درست و قوی‌ترین رقیب را می‌سنجد و
در بازه $[-1,1]$ قرار دارد. اگر مثبت باشد، تصمیم
حتماً درست و اگر منفی باشد، حتماً نادرست است؛
در مقدار صفر، نتیجه به قرارداد رفع تساوی بستگی دارد.
ازاین‌رو برای هر نمونه
$\mathbf1\{M_i^{\mathrm{JS}}<0\}
\le\mathbf1\{\widehat y_i\ne y_i\}
\le\mathbf1\{M_i^{\mathrm{JS}}\le0\}$.
با قرارداد محافظه‌کارانه شمردن تساوی به‌عنوان خطا،
صورت زیر را داریم:
\begin{equation}\label{cmp:eq:ch3-js-hard-margin-error}
\widehat E_{\mathrm{bal}}^{\,\le0}
=\frac1K\sum_{c=1}^{K}\frac1{|C_c|}
\sum_{i\in C_c}
\mathbf1\{M_i^{\mathrm{JS}}\le0\}.
\end{equation}
در نبود تساوی، این مقدار دقیقاً برابر
$\widehat E_{\mathrm{bal}}$ است و در حضور تساوی
کران بالای آن خواهد بود.
پس حاشیه، اطلاعاتی درباره \emph{میزان برتری یا شکست}
تصمیم دارد که در خطای صفر--یک حذف می‌شود.
بااین‌حال، تابع نشانگر خطا نسبت به حاشیه یک تابع
پله‌ای است: بیرون مرز تقریباً همه‌جا مشتق صفر
دارد و روی مرز ناپیوسته است؛ بنابراین
برای هدایت تدریجی بهینه‌سازی گرادیانی مناسب نیست.
\emph{خود حاشیه} را نباید به‌طور کلی نامشتق‌پذیر
خواند: برای امتیازهای هموار، حاشیه در نواحی
دارای رقیب برنده یکتا مشتق‌پذیر است،
ولی عملگر $\max$ در محل تساوی رقبا ممکن است
نامشتق‌پذیر باشد.



\subsection{حاشیه هندسی جنسن--شنون و نقش تراکم و جدایی مراکز}
\label{cmp:subsec:ch3-js-metric-classification}
در بخش~\ref{cmp:subsec:ch3-probability-dissimilarity-families}، فاصله متریک جنسن--شنون و شباهت نرمال‌شده متناظر با آن را تعریف کردیم؛ سپس در بخش~\ref{cmp:subsec:ch3-compactness-separation}، نقش هم‌زمان فشردگی نمونه‌ها و جدایی مراکز را توضیح دادیم. اکنون این دو مفهوم را مستقیماً در قاعده تصمیم نزدیک‌ترین مرکز به کار می‌بریم. این قاعده به نرمال‌سازی با بیشینهٔ هموار یا تفسیر احتمال پسین نیازی ندارد و مستقیماً اختلاف امتیاز طبقه درست و نزدیک‌ترین رقیب را معیار قرار می‌دهد.

\paragraph{حفظ ترتیب احتمال‌ها و قاعده تصمیم.}
برای بردار یک‌داغ $e_c$ و توزیع $P$، محاسبه شباهت تصاویر به رابطه زیر منتهی می‌شود:
\begin{equation}\label{cmp:eq:ch3-js-onehot-order}
\begin{aligned}
s(e_c,P)
&=\frac{(1+p_c)\log(1+p_c)-p_c\log p_c}{2\log2},\\[4pt]
\frac{\partial s(e_c,P)}{\partial p_c}
&=\frac{\log[(1+p_c)/p_c]}{2\log2}>0,
\qquad 0<p_c<1.
\end{aligned}
\end{equation}
پس در مقایسه با بردارهای یک‌داغ، هر دو معیار فاصله و شباهت ترتیب احتمال‌های طبقه‌ها را حفظ می‌کنند. همچنین از رابطه $s(P,Q)=1-d_{\mathrm{JS}}(P,Q)^2/\log2$ نتیجه می‌شود قاعده‌های تصمیم زیر دقیقاً یکسان‌اند:
\begin{equation}\label{cmp:eq:ch3-js-decision}
\boxed{
\widehat y_i
=\arg\min_c d_{\mathrm{JS}}(\mathscr Y_i,\mathcal V_c)
=\arg\max_c s(\mathscr Y_i,\mathcal V_c).
}
\end{equation}
این قاعده به نرمال‌سازی با بیشینهٔ هموار یا نرمال‌سازی مجدد امتیازها نیاز ندارد. البته حفظ ترتیب در مقایسه با $e_c$ به معنای برابری شباهت با احتمال پسین واقعی نیست؛ در مدل، بازنمایی همسایگی $\mathscr Y_i$ با مراکز طبقه‌ها مقایسه می‌شود. جزئیات اثبات خواص نگاشت، محاسبه فاصله و شباهت و حفظ ترتیب در بخش~\ref{app:ch3-js-error-proof} از پیوست~\ref{app:ch3-proofs} ارائه شده‌اند.

\paragraph{حاشیه و کران هندسی خطای طبقه‌بندی.}
برای نمونه‌ای از طبقه واقعی $c$، حاشیه تصمیم $M_i^{\mathrm{JS}}$ را مطابق رابطه~\eqref{cmp:eq:ch3-js-margin-distance-difference} در نظر می‌گیریم. مثبت بودن آن به معنای پیشی‌گرفتن طبقه واقعی از تمام رقباست. برخلاف تفاوت‌های نامحدود KL، این حاشیه همواره در بازه $[-1,1]$ قرار دارد. فاصله نزدیک‌ترین مرکز رقیب $\delta_c$ و پراکندگی $W_c^{\mathrm{JS}}$ پیش‌تر در روابط~\eqref{cmp:eq:ch3-js-center-separation} و \eqref{cmp:eq:ch3-js-within-dispersion} معرفی شدند. با نامساوی مثلثی، اگر $d_{\mathrm{JS}}(\mathscr Y_i,\mathcal V_c)<\delta_c/2$ باشد، فاصله نمونه از مرکز واقعی اکیداً کمتر از تمام مراکز رقیب است و تصمیم حتماً درست خواهد بود. بنابراین برای $\delta_c>0$، با نمونه‌گیری یکنواخت از طبقه $c$ و کاربرد نامساوی مارکوف:
\begin{equation}\label{cmp:eq:ch3-js-class-error-bound}
\boxed{
\begin{aligned}
\widehat\varepsilon_c
&\le\min\left\{1,\,
\frac{4W_c^{\mathrm{JS}}}{\delta_c^2}\right\}\\
&=\min\left\{1,\,
\frac{4\left(1-\overline S_c\right)}
{1-\max_{r\ne c}s(\mathcal V_c,\mathcal V_r)}\right\},
\end{aligned}}
\end{equation}
که در آن $\overline S_c=|C_c|^{-1}\sum_{i\in C_c}s(\mathscr Y_i,\mathcal V_c)$ و $\widehat\varepsilon_c$ نسبت خطای طبقه‌بندی تجربی طبقه $c$ است. اگر $\delta_c=0$، از این استدلال کران مفیدی به‌دست نمی‌آید و تنها کران بدیهی $\widehat\varepsilon_c\le1$ را به کار می‌بریم. اثبات کران در بخش~\ref{app:ch3-js-error-proof} از پیوست~\ref{app:ch3-proofs} آمده است.

این کران، برخلاف معیار مبتنی بر مرکز به‌تنهایی، هر دو عامل \emph{جدایی مراکز} و \emph{پراکندگی نمونه‌ها در طبقه خودشان} را هم‌زمان در نظر می‌گیرد و مستقیماً خطای تصمیم نزدیک‌ترین مرکز را هدف می‌گیرد. البته این کران اولیه ممکن است شل باشد و تعمیم آن به خطای نمونه‌های دیده‌نشده نیازمند تحلیل آماری مستقل است. بررسی کران‌های دقیق‌تر بر اساس حاشیه‌های نمونه‌ها و استخراج هدف مناسب یادگیری سیاست، ادامه این مسیر خواهد بود.

کران هندسی بالا به‌صورت صریح نشان می‌دهد
که فشردگی نمونه‌ها و دوری مراکز در کنار هم
برای تصمیم مهم‌اند، اما حتی کوچک‌شدن این کران
لزوماً بهبود قدم‌به‌قدم دقت را تضمین نمی‌کند.
همچنین خودِ خطای صفر--یک برای بهینه‌سازی
گرادیانی هموار نیست. بنابراین معیار بعدی را
مستقیماً از \emph{علامت و اندازه حاشیه هر نمونه}
می‌سازیم و یک تقریب هموار از خطای متوازن ارائه می‌کنیم.


\paragraph{تمایز کاهش شعاع و حرکت در جهت بهبود حاشیه.}
کاهش فاصله نمونه از مرکز طبقه خودش، به‌تنهایی مشخص نمی‌کند
فاصله آن از مراکز رقیب چگونه تغییر کرده است. این موضوع را
می‌توان در حالت یک بازنمایی مشترک برای هر نمونه، با
نگاشت هیلبرتی $\Phi$ متناظر با متریک
$d_{\mathrm{JS}}=\sqrt{\operatorname{JSD}}$ صورت‌بندی کرد.
اگر $z=\Phi(\mathscr Y_i)$، $v_c=\Phi(\mathcal V_c)$ و
$v_r=\Phi(\mathcal V_r)$ باشند، حاشیه نسبت به رقیب $r$
برابر است با
\begin{equation}\label{cmp:eq:ch3-fixed-center-pairwise-js-margin}
m_{i,c,r}
=\frac{\|z-v_r\|^2-\|z-v_c\|^2}{\ln2}
=\frac{\|v_c-v_r\|^2+2\langle z-v_c,v_c-v_r\rangle}{\ln2}.
\end{equation}
پس اگر مراکز ثابت بمانند و بازنمایی نمونه از $z$ به
$z+\Delta z$ تغییر کند، خواهیم داشت:
\begin{equation}\label{cmp:eq:ch3-fixed-center-pairwise-js-change}
\Delta m_{i,c,r}
=\frac{2\langle\Delta z,v_c-v_r\rangle}{\ln2}.
\end{equation}
بنابراین جهت جابه‌جایی نسبت به بردار اتصال مراکز برای
بهبود یا تضعیف حاشیه تعیین‌کننده است؛ کم‌شدن شعاع
$\|z-v_c\|$ بدون اطلاع از این جهت، افزایش حاشیه را
تضمین نمی‌کند.


\paragraph{تعبیر شماتیک در سیمپلکس سه برچسب.}
برای روشن‌شدن تفاوت تراکم و حاشیه،
سه مرکز طبقه را در دو آرایش
\emph{کاملاً یکسان} در نظر می‌گیریم.
در آرایش نخست، دو بازنمایی
متعلق به طبقه اول از مرکز
خود پراکندگی بیشتری دارند،
اما هر دو در ناحیه تصمیم
همان طبقه قرار می‌گیرند.
در آرایش دوم، بازنمایی‌ها
به مرکز طبقه خود نزدیک‌ترند
و تراکم ظاهراً بهتر است،
بااین‌حال یکی از نمونه‌ها
از مرز تصمیم عبور کرده و
به طبقه رقیب نسبت داده
می‌شود. این مثال
\emph{تعبیری هندسی و شماتیک}
است، نه نتیجه یک آزمایش
روی داده‌ها؛ پیام آن
این است که فاصله
شعاعی تا مرکز،
جهت مناسب جابه‌جایی
نسبت به مراکز رقیب
را مشخص نمی‌کند.
ازاین‌رو حتی با
مرکزهای ثابت و
معیار جدایی ثابت،
بهبود تراکم به‌تنهایی
بهبود همه حاشیه‌ها
و تصمیم‌ها را
تضمین نمی‌کند.

\begin{figure}[htbp]
\centering
\resizebox{0.98\linewidth}{!}{%
% نمایش برداری دو آرایش از توزیع‌های سه‌برچسبی
% مراکز در دو پنل ثابت‌اند؛ تنها موقعیت نمونه‌های طبقه اول تفاوت دارد.
% تصویر توضیحی است؛ اعداد و نسبت‌های شکل نتیجه آزمایش نیستند.
\begingroup
\definecolor{marginedge}{HTML}{61748C}
\definecolor{marginpurple}{HTML}{EEEAFB}
\definecolor{marginorange}{HTML}{FFF0E6}
\definecolor{margingreen}{HTML}{DFF4EE}
\definecolor{marginblue}{HTML}{26829A}
\definecolor{marginred}{HTML}{D16C69}
\definecolor{marginviolet}{HTML}{8168BC}
\begin{tikzpicture}[x=1.30cm,y=1.30cm,line cap=round,line join=round]
% پنل چپ: مراکز ثابت، دو نمونه با پراکندگی بیشتر و تصمیم صحیح.
\begin{scope}
  \begin{scope}
    \clip (0,0)--(4.8,0)--(2.4,4.16)--cycle;
    \fill[marginpurple] (0,0) rectangle (4.8,4.16);
    \fill[marginorange] (0,0)--(1.53,2.65)
        .. controls (1.50,1.70) and (1.68,0.60) .. (2.07,0)--cycle;
    \fill[margingreen] (2.07,0)
        .. controls (1.78,0.50) and (1.67,1.66) .. (1.76,2.17)
        .. controls (2.80,2.41) and (4.04,1.73) .. (4.80,0)--cycle;
    \draw[marginred!65,dashed,line width=.6pt] (2.07,0)
       .. controls (1.68,.73) and (1.48,1.65) .. (1.53,2.65);
  \end{scope}
  \draw[marginedge,line width=1.3pt] (0,0)--(4.8,0)--(2.4,4.16)--cycle;
  \node[below left, text=marginedge] at (0,0) {$e_1$};
  \node[below right, text=marginedge] at (4.8,0) {$e_2$};
  \node[above, text=marginedge] at (2.4,4.16) {$e_3$};
  % مراکز طبقات
  \draw[marginred,line width=1.55pt,fill=white] (1.38,.76) circle (.115);
  \node[below right, text=marginred] at (1.38,.76) {$\mathcal V_2$};
  \draw[marginviolet,line width=1.55pt,fill=white] (2.53,3.50) circle (.115);
  \node[below right, text=marginviolet] at (2.53,3.50) {$\mathcal V_3$};
  \draw[marginblue,line width=1.7pt,fill=white] (2.23,.96) circle (.13);
  \node[below left,text=marginblue] at (2.23,.96) {$\mathcal V_1$};
  % دو نمونه طبقه اول با مرکز میانگین V1
  \draw[marginblue,dashed,line width=1.15pt] (1.87,1.73)--(2.23,.96)--(2.59,.19);
  \filldraw[draw=white,fill=marginblue,line width=1pt] (1.87,1.73) circle (.105);
  \node[above left] at (1.87,1.73) {$P_1$};
  \filldraw[draw=white,fill=marginblue,line width=1pt] (2.59,.19) circle (.105);
  \node[below] at (2.59,.19) {$P_2$};
\end{scope}
% پنل راست: مراکز دقیقاً همان‌اند، فاصله نمونه‌ها از مرکز کمتر است،
% اما نمونه اول به ناحیه تصمیم طبقه دوم منتقل می‌شود.
\begin{scope}[xshift=5.75cm]
  \begin{scope}
    \clip (0,0)--(4.8,0)--(2.4,4.16)--cycle;
    \fill[marginpurple] (0,0) rectangle (4.8,4.16);
    \fill[marginorange] (0,0)--(1.53,2.65)
        .. controls (1.50,1.70) and (1.68,0.60) .. (2.07,0)--cycle;
    \fill[margingreen] (2.07,0)
        .. controls (1.78,0.50) and (1.67,1.66) .. (1.76,2.17)
        .. controls (2.80,2.41) and (4.04,1.73) .. (4.80,0)--cycle;
    \draw[marginred!65,dashed,line width=.6pt] (2.07,0)
       .. controls (1.68,.73) and (1.48,1.65) .. (1.53,2.65);
  \end{scope}
  \draw[marginedge,line width=1.3pt] (0,0)--(4.8,0)--(2.4,4.16)--cycle;
  \node[below left, text=marginedge] at (0,0) {$e_1$};
  \node[below right, text=marginedge] at (4.8,0) {$e_2$};
  \node[above, text=marginedge] at (2.4,4.16) {$e_3$};
  \draw[marginred,line width=1.55pt,fill=white] (1.38,.76) circle (.115);
  \node[below right, text=marginred] at (1.38,.76) {$\mathcal V_2$};
  \draw[marginviolet,line width=1.55pt,fill=white] (2.53,3.50) circle (.115);
  \node[below right, text=marginviolet] at (2.53,3.50) {$\mathcal V_3$};
  \draw[marginblue,line width=1.7pt,fill=white] (2.23,.96) circle (.13);
  \node[below left,text=marginblue] at (2.23,.96) {$\mathcal V_1$};
  \draw[marginblue,dashed,line width=1.15pt] (1.65,.82)--(2.23,.96)--(2.81,1.10);
  % نمونه سمت چپ متعلق به طبقه اول است، اما در ناحیه رقیب قرار دارد.
  \node[marginred,font=\bfseries\large] at (1.65,.82) {$\times$};
  \node[above left] at (1.65,.82) {$P_1$};
  \filldraw[draw=white,fill=marginblue,line width=1pt] (2.81,1.10) circle (.105);
  \node[below right] at (2.81,1.10) {$P_2$};
\end{scope}
\end{tikzpicture}
\endgroup
%
}
\caption{دو آرایش شماتیک از بازنمایی‌های متعلق به یک طبقه در سیمپلکس سه‌برچسبی با مراکز ثابت. در سمت چپ، پراکندگی نمونه‌ها بیشتر است، اما هر دو در ناحیه تصمیم درست قرار دارند؛ در سمت راست، نمونه‌ها به مرکز خود نزدیک‌ترند، ولی یکی از آن‌ها وارد ناحیه طبقه رقیب شده است. خطوط گسسته فاصله‌های نمونه تا مرکز و مرزهای رنگی ناحیه تصمیم نزدیک‌ترین مرکز را نشان می‌دهند.}
\label{cmp:fig:ch3-margin-vs-compactness}
\end{figure}

شکل~\ref{cmp:fig:ch3-margin-vs-compactness}
تفاوت میان \emph{فاصله تا مرکز واقعی} و
\emph{موقعیت نسبت به مرزهای تصمیم} را
به‌صورت بصری بیان می‌کند. این طرح
نتیجه یک اجرای تجربی نیست و صرفاً
برای توضیح هندسیِ کفایت‌نداشتن
تراکم و جدایی مراکز به‌تنهایی
ارائه شده است.

 روابط بالا برای \emph{بازنمایی مشترک}
نوشته شده‌اند. در استنتاجِ شرطی طبقه، که برای هر نامزد
$r$ بازنمایی جداگانه $\mathscr Y_i^{(r)}$ به‌دست می‌آید،
نمی‌توان بدون فرض اضافی همه امتیازها را به حرکت یک
بردار $z$ نسبت داد. تحلیل و سنجش تجربی این تمایز در
فصل~\ref{chap:results} انجام می‌شود.





\subsection{تحلیل آماری حاشیه KL: جدایی میانگین و پراکندگی تصمیم}
\label{cmp:subsec:ch3-margin-statistics}
معیار اطلاعات متقابل $I$ عمدتاً بر توزیع‌های مرکزی طبقه‌ها تکیه دارد، اما طبقه‌بندی یک نمونه زمانی درست است که فاصله آن از \emph{مرکز طبقه خودش} کمتر از فاصله آن از تمام مراکز رقیب باشد. به همین دلیل علاوه بر موقعیت مراکز، پراکندگی نمونه‌ها نسبت به مرزهای تصمیم نیز اهمیت دارد.

برای نمونه $X_i$ متعلق به طبقه $c$ و طبقه رقیب $r\ne c$، حاشیه جدایی را تعریف می‌کنیم:
\begin{equation}\label{cmp:eq:ch3-margin-definition}
\boxed{M_{c,r}(X_i)=d_{i,r}-d_{i,c}
=\sum_{j=1}^{K}\mathscr Y_i(j)
\log\frac{\mathcal V_c(j)}{\mathcal V_r(j)}.}
\end{equation}
حاشیه مثبت به معنای نزدیک‌تر بودن نمونه به مرکز طبقه واقعی نسبت به مرکز رقیب است و حاشیه نامثبت به معنی وجود رقابت یا تساوی در تصمیم است. برای محاسبه آمار این حاشیه، بردار ستونی و ماتریس کوواریانس تجربی طبقه $c$ را به شکل زیر تعریف می‌کنیم:
\begin{equation}\label{cmp:eq:ch3-margin-coeff-cov}
a_{c,r}=\left(\log\frac{\mathcal V_c(j)}{\mathcal V_r(j)}\right)_{j=1}^{K},
\qquad
\Sigma_c=\frac1{|C_c|}\sum_{i\in C_c}
(\mathscr Y_i-\mathcal V_c)^\top
(\mathscr Y_i-\mathcal V_c).
\end{equation}
در اینجا تقسیم و لگاریتم روی مؤلفه‌های مراکز انجام می‌شود و مثبت بودن مؤلفه‌های مراکز فرض شده است. با توجه به میانگین بودن $\mathcal V_c$، میانگین حاشیه نمونه‌های طبقه $c$ نسبت به طبقه رقیب $r$ دقیقاً برابر واگرایی دو مرکز است:
\begin{equation}\label{cmp:eq:ch3-margin-mean}
\boxed{
\mu_{c,r}=\frac1{|C_c|}\sum_{i\in C_c}M_{c,r}(X_i)
=\mathcal V_c a_{c,r}
=D_{\mathrm{KL}}(\mathcal V_c\|\mathcal V_r).}
\end{equation}
اما واریانس این حاشیه برابر است با:
\begin{equation}\label{cmp:eq:ch3-margin-variance}
\boxed{
\nu_{c,r}=\frac1{|C_c|}\sum_{i\in C_c}
\bigl(M_{c,r}(X_i)-\mu_{c,r}\bigr)^2
=a_{c,r}^{\top}\Sigma_c a_{c,r}.}
\end{equation}
پس واگرایی میان مراکز، \emph{میانگین} برتری طبقه واقعی در برابر یک رقیب را اندازه می‌گیرد، ولی درباره \emph{پراکندگی این برتری در میان نمونه‌ها} چیزی نمی‌گوید. واریانس $\nu_{c,r}$ تنها تغییرات بازنمایی در جهتی را می‌سنجد که بر تصمیم میان دو طبقه اثر دارد؛ نه تمام پراکندگی درون‌طبقه‌ای را.

\paragraph{دامنه اعتبار حاشیه KL.}
هویت میانگین و کوواریانس این قسمت
برای \emph{بازنمایی مشترک} نمونه
و مراکز ساخته‌شده از میانگین
همان نمونه‌های طبقه برقرار است.
برای جلوگیری از KL نامتناهی
می‌توان از هموارسازی مشترکِ نمونه‌ها و مراکز
استفاده کرد. صورت دقیق این قرارداد
در بخش نظریِ
\ref{cmp:sec:ch3-score-normalization-theory}
معرفی می‌شود؛ در آن صورت تمام
نمادهای $\mathscr Y$،
$\mathcal V$، $\mu$ و $\nu$ باید
به نسخهٔ هموارشده تفسیر شوند.
اگر در استنتاج شرطی برای هر رقیب
بازنمایی متفاوتی ساخته شود،
رابطه $\mu_{c,r}
=D_{\mathrm{KL}}(\mathcal V_c\|\mathcal V_r)$
لزومـاً برقرار نمی‌ماند
و به یک برهان مستقل نیاز دارد.




\subsection{کران‌های آماری خطای حاشیه: کانتلی و چرنوف}
\label{cmp:subsec:ch3-margin-criterion}
برای سنجش جدایی یک طبقه از رقیب، فقط بزرگ بودن میانگین حاشیه کافی نیست؛ مطلوب است که این میانگین در مقایسه با انحراف معیار حاشیه نیز بزرگ باشد. برای هر زوج مرتب $(c,r)$، معیار زیر را تعریف می‌کنیم:
\begin{equation}\label{cmp:eq:ch3-margin-pairwise}
\boxed{
G_{c,r}=\frac{\nu_{c,r}}{\nu_{c,r}+\mu_{c,r}^{\,2}},
\qquad c\ne r.}
\end{equation}
اگر دو مرکز یکسان باشند، $\mu_{c,r}=\nu_{c,r}=0$ و مقدار $G_{c,r}$ را برابر یک در نظر می‌گیریم. در دیگر حالت‌ها $0\le G_{c,r}\le1$. مقدار کوچک‌تر نشان می‌دهد میانگین جدایی نسبت به پراکندگی حاشیه بزرگ‌تر است. این معیار در حالت کلی جهت‌دار است؛ زیرا $\mu_{c,r}$ و $\nu_{c,r}$ برای جابه‌جایی طبقه مبدأ و رقیب الزاماً ثابت نمی‌مانند.

معیار کل جدایی حاشیه‌ای را با میانگین متوازن روی طبقه‌های مبدأ تعریف می‌کنیم:
\begin{equation}\label{cmp:eq:ch3-margin-criterion}
\boxed{
J_{\mathrm{margin}}
=\frac1K\sum_{c=1}^{K}\sum_{\substack{r=1\\r\ne c}}^{K}
G_{c,r}.}
\end{equation}
هدف کمینه کردن این معیار است. برخلاف معیار اطلاعات متقابل $I$ که تنها از مراکز $\mathcal V_c$ محاسبه می‌شود، $J_{\mathrm{margin}}$ علاوه بر مراکز به کوواریانس‌های $\Sigma_c$ نیز وابسته است؛ بنابراین شکل توزیع نمونه‌ها پیرامون مراکز را در جهت‌های مؤثر بر طبقه‌بندی لحاظ می‌کند.

\paragraph{ارتباط با خطای طبقه‌بندی.}
نمونه‌ای را به‌طور یکنواخت از طبقه $c$ انتخاب کنیم. متغیر حاشیه $M_{c,r}(X)$ میانگین $\mu_{c,r}$ و واریانس $\nu_{c,r}$ دارد. نامساوی یک‌طرفه چبیشف یا کانتلی نتیجه می‌دهد:
\begin{equation}\label{cmp:eq:ch3-margin-Cantelli}
\mathrm P\bigl(M_{c,r}(X)\le0\mid y_X=c\bigr)
\le G_{c,r}.
\end{equation}
اگر نمونه اشتباه طبقه‌بندی شود، دست‌کم یکی از طبقه‌های رقیب باید حاشیه‌ای نامثبت داشته باشد. از کران اجتماع رویدادها و میانگین‌گیری متوازن نتیجه می‌شود:
\begin{equation}\label{cmp:eq:ch3-margin-error-bound}
\boxed{
E_{\mathrm{bal}}
=\frac1K\sum_{c=1}^{K}\frac1{|C_c|}
\sum_{i\in C_c}\mathbf1_{\{\widehat y_i\ne c\}}
\le\min\{1,J_{\mathrm{margin}}\}.}
\end{equation}
اثبات دقیق این کران در بخش~\ref{app:ch3-margin-error} از پیوست~\ref{app:ch3-proofs} آمده است. این نتیجه درباره خطای تجربی روی همین داده‌ها و همین مراکز است؛ تضمینی برای خطای داده‌های دیده‌نشده بدون فرض‌های تعمیم‌پذیری نیست.

این کران محدودیت معیار صرفاً مبتنی بر فاصله مراکز را روشن می‌کند. حتی اگر میانگین حاشیه دو طبقه افزایش یابد، ممکن است پراکندگی حاشیه نیز چنان بزرگ شود که احتمال نمونه‌های واقع در سمت نادرست مرز تصمیم کاهش نیابد. معیار $J_{\mathrm{margin}}$ جدایی میانگین را نسبت به همین پراکندگی می‌سنجد و در برابر ضرب مشترک تمام حاشیه‌های یک زوج طبقه در ضریب مثبت ثابت می‌ماند.


\paragraph{کران نمایی چرنوف بر اساس حاشیه‌های تجربی.}
کران کانتلی تنها از میانگین و
واریانس حاشیه بهره می‌برد و ممکن
است به‌ویژه در توزیع‌های نامتقارن
یا چندقله‌ای شل باشد. وقتی
حاشیه‌ها متناهی‌اند، مثلاً با
استفاده از KL هموارشده،
برای هر $t>0$ نامساوی مارکوف
روی $e^{-tM_{c,r}}$ می‌دهد:
\begin{equation}\label{cmp:eq:ch3-margin-Chernoff}
\boxed{
\begin{aligned}
\mathrm P(M_{c,r}(X)\le0\mid y_X=c)
&\le G_{c,r}^{\mathrm{Ch}},\\
G_{c,r}^{\mathrm{Ch}}
&:=\inf_{t>0}
\frac1{|C_c|}\sum_{i\in C_c}
\exp[-t M_{c,r}(X_i)].
\end{aligned}
}
\end{equation}
نمونه $X$ در این رابطه
به‌طور یکنواخت از همان طبقه
و همان مجموعه تجربی انتخاب می‌شود.
پس با کران اجتماع رویدادها داریم:
\begin{equation}\label{cmp:eq:ch3-margin-Chernoff-balanced}
\boxed{
E_{\mathrm{bal}}\le
\frac1K\sum_{c=1}^{K}
\min\!\left\{1,\,
\sum_{r\ne c}G_{c,r}^{\mathrm{Ch}}\right\}.
}
\end{equation}
کران چرنوف و کانتلی
در حالت کلی ترتیب برتری
ثابتی نسبت به هم ندارند؛
برای ارزیابی می‌توان کمینه
دو کران معتبر را در نظر گرفت.
پارامتر $t$ مربوط به کران
نمایی است و با دمای نرمال‌سازی با بیشینهٔ هموار
و پهنای زیان سیگموئیدی
تفاوت دارد. کمینه‌گیری
داخلی روی $t$ به‌خودی‌خود
یک زیان هموار برای آموزش
سیاست تعریف نمی‌کند؛
کاربرد اصلی این کران،
تحلیل توزیع تجربی حاشیه‌هاست.
اثبات آن در بخش
\ref{app:ch3-margin-error}
از پیوست آمده است.



\subsection{خطای صفر--یک و پایداری هندسی حاشیه}
\label{cmp:subsec:ch3-margin-robustness}

\paragraph{کفایت برای ارزیابی، نه الزاماً برای یادگیری.}
برای مجموعه‌داده و قرارداد رفع تساوی ثابت،
خطای صفر--یک یا دقت متوازن دقیقاً
کیفیت تصمیم‌های نهایی را روی همان مجموعه مشخص می‌کند.
بنابراین برای \emph{شمارش تصمیم‌های درست} نیازی
به معیار مکمل نیست. بااین‌حال، این شاخص میان
نمونه درستِ نزدیک مرز و نمونه درستِ دور از مرز
تمایز نمی‌گذارد و در نواحی بدون تغییر برچسب،
نسبت به پارامترهای سیاست مشتق صفر دارد.
این مسئله هم محدودیت بهینه‌سازی مستقیم آن است
و هم نشان می‌دهد که درستی فعلی تصمیم،
پایداری آن در برابر اختلال را تضمین نمی‌کند.

\paragraph{تعبیر هندسی حاشیه: جدایی مراکز در برابر موقعیت نمونه.}
در حالت \emph{بازنمایی مشترک}،
بگذارید $z_i=\Phi(\mathscr Y_i)$ و
$v_c=\Phi(\mathcal V_c)$ باشند؛
تمام این تصاویر بر کره هیلبرتی
با نرم ثابت $\sqrt{\log2/2}$ قرار دارند.
برای یک رقیب $r$ با
$\delta_{c,r}=\|v_c-v_r\|_{\mathcal H}>0$،
بردار یکه اتصال مراکز و فاصله علامت‌دار
نمونه از صفحه عمودمنصف را با
\[
e_{c,r}=\frac{v_c-v_r}{\delta_{c,r}},
\qquad
h_{i,c,r}
=\left\langle
z_i-\frac{v_c+v_r}{2},e_{c,r}
\right\rangle_{\mathcal H}
\]
تعریف می‌کنیم. از هم‌نرم‌بودن مراکز
و رابطه~\eqref{cmp:eq:ch3-fixed-center-pairwise-js-margin}
نتیجه می‌شود:
\begin{equation}\label{cmp:eq:ch3-js-margin-separation-factorization}
\boxed{
m_{i,c,r}
=s(\mathscr Y_i,\mathcal V_c)
-s(\mathscr Y_i,\mathcal V_r)
=\frac{2\delta_{c,r}}{\log2}\,h_{i,c,r},
\qquad
|m_{i,c,r}|
\le\sqrt{\frac{2}{\log2}}\,\delta_{c,r}.
}
\end{equation}
این تجزیه نشان می‌دهد اندازه حاشیه شباهت
هم از \emph{جدایی مراکز}
و هم از \emph{موقعیت جهت‌دار نمونه}
اثر می‌پذیرد. بنابراین حاشیه کوچک
ممکن است ناشی از نزدیک‌بودن مراکز باشد،
حتی وقتی نمونه در سمت صحیح و نسبتاً دور
از صفحه تصمیم قرار گرفته است.
اگر دو مرکز یکسان باشند،
امتیازهای آن‌ها برای تمام نمونه‌ها برابر
و صفحه تصمیمِ یکتای میان آن دو نامعین است.
توجه کنیم این صفحه در \emph{فضای هیلبرتی}
خطی است؛ مرز متناظر در فضای توزیع‌های
پیش از نگاشت لزوماً خطی نیست.

\paragraph{گواه پایداری با حاشیه متریک.}
برای بررسی مقاومت تصمیم، صرف اندازه
حاشیه شباهت کافی نیست؛ حاشیه اختلاف
فاصله‌های متریک را نیز تعریف می‌کنیم:
\begin{equation}\label{cmp:eq:ch3-metric-robust-margin}
\boxed{
\gamma_{i,c,r}
=d_{\mathrm{JS}}(\mathscr Y_i,\mathcal V_r)
-d_{\mathrm{JS}}(\mathscr Y_i,\mathcal V_c),
\qquad
\gamma_i=\min_{r\ne y_i}\gamma_{i,y_i,r}.
}
\end{equation}
این حاشیه و $m_{i,c,r}$ علامت
یکسان دارند، ولی در برابر اختلال
رفتارشان باید با توجه به متریک سنجیده شود.
اگر نمونه به $\mathscr Y_i'$ و مراکز
به $\mathcal V_c'$ تغییر کنند و
\[
\begin{aligned}
d_{\mathrm{JS}}(\mathscr Y_i,\mathscr Y_i')&\le\varepsilon_i,\\
d_{\mathrm{JS}}(\mathcal V_c,\mathcal V_c')&\le\varepsilon_c,
\end{aligned}
\]
آنگاه برای هر رقیب $r$، نامساوی مثلثی
کران زیر را به‌دست می‌دهد:
\begin{equation}\label{cmp:eq:ch3-margin-perturbation-bound}
\boxed{
\gamma_{i,c,r}'
\ge \gamma_{i,c,r}
-2\varepsilon_i-\varepsilon_c-\varepsilon_r.
}
\end{equation}
در نتیجه اگر برای تمام رقبا
$\gamma_{i,y_i,r}>
2\varepsilon_i+\varepsilon_{y_i}+\varepsilon_r$،
پیش‌بینی صحیح نمونه تحت این
اختلال‌ها حفظ می‌شود.
این یک \emph{شرط کافیِ پایداری} است؛
نقض آن به معنای وقوع قطعی خطا نیست.
در حالت مراکز ثابت، تنها جمله
$2\varepsilon_i$ باقی می‌ماند؛
اگر نمونه و همه مراکز حداکثر
$\varepsilon$ تغییر کنند،
شرط ساده‌تر $\gamma_i>4\varepsilon$
کفایت می‌کند.

از نامساوی مثلثی همچنین
$\gamma_{i,c,r}\le\delta_{c,r}$
به‌دست می‌آید. پس وقتی مراکز نزدیک‌اند،
گواه پایداری مبتنی بر حاشیه متریک
در برابر بودجه اختلال ثابت
محدود می‌شود؛ بااین‌حال،
\emph{نزدیکی مراکز به‌تنهایی ناپایداری
واقعی هر نمونه را اثبات نمی‌کند}.
نتیجه حاضر به اندازه و نوع
اختلال مجاز نیز وابسته است.
اگر در استنتاج شرطی برای
هر نامزد بازنمایی متفاوتی ساخته شود،
هویت صفحه عمودمنصف یک بازنمایی مشترک
دیگر قابل استفاده نیست و کران پایداری
باید جداگانه برای تغییر هر امتیاز
شرطی محاسبه شود.
اثبات این نتایج و حدود آن‌ها در
پیوست~\ref{app:ch3-margin-robustness}
ارائه شده است.


در نتیجه، حاشیه علاوه بر شمارش خطا،
موقعیت نمونه نسبت به مرز و امکان پایداری تصمیم
را نیز توصیف می‌کند. ساخت زیان‌های آموزشی از حاشیه
و بررسی رابطهٔ آن‌ها با معیارهای احتمالاتی و هندسی
در بخش بعد انجام می‌شود.


\section{توابع زیان برای پیش‌بینی برچسب و حاشیهٔ تصمیم}
\label{cmp:sec:ch3-learning-objectives}
در چارچوب طبقه‌بند معرفی‌شده، دو کمیت پیوسته
برای یادگیری در دسترس‌اند: احتمال تخمینیِ برچسب
و حاشیهٔ تصمیم. بر این اساس، زیان
لگاریتمیِ مبتنی بر احتمال، زیان هموار
حاشیه‌ای و صورت Softplus--JS را تعریف می‌کنیم.
این تعریف‌ها به‌خودی‌خود الگوریتم آموزش نیستند؛
نقش مقیاس و پایداری عددی، تضمین‌های نظری
و مشتق‌گیری سیاست به‌ترتیب در بخش‌های بعد بررسی می‌شوند.

\subsection{زیان لگاریتمیِ پیش‌بینی برچسب}
\label{cmp:sec:ch3-NLL}
در بخش~\ref{cmp:subsec:ch3-class-centers}،
از ناهمانندی نمونه تا مراکز، احتمال‌های تخمینی
$\widehat p_{i,c}$ ساخته شدند. اکنون تابع زیان
مناسب برای ارزیابی و یادگیری این احتمال‌ها را
مستقل از انتخاب اولیهٔ KL یا JSD تعریف می‌کنیم.
این زیان مربوط به \emph{برچسب واقعی نمونه} است
و با هزینهٔ درست‌نماییِ مدل مشترکِ مشاهدات همسایگی
در بخش~\ref{cmp:subsec:ch3-mean-center-justification}
یکسان نیست.

\paragraph{زیان نمونه و میانگین وزن‌دار آن.}
برای نمونهٔ دارای برچسب واقعی $y_i$، بر اساس احتمال تخمینی عمومی، زیان لگاریتمی منفی به شکل زیر است؛ این رابطه برای هر سنجهٔ ناهمانندی با امتیاز $-d_{i,c}$ و فاصله‌های متناهی برقرار است:
\begin{equation}\label{cmp:eq:ch3-NLL-sample}
\begin{aligned}
\ell_i
&=-\log\widehat p_{i,y_i}\\
&=d_{i,y_i}+\log\sum_{r=1}^{K}e^{-d_{i,r}}\\
&=\log\left(1+\sum_{r\ne y_i}e^{d_{i,y_i}-d_{i,r}}\right).
\end{aligned}
\end{equation}
با استفاده از وزن طبقات $\Omega=(\omega_1,\ldots,\omega_K)$، به شرط $\omega_c\ge0$ و $\sum_c\omega_c=1$، معیار را به‌صورت میانگین وزن‌دار خطاهای درون هر طبقه تعریف می‌کنیم:
\begin{equation}\label{cmp:eq:ch3-NLL-balanced}
\boxed{
NLL_{\Omega}=\sum_{c=1}^{K}\frac{\omega_c}{|C_c|}
\sum_{i\in C_c}\ell_i.}
\end{equation}
این تعریف اجازه می‌دهد اهمیت طبقات مستقل از تعداد نمونه‌هایشان تعیین شود. اگر $\omega_c=|C_c|/n$ باشد، معیار به میانگین خطای تمام نمونه‌ها تبدیل می‌شود؛ در مقابل، با $\omega_c=1/K$ همه طبقات بدون توجه به فراوانی نمونه‌هایشان سهم یکسان دارند. حالت دوم به‌ویژه برای جلوگیری از غلبه طبقات پرتعداد و بررسی کیفیت تفکیک همه طبقات مناسب است. در ادامه، حالت یکنواخت را با $NLL_U$ نشان می‌دهیم، درحالی‌که تعریف اصلی معیار همان $NLL_\Omega$ است.

\phantomsection\label{cmp:subsec:ch3-NLL-definition}
\paragraph{صورت KL و ارتباط آن با پسین واقعی.}
برای روشن شدن انگیزه استفاده از KL و نرمال‌سازی با بیشینهٔ هموار، ابتدا حالتی فرضی را در نظر بگیریم که توزیع پسین واقعی برچسب نمونه، با نماد $p_i(c)=\mathrm P(y_i=c\mid X_i)$، معلوم باشد. اگر $e_c$ بردار یک‌داغ طبقه $c$ باشد، از تعریف واگرایی KL داریم:
\begin{equation}\label{cmp:eq:ch3-NLL-true-posterior-KL}
D_{\mathrm{KL}}(e_c\|p_i)=-\log p_i(c),
\qquad
\exp\bigl(-D_{\mathrm{KL}}(e_c\|p_i)\bigr)=p_i(c).
\end{equation}
در این رابطه، برای $p_i(c)=0$ قرارداد $D_{\mathrm{KL}}(e_c\|p_i)=+\infty$ و $e^{-\infty}=0$ را به کار می‌بریم. چون مجموع مؤلفه‌های $p_i$ یک است، نرمال‌سازی نمایی منفی این واگرایی‌ها همان توزیع اولیه را بازمی‌گرداند:
\begin{equation}\label{cmp:eq:ch3-NLL-true-posterior-softmax}
p_i(c)=
\frac{\exp[-D_{\mathrm{KL}}(e_c\|p_i)]}
{\sum_{r=1}^{K}\exp[-D_{\mathrm{KL}}(e_r\|p_i)]}.
\end{equation}
این برابری توضیح می‌دهد که چگونه نرمال‌سازی با بیشینهٔ هموار می‌تواند امتیازهای مبتنی بر KL را به احتمال تبدیل کند؛ اما در مسئله ما $p_i$ مجهول است و رابطه بالا به‌تنهایی روشی برای محاسبه آن فراهم نمی‌کند.

در صورت مبتنی بر KL، سنجهٔ عمومی $d$ را با
واگرایی $D_{\mathrm{KL}}$ جایگزین می‌کنیم:
\begin{equation}\label{cmp:eq:ch3-NLL-distance}
d_{i,c}=D_{\mathrm{KL}}(\mathscr Y_i\|\mathcal V_c).
\end{equation}
با جایگذاری این انتخاب در نرمال‌سازی
رابطه~\eqref{cmp:eq:ch3-NLL-softmax}،
احتمال‌های تخمینی KL به‌دست می‌آیند.
جهت واگرایی در اینجا $\mathscr Y_i\|\mathcal V_c$ است؛
برابری صوری نرمال‌سازی نمایی با پسینِ معلوم در
رابطه~\eqref{cmp:eq:ch3-NLL-true-posterior-softmax}
به معنای آن نیست که خروجی مدل، پسین واقعی باشد.
اگر دست‌کم یک فاصله متناهی باشد، قاعدهٔ تصمیم برابر است با:
\begin{equation}\label{cmp:eq:ch3-estimated-label-from-neighborhood}
\boxed{
\widehat y_i
=\arg\max_{c}\widehat p_{i,c}
=\arg\min_{c}D_{\mathrm{KL}}(\mathscr Y_i\|\mathcal V_c).
}
\end{equation}
این برابری حالت خاص قاعدهٔ عمومیِ نزدیک‌ترین مرکز است.
در صورت برابری چند امتیاز، قرارداد رفع تساوی باید مشخص شود.
نرمال‌سازی با بیشینهٔ هموار علاوه بر تصمیم، توزیع احتمال
تخمینی تولید می‌کند؛ صحت کالیبراسیون آن
به ارزیابی یا استدلال جداگانه نیاز دارد.

\subsection{صورت‌های زیان مبتنی بر JSD و جذر آن}
\label{cmp:subsec:ch3-NLL-JS-sqrt-variants}

KL ممکن است نامتناهی شود و متقارن نیست.
در مقابل JS متقارن و کراندار است و
جذر آن، برخلاف خود واگرایی JS،
متریک فاصله با نامساوی مثلثی است.
این ویژگی‌های هندسی مستقل از مسئله
کالیبراسیون احتمال‌اند: کرانداربودن
فاصله، خروجی نرمال‌سازی با بیشینهٔ هموار را به پسین واقعی تبدیل
نمی‌کند و مسئله انتخاب دما را نیز
به‌تنهایی حل نمی‌کند.

یکی از محدودیت‌های فاصله KL آن است که با نامنطبق‌شدن
تکیه‌گاه‌ها می‌تواند بی‌نهایت شود. با تعریف متریک
جنسن--شنون، می‌توان به‌جای آن از دو فاصله کراندار
استفاده کرد. برای هر نمونه و طبقه نامزد می‌نویسیم:
\begin{equation}\label{cmp:eq:ch3-NLL-JS-distances}
\begin{aligned}
d^{\mathrm{JS}}_{i,c}
&=\operatorname{JSD}(\mathscr Y_i,\mathcal V_c),\\
d^{\sqrt{\mathrm{JS}}}_{i,c}
&=\sqrt{\operatorname{JSD}(\mathscr Y_i,\mathcal V_c)}.
\end{aligned}
\end{equation}
برای هر انتخاب $D\in\{\mathrm{KL},\mathrm{JS},
\sqrt{\mathrm{JS}}\}$ و مقیاس مثبت
$\tau_D$، خطای نرمال‌سازی با بیشینهٔ هموار مبتنی بر فاصله را به‌صورت
\begin{equation}\label{cmp:eq:ch3-NLL-distance-family}
\boxed{
\begin{aligned}
\ell_i^{D,\tau_D}
&=-\log\frac{\exp[-d^D_{i,y_i}/\tau_D]}
{\sum_{r=1}^{K}\exp[-d^D_{i,r}/\tau_D]},\\
NLL^{D,\tau_D}_\Omega
&=\sum_{c=1}^{K}\frac{\omega_c}{|C_c|}
\sum_{i\in C_c}\ell_i^{D,\tau_D}.
\end{aligned}}
\end{equation}
تعریف می‌کنیم. مدل پیشین KL متناظر با $\tau_{\mathrm{KL}}=1$
است. برای JS، اگر $\tau_{\mathrm{JS}}=\log2$ قرار دهیم،
نرمال‌سازی با بیشینهٔ هموار روی $-d^{\mathrm{JS}}/\log2$ به دلیل
ناوردایی نسبت به افزودن ثابت دقیقاً برابر نرمال‌سازی با بیشینهٔ هموار
روی شباهت $s=1-d^{\mathrm{JS}}/\log2$ است؛
این همان تعریف \emph{NLL شباهت JS} خواهد بود.
در مقابل، نرمال‌سازی با بیشینهٔ هموار روی $-d^{\mathrm{JS}}$
به دمای $\tau_{\mathrm{JS}}=1$ مربوط است و
لزوماً احتمال‌های یکسانی تولید نمی‌کند.
برای جذر JS نیز باید $\tau_{\sqrt{\mathrm{JS}}}$
به‌طور مستقل تعیین و گزارش شود؛ زیرا واحد
فاصله آن با JSD متفاوت است.

ازآنجاکه ریشه دوم صعودی است،
فاصله‌های JS و $\sqrt{\mathrm{JS}}$ در
\emph{سیاست و مراکز ثابت} طبقه برنده یکسانی
می‌دهند، ولی اندازه شکاف بین امتیازهای رقبا
و در نتیجه احتمال‌های حاصل از نرمال‌سازی با بیشینهٔ هموار، NLL و گرادیان‌های آموزش
آن‌ها یکسان نیست. هر دو فاصله متناهی و کراندارند؛
در عوض، این کرانداری به‌تنهایی احتمال‌های
خروجی نرمال‌سازی با بیشینهٔ هموار را به پسین بیزی تبدیل نمی‌کند.

تعریف این خانواده‌ها به معنای اثبات بهتر بودن
یک صورت NLL نیست. برای سیاست شرطی طبقه، در
تمام روابط بالا باید در امتیاز نامزد $c$،
$\mathscr Y_i$ را با
$\mathscr Y_i^{(c)}$ جایگزین کرد؛ برچسب واقعی
فقط برای محاسبه زیان نظارت‌شده وارد می‌شود.
مقایسه عددی این انتخاب‌ها و محدودیت‌های
یادگیری، موضوع فصل~\ref{chap:results} است.

\noindent\textbf{صورت تصمیم در سیاست شرطی طبقه.}
اگر برای هر نامزد $c$ بازنمایی
$\mathscr Y_i^{(c)}$ حاصل شود، باید
امتیاز همان نامزد را از مقایسه
بازنمایی شرطی با مرکز طبقه به‌دست آورد:
\[
d_{i,c}^{D,\mathrm{cond}}
=D\bigl(\mathscr Y_i^{(c)},\mathcal V_c\bigr),
\qquad
\widehat y_i=\arg\min_c d_{i,c}^{D,\mathrm{cond}}.
\]
در صورت استفاده از نرمال‌سازی با بیشینهٔ هموار، تمام امتیازهای
نامزدها باید با یک قرارداد مشترک
برای دما یا نرمال‌سازی مقایسه شوند.
برچسب واقعی در زمان آزمون در انتخاب
سیاست نامزد نقش ندارد و فقط هنگام
محاسبه زیان نظارت‌شده وارد می‌شود.
نتایج کران‌بندی NLL بر پایهٔ جدایی مراکز،
که در بخش تحلیل نظریِ
\ref{cmp:sec:ch3-loss-quality-theory}
برای بازنمایی مشترک بررسی می‌شوند،
بدون استدلال تازه به این حالت تعمیم نمی‌یابند.



\subsection{زیان Softplus--JS برای امتیازهای طبقات}
\label{cmp:subsec:ch3-softplus-js-objective}
علاوه بر NLL و زیان سیگموئیدیِ حاشیه،
صورت \emph{Softplus--JS} را برای یادگیری
از امتیازهای تمام طبقات نامزد در نظر می‌گیریم.
برای هر نمونهٔ آموزشی $i$، امتیاز نامزد $r$
را با قرارداد شباهت جنسن--شنون تعریف می‌کنیم:
\[
S_{i,r}:=
1-\frac{\operatorname{JSD}(\mathscr Y_i^{[r]},\mathcal V_r)}{\log2}.
\]
در سیاست مستقل از فرضیهٔ طبقه،
$\mathscr Y_i^{[r]}$ برای همهٔ نامزدها یکی است؛
در سیاست شرطی طبقه، بازنمایی هر نامزد با
\emph{فرضیهٔ همان طبقه} ساخته می‌شود
و در هیچ‌کدام برچسب واقعی نمونه برای تعیین
سیاست نامزد به کار نمی‌رود. تنها در زیان نظارت‌شده
از $y_i$ استفاده می‌کنیم.
با مقیاس $\tau_{\mathrm{SP}}>0$، زیان نمونه
عبارت است از:
\begin{equation}\label{cmp:eq:ch3-softplus-js-training-objective}
\boxed{
\ell_i^{\mathrm{SP}}
=\tau_{\mathrm{SP}}
\log\sum_{r=1}^{K}
\exp\left(\frac{S_{i,r}-S_{i,y_i}}{\tau_{\mathrm{SP}}}\right).
}
\end{equation}
این تعریف در پیاده‌سازی با
$\operatorname{LogSumExp}$ پایدار محاسبه می‌شود.
$S_{i,r}$ امتیاز \emph{طبقهٔ نامزد} است
و با logits انتخاب ویژگی اشتباه گرفته نمی‌شود.
مقدار $\tau_{\mathrm{SP}}$ با پهنای
$\tau$ در زیان سیگموئیدی تفاوت دارد.
مانند دیگر اهداف، زیان متوازن از میانگین
وزن‌دار زیان نمونه‌ها در طبقات آموزش حاصل می‌شود.
این تابع هدف لزوماً با زیان سیگموئیدی
یا NLL هم‌ارز نیست؛ مشتق‌های آن در
بخش~\ref{cmp:subsec:ch3-softplus-js-policy-gradient}
بررسی می‌شوند.

\subsection{زیان هموار حاشیه و حساسیت نزدیک مرز تصمیم}
\label{cmp:subsec:ch3-sigmoid-balanced-error}
در بخش حاشیه، علامت
رابطه~\eqref{cmp:eq:ch3-general-decision-margin}
معیار درست یا نادرست بودن تصمیم بود.
اکنون می‌خواهیم از \emph{اندازهٔ}
حاشیه نیز برای ساخت تابع هدفی پیوسته
استفاده کنیم تا تغییرات سیاست
در نزدیکی مرز تصمیم قابل‌پایش باشند.

\paragraph{چرا زیان سیگموئیدی؟ تقریب مشتق‌پذیرِ خطای صفر--یک.}
به‌جای کمینه‌کردن مستقیم نشانگر ناپیوسته خطای صفر--یک،
یک \emph{جانشین هموار و کراندار} از حاشیه می‌سازیم.
سیگموئید نزولی انتخاب مناسبی است چون
علامت حاشیه را به جریمه نزدیک به صفر برای
پیش‌بینی‌های درستِ مطمئن و نزدیک به یک برای
پیش‌بینی‌های نادرستِ مطمئن تبدیل می‌کند،
پارامتر $\tau$ پهنای ناحیه حساس را کنترل می‌کند،
و مشتق آن نسبت به حاشیه نزدیک مرز بیشترین اندازه را دارد.
این مزایا، تضمین بهینه سراسری یا بهبود یکنواخت دقت نیستند.
برای $\tau>0$ تعریف می‌کنیم:
\begin{equation}\label{cmp:eq:ch3-js-sigmoid-loss}
\boxed{
\ell_\tau(m)=\frac{1}{1+\exp(m/\tau)},
\qquad 0<\ell_\tau(m)<1.
}
\end{equation}
در این رابطه $\tau$ پارامتر پهنای ناحیه حساس اطراف مرز تصمیم است و با اندیس تصادفی گام مشاهده $\xi$ و پارامتر توقف هندسی $\beta$ تفاوت دارد. برای هر حاشیه ثابت $m\ne0$، با میل‌کردن $\tau$ به صفر داریم
$\ell_\tau(m)\to\mathbf1\{m<0\}$؛ اما برای $m=0$، مقدار تابع همواره $1/2$ است. بنابراین در صورت وجود تساوی، حد تابع هدف به هر تساوی وزن نصف می‌دهد و لزوماً با قرارداد محافظه‌کارانه رابطه~\eqref{cmp:eq:ch3-js-hard-margin-error} برابر نیست.

تابع هدف \emph{هموارِ خطای متوازن} را با میانگین‌گیری برابر از طبقه‌ها تعریف می‌کنیم:
\begin{equation}\label{cmp:eq:ch3-js-sigmoid-balanced-objective}
\boxed{
\mathcal L_{\mathrm{sig}}(\tau)
=\frac1K\sum_{c=1}^{K}\frac1{|C_c|}
\sum_{i\in C_c}\ell_\tau(M_i^{\mathrm{JS}}).
}
\end{equation}
در صورت وابستگی بازنمایی‌ها و مراکز به سیاست $\pi_\theta$، این مقدار به تابعی از پارامترهای سیاست، یعنی $\mathcal L_{\mathrm{sig}}(\theta;\tau)$، تبدیل می‌شود. هدف کمینه‌کردن آن است؛ انتخاب مقدار $\tau$ باید با داده آموزش یا اعتبارسنجی و مستقل از داده آزمون انجام گیرد.

\paragraph{مزیت زیان سیگموئیدی: حساسیت گرادیانی نزدیک مرز تصمیم.}
مشتق زیان نسبت به حاشیه عبارت است از:
\begin{equation}\label{cmp:eq:ch3-js-sigmoid-derivative}
\boxed{
\frac{\mathrm d\ell_\tau(m)}{\mathrm dm}
=-\frac{1}{\tau}\ell_\tau(m)
\bigl(1-\ell_\tau(m)\bigr)
=-\frac{1}{4\tau}
\operatorname{sech}^{2}\!\left(\frac{m}{2\tau}\right).
}
\end{equation}
برای مقایسهٔ شکل حساسیت در مقادیر مختلف
$\tau$، \emph{حساسیت نرمال‌شده نسبت به خود حاشیه} را
به‌صورت
\begin{equation}\label{cmp:eq:ch3-js-sigmoid-normalized-sensitivity}
\boxed{
g_\tau(m):=4\tau\left|\ell'_\tau(m)\right|
=4\ell_\tau(m)\bigl(1-\ell_\tau(m)\bigr)
=\operatorname{sech}^{2}\!\left(\frac{m}{2\tau}\right),
\qquad 0<g_\tau(m)\le1
}
\end{equation}
تعریف می‌کنیم؛ بیشینهٔ یک دقیقاً در $m=0$ رخ می‌دهد.
پارامتر $\tau$ در این قرارداد \emph{پهنای گذار}
است، زیرا $\ell_\tau(m)=\sigma(-m/\tau)$،
نه $\sigma(-\tau m)$؛ ازاین‌رو
$g_\tau$ گرادیان سیاست یا بزرگی
$\|\nabla_\theta\ell_\tau(M_i(\theta))\|$
نیست. در نقاط مشتق‌پذیر،
$\nabla_\theta\ell_\tau(M_i(\theta))
=\ell'_\tau(M_i(\theta))\nabla_\theta M_i(\theta)$
و عامل دوم می‌تواند میان نمونه‌ها تفاوت زیادی داشته باشد.

در نتیجه، بزرگی مشتق در $m=0$ برابر $1/(4\tau)$ و بیشینه است؛ با افزایش $|m|$ کاهش یافته و به صفر میل می‌کند. بنابراین نمونه‌های اشتباهِ نزدیک مرز، نسبت به نمونه‌های اشتباهِ بسیار دور از آن، سهم گرادیانی بیشتری از ناحیه تغییر حاشیه دارند. نمونه‌های درست اما شکننده، یعنی نمونه‌هایی با حاشیه مثبت کوچک، نیز در تابع هدف مؤثر باقی می‌مانند. در مقابل، جریمه نمونه‌های بسیار دشوار یا نویزی حداکثر به یک نزدیک می‌شود و مانند زیان مربعی حاشیه، با شدت خطا بدون کران افزایش نمی‌یابد.

این ویژگی به معنای کم‌هزینه‌بودن قطعی اصلاح نمونه‌های نزدیک مرز نیست؛ زیرا تغییر حاشیه هر نمونه به میزان حساسیت آن به پارامترهای سیاست نیز وابسته است. در نقاطی که مشتق حاشیه نسبت به $\theta$ وجود دارد، سهم آن نمونه در گرادیان کل از ضرب $\ell_\tau'(M_i^{\mathrm{JS}})$ در $\nabla_\theta M_i^{\mathrm{JS}}$ به‌دست می‌آید. تغییر یک سیاست مشترک می‌تواند هم‌زمان حاشیه چند نمونه را بهبود دهد یا تضعیف کند؛ ازاین‌رو هیچ تضمینی وجود ندارد که هر گام کاهش این تابع الزاماً تعداد خطاهای سخت را کاهش دهد.


شکل~\ref{cmp:fig:ch3-sigmoid-loss-formula} رفتار زیان سیگموئیدی و شکل~\ref{cmp:fig:ch3-sigmoid-sensitivity-formula} حساسیت نرمال‌شده آن را برای چهار مقدار مختلف $\tau$ نشان می‌دهند. هرچه $\tau$ کوچک‌تر شود، ناحیه گذار زیان و ناحیه اثرگذاری گرادیان در اطراف مرز تصمیم $m=0$ باریک‌تر می‌شوند. مقیاس عمودی شکل دوم برابر $4\tau|\ell_\tau'(m)|$ است تا بیشینه حساسیت برای همه منحنی‌ها به یک نرمال شود؛ ازاین‌رو ارتفاع‌های شکل دوم بیانگر اندازه مطلق گرادیان برای $\tau$های مختلف نیستند.

\begin{figure}[htbp]
\centering
\includegraphics[width=0.86\linewidth]{chap3/sigmoid_loss_formula_bw.pdf}
\caption{زیان سیگموئیدی بر حسب حاشیه تصمیم $m$ برای چهار مقدار پارامتر $\tau$.}
\label{cmp:fig:ch3-sigmoid-loss-formula}
\end{figure}

\begin{figure}[htbp]
\centering
\includegraphics[width=0.86\linewidth]{chap3/sigmoid_sensitivity_formula_bw.pdf}
\caption{حساسیت نرمال‌شده $4\tau|\ell_\tau'(m)|$ نسبت به حاشیه $m$ برای چهار مقدار $\tau$؛ بیشینه همه منحنی‌ها در مرز تصمیم رخ می‌دهد.}
\label{cmp:fig:ch3-sigmoid-sensitivity-formula}
\end{figure}

\paragraph{مسئله همواری نزدیک‌ترین رقیب و حالت سیاست شرطی طبقه.}
خود تابع $\ell_\tau$ نسبت به حاشیه هموار است، اما عملگر بیشینه یا کمینه در تعیین نزدیک‌ترین رقیب ممکن است در نقاط تساوی رقبا مشتق‌پذیر نباشد. در صورت نیاز به تابعی کاملاً هموار بر حسب \emph{امتیازهای} طبقه‌ها، می‌توان بیشینه امتیاز رقیب را با تقریب $\kappa\log\sum_{r\ne y_i}\exp(s_{i,r}/\kappa)$ برای $\kappa>0$ جایگزین کرد؛ این تقریب با میل $\kappa$ به صفر به بیشینه واقعی نزدیک می‌شود. چنین هموارسازی‌ای جدا از پارامتر $\tau$ است و ممکن است مقدار حاشیه را به اندازه‌ای محدود تغییر دهد.

در روش اصلی با یک بازنمایی مشترک برای هر نمونه، $s_{i,r}=s(\mathscr Y_i,\mathcal V_r)$ است. اگر در یک روش دیگر برای هر \emph{فرض طبقه} بازنمایی جداگانه $\mathscr Y_i^{(r)}$ ساخته شود، باید به جای آن $s_{i,r}=s(\mathscr Y_i^{(r)},\mathcal V_r)$ قرار گیرد تا حاشیه با قاعده تصمیم همان روش سازگار بماند. برچسب واقعی نمونه در این حالت فقط برای محاسبه زیان نظارت‌شده به کار می‌رود و نباید در ساخت پیش‌بینی نمونه آزمون استفاده شود.

این معیار بر خلاف $I$ یا شعاع درون‌طبقه‌ای $R$، مستقیماً علامت حاشیه تصمیم نمونه‌ها را هدف می‌گیرد. بااین‌حال، تابع زیان سیگموئیدی در کل محدب نیست و کمینه‌سازی آن نیز به‌تنهایی تضمین بهینه سراسری یا تعمیم به داده‌های دیده‌نشده را فراهم نمی‌کند. ازاین‌رو رابطه این معیار با دقت طبقه‌بندی و کیفیت آموزش سیاست، مسئله‌ای تجربی و تحلیلی برای فصل~\ref{chap:results} خواهد بود.






\subsection{کران زیان حاشیه‌ای و خطای تجربی}
زیان سیگموئیدی رابطه
\eqref{cmp:eq:ch3-js-sigmoid-loss}
به علامت و اندازه حاشیه حساس است؛
اما میانگین این زیان نیز با خطای صفر--یک
\emph{هم‌ارز ترتیبی} نیست:
ممکن است هنگام کاهش میانگین زیان نرم،
تعداد نمونه‌های اشتباه بیشتر شود.
با وجود این، چون
$\ell_\tau(m)\ge\tfrac12$ برای $m\le0$،
برای هر $\tau>0$ کران عمومی
\begin{equation}\label{cmp:eq:ch3-margin-sigmoid-error-bound}
\boxed{
\widehat E_{\mathrm{bal}}
\le \widehat E_{\mathrm{bal}}^{\,\le0}
\le 2\mathcal L_{\mathrm{sig}}(\tau)
}
\end{equation}
برقرار است. پس زیان نرم یک \emph{کران قابل کمینه‌سازی}
برای خطای تجربی فراهم می‌کند، نه تضمین
بهبود یکنواخت آن در هر گام.
برای تقریب هموار نزدیک‌ترین رقیب نیز
$\widetilde M_{i,\kappa}^{\mathrm{JS}}\le M_i^{\mathrm{JS}}$
و نزولی‌بودن $\ell_\tau$
همین کران را با زیان تقریب‌زده حفظ می‌کند.
یک مثال نقض برای هم‌ارزی ترتیبی
و اثبات کران در پیوست~\ref{app:ch3-margin-robustness}
آمده است.


\section{مقیاس امتیازها و پایداری عددی طبقه‌بندی}
\label{cmp:sec:ch3-score-normalization-analysis}
احتمال‌های Softmax و زیان لگاریتمی علاوه بر ترتیب
طبقات، به مقیاس امتیازها وابسته‌اند. در این بخش،
اثر دما و تغییر مقیاس را از مسئلهٔ مؤلفه‌های صفر
و نامتناهی‌شدن KL تفکیک می‌کنیم و قراردادهای
پایهٔ محاسباتی را توضیح می‌دهیم.
صورت‌های تخصصی نرمال‌سازی و کران‌های حاصل
در بخش تکمیلیِ
\ref{cmp:sec:ch3-score-normalization-theory}
تحلیل می‌شوند.

\subsection{دما، مقیاس و نرمال‌سازی امتیازهای Softmax}
\label{cmp:subsec:ch3-NLL-scaling}

دو محدودیت مستقل باید از هم جدا شوند:
تغییر مقیاس امتیازها، با وجود حفظ برچسب برنده،
مقدار NLL را عوض می‌کند؛ اما نامتناهی‌شدن KL
از ناسازگاری تکیه‌گاه توزیع‌ها ناشی می‌شود.
نرمال‌سازی مقیاس لزوماً مسئله دوم را حل نمی‌کند
و جایگزینی KL با JS نیز وابستگی نرمال‌سازی با بیشینهٔ هموار
به دما را به‌خودی‌خود حذف نمی‌کند.

برای روشن شدن تفاوت کیفیت تصمیم با مقدار NLL، ابتدا چند خاصیت نرمال‌سازی با بیشینهٔ هموار را بررسی می‌کنیم. اگر $S(v)_c=e^{v_c}/\sum_re^{v_r}$، افزودن \emph{یک ثابت مشترک} به تمام مؤلفه‌های بردار امتیاز، احتمال‌ها را تغییر نمی‌دهد؛ اما ضرب بردار امتیاز در ضریب مثبت، در عین حفظ طبقه برنده، می‌تواند احتمال‌ها را تغییر دهد:
\begin{equation}\label{cmp:eq:ch3-NLL-softmax-invariance}
S(v+b\mathbf1)=S(v),
\qquad
\arg\max_c S(av)_c=\arg\max_c v_c,\quad a>0.
\end{equation}
در رابطه دوم، برابری بردار احتمال‌ها ادعا نشده است. همچنین افزودن یک \emph{بردار دلخواه} به امتیازها، برخلاف افزودن ثابت مشترک، ممکن است حتی طبقه پیش‌بینی‌شده را تغییر دهد. اثبات خواص نرمال‌سازی با بیشینهٔ هموار در پیوست~\ref{app:ch3-softmax-properties} آمده است.

در مدل ما امتیازهای ورودی نرمال‌سازی با بیشینهٔ هموار برابر $v_{i,c}=-d_{i,c}$ هستند. برای بررسی وابستگی به مقیاس، آن‌ها را در ضریب مثبت $\gamma$ ضرب می‌کنیم:
\begin{equation}\label{cmp:eq:ch3-NLL-scaled-softmax}
\widehat p_{i,c}^{(\gamma)}
=\frac{e^{-\gamma d_{i,c}}}
{\sum_{r=1}^{K}e^{-\gamma d_{i,r}}},\qquad\gamma>0.
\end{equation}
ضریب $\gamma$ معادل معکوس دمای Softmax برای
خانوادهٔ ناهمانندی مورد استفاده است؛
با نماد مقیاس $\tau_D>0$ در بخش توابع زیان،
$\gamma=1/\tau_D$ خواهد بود.
این پارامتر با پهنای $\tau$ در زیان سیگموئیدی
متفاوت است.
برای تمام $\gamma>0$، طبقه نزدیک‌ترین مرکز تغییر نمی‌کند،
ولی احتمال طبقه برنده و مقدار NLL تغییر می‌کند. برای نمونه دودویی با برچسب واقعی طبقه اول و فاصله‌های $(1,2)$، خطا برابر $\log(1+e^{-\gamma})$ است؛ این مقدار با افزایش $\gamma$ به صفر میل می‌کند، بی‌آنکه تصمیم نمونه تغییر کرده باشد. در مقابل، خطای یک نمونه اشتباه طبقه‌بندی‌شده می‌تواند با افزایش مقیاس بزرگ‌تر شود. پس کاهش NLL به‌تنهایی تضمین‌کننده افزایش دقت نیست. این تحلیل درباره \emph{امتیازهای ورودی نرمال‌سازی با بیشینهٔ هموار} است و ادعا نمی‌کند هر سیاست انتخاب ویژگی بتواند تمام فاصله‌های KL را آزادانه در یک ضریب مشترک ضرب کند.

در شبکه‌های عصبی نیز افزایش نرم بردار امتیازهای پیش از نرمال‌سازی با بیشینهٔ هموار می‌تواند موجب افزایش اطمینان بدون بهبود دقت شود. نرمال‌سازی این امتیازها در روش LogitNorm پیشنهاد شده است \cite{Wei2022LogitNorm}؛ تنظیم دمای نرمال‌سازی با بیشینهٔ هموار نیز برای کالیبراسیون خروجی به کار رفته است \cite{Guo2017Calibration}. هیچ‌کدام را نباید با Batch Normalization فعال‌سازی‌های لایه‌های میانی یکی دانست؛ نرمال‌کردن امتیازهای خروجی پیش از اعمال نرمال‌سازی با بیشینهٔ هموار روش اجباری همه شبکه‌ها نیست.

\paragraph{ثابت نگه‌داشتن اندازه امتیازها.}
چون افزودن یک ثابت مشترک به امتیازها بی‌اثر است،
می‌توان آن‌ها را حول میانگین مرکز کرد. برای یک نمونه
با فاصله‌های متناهی، میانگین و انحراف معیار فاصله‌های
طبقات را ابتدا چنین تعریف می‌کنیم:
\begin{equation}\label{cmp:eq:ch3-distance-statistics}
\overline d_i=\frac1K\sum_{c=1}^{K}d_{i,c},
\qquad
\sigma_i=\sqrt{\frac1K\sum_{c=1}^{K}
(d_{i,c}-\overline d_i)^2}.
\end{equation}
اگر $v_{i,c}=\overline d_i-d_{i,c}$ باشد،
آنگاه $\sum_cv_{i,c}=0$ و مربع نرم اقلیدسی
بردار امتیاز مرکز‌شده با واریانس فاصله‌ها رابطهٔ دقیق زیر را دارد:
\begin{equation}\label{cmp:eq:ch3-NLL-centered-L2}
\|v_i\|_2^2=\sum_{c=1}^{K}(d_{i,c}-\overline d_i)^2
=K\,\sigma_i^2.
\end{equation}
پس ثابت‌کردن نرم ۲ بردار مرکز‌شده، معادل ثابت‌کردن واریانس مؤلفه‌های آن است؛ به‌طور مشخص، تقسیم بر $\sigma_i$ نرم بردار را برابر $\sqrt K$ می‌کند. این انتخاب، اثر ضرب مشترک امتیازها را حذف می‌کند.

\paragraph{گزینه مثبت‌سازی و نرم یک ثابت.}
راه دیگر آن است که ابتدا از بردار فاصله‌ها یک بردار امتیاز نامنفی بسازیم و سپس مجموع مؤلفه‌های آن را ثابت نگه داریم. برای نمونه، اگر همه فاصله‌ها متناهی باشند و حداقل دو فاصله متفاوت وجود داشته باشد، می‌توان نوشت:
\begin{equation}\label{cmp:eq:ch3-NLL-positive-L1-example}
u_{i,c}=\max_{r}d_{i,r}-d_{i,c},
\qquad
\widetilde u_{i,c}=\frac{u_{i,c}}{\sum_{r=1}^{K}u_{i,r}},
\qquad \|\widetilde u_i\|_1=1.
\end{equation}
در این نمونه از مثبت‌سازی، ترتیب طبقات حفظ می‌شود و $\widetilde u_i$ نسبت به تغییر $d'_i=a_i d_i+b_i\mathbf1$ با $a_i>0$ ناورداست. بنابراین می‌توان مقیاس را با نرم یک ثابت کنترل کرد و نیازی به تحمیل قید جداگانه بر نرم دو نیست. این روش با نرمال‌سازی مرکز‌شده مبتنی بر انحراف معیار یکسان نیست؛ در بخش~\ref{cmp:subsec:ch3-normalized-NLL} صورت‌بندی مبتنی بر انحراف معیار را بررسی می‌کنیم، و به‌کارگیری صورت مثبت و نرم یک به تحلیل متناظر خود نیاز دارد. اگر همه فاصله‌ها برابر باشند، مخرج رابطه بالا صفر می‌شود و باید حالت تساوی جداگانه مدیریت شود.

\paragraph{واگرایی بی‌نهایت و مسئله محاسباتی.}
واگرایی $D_{\mathrm{KL}}(\mathscr Y_i\|\mathcal V_c)$ زمانی بی‌نهایت است که مؤلفه‌ای از $\mathscr Y_i$ مثبت، اما مؤلفه متناظر مرکز $\mathcal V_c$ صفر باشد. در نرمال‌سازی با بیشینهٔ هموارِ رابطه~\eqref{cmp:eq:ch3-NLL-softmax}، چنین فاصله‌ای فقط وزن نمایی صفر به آن طبقه می‌دهد و \emph{تا وقتی دست‌کم یک فاصله متناهی باشد} نرمال‌سازی نمایی تعریف می‌شود. برای نمونه‌های آموزش، با مراکزی که دقیقاً میانگین توزیع‌های همان طبقه هستند، فاصله هر نمونه از مرکز طبقه خودش متناهی است. اما برای یک نمونه جدید ممکن است تمام فاصله‌ها بی‌نهایت شوند و در آن صورت حتی مخرج نرمال‌سازی با بیشینهٔ هموار نیز صفر خواهد بود.

برخلاف محاسبهٔ احتمال با نرمال‌سازی با بیشینهٔ هموار در حضور دست‌کم یک امتیاز متناهی، مرکزکردن و محاسبه واریانس یا مثبت‌سازی بردار فاصله‌هایی که مؤلفه بی‌نهایت دارند می‌تواند به عبارت‌های تعریف‌نشده مانند $\infty-\infty$ بینجامد. بنابراین نرمال‌سازی مستقیم فاصله‌های KL به یک قرارداد محاسباتی روشن نیاز دارد؛ مثلاً می‌توان مراکز را به‌صورت سازگار هموارسازی کرد تا مؤلفه‌های آن‌ها مثبت شوند. چنین تغییری ممکن است خواص روابط مبتنی بر میانگین دقیق مراکز را تحت تأثیر قرار دهد و باید در تحلیل‌های بعدی مشخص شود. این موضوع یکی از محدودیت‌های عملی نرمال‌سازی امتیازهای KL است.



\section{تحلیل نظری نرمال‌سازی امتیازها و کران‌های آن}
\label{cmp:sec:ch3-score-normalization-theory}
این بخش صورت‌های نرمال‌سازی امتیازها،
کران‌های احتمال و زیان نرمال‌شده و
رفتار حالت‌های مرزی را تحلیل می‌کند.
این جزئیات برای فهم اصل گشت،
بازنمایی و قاعدهٔ پیش‌بینی ضروری نیستند.
خواننده‌ای که تنها سازوکار اجرایی اصلی
را دنبال می‌کند، می‌تواند از پایان بخش
مقیاس امتیازها به
بخش~\ref{cmp:sec:ch3-policy-learning-algorithm}
برود؛ بخش~\ref{cmp:sec:ch3-loss-quality-theory}
نیز ادامهٔ تحلیل نظریِ اختیاری است.

\subsection{نرمال‌سازی امتیازها و پیامدهای آن}
\label{cmp:subsec:ch3-normalized-NLL}

در اینجا نرمال‌سازی مبتنی بر انحراف معیار
را تحلیل می‌کنیم؛ روش مثبت‌سازی و تثبیت نرم یک،
انتخابی متمایز است. حفظ ترتیب برچسب‌ها
در این دو روش، به معنای یکسان‌بودن NLL
یا گرادیان‌های آن‌ها نیست.

در این زیربخش، برای تحلیل صورت نرمال‌شده مبتنی بر انحراف معیار، متناهی بودن تمام فاصله‌های مورد استفاده را فرض می‌کنیم. این فرض برای تعریف میانگین و واریانس فاصله‌ها ضروری است؛ حالت‌های دارای KL بی‌نهایت در بخش~\ref{cmp:subsec:ch3-NLL-scaling} بررسی شدند. روش مثبت‌سازی و نرم یک ثابت نیز در همان بخش به‌عنوان یک انتخاب متفاوت معرفی شد؛ کران‌های این بخش مخصوص نرمال‌سازی مبتنی بر انحراف معیار هستند.

برای حذف اثر ضرب مشترک فاصله‌های یک نمونه
در یک ضریب مثبت، به‌جای مقدار خام فاصله‌ها
از موقعیت نسبی آن‌ها استفاده می‌کنیم.
میانگین و انحراف معیار در
رابطه~\eqref{cmp:eq:ch3-distance-statistics}
تعریف شده‌اند.
برای $\sigma_i>0$، امتیاز نرمال‌شده طبقه $c$ برابر است با:
\begin{equation}\label{cmp:eq:ch3-normalized-score}
\boxed{z_{i,c}=\frac{\overline d_i-d_{i,c}}{\sigma_i}.}
\end{equation}
در نتیجه $\sum_cz_{i,c}=0$ و $\sum_cz_{i,c}^2=K$. بنابراین بردار امتیازها میانگین صفر و نرم اقلیدسی ثابت $\sqrt K$ دارد؛ اما ترتیب طبقات تغییر نمی‌کند. با استفاده از نرمال‌سازی با بیشینهٔ هموار روی همین امتیازها، داریم:
\begin{equation}\label{cmp:eq:ch3-normalized-softmax}
\widehat p_{i,c}^{\,\mathrm{norm}}
=\frac{e^{z_{i,c}}}{\sum_{r=1}^{K}e^{z_{i,r}}},
\qquad
\widehat y_i=\arg\max_c z_{i,c}=\arg\min_c d_{i,c}.
\end{equation}
خطای نمونه و NLL نرمال‌شده وزن‌دار را مشابه تعریف پیشین به‌صورت زیر در نظر می‌گیریم:
\begin{equation}\label{cmp:eq:ch3-normalized-NLL}
\boxed{
\begin{aligned}
\ell_i^{\mathrm{norm}}
&=-\log\widehat p_{i,y_i}^{\,\mathrm{norm}}
=\log\left(1+\sum_{r\ne y_i}
e^{(d_{i,y_i}-d_{i,r})/\sigma_i}\right),\\
NLL_{\mathrm{norm},\Omega}
&=\sum_{c=1}^{K}\frac{\omega_c}{|C_c|}
\sum_{i\in C_c}\ell_i^{\mathrm{norm}}.
\end{aligned}}
\end{equation}
در ادامه، هنگام ارتباط دادن این خطا با معیارهای تفکیک، حالت وزن یکنواخت $\omega_c=1/K$ را با $NLL_{\mathrm{norm},U}$ نشان می‌دهیم.
در حالت $\sigma_i=0$ همه فاصله‌ها برابرند و هیچ طبقه بر اساس این بردار برتری ندارد؛ در این حالت امتیازها را صفر و احتمال تمام طبقات را $1/K$ در نظر می‌گیریم. این قرارداد از تقسیم بر صفر جلوگیری می‌کند و خطای نمونه را برابر $\log K$ قرار می‌دهد.

این تعریف در نقطه تساوی فاصله‌ها عموماً
امتداد پیوسته نگاشت نرمال‌سازی نیست:
حد امتیازها هنگام نزدیک‌شدن به این نقطه
می‌تواند به جهت نزدیک‌شدن بستگی داشته باشد.
ازاین‌رو، پایداری و مشتق‌پذیری زیان
در اطراف این نقاط به بررسی جداگانه نیاز دارد.

\paragraph{نرمال‌سازی برای خانواده‌های گوناگون ناهمانندی.}
برای هر سه خانوادهٔ KL، JSD و جذر JSD،
می‌توان فاصله‌ها را به‌طور مستقل نرمال کرد:
\begin{equation}\label{cmp:eq:ch3-normalized-distance-family}
\boxed{
z_{i,c}^{D}
=\frac{\overline d_i^{D}-d_{i,c}^{D}}
{\sigma_i^D},
\quad
(\sigma_i^D)^2
=\frac1K\sum_r(d_{i,r}^{D}-\overline d_i^D)^2,
\quad
\overline d_i^D=\frac1K\sum_r d_{i,r}^D.
}
\end{equation}
در این رابطه $\sigma_i^D>0$ فرض شده است.
اگر فاصله‌ها برابر باشند، قرارداد امتیاز صفر
اعمال می‌شود. نرمال‌سازی نسبت به تغییر مقیاس
مشترک فاصله‌های همان نمونه ناوردا و تصمیم
نزدیک‌ترین مرکز را حفظ می‌کند، اما اطلاعات
\emph{اندازه حاشیه} را تا حدی حذف می‌کند.
به‌ویژه برای $K=2$، بردار امتیازهای
نرمال‌شده برای هر سه معیار، در حالت بدون تساوی،
صرفاً یکی از دو بردار $(1,-1)$ و $(-1,1)$
است. این ویژگی ممکن است به کاهش حساسیت
گرادیان در مسئله دودویی منجر شود.





\paragraph{کران تیز حاشیه امتیاز، سقف احتمال و کف NLL.}
ثابت‌بودن نرم امتیازهای مرکز‌شده
تنها اثر مقیاس را حذف نمی‌کند؛
بلکه \emph{بیشترین حاشیه امتیازی}
و بنابراین بزرگ‌ترین احتمال قابل
تولید توسط نرمال‌سازی با بیشینهٔ هموار را نیز محدود می‌کند.
برای $K\ge2$ و هر نمونه با $\sigma_i>0$،
از $\sum_c z_{i,c}=0$ و
$\|z_i\|_2^2=K$ نتیجه می‌شود
که اختلاف امتیاز هر دو طبقه و
حاشیه امتیازی طبقه نامزد $c$
از کران‌های زیر پیروی می‌کنند:
\begin{equation}\label{cmp:eq:ch3-normalized-score-margin-ceilings}
\boxed{
|z_{i,c}-z_{i,r}|\le\sqrt{2K},
\qquad
z_{i,c}-\max_{r\ne c}z_{i,r}
\le\frac{K}{\sqrt{K-1}}.
}
\end{equation}
کران دوم می‌گوید حتی اگر طبقه $c$
برنده باشد، فاصله امتیاز آن از
\emph{نزدیک‌ترین رقیب} نمی‌تواند
از مقدار مشخصی بیشتر شود.
نتیجه قوی‌تر برای احتمال خروجی،
کران تیز و قابل‌دستیابی زیر روی
تمام بردارهای امتیاز با این قیود است:
\begin{equation}\label{cmp:eq:ch3-normalized-softmax-pmax}
\boxed{
\max_{1\le c\le K}\widehat p_{i,c}^{\,\mathrm{norm}}
\le p_K^\star
:=\frac{1}
{1+(K-1)\exp(-K/\sqrt{K-1})}
<1.
}
\end{equation}
در نتیجه، برای برچسب واقعی نمونه،
زیان نرمال‌شده و میانگین وزن‌دار آن
دارای \emph{کران پایین مثبت} هستند:
\begin{equation}\label{cmp:eq:ch3-normalized-NLL-positive-floor}
\boxed{
\ell_i^{\mathrm{norm}}\ge L_K,
\qquad
NLL_{\mathrm{norm},\Omega}\ge L_K,
\qquad
L_K:=-\log p_K^\star
=\log\bigl(1+(K-1)e^{-K/\sqrt{K-1}}\bigr)>0.
}
\end{equation}
برابری کران احتمال در فضای
بردارهای امتیاز نرمال‌شده
دقیقاً هنگامی حاصل می‌شود که
برای یک طبقه $c$ داشته باشیم
$z_{i,c}=\sqrt{K-1}$ و
$z_{i,r}=-1/\sqrt{K-1}$
برای همه $r\ne c$.
این تیزی درباره \emph{فضای امتیازهای مجاز}
است و ادعا نمی‌کند هر سیاست
گام‌برداری الزاماً چنین برداری می‌سازد.
برای $K=2$، کف زیان
$L_2=\log(1+e^{-2})$ است.
ارتباط دقیق این کف با خطای
متوازن دودویی در
بخش~\ref{cmp:subsec:ch3-margin-normalized-bound}
صورت‌بندی خواهد شد.
به این ترتیب سقف احتمال و کف NLL
با \emph{کران بالای NLL برای نمونه
درست‌طبقه‌بندی‌شده} متفاوت‌اند:
اولی محدودیت دامنه اطمینان
پس از نرمال‌سازی است و دومی
برای مطالعه خطای تصمیم به کار می‌رود.
اثبات تیزبودن، شرایط برابری و
تعمیم به نرم کوچک‌تر در
پیوست~\ref{app:ch3-normalized-softmax-ceiling}
آمده است.

\paragraph{ناوردایی نسبت به تغییر مقیاس.}
فرض کنیم برای یک نمونه تمام فاصله‌ها طبق $d'_{i,c}=a_i d_{i,c}+b_i$ تغییر کنند، که در آن $a_i>0$ و $b_i$ برای همه طبقات همان نمونه یکسان است. در این صورت $\overline d'_i=a_i\overline d_i+b_i$ و $\sigma'_i=a_i\sigma_i$؛ پس $z'_{i,c}=z_{i,c}$. بنابراین احتمال‌های نرمال‌شده و NLL نرمال‌شده در برابر این تبدیل ثابت می‌مانند. این خاصیت وابستگی تابع هدف به ضریب مقیاس مشترک فاصله‌های نمونه را حذف می‌کند.

بااین‌حال، نرمال‌سازی تضمین نمی‌کند که هر کاهش $NLL_{\mathrm{norm},U}$ به افزایش دقت بینجامد؛ زیرا شکل نسبی بردار فاصله‌ها هنوز می‌تواند بدون عوض شدن طبقه برنده تغییر کند. همچنین برای $K=2$ و $\sigma_i>0$، بردار امتیازهای نرمال‌شده فقط دو آرایش $(1,-1)$ و $(-1,1)$ دارد. بنابراین مقدار خطای هر نمونه دودویی صرفاً به درست یا نادرست بودن تصمیم وابسته است و بزرگی حاشیه فاصله را نشان نمی‌دهد. این محدودیت، انگیزه‌ای برای بررسی مستقیم حاشیه‌های تفکیک و پراکندگی آن‌ها فراهم می‌کند.

به‌طور مشخص برای $K=2$ و دو فاصله نابرابر،
امتیازهای نرمال‌شده در هر ناحیه با ترتیب ثابت
فاصله‌ها، ثابت‌اند؛ بنابراین NLL نمونه
برای تصمیم درست و نادرست به‌ترتیب
$\log(1+e^{-2})$ و $\log(1+e^2)$ است
و گرادیان معمول آن نسبت به فاصله‌ها
در درون این نواحی صفر خواهد بود.
این محدودیتِ یادگیری از حاشیه مستقل
از مشکل KL نامتناهی است.






\subsection{صورت پایدار KL و نرمال‌سازی امتیازها در حالت‌های مرزی}
\label{cmp:subsec:ch3-stable-KL-normalization}
دو دشواری پیشین منشأ یکسانی ندارند: صفرشدن
انحراف معیار زمانی رخ می‌دهد که فاصله نمونه
از تمام مراکز برابر باشد؛ در مقابل،
نامتناهی‌شدن KL ناشی از صفر بودن
احتمال مرجع در مؤلفه‌ای با جرم مثبت است.
برای ارائه صورت محاسباتی تعریف‌شده و
مشتق‌پذیر در هر دو وضعیت، ابتدا
هموارسازی توزیع‌ها و سپس یک نرمال‌سازی
با مخرج منظم‌شده را معرفی می‌کنیم.
این صورت \emph{مکمل} تعریف دقیق قبلی است
و نتایج آن جایگزین خودکار کران‌های
صورت قبلی نیستند.

\paragraph{هموارسازی سازگار نمونه‌ها و مراکز.}
بردار یکنواخت $U=(1/K,\ldots,1/K)$ و
پارامتر ثابت $0<\eta<1$ را انتخاب می‌کنیم.
برای هر نمونه و هر طبقه تعریف می‌کنیم:
\begin{equation}\label{cmp:eq:ch3-stable-smoothed-representations}
\boxed{
\mathscr Y_i^{(\eta)}=(1-\eta)\mathscr Y_i+\eta U,
\qquad
\mathcal V_c^{(\eta)}=(1-\eta)\mathcal V_c+\eta U
=\frac1{|C_c|}\sum_{i\in C_c}\mathscr Y_i^{(\eta)}.
}
\end{equation}
بنابراین مراکز هموارشده همچنان
\emph{میانگین دقیق} بازنمایی‌های هموارشده
اعضای همان طبقه هستند و همه مؤلفه‌های
آن‌ها حداقل $\eta/K$ است. با تعریف
\begin{equation}\label{cmp:eq:ch3-stable-KL-distance}
d_{i,c}^{(\eta)}
=D_{\mathrm{KL}}\!\left(
\mathscr Y_i^{(\eta)}\middle\|\mathcal V_c^{(\eta)}
\right),
\qquad
0\le d_{i,c}^{(\eta)}\le\log\frac K\eta<\infty,
\end{equation}
هیچ فاصله‌ای نامتناهی نمی‌شود؛
از سوی دیگر، مثبت‌شدن تمام مؤلفه‌ها
مشتق‌گیری نسبت به بازنمایی‌ها را
در مدل هموارشده ممکن می‌کند.
هموارسازی، خودِ بازنمایی و فاصله‌ها
را تغییر می‌دهد و اندازه $\eta$
باید در آزمایش‌ها ثابت و گزارش شود.
در حد $\eta\downarrow0$، هر فاصله
به مقدار KL اولیه، با امکان واگرایی
به بی‌نهایت، میل می‌کند.

\paragraph{منظم‌سازی مخرج واریانس.}
حال برای فاصله‌های متناهی فوق،
$\overline d_i^{(\eta)}$ و
$(\sigma_i^{(\eta)})^2$
را مانند رابطه~\eqref{cmp:eq:ch3-distance-statistics}
تعریف می‌کنیم. برای یک پارامتر ثابت
$\varepsilon>0$، امتیاز پایدار را
به‌صورت زیر قرار می‌دهیم:
\begin{equation}\label{cmp:eq:ch3-stable-normalized-score}
\boxed{
z_{i,c}^{(\eta,\varepsilon)}
=\frac{\overline d_i^{(\eta)}-d_{i,c}^{(\eta)}}
{\sqrt{(\sigma_i^{(\eta)})^2+\varepsilon^2}},
\qquad
\widehat p_{i,c}^{(\eta,\varepsilon)}
=\frac{\exp z_{i,c}^{(\eta,\varepsilon)}}
{\sum_r\exp z_{i,r}^{(\eta,\varepsilon)}}.
}
\end{equation}
اگر واریانس صفر باشد، تمام فاصله‌ها
برابرند، همه امتیازها دقیقاً صفر می‌شوند
و توزیع خروجی یکنواخت است؛
اما برخلاف قرارداد نقطه‌ای
در صورت قبلی، تابع امتیازها
در اطراف این وضعیت نیز پیوسته و
مشتق‌پذیر است. برای هر $\varepsilon>0$
و هر نمونه داریم:
\[
\sum_cz_{i,c}^{(\eta,\varepsilon)}=0,
\qquad
\sum_c\bigl(z_{i,c}^{(\eta,\varepsilon)}\bigr)^2
=K\,\frac{(\sigma_i^{(\eta)})^2}
{(\sigma_i^{(\eta)})^2+\varepsilon^2}\le K.
\]
ترتیب فاصله‌ها و در نتیجه
قاعده نزدیک‌ترین مرکز تغییر نمی‌کند.
هرگاه $\sigma_i^{(\eta)}>0$،
در حد $\varepsilon\downarrow0$،
امتیازهای نرمال‌شده دقیق قبلی
بازیابی می‌شوند. بااین‌حال،
برای $\varepsilon$ ثابت و مثبت،
\emph{ناوردایی دقیق در برابر ضرب همه
فاصله‌ها در یک ضریب مثبت}
دیگر برقرار نیست؛ این محدودیت
بهای رفع ناپیوستگی در واریانس صفر است.

در واقع نمی‌توان برای یک نگاشت امتیاز
\emph{غیرثابت}، هم ناوردایی دقیق نسبت
به تمام ضرایب مثبت مقیاس و هم
پیوستگی در بردار فاصله‌های برابر را
داشت؛ برهان این ناسازگاری در پیوست آمده است.
پارامترهای $\eta$ و $\varepsilon$
باید پیش از ارزیابی نهایی تعیین شوند؛
همچنین برای خانواده‌های فاصله متفاوت،
مقیاس مناسب $\varepsilon$ ممکن است
یکسان نباشد.


\paragraph{حفظ حساسیت حاشیه در حالت دودویی.}
برای $K=2$ و حاشیه خام
$m=d_{i,r}^{(\eta)}-d_{i,c}^{(\eta)}$
در مقایسه طبقه درست $c$ با رقیب $r$،
شکاف امتیازهای پایدار برابر است با:
\begin{equation}\label{cmp:eq:ch3-stable-binary-margin}
g_\varepsilon(m)
=\frac{m}{\sqrt{m^2/4+\varepsilon^2}},
\qquad
g_\varepsilon'(m)
=\frac{\varepsilon^2}
{(m^2/4+\varepsilon^2)^{3/2}}>0.
\end{equation}
ازاین‌رو
$\ell_i^{(\eta,\varepsilon)}
=\log(1+e^{-g_\varepsilon(m)})$
در تمام مقادیر متناهی حاشیه مشتق‌پذیر
است و با افزایش حاشیه کاهش می‌یابد.
در مقابل، نرمال‌سازی دقیق با
$\varepsilon=0$ در هر نیم‌خط
$m>0$ یا $m<0$ امتیاز ثابتی تولید
می‌کند. با وجود این، گرادیانِ
صورت پایدار برای حاشیه‌های
بسیار بزرگ ممکن است کوچک شود
و انتخاب $\varepsilon$ بر
حساسیت به نواحی نزدیک مرز اثر دارد.

\paragraph{حد کران‌های قبلی و نسخه پایدار.}
چون نرم امتیاز پایدار حداکثر $\sqrt K$
است، زیان تک‌نمونه‌ای نیز
\emph{از بالا کراندار} می‌ماند.
صورت عددی کران عمومی و اثبات آن
در بخش~\ref{cmp:subsec:ch3-margin-normalized-bound}
خواهد آمد؛ این کران برای
نسخهٔ پایدار نیز معتبر است.
اما کران قوی‌ترِ مخصوص نمونهٔ
درست‌طبقه‌بندی‌شده در حالت
نرم دقیق $\sqrt K$ را
نباید بدون اثبات به صورت پایدار منتقل کرد.
برای هر نمونه درست طبقه‌بندی‌شده،
احتمال طبقه درست دست‌کم $1/K$
و در نتیجه زیان حداکثر $\log K$ است.
پس با تعریف خطای تجربی وزن‌دار
$E_\Omega=\sum_c(\omega_c/|C_c|)
\sum_{i\in C_c}\mathbf1_{\{\widehat y_i\ne c\}}$
داریم:
\begin{equation}\label{cmp:eq:ch3-stable-normalized-error-bound}
\boxed{
NLL_{\mathrm{norm},\Omega}^{(\eta,\varepsilon)}
\le\log K+(B_K-\log K)E_\Omega,
\quad B_K=\log(1+(K-1)e^{\sqrt{2K}}).
}
\end{equation}
این کران از نسخه با نرم ثابت ضعیف‌تر است،
اما حالت صفرشدن واریانس را نیز
بدون استثنا پوشش می‌دهد.
در سیاست شرطی طبقه می‌توان همین
هموارسازی و منظم‌سازی را بر
امتیازهای تمام نامزدها اعمال کرد،
اما تعمیم کران‌های اطلاعاتیِ مبتنی بر
جدایی مراکز به این مدل نیازمند اثبات جداگانه است.
استدلال‌های دقیق این قسمت
در پیوست~\ref{app:ch3-stable-KL-normalization}
آمده‌اند.


\section{روابط نظری توابع زیان با کیفیت بازنمایی و تصمیم}
\label{cmp:sec:ch3-loss-quality-theory}
این بخش نیز تحلیلی تکمیلی است و برای اجرای
الگوریتم اصلی آموزش ضرورت ندارد؛
خواننده‌ای که صرفاً چرخهٔ یادگیری سیاست
را دنبال می‌کند، می‌تواند مستقیماً به
بخش~\ref{cmp:sec:ch3-policy-learning-algorithm}
مراجعه کند.
پس از تعریف توابع زیان و قراردادهای نرمال‌سازی،
می‌توان روابط آن‌ها را با معیارهای تراکم،
جدایی و خطای تصمیم بررسی کرد.
در این بخش، ابتدا کران‌های ویژهٔ KL و سپس
کران‌های مربوط به خطای طبقه‌بندی و حاشیه را
با فرض‌های معتبر هر یک تحلیل می‌کنیم؛
این روابط به معنای تضمین بهبود گام‌به‌گام
دقت در فرایند یادگیری نیستند.

\subsection{پیوند معیارهای تراکم و جدایی با NLL}
\label{cmp:subsec:ch3-density-separation-NLL-bound}
از تعریف شعاع درون‌طبقه‌ای~\eqref{cmp:eq:ch3-R-average} و تجزیه رابطه~\eqref{cmp:eq:ch3-NLL-sample} داریم:
\begin{equation}\label{cmp:eq:ch3-NLL-R}
NLL_\Omega=R+
\sum_{c=1}^{K}\frac{\omega_c}{|C_c|}\sum_{i\in C_c}
\log\left(\sum_{r=1}^{K}
e^{-D_{\mathrm{KL}}(\mathscr Y_i\|\mathcal V_r)}\right).
\end{equation}
در این تساوی، $R$ نیز با \emph{همان وزن‌های} $\Omega$ محاسبه می‌شود. پس NLL علاوه بر فاصله نمونه از مرکز طبقه واقعی، به فاصله آن از همه مراکز رقیب وابسته است.

\paragraph{کران بالا در حالت وزن‌های یکنواخت.}
برای توزیع احتمال ثابت $U$، تابع $f_U(V)=e^{-D_{\mathrm{KL}}(U\|V)}$ نسبت به $V$ مقعر است. بنابراین نامساوی ینسن، با میانگین یکنواخت مراکز $\overline{\mathcal V}=\frac1K\sum_c\mathcal V_c$، نتیجه می‌دهد:
\begin{equation}\label{cmp:eq:ch3-NLL-Jensen}
\sum_{r=1}^{K}e^{-D_{\mathrm{KL}}(\mathscr Y_i\|\mathcal V_r)}
\le K e^{-D_{\mathrm{KL}}(\mathscr Y_i\|\overline{\mathcal V})}.
\end{equation}
با گرفتن لگاریتم و جایگذاری در رابطه~\eqref{cmp:eq:ch3-NLL-R} برای $\omega_c=1/K$، خواهیم داشت:
\begin{equation}\label{cmp:eq:ch3-NLL-pre-I}
NLL_U\le R+\log K-\frac1K\sum_{c=1}^{K}
\frac1{|C_c|}\sum_{i\in C_c}
D_{\mathrm{KL}}(\mathscr Y_i\|\overline{\mathcal V}).
\end{equation}
طبق تجزیه~\eqref{cmp:eq:ch3-R-I-decomposition}، میانگین واگرایی عبارت آخر برابر $R+I(\mathcal V)$ است. در نتیجه، $R$ حذف می‌شود و کران زیر به‌دست می‌آید:
\begin{equation}\label{cmp:eq:ch3-NLL-I-bound}
\boxed{0\le NLL_U\le\log K-I(\mathcal V),
\qquad \omega_c=\frac1K.}
\end{equation}
این کران برای وزن‌های یکنواخت، مستقل از شعاع درون‌طبقه‌ای است. نکته مهم آن است که \emph{همین نرمال‌سازی با بیشینهٔ هموار بدون وزن طبقه‌ای} امکان تعمیم به وزن‌های دلخواه $\Omega$ را نیز دارد. برای این منظور میانگین یکنواخت مراکز را با $\overline{\mathcal V}_U=\frac1K\sum_{r=1}^{K}\mathcal V_r$ و میانگین وزن‌دار آن‌ها را مطابق تعریف قبلی با $\overline{\mathcal V}=\sum_c\omega_c\mathcal V_c$ نشان می‌دهیم. با اعمال همان نامساوی ینسن و تجزیه KL، نتیجه دقیق‌تر زیر حاصل می‌شود:
\begin{equation}\label{cmp:eq:ch3-NLL-weighted-I-bound}
\boxed{
\begin{aligned}
0\le NLL_\Omega
&\le\log K-I(\mathcal V)
-D_{\mathrm{KL}}(\overline{\mathcal V}\|\overline{\mathcal V}_U)\\
&\le\log K-I(\mathcal V).
\end{aligned}}
\end{equation}
در این رابطه $I(\mathcal V)$ با همان وزن‌های $\Omega$ محاسبه می‌شود، اما $\overline{\mathcal V}_U$ میانگین \emph{بدون وزن} مراکز در مخرج نرمال‌سازی با بیشینهٔ هموار است. اگر $\omega_c=1/K$ باشد، دو میانگین برابر می‌شوند و جمله KL اضافی صفر خواهد شد؛ بنابراین رابطه~\eqref{cmp:eq:ch3-NLL-I-bound} حالت خاصی از این کران است. اثبات تقعر و تعمیم وزن‌دار در بخش~\ref{app:ch3-NLL-concavity} از پیوست~\ref{app:ch3-proofs} آمده است.

این رابطه یک \textbf{هم‌سویی نظری میان دو هدف} را نشان می‌دهد:
افزایش تمایز اطلاعاتی مراکز با معیار $I$، سقف قابل‌اثبات
برای NLL مبتنی بر KL را کاهش می‌دهد. به بیان دیگر،
بیشینه‌سازی $I$ و کوچک‌کردن \emph{این کران بالای} NLL
در یک جهت مطلوب حرکت می‌کنند. اهمیت حذف $R$ نیز
آن است که این پیوند بدون افزودن فرضِ فشردگی درون‌طبقه‌ای
به‌دست می‌آید. بااین‌حال، از کاهش کران نمی‌توان
کاهش یکنواخت \emph{مقدار واقعی} NLL در همه مراحل آموزش
یا افزایش قطعی دقت طبقه‌بندی را نتیجه گرفت.
این استدلال برای بازنمایی \emph{مشترک} نمونه‌ها،
فاصله KL، مراکز میانگین و نرمال‌سازی با بیشینهٔ هموار با امتیاز
$-d_{i,c}$ است؛ تعمیم آن به بازنمایی شرطی طبقه،
امتیازهای نرمال‌شده یا فاصله JS به استدلال جداگانه نیاز دارد.

\paragraph{کران متناظر با بازنمایی‌های هموارشده.}
برای بازنمایی‌ها و مراکز هموارشدهٔ
رابطه~\eqref{cmp:eq:ch3-stable-smoothed-representations}،
از آنجا که مراکز همچنان میانگین بازنمایی‌ها هستند، برهان کران NLL
مبتنی بر تقعر نمایی منفی KL
بدون تغییر ساختاری روی این خانواده
هم برقرار است. اگر
$I^{(\eta)}(\mathcal V^{(\eta)})$
همان معیار اطلاعات متقابل با
مراکز هموارشده و وزن‌های $\Omega$ باشد و
$\overline{\mathcal V}^{(\eta)}_{\Omega}$
و $\overline{\mathcal V}^{(\eta)}_U$
به‌ترتیب میانگین وزن‌دار و یکنواخت
همین مراکز باشند، برای
\emph{نرمال‌سازی با بیشینهٔ هموار روی امتیازهای نرمال‌نشده}
با امتیاز $-d_{i,c}^{(\eta)}$ داریم:
\begin{equation}\label{cmp:eq:ch3-stable-NLL-I-bound}
\boxed{
0\le NLL_\Omega^{(\eta)}
\le \log K-I^{(\eta)}(\mathcal V^{(\eta)})
-D_{\mathrm{KL}}\!\left(
\overline{\mathcal V}^{(\eta)}_{\Omega}
\middle\|\overline{\mathcal V}^{(\eta)}_U
\right)
\le\log K-I^{(\eta)}(\mathcal V^{(\eta)}).
}
\end{equation}
این کران درباره کمیت‌های \emph{هموارشده}
است و نباید $I^{(\eta)}$ را با $I$
اولیه یکی دانست.

\subsection{کران‌های خطای تصمیم بر حسب NLL نرمال‌شده}
\label{cmp:subsec:ch3-margin-normalized-bound}
تحلیل حاشیه KL و کران خطای طبقه‌بندی در زیربخش‌های قبل مستقل از NLL به‌دست آمد. برای مقایسه تکمیلی با معیار احتمالاتی بخش~\ref{cmp:sec:ch3-NLL}، اکنون حالت $\sigma_i>0$ را برای تمام نمونه‌ها فرض می‌کنیم. طبق تعریف~\eqref{cmp:eq:ch3-normalized-score}، بردار امتیازهای نرمال‌شده روی مجموعه $\sum_cz_{i,c}=0$ و $\sum_cz_{i,c}^2=K$ قرار دارد. همین قیدها اجازه می‌دهند برای مقدار NLL هر نمونه کران‌های مستقل از مقیاس فاصله‌ها به‌دست آوریم.

اگر نمونه درست طبقه‌بندی شده باشد، از بیشینه بودن امتیاز طبقه واقعی نتیجه می‌شود:
\begin{equation}\label{cmp:eq:ch3-normalized-correct-bound}
\ell_i^{\mathrm{norm}}\le
C_K=\log\left(K-1+e^{-K/\sqrt{K-1}}\right).
\end{equation}
برای نمونه دلخواه نیز، با توجه به کرانداری اختلاف دو مؤلفه از بردار نرمال‌شده، کران عمومی زیر برقرار است:
\begin{equation}\label{cmp:eq:ch3-normalized-global-bound}
\ell_i^{\mathrm{norm}}\le
B_K=\log\left(1+(K-1)e^{\sqrt{2K}}\right).
\end{equation}
در این روابط $K\ge2$ و $C_K<B_K$ است. بنابراین خطای هر نمونه حداکثر $C_K$ است، مگر آنکه نمونه اشتباه طبقه‌بندی شده باشد؛ در آن صورت کران $B_K$ را در نظر می‌گیریم. با میانگین‌گیری متوازن و استفاده از رابطه~\eqref{cmp:eq:ch3-margin-error-bound} خواهیم داشت:
\begin{equation}\label{cmp:eq:ch3-normalized-NLL-margin-bound}
\boxed{
NLL_{\mathrm{norm},U}
\le C_K+(B_K-C_K)E_{\mathrm{bal}}
\le C_K+(B_K-C_K)\min\{1,J_{\mathrm{margin}}\}.}
\end{equation}
اثبات کران‌ها در بخش~\ref{app:ch3-normalized-NLL-bound} از پیوست~\ref{app:ch3-proofs} آمده است. کران عمومی $B_K$ برای سادگی انتخاب شده و ادعا نمی‌کنیم که برای هر $K$ تیز است.


\paragraph{کران مکمل خطای طبقه‌بندی بر اساس NLL.}
در بخش~\ref{cmp:subsec:ch3-normalized-NLL}
نشان دادیم که حتی برای یک
بردار امتیاز نرمال‌شده با تصمیم درست،
زیان نمونه نمی‌تواند از کف مثبت
$L_K$ در
رابطه~\eqref{cmp:eq:ch3-normalized-NLL-positive-floor}
کمتر شود. از سوی دیگر،
اگر برچسب پیش‌بینی‌شده اشتباه باشد،
دست‌کم یک رقیب امتیازی برابر یا
بزرگ‌تر از طبقه واقعی دارد؛
در نتیجه احتمال طبقه واقعی حداکثر
$1/2$ و زیان آن دست‌کم $\log2$ است.
بنابراین برای همه نمونه‌ها:
\begin{equation}\label{cmp:eq:ch3-normalized-NLL-error-lower-bound}
\boxed{
\ell_i^{\mathrm{norm}}
\ge L_K+
(\log2-L_K)\mathbf1_{\{\widehat y_i\ne y_i\}},
\qquad
E_{\mathrm{bal}}\le
\frac{NLL_{\mathrm{norm},U}-L_K}{\log2-L_K}.
}
\end{equation}
این رابطه و کران بالای پیشین
\emph{دو جهت مکمل} دارند:
کران کانتلی از آمار حاشیه
برای محدودکردن NLL استفاده می‌کند،
درحالی‌که رابطه اخیر از NLL
برای کنترل خطای تصمیم استفاده
می‌کند. هیچ‌کدام برابری یکنواخت
میان کمینه‌سازی NLL و
بیشینه‌سازی دقت را اثبات نمی‌کنند.
همین استدلال برای نرمال‌سازی
پایدار بخش~\ref{cmp:subsec:ch3-stable-KL-normalization}
نیز معتبر است، زیرا در آن
$\|z_i\|_2\le\sqrt K$ و
کران پایین $L_K$ حفظ می‌شود؛
ولی کران بالای مخصوص نرم
دقیق باید با نسخه پایدار
جایگزین شود. اثبات کران مکمل
در پیوست
\ref{app:ch3-normalized-NLL-bound}
آمده است.

این رابطه نشان می‌دهد کوچک‌تر شدن $J_{\mathrm{margin}}$، کران بالای قابل محاسبه‌ای برای $NLL_{\mathrm{norm},U}$ را کاهش می‌دهد. اما نه هر کاهش این کران الزاماً مقدار واقعی NLL را کم می‌کند و نه هر کاهش NLL لزوماً خطای طبقه‌بندی را کاهش می‌دهد. در حالت خاص دو طبقه، اگر همه $\sigma_i>0$ باشند، رابطه دقیق‌تری نیز داریم:
\begin{equation}\label{cmp:eq:ch3-normalized-NLL-binary}
\boxed{
K=2:\quad
NLL_{\mathrm{norm},U}
=\log(1+e^{-2})+2E_{\mathrm{bal}}.
}
\end{equation}
این برابری حاصل آن است که نرمال‌سازی دو فاصله متفاوت فقط علامت حاشیه را نگه می‌دارد و بزرگی آن را حذف می‌کند. ازاین‌رو نسخه نرمال‌شده NLL، به‌ویژه برای دو طبقه، به‌تنهایی هدف نرمی برای بهبود تدریجی حاشیه‌ها نیست. معیار حاشیه‌ای مکملی است که اطلاعات میانگین و پراکندگی جدایی را حفظ می‌کند.

این روابط نشان می‌دهند که معیارهای جدایی و پراکندگی
با خطای تصمیم بی‌ارتباط نیستند؛ برای نمونه،
کاهش $J_{\mathrm{margin}}$ می‌تواند کران بالای
خطای تجربی و کران بالای NLL نرمال‌شده را
کوچک‌تر کند. بااین‌حال، این کران‌ها ممکن است
شل باشند و کاهش آن‌ها، کاهش واقعی خطای سخت را
در هر مرحله تضمین نمی‌کند. ازاین‌رو،
حاشیه مبتنی بر JS را نیز مستقیماً بررسی می‌کنیم.



\subsection{ناهم‌راستایی NLL و دقت طبقه‌بندی}

بهبود معیارهای $I(\mathcal V)$ و $NLL$ الزاماً افزایش دقت طبقه‌بندی را نتیجه نمی‌دهد. توضیح این پدیده صرفاً ضرب مصنوعی فاصله‌ها در یک ضریب ثابت نیست؛ خودِ تغییر سیاست انتخاب ویژگی می‌تواند هم‌زمان موجب افزایش فاصله مراکز طبقات و افزایش پراکندگی توزیع‌های همسایگی درون هر طبقه شود. افزایش فاصله مراکز ممکن است نمونه‌های بیشتری را قابل تفکیک کند، اما افزایش شعاع درون‌طبقه‌ای می‌تواند برخی نمونه‌های نزدیک مرز تصمیم را به سمت مرکز رقیب ببرد. در همان زمان ممکن است فاصله نسبی برخی نمونه‌های دیگر از مراکز رقیب بسیار بیشتر شود.

برای دیدن این اثر، در حالت دودویی حاشیه نمونه را با $m_i=d_{i,r}-d_{i,y_i}$، که $r$ طبقه رقیب است، نشان می‌دهیم. در این صورت:
\begin{equation}\label{cmp:eq:ch3-NLL-binary-margin-contrast}
\ell_i=\log(1+e^{-m_i}),
\qquad
\widehat y_i=y_i\ \Longleftrightarrow\ m_i>0
\end{equation}
(با صرف‌نظر از حالت تساوی). بنابراین، اندازه حاشیه در مقدار $NLL$ اثر دارد، ولی صحت تصمیم فقط به علامت آن وابسته است. برای نمونه، در ده مشاهده دودویی با حاشیه اولیه $0.2$ برای همه نمونه‌ها، دقت $100\%$ و میانگین NLL حدود $0.598$ است؛ اگر در وضعیت بعدی نه نمونه حاشیه $5$ و یک نمونه حاشیه $-4$ داشته باشند، دقت به $90\%$ کاهش می‌یابد، اما میانگین NLL به حدود $0.408$ می‌رسد. این مثال در سطح حاشیه‌های تصمیم است و ادعا نمی‌کند که هر بردار دلخواه از این حاشیه‌ها الزاماً با یک سیاست گام‌برداری معین تولید می‌شود؛ بلکه سازوکار ناسازگاری دو معیار را نشان می‌دهد.

معیارهای مرکز‌محور بخش~\ref{cmp:sec:ch3-quality-criteria}، مانند $T$، $BT$، $TNL$، $BTNL$ و $I$، تنها از ماتریس مراکز $\mathcal V$ محاسبه می‌شوند؛ در نتیجه، تغییر پراکندگی نمونه‌ها در اطراف مراکز و عبور آن‌ها از مرزهای تصمیم را مستقیماً منعکس نمی‌کنند. معیار $I$ علاوه بر این، کران بالای $H(\Omega)$ دارد و نزدیک‌شدن آن به بیشینه به‌تنهایی تضمینی برای اصلاح نمونه‌های نزدیک مرز یا اشتباه‌طبقه‌بندی‌شده فراهم نمی‌کند. معیار شعاع $R$ پراکندگی درون‌طبقه‌ای را اندازه می‌گیرد، اما به‌تنهایی برتری فاصله از مرکز واقعی نسبت به مراکز رقیب را نمی‌سنجد. از سوی دیگر، حتی کران‌های اثبات‌شده $NLL$ بر حسب $I$ نیز کنترل یکنواخت خطای صفر--یک نیستند. پس این معیارها، با وجود ارزش نظری و اطلاعاتی،
به‌تنهایی هدف کافی برای یادگیری سیاستی
که مقصود اصلی آن افزایش دقت طبقه‌بندی است
محسوب نمی‌شوند.
در مجموع، مثال دودویی بالا روشن می‌کند
که در نگاشت توزیع برچسب همسایگی
به برچسب تخمینی، \emph{درستی تصمیم}
به علامت برتری مرکز درست نسبت به
نزدیک‌ترین رقیب وابسته است؛
درحالی‌که زیان NLL به \emph{اندازه
همان برتری} نیز حساس است.
این تمایز همچنین نشان می‌دهد چرا
حاشیهٔ تصمیم باید جدا از NLL پایش شود.

بنابراین کاهش NLL، حتی هنگام ثابت‌ماندن
روش تصمیم‌گیری، لزوماً معادل کاهش خطای سخت
نیست. برای ارزیابی نتیجهٔ یادگیری،
حاشیهٔ نمونه‌ها، تراکم و جدایی مراکز
باید در کنار دقت متوازن پایش شوند.


\paragraph{پیامد نظری برای انتخاب تابع هدف.}
هر یک از کمیت‌های $R$، $I$، حاشیه، NLL و
خطای صفر--یک جنبه‌ای متفاوت از مدل را می‌سنجد:
کران~\eqref{cmp:eq:ch3-NLL-I-bound} هم‌سویی
$I$ با \emph{سقف} NLL در حالت KL را نشان می‌دهد،
درحالی‌که کران‌های~\eqref{cmp:eq:ch3-margin-error-bound}
و~\eqref{cmp:eq:ch3-js-class-error-bound}
به پراکندگی نمونه‌ها و فاصلهٔ تصمیم وابسته‌اند.
هیچ‌یک کاهش یکنواخت خطای سخت در گام‌های
آموزش را تضمین نمی‌کند. پس انتخاب تابع هدف باید
با قاعدهٔ تصمیم سازگار باشد و با دقت متوازن،
حاشیه و معیارهای هندسی به‌طور مستقل ارزیابی شود؛
نتایج تجربی در فصل~\ref{chap:results} گزارش می‌شوند.






\section{آموزش سیاست انتخاب ویژگی و استنتاج برون‌نمونه‌ای}
\label{cmp:sec:ch3-policy-learning-algorithm}
پس از تعریف توابع هدف، باید روشن کنیم
سیاست چگونه آموزش می‌بیند و چگونه
برای نمونهٔ جدید به کار می‌رود.
پارامترهای سیاست، احتمال‌های انتخاب
ویژگی در گشت را تعیین می‌کنند؛
این احتمال‌ها بازنمایی‌های همسایگی
و در نتیجه مراکز طبقات و امتیازهای
پیش‌بینی را تغییر می‌دهند. بنابراین
در یادگیری مبتنی بر گرادیان، وابستگی
\emph{هر دو} جزءِ بازنمایی نمونه
و مرکز طبقه به سیاست باید لحاظ شود؛
نمی‌توان مراکز را در محاسبهٔ گرادیان
به‌طور ضمنی ثابت فرض کرد.
در این بخش، بدون ورود به مشتق‌های
مؤلفه‌ای JSD و تحلیل تکیه‌گاه صفر،
چرخهٔ آموزش و قرارداد استنتاج
و اعتبارسنجی را به شکل الگوریتم
شرح می‌دهیم. صورت‌بندی دقیق مشتق‌ها،
شرایط اعتبار آن‌ها و رابطهٔ گرادیان
با به‌روزرسانی بهینه‌ساز در
بخش~\ref{cmp:sec:ch3-policy-gradient-theory}
آمده است.

\subsection{شبه‌کد یکپارچهٔ آموزش سیاست}
\label{cmp:subsec:ch3-policy-training-pseudocode}
سه زیرروال
\ref{cmp:alg:ch3-capped-walk-weights}،
\ref{cmp:alg:ch3-advance-walk-state} و
\ref{cmp:alg:ch3-exact-walk-expectation}
به‌ترتیب ضرایب گشت، توزیع حالت و امید دقیق
بازنمایی را برای \emph{سیاست و گراف ثابت} محاسبه
می‌کنند. الگوریتم حاضر این محاسبهٔ پیشرو را
درون چرخهٔ یادگیری پارامترهای سیاست به کار می‌گیرد؛
جزئیات گشت در اینجا تکرار نمی‌شوند.
$\mathcal I_{\mathrm{tr}}$ اندیس‌های مجموعهٔ آموزش و
$C_c=\{i\in\mathcal I_{\mathrm{tr}}:y_i=c\}$ است؛
تمام طبقاتِ مورد استفاده در آموزش باید $|C_c|>0$
داشته باشند. دادهٔ مرجع، برچسب‌ها و ماسک‌های
ویژگی فقط از مجموعهٔ آموزش ساخته می‌شوند.

در این الگوریتم، فراخوانی کوتاه
$\operatorname{ExactWalk}_{\mathrm{tr}}(i,c;\theta)$
به‌معنای اجرای الگوریتم
\ref{cmp:alg:ch3-exact-walk-expectation}
برای مبدأ آموزشی $i$، فرضیهٔ طبقهٔ $c$،
گراف و کرنل مرجع آموزش، سیاست $\pi_\theta$
و پارامترهای ثابت $H,\beta,\lambda$ است.
مبدأ $i$ از مقصدهای تمام گام‌ها حذف می‌شود
و سیاست افقی فقط روی ویژگی‌های موجود و مجاز
همان وضعیت نرمال می‌شود.
$\mathscr Y_i^{[c]}$ بازنمایی این اجراست؛
برای خانواده‌های مستقل از فرضیهٔ طبقهٔ مبدأ
(\lr{shared} و \lr{node}) محاسبهٔ آن تنها یک‌بار
لازم است. برای خانواده‌های وابسته به فرضیه
(\lr{class} و \lr{node\_class}) هر نامزد $c$
جداگانه ارزیابی می‌شود. حتی در آموزش نیز
استفاده از $c=y_i$ برای ساخت \emph{تمام}
امتیازهای پیش‌بینی مجاز نیست؛ برچسب واقعیِ مبدأ برای تعیین
\emph{سیاست فرضیهٔ نامزد} به کار نمی‌رود؛
برچسب‌های آموزش برای تعیین مجموعه‌های $C_c$،
تجمیع برچسبِ مقصد در گشت‌ها و زیان نظارت‌شده
استفاده می‌شوند.

زیان Softplus--JS در بخش~\ref{cmp:subsec:ch3-softplus-js-objective}
تعریف شد و مشتق آن در بخش~\ref{cmp:subsec:ch3-softplus-js-policy-gradient}
آمده است. در آزمایش حساسیت نویز فصل چهارم،
$K=3$ و $\tau_{\mathrm{SP}}=0.005$ به‌کار رفته‌اند؛
این مقادیر، تعریف عمومی تابع هدف را محدود نمی‌کنند.

\begin{algorithm}[!htbp]
\caption{آموزش سیاست انتخاب ویژگی با گشت دقیق و مراکز وابسته به سیاست}
\label{cmp:alg:ch3-learn-feature-policy}
\begin{algorithmic}[1]
\REQUIRE داده، برچسب و ماسک مجموعهٔ آموزش
$D_{\mathrm{tr}},Y_{\mathrm{tr}},M_{\mathrm{tr}}$؛
خانوادهٔ سیاست، $\theta^{(0)}$، $H,\beta,\lambda$،
وزن‌های $\Omega$ و قرارداد $\mathsf{cfg}$
شامل $D,\tau_D,\tau,\tau_{\mathrm{SP}}$،
روش نرمال‌سازی و پایداری و
$\mathsf{loss}\in\{\mathrm{NLL},\mathrm{sig},
\mathrm{SP\!-\!JS}\}$؛ بهینه‌ساز و شمار گام $T$
\ENSURE مدل آموزش‌دیده
$\mathcal M_{\mathrm{tr}}$ و در صورت درخواست،
حالت بهینه‌ساز برای ادامهٔ آموزش
\STATE ساخت گراف، کرنل‌ها، مجموعه‌های مجاز و
$C_1,\ldots,C_K$ \textbf{فقط} از $D_{\mathrm{tr}}$
\STATE $\theta\gets\theta^{(0)}$؛
$\xi\gets\operatorname{InitOptimizer}(\theta)$
\FOR{$t=0,\ldots,T-1$}
    \FORALL{$i\in\mathcal I_{\mathrm{tr}}$}
        \IF{سیاست مستقل از فرضیهٔ طبقهٔ مبدأ است}
            \STATE $\mathscr Y_i^{[\cdot]}
            \gets\operatorname{ExactWalk}_{\mathrm{tr}}(i,\cdot;\theta)$
            \STATE $\mathscr Y_i^{[c]}\gets\mathscr Y_i^{[\cdot]}$
            برای همهٔ $c=1,\ldots,K$
        \ELSE
            \FOR{$c=1,\ldots,K$}
                \STATE $\mathscr Y_i^{[c]}
                \gets\operatorname{ExactWalk}_{\mathrm{tr}}(i,c;\theta)$
            \ENDFOR
        \ENDIF
    \ENDFOR
    \FOR{$c=1,\ldots,K$}
        \STATE $\mathcal V_c(\theta)\gets
        |C_c|^{-1}\sum_{i\in C_c}\mathscr Y_i^{[c]}$
    \ENDFOR
    \FORALL{$i\in\mathcal I_{\mathrm{tr}}$}
        \FOR{$c=1,\ldots,K$}
            \STATE $d_{i,c}\gets
            D(\mathscr Y_i^{[c]},\mathcal V_c(\theta))$
        \ENDFOR
        \IF{$\mathsf{loss}=\mathrm{NLL}$}
            \STATE $s_{i,:}\gets
            \operatorname{Score}_{\mathsf{cfg}}(d_{i,:})$
            \STATE $\ell_i\gets
            -s_{i,y_i}+\operatorname{LogSumExp}_r(s_{i,r})$
        \ELSE
            \STATE $S_{i,c}\gets 1-
            \operatorname{JSD}(\mathscr Y_i^{[c]},\mathcal V_c)/\log2$
            برای تمام $c$
            \IF{$\mathsf{loss}=\mathrm{sig}$}
                \STATE $m_i\gets S_{i,y_i}-\max_{r\ne y_i}S_{i,r}$
                \STATE $\ell_i\gets \ell_\tau(m_i)$
            \ELSE
                \STATE $\ell_i\gets\tau_{\mathrm{SP}}\,
                \operatorname{LogSumExp}_r
                ((S_{i,r}-S_{i,y_i})/\tau_{\mathrm{SP}})$
            \ENDIF
        \ENDIF
    \ENDFOR
    \STATE $\mathcal L_{\mathrm{tr}}(\theta)\gets
    \sum_{c=1}^{K}\frac{\omega_c}{|C_c|}
    \sum_{i\in C_c}\ell_i$
    \STATE $g\gets\nabla_\theta\mathcal L_{\mathrm{tr}}(\theta)$
    با مشتق‌گیری از \textbf{گشت‌ها و مراکز}
    \IF{$\mathcal L_{\mathrm{tr}}$ یا $g$ متناهی نیست}
        \STATE \textbf{توقف با گزارش خطای عددی}
    \ENDIF
    \STATE $(\theta,\xi)\gets
    \operatorname{OptimizerStep}(\theta,g,\xi)$
\ENDFOR
\STATE $\theta^{\mathrm{fit}}\gets\theta$
\STATE محاسبهٔ مجدد
$\mathscr Y_i^{[c]}(\theta^{\mathrm{fit}})$
برای همهٔ مبدأها و نامزدهای لازم با
$\operatorname{ExactWalk}_{\mathrm{tr}}$
\STATE $\mathcal V_c^{\mathrm{fit}}\gets
|C_c|^{-1}\sum_{i\in C_c}
\mathscr Y_i^{[c]}(\theta^{\mathrm{fit}})$
برای تمام $c$
\STATE $\mathcal M_{\mathrm{tr}}\gets
(\theta^{\mathrm{fit}},\mathcal V^{\mathrm{fit}},
D_{\mathrm{tr}},Y_{\mathrm{tr}},M_{\mathrm{tr}},
\mathsf{cfg},H,\beta,\lambda,\text{کرنل و قرارداد سیاست})$
\STATE \textbf{بازگردان}
$\mathcal M_{\mathrm{tr}}$ و، در صورت نیاز، $\xi$
\end{algorithmic}
\end{algorithm}

در شاخهٔ NLL، نگاشت
$\operatorname{Score}_{\mathsf{cfg}}$
دقیقاً قرارداد انتخاب‌شدهٔ
$-d/\tau_D$ یا نرمال‌سازی مجاز بخش
\ref{cmp:sec:ch3-score-normalization-analysis}
را اعمال می‌کند؛ این انتخاب باید در همهٔ
نامزدها یکسان و در بستهٔ مدل ذخیره شود.
در شاخه‌های سیگموئید و Softplus--JS،
$D=\operatorname{JSD}$ و امتیازها با
شباهت نرمال‌شدهٔ JS ساخته می‌شوند.
نقاط تساوی رقبا در حاشیهٔ سخت ممکن است
مشتق‌پذیر نباشند؛ تقریب هموار
رابطه~\eqref{cmp:eq:ch3-margin-smooth-competitor}
یک قرارداد اختیاریِ مستقل است.
برای KL نامتناهی و نرمال‌سازی تکین نیز
فرض‌ها یا روش پایدارسازی باید به‌صراحت
اعمال شوند؛ رابطهٔ گرادیانِ مشتق‌پذیر
به‌طور خودکار این وضعیت‌ها را حل نمی‌کند.

حالت بهینه‌ساز $\xi$ در گام‌های آموزش
نگهداری می‌شود: در SGD می‌تواند تهی و
در Adam شامل ممان‌ها و شمار گام باشد.
در نتیجهٔ ادامهٔ آموزش از checkpoint،
صرفاً ذخیرهٔ $\theta$ برای بازسازی همان
مسیر به‌روزرسانی کافی نیست.
به‌عنوان کنترل سازگاری، در پایان هر محاسبهٔ
معتبر باید $\sum_u\mathscr Y_i^{[c]}(u)=1$
و $\sum_u\mathcal V_c^{\mathrm{fit}}(u)=1$
برقرار باشد و مرکز نهایی با میانگین
بازنمایی‌های \emph{سیاست نهایی} مطابقت کند.
در مقابل، برای \emph{استنتاج}، پارامترها و
مراکزِ بازحساب‌شدهٔ نهایی همراه با
مرجع آموزش و قرارداد امتیاز لازم‌اند؛
مراکز محاسبه‌شده \emph{پیش از آخرین}
به‌روزرسانی عموماً با $\theta^{\mathrm{fit}}$
سازگار نیستند. در این شبه‌کد دادهٔ
اعتبارسنجی در محاسبهٔ گرادیان یا
بازبرآورد مراکز وارد نمی‌شود؛
پایش و انتخاب checkpoint با اعتبارسنجی
در الگوریتم جداگانه توضیح داده خواهد شد.

خانواده‌های گره‌محور
(\lr{node} و \lr{node\_class})
در مجموعهٔ \emph{آموزشِ ثابت} قابل
صورت‌بندی‌اند، ولی برای استفاده از
$\mathcal M_{\mathrm{tr}}$ روی نمونهٔ جدید،
به قاعدهٔ آغاز برون‌گراف نیاز دارند؛
مطابق بخش~\ref{cmp:subsec:ch3-node-policy-families}،
تا تعیین این قاعده، اعتبارسنجی برون‌نمونه‌ای
آن‌ها تعریف کامل ندارد.
همچنین اگر سیاست‌های شرطی طبقه به
\emph{سیاست میانگین} تبدیل شوند،
مراکز باید با همان سیاست میانگین
بازمحاسبه گردند؛ مراکز
$\mathcal V_c^{\mathrm{fit}}$ در الگوریتم حاضر
برای استنتاج \emph{فرضیه‌محور} ساخته شده‌اند
و بدون این بازبرآورد قابل جایگزینی نیستند.

\subsection{استنتاج برون‌نمونه‌ای و الگوریتم اعتبارسنجی}
\label{cmp:subsec:ch3-heldout-validation-algorithm}
الگوریتم~\ref{cmp:alg:ch3-learn-feature-policy} یک بستهٔ
مدل $\mathcal M_{\mathrm{tr}}$ تولید می‌کند که
پارامترهای سیاست، مراکز محاسبه‌شده با سیاست نهایی،
مرجع آموزشی و قرارداد امتیازدهی را در بر دارد.
برای ارزیابی نمونهٔ جدید $x$، \emph{نه} پارامترهای
سیاست و \emph{نه} مراکز آموزش دوباره برازش می‌شوند.
نمونهٔ پرس‌وجو، حتی اگر برچسبش در مجموعهٔ
اعتبارسنجی موجود باشد، عضو گراف مرجع برچسب‌دار
نیست؛ برچسب آن در ساخت بازنمایی به کار نمی‌رود.
این قرارداد با تفکیک مجموعه‌ها در
بخش~\ref{subsec:ch4-data-splits} نیز سازگار است.

\paragraph{گام نخست پرس‌وجو در برابر گشت آموزشی.}
در این زیربخش، $x$ شناسهٔ یک مشاهدهٔ پرس‌وجو
به‌همراه بردار ویژگی آن است؛ دو مشاهده با
مقادیر یکسان همچنان دو شناسهٔ متمایزند.
برای هر ویژگی مشاهده‌شدهٔ $b$ از نمونهٔ جدید،
$D_b^{\mathrm{tr}}=\{X_j\in D_{\mathrm{tr}}:
b\in A_{X_j}\}$ مجموعهٔ مقصدهای آموزشیِ
دارای آن ویژگی است. با همان وزن کرنل
رابطه~\eqref{cmp:eq:ch3-edge-weight}،
توزیع \emph{گام نخست} را به‌شکل زیر می‌سازیم:
\begin{equation}\label{cmp:eq:ch3-heldout-initial-transition}
q_b^{(x)}(j\mid x)=
\frac{w_b(x,X_j)}
{\sum_{X_k\in D_b^{\mathrm{tr}}}w_b(x,X_k)},
\qquad X_j\in D_b^{\mathrm{tr}}.
\end{equation}
این رابطه فقط برای ویژگی دارای مقدار مشاهده‌شده
در $x$ و مخرج مثبت به کار می‌رود؛ ویژگی‌های
ناموجود یا بدون مقصد معتبر از انتخاب نخست
حذف می‌شوند. اگر هیچ ویژگی قابل‌انتخابی
باقی نماند، پیش‌بینی در این قرارداد
تعریف نشده است و باید به‌جای ساخت
احتمال مصنوعی، وضعیت نامعتبر گزارش شود.

با سیاست آغازِ مجاز
$\pi_{\star}^{[c]}(b\mid x,\bot)$،
توزیع حالت پس از گام اول برابر است با
\begin{equation}\label{cmp:eq:ch3-heldout-first-state}
\nu_{x,1}^{[c]}(j,b)
=\pi_{\star}^{[c]}(b\mid x,\bot)\,
q_b^{(x)}(j\mid x).
\end{equation}
از گام دوم به بعد، گشت فقط روی
نمونه--ویژگی‌های مجموعهٔ آموزش ادامه
می‌یابد؛ توزیع‌های عمودی از
$D_b^{\mathrm{tr}}\setminus\{X_i\}$ برای
گرهٔ جاری $X_i$ ساخته می‌شوند.
چون $x\notin D_{\mathrm{tr}}$ است،
نیازی به حذف مبدأ پرس‌وجو از گره‌های
آموزشی نیست. همین تفاوت، استفادهٔ
بدون تغییر از فراخوانی گشت
\emph{مبدأ آموزشی} را نامعتبر می‌کند.

برای اختصار، $\mathfrak a$ را قرارداد آغاز
برون‌گراف می‌نامیم؛ در خانواده‌های
\lr{shared} و \lr{class}، این قرارداد
از سطر آغاز سیاست آموزش‌دیده به‌دست می‌آید،
اما در خانواده‌های گره‌محور آزاد باید
جداگانه مشخص شود.
$\operatorname{QueryExactWalk}(x,c;
\mathcal M_{\mathrm{tr}},\mathfrak a)$
به این قرارداد اشاره دارد: حالت اولیهٔ
رابطه~\eqref{cmp:eq:ch3-heldout-first-state}،
ادامهٔ گشت روی مرجع آموزش با سازوکار
$\operatorname{Advance}$ در
الگوریتم~\ref{cmp:alg:ch3-advance-walk-state}،
ضرایب $\alpha_t^{(H)}$ از
الگوریتم~\ref{cmp:alg:ch3-capped-walk-weights}
و جمع وزن‌دار برچسب \emph{مقصدهای آموزش}
مانند الگوریتم~\ref{cmp:alg:ch3-exact-walk-expectation}.
این نام، یک \emph{فراخوانی برون‌نمونه‌ای}
است، نه الگوریتم آموزشی چهارم یا روشی
مبتنی بر مونت‌کارلو. برای هر فرضیهٔ $c$،
خروجی آن $\mathscr Y_x^{[c]}\in\Delta_K$
است؛ شرط نرمال‌بودن فقط در صورت اعتبار
همهٔ انتقال‌های لازم برقرار خواهد بود.

\paragraph{سیاست شرطی و قابلیت پیش‌بینی برون‌نمونه‌ای.}
در خانوادهٔ \lr{shared} بازنمایی پرس‌وجو
بین همهٔ طبقات نامزد مشترک است؛ در خانوادهٔ
\lr{class} برای هر فرضیهٔ طبقهٔ مبدأ $c$
گشت جداگانه اجرا و فقط با مرکز سازگارِ
$\mathcal V_c^{\mathrm{fit}}$ مقایسه می‌شود.
این همان روش فرضیه‌محور
رابطه~\eqref{cmp:eq:ch3-class-hypothesis-inference}
است. برای خانواده‌های آزاد \lr{node}
و \lr{node\_class}، سطر آغاز مخصوص
نمونهٔ جدید در پارامترهای آموزش وجود ندارد؛
ازاین‌رو تنها در صورت ارائهٔ \emph{قاعدهٔ آغاز
برون‌گرافِ ازپیش‌تعیین‌شده} می‌توان از
الگوریتم زیر استفاده کرد.
راه‌حل مبتنی بر سیاست میانگینِ
رابطه~\eqref{cmp:eq:ch3-mean-policy-inference}
قرارداد متفاوتی است: در آن باید مراکز
نیز با همان سیاست میانگین ساخته شده باشند،
نه با مراکز فرضیه‌محور فعلی.

\paragraph{تفکیک پیش‌بینی یک نمونه از سنجش عملکرد.}
محاسبهٔ امتیازهای نامزد و برچسب پیش‌بینی‌شده برای
یک پرس‌وجوی برون‌نمونه‌ای، به برچسب واقعی آن نیاز ندارد.
ازاین‌رو زیرروال
$\operatorname{PredictOne}(x;\mathcal M_{\mathrm{tr}},\mathfrak a)$
را فقط یک‌بار تعریف می‌کنیم. ورودی آن یک \emph{نمونهٔ
دارای شناسه و ماسک} است و خروجی آن، وضعیت اعتبار
پیش‌بینی و در صورت اعتبار، همهٔ امتیازها،
احتمال‌های تخمینی و برچسب برنده است.
این زیرروال نه پارامتر سیاست را به‌روزرسانی می‌کند،
نه مرکز جدید می‌سازد و نه به برچسب واقعی
پرس‌وجو دسترسی دارد. الگوریتم اعتبارسنجی
\ref{cmp:alg:ch3-heldout-validation} تنها همین
زیرروال را برای همهٔ نمونه‌ها تکرار می‌کند.
این تفکیک برای استنتاج یک نمونهٔ جدید، ارزیابی
مجموعهٔ اعتبارسنجی و استنتاج مجموعهٔ آزمونِ
پس از قفل‌شدن مدل \emph{یک قرارداد مشترک} می‌سازد.

\begin{algorithm}[!htbp]
\caption{زیرروال $\operatorname{PredictOne}$ برای یک نمونهٔ برون‌نمونه‌ای}
\label{cmp:alg:ch3-heldout-predict-one}
\begin{algorithmic}[1]
\REQUIRE نمونهٔ پرس‌وجو $x$ با ماسک مشاهده و شناسهٔ مستقل؛
مدل منجمد $\mathcal M_{\mathrm{tr}}$؛
قاعدهٔ آغاز $\mathfrak a$، در صورت نیازِ خانوادهٔ سیاست
\ENSURE $(v_x,s_{x,:},\widehat p_{x,:},\widehat y_x)$؛
$v_x\in\{0,1\}$ نشانگر پیش‌بینی معتبر است
\STATE بازیابی سیاست، مراکز $\mathcal V^{\mathrm{fit}}$،
مرجع برچسب‌دار Train و $\mathsf{cfg}$ از مدل
\IF{خانوادهٔ گره‌محور است و قاعدهٔ آغاز برون‌گراف موجود نیست}
    \STATE \textbf{توقف: قرارداد استنتاج این خانواده تعریف نشده است}
\ENDIF
\STATE $B_x\gets\{b\in A_x:
\sum_{X_j\in D_b^{\mathrm{tr}}}w_b(x,X_j)>0\}$
\IF{$B_x=\varnothing$}
    \STATE \textbf{بازگردان} $(0,\bot,\bot,\bot)$
\ENDIF
\STATE ساخت انتقال آغازین روی $B_x$ از
رابطه~\eqref{cmp:eq:ch3-heldout-initial-transition}
و نرمال‌سازی سیاست آغاز فقط روی گزینه‌های مجاز
\IF{سیاست مستقل از فرضیهٔ طبقهٔ مبدأ است}
    \STATE $\mathscr Y_x^{[\cdot]}\gets
    \operatorname{QueryExactWalk}(x,\cdot;
    \mathcal M_{\mathrm{tr}},\mathfrak a)$
    \STATE $\mathscr Y_x^{[c]}\gets\mathscr Y_x^{[\cdot]}$
    برای همهٔ $c=1,\ldots,K$
\ELSE
    \FOR{$c=1,\ldots,K$}
        \STATE $\mathscr Y_x^{[c]}\gets
        \operatorname{QueryExactWalk}(x,c;
        \mathcal M_{\mathrm{tr}},\mathfrak a)$
    \ENDFOR
\ENDIF
\IF{یکی از بازنمایی‌های لازم نامعتبر یا نانرمال است}
    \STATE \textbf{بازگردان} $(0,\bot,\bot,\bot)$
\ENDIF
\IF{$\mathsf{loss}=\mathrm{NLL}$}
    \FOR{$c=1,\ldots,K$}
        \STATE $d_{x,c}\gets
        D(\mathscr Y_x^{[c]},\mathcal V_c^{\mathrm{fit}})$
    \ENDFOR
    \STATE $s_{x,:}\gets
    \operatorname{Score}_{\mathsf{cfg}}(d_{x,:})$
\ELSE
    \FOR{$c=1,\ldots,K$}
        \STATE $s_{x,c}\gets
        1-\operatorname{JSD}(\mathscr Y_x^{[c]},
        \mathcal V_c^{\mathrm{fit}})/\log2$
    \ENDFOR
\ENDIF
\IF{هر یک از $K$ امتیاز نامعتبر یا نامتناهی است}
    \STATE \textbf{بازگردان} $(0,\bot,\bot,\bot)$
\ENDIF
\STATE $\widehat y_x\gets
\min\{c:s_{x,c}=\max_{r=1,\ldots,K}s_{x,r}\}$
\STATE $\widehat p_{x,:}\gets
\operatorname{Softmax}(s_{x,:})$
بدون تکرار مقیاس‌دهی قرارداد $\mathsf{cfg}$
\STATE \textbf{بازگردان}
$(1,s_{x,:},\widehat p_{x,:},\widehat y_x)$
\end{algorithmic}
\end{algorithm}

\begin{algorithm}[!htbp]
\caption{استنتاج و اعتبارسنجی با فراخوانی مشترک پیش‌بینیِ یک نمونه}
\label{cmp:alg:ch3-heldout-validation}
\begin{algorithmic}[1]
\REQUIRE مدل منجمد $\mathcal M_{\mathrm{tr}}$؛
نمونه‌های برون‌نمونه‌ای و ماسک‌هایشان
$X_{\mathrm{val}},M_{\mathrm{val}}$ با شناسه‌های متمایز؛
$Y_{\mathrm{val}}$ \textbf{اختیاری و فقط برای ارزیابی}؛
قاعدهٔ آغاز برون‌گراف $\mathfrak a$ در صورت نیاز
\ENSURE وضعیت، امتیاز، احتمال و برچسب پیش‌بینی
هر پرس‌وجو؛ گزارش $\mathcal R_{\mathrm{val}}$
\STATE $\mathcal Q_{\mathrm{ok}}\gets\varnothing$؛
$\mathcal Q_{\mathrm{bad}}\gets\varnothing$؛
$\mathcal R_{\mathrm{val}}\gets\varnothing$
\STATE نگه‌داشتن تمام مؤلفه‌های
$\mathcal M_{\mathrm{tr}}$ به‌صورت ثابت
\FORALL{$x\in X_{\mathrm{val}}$}
    \STATE $(v_x,s_{x,:},\widehat p_{x,:},\widehat y_x)
    \gets\operatorname{PredictOne}
    (x;\mathcal M_{\mathrm{tr}},\mathfrak a)$
    با الگوریتم~\ref{cmp:alg:ch3-heldout-predict-one}
    \IF{$v_x=1$}
        \STATE $\mathcal Q_{\mathrm{ok}}\gets
        \mathcal Q_{\mathrm{ok}}\cup\{\operatorname{id}(x)\}$
    \ELSE
        \STATE $\mathcal Q_{\mathrm{bad}}\gets
        \mathcal Q_{\mathrm{bad}}\cup\{\operatorname{id}(x)\}$
    \ENDIF
\ENDFOR
\STATE ثبت $\mathrm{Coverage}\gets
|\mathcal Q_{\mathrm{ok}}|/|X_{\mathrm{val}}|$
و شمار پرس‌وجوهای نامعتبر؛ فرض $|X_{\mathrm{val}}|>0$
\IF{$Y_{\mathrm{val}}$ برای ارزیابی موجود است}
    \STATE \textbf{اکنون و فقط اکنون} استفاده از
    $Y_{\mathrm{val}}$ برای مقایسه با پیش‌بینی‌ها
    \STATE $\mathsf{Conf}_{a,b}^{\mathrm{val}}\gets
    \sum_{x\in X_{\mathrm{val}}:\,v_x=1}
    \mathbf1\{y_x=a,\widehat y_x=b\}$
    برای $a,b=1,\ldots,K$
    \STATE محاسبهٔ Accuracy، Recall هر طبقه،
    BA و Precision از ماتریس اغتشاش، در صورت تعریف
    \STATE محاسبهٔ حاشیهٔ صحیح و زیان
    برچسب‌دار با همان $\mathsf{cfg}$، در صورت تعریف
    \STATE $\mathcal R_{\mathrm{val}}\gets
    (\mathsf{Conf}^{\mathrm{val}},\mathrm{Accuracy},
    \mathrm{Recall},\mathrm{BA},\mathrm{Precision},
    \mathrm{Coverage},|\mathcal Q_{\mathrm{bad}}|,
    \text{سایر سنجه‌های معتبر})$
\ELSE
    \STATE $\mathcal R_{\mathrm{val}}\gets
    (\mathrm{Coverage},|\mathcal Q_{\mathrm{bad}}|)$
\ENDIF
\STATE \textbf{بازگردان}
$(v,s,\widehat p,\widehat y,\mathcal R_{\mathrm{val}})$
\end{algorithmic}
\end{algorithm}

در الگوریتم~\ref{cmp:alg:ch3-heldout-predict-one}،
نماد $\bot$ به معنی \emph{نبود خروجی معتبر} است،
نه یک امتیاز عددی یا برچسب طبقه.
بنابراین تنها پرس‌وجوهای دارای $v_x=1$
در ماتریس اغتشاشِ اعتبارسنجی وارد می‌شوند.
در حالت یک نمونهٔ بدون برچسب، همان زیرروال
بدون فراخوانی الگوریتم سنجش مجموعه‌ای قابل‌استفاده است.
این زیرروال تنها یک لایهٔ امتیازدهی و تصمیم
بر فراخوانی محاسباتی
$\operatorname{QueryExactWalk}$ است؛
نسخهٔ تازه‌ای از موتور گشت، $\operatorname{ExactWalk}_{\mathrm{tr}}$
یا الگوریتم~\ref{cmp:alg:ch3-exact-walk-expectation}
تعریف نمی‌کند. ثبت شناسهٔ هر پرس‌وجو مانع
یکی‌گرفتن دو مشاهدهٔ هم‌مقدار در شمارش
پوشش و ماتریس اغتشاش می‌شود.

منظور از نامعتبربودن بازنمایی یا امتیاز
در الگوریتم، وجود مقدار نامتناهی،
جرم احتمال نامتعارف یا نقض شرط
نرمال‌بودن بازنمایی است؛ چنین خطایی
با دست‌کاری برچسب یا جایگذاری یک
برچسب دلخواه رفع نمی‌شود.
در ارزیابی، پوشش پیش‌بینی برابر
$|\mathcal Q_{\mathrm{ok}}|/|X_{\mathrm{val}}|$
و Accuracy نسبت نمونه‌های درست
به شمار نمونه‌های دارای پیش‌بینی معتبر است.
Recall طبقهٔ $c$
برابر
$\mathsf{Conf}_{c,c}^{\mathrm{val}}/
\sum_b\mathsf{Conf}_{c,b}^{\mathrm{val}}$
است. اگر تمام $K$ طبقه در
$\mathcal Q_{\mathrm{ok}}$ حضور داشته باشند،
$\mathrm{BA}=K^{-1}\sum_c\mathrm{Recall}_c$
خواهد بود. Precision طبقهٔ $c$
نیز در صورت مثبت بودن شمار پیش‌بینی‌شده
برای آن، برابر
$\mathsf{Conf}_{c,c}^{\mathrm{val}}/
\sum_a\mathsf{Conf}_{a,c}^{\mathrm{val}}$
است. در نبودِ نمونهٔ واقعی برای یک طبقه،
BA تعریف‌شده روی تمام $K$ طبقه
قابل محاسبه نیست و باید وضعیت آن گزارش
شود؛ Precision با مخرج صفر نیز بدون
تعیین قرارداد، \emph{تعریف‌نشده} است.
اگر پرس‌وجویی فاقد پیش‌بینی معتبر باشد،
پوشش پیش‌بینی و تعداد موارد حذف‌شده باید
گزارش شوند و نباید BA/Accuracy نمونه‌های
پوشش‌داده‌شده را به‌جای معیار کل مجموعه
معرفی کرد.

برچسب واقعی $y_x$ برای محاسبهٔ
حاشیهٔ نظارت‌شده
$m_x=s_{x,y_x}-\max_{r\ne y_x}s_{x,r}$
و زیان اعتبارسنجی، \emph{پس از}
تکمیل همهٔ امتیازهای نامزد به کار می‌رود.
در صورت انتخاب NLL، زیان با همان احتمال
تخمینی تعریف می‌شود؛ برای سیگموئید و
Softplus--JS باید به‌ترتیب از زیان
رابطه~\eqref{cmp:eq:ch3-js-sigmoid-loss}
و~\eqref{cmp:eq:ch3-softplus-js-training-objective}
استفاده کرد. خروجی Softmax در اینجا
\emph{احتمال تخمینی مدل} است، نه
پسین واقعیِ معلوم از مولد داده.
برای مقایسهٔ checkpointهای مختلف،
این الگوریتم می‌تواند برای هر مدلِ حاصل از
\emph{گام تکمیل‌شدهٔ آموزش} جداگانه اجرا شود؛
در هر بار، مراکز باید از سیاست همان checkpoint
و \emph{فقط} نمونه‌های Train محاسبه شده باشند.
انتخاب بهترین checkpoint یا فراپارامتر
از روی معیار مشخصِ اعتبارسنجی به‌معنای
پس‌انتشار گرادیان Validation نیست.
از همین سازوکار استنتاج می‌توان برای
نمونهٔ آزمون مستقل \emph{پس از تثبیت مدل}
استفاده کرد؛ انتخاب پارامتر، checkpoint
یا آستانه بر اساس Test مجاز نیست.

\section{تحلیل نظری مشتق‌ها و گرادیان یادگیری سیاست}
\label{cmp:sec:ch3-policy-gradient-theory}
این بخش جزئیات تحلیلی
مشتق‌گیری از گشت، بازنمایی‌ها،
مراکز و توابع زیان را تکمیل می‌کند.
مفاهیم لازم برای درک چرخهٔ آموزش
و پیش‌بینی در
بخش~\ref{cmp:sec:ch3-policy-learning-algorithm}
معرفی شدند و مطالعهٔ این مشتق‌های
مؤلفه‌ای برای اجرای شبه‌کد آن بخش
ضروری نیست. روابط زیر روشن می‌کنند
چگونه گرادیان به سیاست بازمی‌گردد،
در چه شرایطی فرمول‌های مشتق‌پذیر
معتبرند و چرا گرادیان و
\emph{تغییر واقعی پارامترها}
یک کمیت نیستند.

\subsection{مشتق‌گیری از بازنمایی‌ها، مراکز و حاشیه}
\label{cmp:subsec:ch3-margin-policy-derivatives}
زیان سیگموئیدی یک تابع
هموار از \emph{مقدار حاشیه}
است، اما مشتق‌گیری نسبت به
پارامترهای سیاست نیازمند
محاسبه اثر سیاست بر
بازنمایی نمونه‌ها و مراکز
نیز هست. نگاشت پیشرو
سیاست به بازنمایی‌ها،
از طریق احتمالات انتقال
و امید گشت‌ها در بخش‌های
پیشین این فصل معرفی شد.
برای یک سیاست پارامتری
$\pi_\theta$، رابطه
\eqref{cmp:eq:ch3-class-center}
نتیجه می‌دهد:
\begin{equation}\label{cmp:eq:ch3-margin-center-policy-derivative}
\boxed{
\nabla_\theta\mathcal V_c(\theta)
=\frac1{|C_c|}\sum_{j\in C_c}
\nabla_\theta\mathscr Y_j(\theta).
}
\end{equation}
این رابطه برای داده و
تعداد اعضای ثابت طبقه برقرار است.
ازاین‌رو در مدل اصلی،
تغییر سیاست هم خودِ
بازنمایی یک نمونه و
هم تمامی مراکز را
تغییر می‌دهد؛ حذف مشتق
مراکز بدون فرضی صریح
در این تحلیل مجاز نیست.

\paragraph{مشتق امتیاز مبتنی بر JSD.}
برای پرهیز از مشکل مشتق
در مؤلفه‌های صفر،
بازنمایی‌ها و مراکز
هموارشده رابطه
\eqref{cmp:eq:ch3-stable-smoothed-representations}
را در نظر می‌گیریم و
برای اختصار
$P_i=\mathscr Y_i^{(\eta)}$،
$V_c=\mathcal V_c^{(\eta)}$
و $s_{i,c}=s(P_i,V_c)$
می‌نویسیم. با بسط
آنتروپی در تعریف JSD،
برای هر مؤلفه $u$ با
$P_i(u),V_c(u)>0$ داریم:
\begin{equation}\label{cmp:eq:ch3-margin-JS-partials}
\boxed{
\begin{aligned}
\frac{\partial s(P_i,V_c)}{\partial P_i(u)}
&=-\frac1{2\log2}
\log\frac{2P_i(u)}{P_i(u)+V_c(u)},\\
\frac{\partial s(P_i,V_c)}{\partial V_c(u)}
&=-\frac1{2\log2}
\log\frac{2V_c(u)}{P_i(u)+V_c(u)}.
\end{aligned}}
\end{equation}
با قاعده زنجیره‌ای،
مشتق امتیاز نسبت به
پارامترهای سیاست برابر است با:
\begin{equation}\label{cmp:eq:ch3-margin-score-policy-gradient}
\boxed{
\begin{aligned}
\nabla_\theta s_{i,c}
=-\frac1{2\log2}\sum_{u=1}^{K}
\Bigg[
&\log\frac{2P_i(u)}{P_i(u)+V_c(u)}
\nabla_\theta P_i(u)\\
+&\log\frac{2V_c(u)}{P_i(u)+V_c(u)}
\nabla_\theta V_c(u)
\Bigg].
\end{aligned}}
\end{equation}
در این فرمول، برای $\eta$
ثابت، مشتق بازنمایی‌های
هموارشده برابر
$(1-\eta)$ برابر مشتق
بازنمایی اولیه است.
رابطه مشتق مرکز در این
حالت نیز با همان میانگین
مشتق بازنمایی‌ها برقرار
می‌ماند. بنابراین هیچ
نیازی به فرض ثابت بودن
مراکز در یادگیری آزاد
سیاست نداریم.

\paragraph{مشتق حاشیه و نزدیک‌ترین رقیب.}
اگر رقیب برنده
$r_i^\star=\arg\max_{r\ne y_i}s_{i,r}$
یکتا باشد، در ناحیه‌ای
که این رقیب تغییر نمی‌کند:
\begin{equation}\label{cmp:eq:ch3-margin-policy-gradient}
\boxed{
\nabla_\theta M_i^{\mathrm{JS}}
=\nabla_\theta s_{i,y_i}
-\nabla_\theta s_{i,r_i^\star}.
}
\end{equation}
برای تساوی امتیاز رقبا،
مشتق طبقه‌ایک حاشیه ممکن
است وجود نداشته باشد.
برای حذف این نقطه نامشتق‌پذیر،
با $\kappa>0$ حاشیه نرم
را تعریف می‌کنیم:
\begin{equation}\label{cmp:eq:ch3-margin-smooth-competitor}
\boxed{
\widetilde M_{i,\kappa}^{\mathrm{JS}}
=s_{i,y_i}
-\kappa\log\sum_{r\ne y_i}
\exp(s_{i,r}/\kappa),
\qquad
0\le M_i^{\mathrm{JS}}
-\widetilde M_{i,\kappa}^{\mathrm{JS}}
\le\kappa\log(K-1).
}
\end{equation}
پس این تقریب همواره
حاشیه را به‌اندازه‌ای
محدود و \emph{به سمت
مقادیر کوچک‌تر} تغییر
می‌دهد و برای $K=2$
دقیقاً با حاشیه اصلی
برابر است. اگر وزن‌های
نرم رقبا را با
\[
w_{i,r}^{(\kappa)}
=\frac{\exp(s_{i,r}/\kappa)}
{\sum_{q\ne y_i}\exp(s_{i,q}/\kappa)}
\]
نشان دهیم، مشتق حاشیه
نرم دقیقاً برابر است با:
\begin{equation}\label{cmp:eq:ch3-margin-smooth-gradient}
\boxed{
\nabla_\theta\widetilde M_{i,\kappa}^{\mathrm{JS}}
=\nabla_\theta s_{i,y_i}
-\sum_{r\ne y_i}
w_{i,r}^{(\kappa)}
\nabla_\theta s_{i,r}.
}
\end{equation}
در نتیجه، برای میانگین
متوازن زیان سیگموئیدی
با حاشیه نرم، گرادیان
قابل محاسبه زیر حاصل
می‌شود:
\begin{equation}\label{cmp:eq:ch3-margin-loss-policy-gradient}
\boxed{
\nabla_\theta\widetilde{\mathcal L}_{\mathrm{sig}}
=\frac1K\sum_{c=1}^{K}\frac1{|C_c|}
\sum_{i\in C_c}
\ell_\tau'\!\left(\widetilde M_{i,\kappa}^{\mathrm{JS}}\right)
\nabla_\theta\widetilde M_{i,\kappa}^{\mathrm{JS}}.
}
\end{equation}
این فرمول،
\emph{روش محاسبه مشتق
هدف حاشیه‌ای} را مشخص
می‌کند، نه تضمین کاهش
خطای صفر--یک در هر
گام بهینه‌سازی.

\paragraph{بازنمایی مشترک، شرطی و مراکز موقتاً ثابت.}
در حالت بازنمایی مشترک،
$P_i$ در امتیاز همه
طبقات یکسان است و
مشتق آن در تمامی امتیازها
تکرار می‌شود. در
استنتاج شرطی طبقه،
امتیاز نامزد $c$
از $P_i^{(c)}$ و
مرکز متناظر به‌دست
می‌آید؛ در نتیجه
مشتق هر امتیاز باید
از همان مسیر محاسبه
شود و دیگر نمی‌توان
مشتق تمام امتیازها را
با یک بردار $P_i$
مشترک نوشت.
اگر مراکز فقط هنگام
محاسبه یک گام گرادیان
\emph{ثابت فرض شوند}،
جمله
$\nabla_\theta V_c$
از رابطه
\eqref{cmp:eq:ch3-margin-score-policy-gradient}
حذف می‌شود؛ اما این
تقریب با \emph{ثابت
ماندن واقعی مراکز
تحت تغییر سیاست}
یکسان نیست. برای
ثابت‌ماندن واقعی
مراکز باید قید
$\mathcal V_c(\theta)
=\mathcal V_c^\star$
برقرار باشد؛ چگونگی
تحلیل فضای سیاست‌های
دارای مرکز مشترک
در بحث نگاشت پیشرو
و مسئله معکوس این فصل
بررسی می‌شود.

شرح مفصل مشتق‌های
JSD، دقت تقریب
بیشینه نرم و شرایط
مرزی آن‌ها در
پیوست~\ref{app:ch3-margin-gradient}
آمده است.




\subsection{مشتق زیان Softplus--JS و منشأ تغییر سیاست}
\label{cmp:subsec:ch3-softplus-js-policy-gradient}
در زیان Softplus--JS رابطه~\eqref{cmp:eq:ch3-softplus-js-training-objective}،
امتیاز $S_{i,r}$ به \emph{طبقهٔ نامزد} $r$ تعلق دارد؛
این امتیاز نباید با $z_{ab}$، یعنی logits انتخاب \emph{ویژگی}، اشتباه شود.
برای هر نمونهٔ آموزشی $i$ و نامزد $r$، بازنمایی
$U_{i,r}(z)=\mathscr Y_i^{[r]}(z)$ و مرکز
$V_r(z)=|C_r|^{-1}\sum_{j\in C_r}U_{j,r}(z)$
را در نظر می‌گیریم که در آن $z=\operatorname{vec}(Z)$
برای خانواده‌های ماتریسی ثابت است.
هم خودِ بازنمایی‌ها و هم مراکز فقط از گشت دقیق
\emph{ماسک‌دار روی Train} حاصل می‌شوند؛ در محاسبهٔ مرکز،
مبدأهای $j\in C_r$ با فرضیهٔ سازگار $r$ ارزیابی می‌شوند.
اما برای امتیازِ نمونهٔ $i$، \emph{تمام} نامزدهای $r$
بدون استفاده از $y_i$ برای انتخاب سیاست فرضیه محاسبه می‌شوند.

با تعریف وزن‌های نرم \emph{روی امتیاز طبقه‌ها}
و مشتق زیان نمونه نسبت به همان امتیازها، داریم:
\begin{equation}\label{cmp:eq:ch3-softplus-js-score-weights}
\boxed{
w_{i,r}^{\mathrm{SP}}
=\frac{\exp(S_{i,r}/\tau_{\mathrm{SP}})}
{\sum_{q=1}^{K}\exp(S_{i,q}/\tau_{\mathrm{SP}})},
\qquad
\frac{\partial\ell_i^{\mathrm{SP}}}{\partial S_{i,r}}
=w_{i,r}^{\mathrm{SP}}-\mathbf1\{r=y_i\}.
}
\end{equation}
وزن $w_{i,r}^{\mathrm{SP}}$ احتمال انتخاب ویژگی نیست
و با $p_0$ و $P$ در تحلیل نویز نیز تفاوت دارد.
زیان متوازن برای وزن‌های طبقات $\Omega$،
مطابق قرارداد آموزش، برابر است با
$\mathcal L_{\mathrm{SP},\Omega}(z)
=\sum_{c=1}^{K}\frac{\omega_c}{|C_c|}
\sum_{i\in C_c}\ell_i^{\mathrm{SP}}(z)$.
در صورت $\omega_c=1/K$، گرادیان آن برابر است با:
\begin{equation}\label{cmp:eq:ch3-softplus-js-policy-gradient}
\boxed{
\nabla_z\mathcal L_{\mathrm{SP},U}
=\frac1K\sum_{c=1}^{K}\frac1{|C_c|}
\sum_{i\in C_c}\sum_{r=1}^{K}
\bigl(w_{i,r}^{\mathrm{SP}}-\mathbf1\{r=c\}\bigr)
\nabla_z S_{i,r}.
}
\end{equation}
برای وزن‌های عمومی $\Omega$، تنها $1/K$ در جمع بیرونی
با $\omega_c$ جایگزین می‌شود؛ مشتق نمونه تغییر نمی‌کند.

در ناحیه‌ای که مؤلفه‌های لازم مثبت و نگاشت گشت
با ماسک‌های ثابت مشتق‌پذیر است، با استفاده از
رابطه~\eqref{cmp:eq:ch3-margin-JS-partials}،
مشتق امتیاز هر نامزد از \emph{دو مسیر} به‌دست می‌آید:
\begin{equation}\label{cmp:eq:ch3-softplus-js-score-chain}
\boxed{
\begin{aligned}
\nabla_z S_{i,r}
=-\frac1{2\log2}\sum_{u=1}^{K}\Bigg[
&\log\frac{2U_{i,r}(u)}{U_{i,r}(u)+V_r(u)}
\,\nabla_z U_{i,r}(u)\\
+&\log\frac{2V_r(u)}{U_{i,r}(u)+V_r(u)}
\,\frac1{|C_r|}\sum_{j\in C_r}\nabla_z U_{j,r}(u)
\Bigg].
\end{aligned}}
\end{equation}
برای اعمال این فرمول به هموارسازی مثبت مشترکِ
رابطه~\eqref{cmp:eq:ch3-stable-smoothed-representations}،
باید $U_{i,r}$ و $V_r$ را با صورت‌های هموارشده
جایگزین و ضرایب مشتق نگاشت هموارسازی را لحاظ کرد.
حذف جملهٔ مشتق مرکز در آموزش آزاد مجاز نیست.
در خانوادهٔ \lr{shared} مسیر بازنمایی نامزدها مشترک است،
اما برای \lr{class} هر $U_{i,r}$ از سیاست فرضیهٔ $r$ می‌آید.

\paragraph{تکیه‌گاه صفر و پیاده‌سازی واقعی.}
رابطه~\eqref{cmp:eq:ch3-softplus-js-score-chain} فرمول
داخلیِ مشتق‌پذیر است، نه مجوز جایگذاری لگاریتم صفر.
در کد گشت دقیق، JSD با قرارداد $0\log0=0$،
تابع $\operatorname{where}$، کف عددی در لگاریتم
و $\operatorname{clamp}_{\min}(0)$ روی نتیجه اجرا می‌شود؛
\emph{هموارسازی مثبت سراسری با $\eta>0$ در آن فعال نیست}.
روی تکیه‌گاه‌های ساختاری ثابت، مؤلفه‌هایی که همواره صفرند
مشتق صفر دارند و باید سهمشان جداگانه تفسیر شود.
در مرز تغییر تکیه‌گاه، مقادیر بسیار کوچک، یا محل فعال‌شدن
\lr{clamp}، مشتق صوری بالا لزوماً با مشتق انتخابی
\lr{autograd} یکسان نیست؛ کنترل عددی این مرزها جداگانه لازم است.
جزئیات مشتق و حدود اعتبار آن در
بخش~\ref{app:ch3-softplus-gradient-proof}
از پیوست~\ref{app:ch3-proofs}
آمده‌اند.

\paragraph{از گرادیان تا تغییر واقعی logits.}
اگر $g^{(t)}=\nabla_z\mathcal L_{\mathrm{SP},U}(z^{(t)})$
باشد، این گرادیان \emph{خودِ آپدیت} نیست.
برای Adam استاندارد، بدون کاهش وزن یا اصلاح اضافی، با
نرخ یادگیری $\eta_{\mathrm A}$ و ممان‌های $m^{(t)},v^{(t)}$ داریم:
\begin{equation}\label{cmp:eq:ch3-adam-logit-update}
\boxed{
\begin{aligned}
m^{(t+1)}&=\beta_1m^{(t)}+(1-\beta_1)g^{(t)},\\
v^{(t+1)}&=\beta_2v^{(t)}+(1-\beta_2)(g^{(t)})^{\odot2},\\
\widehat m^{(t+1)}&=m^{(t+1)}/(1-\beta_1^{t+1}),\quad
\widehat v^{(t+1)}=v^{(t+1)}/(1-\beta_2^{t+1}),\\
\Delta z^{(t)}&=z^{(t+1)}-z^{(t)}
=-\eta_{\mathrm A}
\frac{\widehat m^{(t+1)}}{\sqrt{\widehat v^{(t+1)}}+
\varepsilon_{\mathrm A}} .
\end{aligned}}
\end{equation}
توان دوم، تقسیم و ریشه در این رابطه مؤلفه‌به‌مؤلفه‌اند؛
$\beta_1,\beta_2$ ممان‌های Adam هستند و
با $\beta$ توقف گشت تفاوت دارند.
اگر logits خروجی تابع پارامتری $Z(\theta)$ باشند،
تغییر واقعی از $Z(\theta^{(t+1)})-Z(\theta^{(t)})$
به‌دست می‌آید، نه از جایگذاری مستقیم گرادیان $z$
در فرمول مخصوص پارامتر مستقل.
حتی در مدل‌های ماتریسی، $\Delta z^{(t)}$ را می‌توان
مستقیماً از دو checkpoint پیاپی به‌دست آورد.
نگاشت این تغییر به تغییر احتمال سیاست
در بخش~\ref{cmp:subsec:ch3-noise-logpolicy-sensitivity}
\emph{مستقل از زیان و بهینه‌ساز} بررسی می‌شود.

\section{رابطهٔ سیاست و مراکز: دستیابی، چندگانگی و اصلاح سیاست}
\label{cmp:sec:ch3-inverse-center-policy}
در بخش‌های پیشین مسیر \emph{پیشرو} را ساختیم:
با داشتن سیاست انتخاب ویژگی، امید گشت،
بازنمایی‌های احتمالاتی و مراکز طبقه‌ها
محاسبه می‌شوند؛ سپس از این مراکز برای
تصمیم و آموزش سیاست استفاده می‌کنیم.
اما موفقیت یک روش بهینه‌سازی به‌تنهایی
پاسخ نمی‌دهد که \emph{هر آرایش مطلوبی
از مراکز} اصولاً از سیاست‌های مجاز
قابل‌دستیابی است یا نه؛ همچنین ممکن
است سیاست‌های متفاوتی دقیقاً مراکز
یکسان بسازند. ازاین‌رو پس از تعریف
آموزش، جهت پرسش را برمی‌گردانیم:
اگر مراکز هدف یا تغییری مطلوب در آن‌ها
داشته باشیم، چه چیزی دربارهٔ
وجود، چندگانگی یا اصلاح سیاست
می‌توان گفت؟

برای پاسخ، ابتدا وجود و یکتایی
بازنمایی‌ها و مراکز حاصل از یک
سیاست ثابت را با خواص مارکوفی
گشت بررسی می‌کنیم؛ سپس
محدودیت دستیابی‌پذیری مراکز
و سیاست‌های هم‌مرکز را توضیح
می‌دهیم و سرانجام چارچوبی
برای اصلاح سیاست با مراکز
هدف ثابت صورت‌بندی می‌کنیم.
این تحلیل، مسئلهٔ معکوس
و حدود توان خانوادهٔ سیاست
را روشن می‌کند و هیچ تضمین
مستقلی برای بهبود دقت یا
همگرایی سراسری بهینه‌ساز
فراهم نمی‌آورد.

در سراسر این بخش، افق $H$، پارامترهای $\beta,\lambda$، گراف مرجع و خانوادهٔ سیاست‌ها ثابت فرض می‌شوند. مراکزِ تابع سیاست با امید دقیق گشت سقف‌دار در رابطه~\eqref{cmp:eq:ch3-capped-alpha} محاسبه می‌شوند. ازاین‌رو، وابستگی مراکز به $H$ نیز در نمادها آشکار خواهد بود.

فرض کنیم خانواده مجاز سیاست‌ها با پارامتر $z=\operatorname{vec}(Z)\in\mathbb R^p$ مشخص شود؛ برای ماتریس امتیاز ثابت رابطه~\eqref{cmp:eq:ch3-fixed-policy-matrix}، تعداد پارامترهای خام $p=d(d+1)$ است. به دلیل ناوردایی نرمال‌سازی با بیشینهٔ هموار نسبت به افزودن یک ثابت به هر سطر، پارامترسازی لزوماً یکتا نیست. با استفاده از رابطه~\eqref{cmp:eq:ch3-exact-matrix-neighborhood} و تعریف مرکز در رابطه~\eqref{cmp:eq:ch3-class-center}، نگاشت پیشرو به‌صورت زیر است:
\begin{equation}\label{cmp:eq:ch3-center-forward-map}
\mathcal F_{\beta,\lambda,H}(z)=\mathcal V(z;\beta,\lambda,H),
\qquad
\mathcal V_c(z;\beta,\lambda,H)=
\frac{1}{|C_c|}\sum_{s\in C_c}\mathscr Y_s(z;\beta,\lambda,H).
\end{equation}
اگر برای هر فرض طبقه، ماتریس امتیاز جداگانه $Z_c$ به کار رود، سطر $c$ام $\mathcal V$ به $Z_c$ وابسته خواهد بود. این خانواده از حالت یک ماتریس $Z$ مشترک بین تمام طبقه‌ها آزادتر است؛ مقایسه آن‌ها تنها با مشخص‌کردن فضای سیاست معتبر خواهد بود.




\subsection{خواص مارکوفی گشت، بازآغاز و همگرایی}

\paragraph{خواص مارکوفی گشت و یکتایی مراکز برای یک سیاست ثابت.}
طبق فرض‌های اول تا پنجمِ بخش
\ref{cmp:subsec:ch3-definitions}
و تعریف گراف در بخش
\ref{cmp:subsec:ch3-graph}،
گراف کامل نمونه--ویژگی پیش از
حذف مبدأ از نظر ساختاری متصل است؛
مثبت بودن وزن‌های عمودی و
احتمال‌های حاصل از نرمال‌سازی با بیشینهٔ هموار برای ویژگی‌های
مجاز نیز امکان طی مسیرهای ساختاری
را فراهم می‌کند. با این حال،
\emph{گشت مورد استفاده در این پژوهش
به مبدأ بازنمی‌گردد}؛ بنابراین
پس از حذف $X_s$، مجموعه حالت
$\mathcal S_s$ ممکن است به
چند مؤلفه تقسیم شود.
فرض پنجم تضمین می‌کند هیچ سطر
قابل استفاده بدون مقصد نماند،
اما به‌تنهایی تضمین‌کننده
تحویل‌ناپذیری و نادوره‌بودن
زنجیره کاهش‌یافته نیست.

برای یک مبدأ $s$ و سیاست
ثابت، ماتریس گذار
$\mathsf P_s$ در رابطه
\eqref{cmp:eq:ch3-block-factorization}
سطری--تصادفی است و
$\rho_s$ در رابطه
\eqref{cmp:eq:ch3-block-first-step}
توزیع شروع معتبری است.
با ثابت‌بودن گراف، سیاست،
$H$ و پارامترهای
$\beta,\lambda$، بازنمایی
احتمالاتی نمونه به‌صورت
\begin{equation}\label{cmp:eq:ch3-policy-center-unique-walk}
\boxed{
\mathscr Y_s(z)
=\rho_s(z)
\left(\sum_{t=1}^{H}
\alpha_t^{(H)}
\mathsf P_s(z)^{t-1}\right)
\mathsf M_s^{(Y)}
}
\end{equation}
\emph{به‌طور یکتا} تعیین
می‌شود؛ در نتیجه میانگین
بازنمایی‌های هر طبقه نیز
مرکز توزیعی یکتایی دارد.
برای نسخه افق نامتناهی،
جمع وزن‌دار متناظر نیز
همگراست؛ زیرا $\alpha_t\ge0$،
$\sum_{t\ge1}\alpha_t=1$
و تمام بردارهای
$\rho_s\mathsf P_s^{t-1}
\mathsf M_s^{(Y)}$
توزیع احتمال هستند.
این \emph{همگرایی سری
وزن‌دار} را نباید با
همگرایی خود توان‌های
$\mathsf P_s^t$
به یک توزیع مانای یکتا
یکی گرفت. همگرایی اخیر
برای یک زنجیره متناهی،
مثلاً تحت تحویل‌ناپذیری
و نادوره‌بودن برقرار است،
ولی از پنج فرض فعلی
به‌تنهایی نتیجه نمی‌شود.
حتی در آن حالت،
بازنمایی وزن‌دار با
افق و توزیع توقف معین
لزومـاً برابر توزیع مانا
نخواهد بود.
اثبات و مثال نقض
در پیوست~\ref{app:ch3-policy-center-markov}
آمده است.


\paragraph{بازآغاز متوالی: یک مسیر پیوسته با گشت‌های پیاپی.}
گشت‌های آغازشده از $X_s$ را می‌توان
نه به‌صورت اجراهای موازی،
بلکه به‌صورت قطعه‌های
\emph{متوالی یک فرایند واحد}
تعبیر کرد: پس از هر مشاهده،
با احتمال $\beta$ گشت جاری
پایان می‌یابد، فرایند به
حالت ویژه آغاز $o_s$
بازمی‌گردد و گشت بعدی
با توزیع $\rho_s$ آغاز می‌شود.
در طول \emph{هر قطعه} بازگشت
به نمونه $X_s$ ممنوع است؛
بازآغاز میان دو قطعه،
این ممنوعیت را نقض نمی‌کند.
در حالت بدون سقف عمق و
برای $0<\beta<1$،
ماتریس گذار این زنجیره
روی $\{o_s\}\cup\mathcal S_s$
برابر است با
\begin{equation}\label{cmp:eq:ch3-sequential-restart-matrix}
\boxed{
\widetilde{\mathsf P}_s=
\begin{pmatrix}
0&\rho_s\\
\beta\mathbf1&(1-\beta)\mathsf P_s
\end{pmatrix}.
}
\end{equation}
این ماتریس نیز سطری--تصادفی
است. از هر حالت داخلی با
احتمال مثبت به $o_s$
می‌رسیم و از $o_s$
تمام حالت‌های دست‌یافتنی
از توزیع آغاز قابل وصول‌اند؛
پس مجموعه حالت‌های
دست‌یافتنی یک طبقه
ارتباطی بازگشتی دارد.
بازگشت‌های دو و سه
انتقالی به $o_s$،
برای $0<\beta<1$
دوره آن را یک می‌کنند.
در نتیجه زنجیره
بازآغازشونده روی
این طبقه \emph{توزیع
مانای یکتا} دارد
و توزیع حالت‌های آن
به توزیع مانا همگرا
می‌شود؛ حتی اگر
$\mathsf P_s^t$
به‌تنهایی همگرا نشود.
به‌طور مشخص، اگر
$\widetilde\pi_s(o_s)$
جرم حالت آغاز و
$\widetilde\pi_s^{\rm int}$
جرم حالت‌های داخلی باشد،
\begin{equation}\label{cmp:eq:ch3-sequential-restart-stationary}
\boxed{
\begin{aligned}
\widetilde\pi_s(o_s)
&=\frac{\beta}{1+\beta},\\
\widetilde\pi_s^{\rm int}
&=\frac{\beta}{1+\beta}\,
\rho_s\bigl(I-(1-\beta)\mathsf P_s\bigr)^{-1},\\
\overline\pi_s^{\rm int}
&=\frac{\widetilde\pi_s^{\rm int}}
{1-\widetilde\pi_s(o_s)}
=\beta\sum_{t=1}^{\infty}
(1-\beta)^{t-1}\rho_s\mathsf P_s^{t-1}.
\end{aligned}
}
\end{equation}
$ \overline\pi_s^{\rm int} $
توزیع مانای \emph{مشروط
به حضور در حالت داخلی}
است؛ حالت آغاز فقط
نقش کنترل بازآغاز
را دارد و در
بازنمایی همسایگی
مشاهده محسوب نمی‌شود.
این صورت‌بندی
با ایده گشت تصادفی
همراه با بازآغاز
ارتباط دارد، اما
وزن‌دهی آماری روش
پیشنهادی را نیز باید
جداگانه در نظر گرفت.

\paragraph{همگرایی میانگین گشت‌های متوالی و تمایز آن از اشغال مانا.}
اگر $L_r$ طول قطعه $r$ام و
$J_{r,t}$ نمونه مشاهده‌شده
در گام $t$ام آن باشد،
بردار مشاهده نرمال‌شده
آن قطعه را با
\begin{equation}\label{cmp:eq:ch3-sequential-episode-average}
\boxed{
B_{s,r}=
\frac{\sum_{t=1}^{L_r}
\lambda^{t-1}Y_{J_{r,t},:}}
{\mathcal N_{L_r}(\lambda)},
\qquad
\widehat{\mathscr Y}_{s,M}
=\frac1M\sum_{r=1}^{M}B_{s,r}
\xrightarrow[M\to\infty]{\mathrm{a.s.}}
\mathscr Y_s=\mathbb E B_{s,1}.
}
\end{equation}
تعریف می‌کنیم.
با سیاست و داده ثابت،
بازآغاز کامل باعث می‌شود
قطعه‌ها هم‌توزیع و،
با تولید تصادفی مستقل
پس از هر بازآغاز،
مستقل باشند.
چون هر مؤلفه $B_{s,r}$
بین صفر و یک قرار دارد،
قانون قوی اعداد بزرگ
همگرایی تقریب تجربی
را تضمین می‌کند.
این بازنمایی دقیقاً همان
امید گشت‌های نرمال‌شده
در روابط
\eqref{cmp:eq:ch3-h-alpha}
و \eqref{cmp:eq:ch3-capped-alpha}
است؛ متوالی یا موازی
اجراشدن قطعه‌ها
مقدار امید را تغییر
نمی‌دهد.

بااین‌حال، \emph{میانگین
روی قطعه‌ها} را
نباید با \emph{میانگین
زمانی تمام حالت‌های
داخلی} یکی دانست:
در اولی سهم کل هر
گشت بعد از نرمال‌سازی
یکسان است؛ در دومی
گشت‌های طولانی‌تر
بازدیدهای بیشتری
به میانگین زمانی
اضافه می‌کنند.
ازاین‌رو به‌طور کلی
\[
\mathscr Y_s
=\sum_{t\ge1}\alpha_t(\beta,\lambda)
\rho_s\mathsf P_s^{t-1}\mathsf M_s^{(Y)}
\ne
\overline\pi_s^{\rm int}\mathsf M_s^{(Y)}.
\]
بنابراین \emph{توزیع
مانای زنجیره بازآغازشونده}
و \emph{بازنمایی همسایگی
مبتنی بر میانگین هر گشت}
هر دو یکتا هستند،
اما عموماً دو کمیت
متفاوت‌اند. اگر
سقف عمق $H$ اعمال شود،
برای بازنمایی همان
$L_r=\min(L,H)$ و
ضرایب $\alpha_t^{(H)}$
به کار می‌روند؛
برای نمایش بازآغاز
سقف‌دار به‌عنوان
زنجیره همگن نیز
باید شمارنده گام
را به فضای حالت
اضافه کنیم.
صورت اثبات و
تفاوت دو میانگین
در پیوست
\ref{app:ch3-policy-center-sequential-restart}
آمده است.


برای مراکز هدف دلخواه $\mathcal V^\star$، مسئله معکوس را می‌توان چنین نوشت:
\begin{equation}\label{cmp:eq:ch3-center-inverse-loss}
\boxed{\min_{z\in\mathcal Z}
\mathcal L_{\mathrm{inv}}(z)=
\frac1K\sum_{c=1}^{K}
\operatorname{JSD}\!\left(
\mathcal V_c(z;\beta,\lambda,H),\mathcal V_c^\star
\right).}
\end{equation}
وجود جواب با خطای صفر تضمین‌شده نیست. مجموعه
\begin{equation}\label{cmp:eq:ch3-attainable-centers}
\mathfrak C_{\beta,\lambda,H}
=\left\{\mathcal V(z;\beta,\lambda,H):z\in\mathcal Z\right\}
\subseteq(\Delta^{K-1})^K
\end{equation}
فضای مراکز قابل تولید تحت گراف، خانواده سیاست و پارامترهای عمق ثابت است؛ رسیدن دقیق به هدف تنها وقتی ممکن است که $\mathcal V^\star$ در این مجموعه قرار گیرد. در مقابل، توزیع پسین برچسب خود نمونه با توزیع \emph{برچسب همسایگی} یکسان نیست؛ بنابراین حتی مراکز حاصل از میانگین پسین بیزیِ طبقه‌ها نیز الزاماً در $\mathfrak C_{\beta,\lambda,H}$ قرار نمی‌گیرند. کران‌های ساختاری این فضا و وابستگی آن به $\beta,\lambda$ در پیوست~\ref{app:ch3-inverse-policy-proof} تشریح شده‌اند.




\subsection{سیاست‌های هم‌مرکز و نامعکوس‌پذیری نگاشت سیاست}

\paragraph{خانواده سیاست‌های هم‌مرکز و آزادی آرایش نمونه‌ها.}
برای یک ماتریس مراکز دست‌یافتنی
$\mathcal V^\star\in\mathfrak C_{\beta,\lambda,H}$،
\emph{مجموعه سیاست‌های هم‌مرکز} را تعریف می‌کنیم:
\begin{equation}\label{cmp:eq:ch3-fixed-center-fiber}
\boxed{
\mathfrak F(\mathcal V^\star)
=\{z\in\mathcal Z:
\mathcal F_{\beta,\lambda,H}(z)=\mathcal V^\star\}.
}
\end{equation}
برای تمام اعضای این مجموعه،
ماتریس مراکز و بنابراین هر معیار
وابسته فقط به مراکز، از جمله $I$،
ثابت می‌ماند:
\begin{equation}\label{cmp:eq:ch3-fixed-center-I}
\boxed{
z_1,z_2\in\mathfrak F(\mathcal V^\star)
\quad\Longrightarrow\quad
I(\mathcal V(z_1))=I(\mathcal V(z_2))
=I(\mathcal V^\star).
}
\end{equation}
اما این برابری هیچ استلزامی برای
یکسان‌بودن بازنمایی تک‌تک نمونه‌ها،
تراکم درون‌طبقه‌ای، حاشیه‌ها
یا تصمیم‌های طبقه‌بندی ندارد.
به بیان دقیق‌تر، اگر یک خانواده سیاست
بتواند دو عضو واقعاً متفاوت
با مراکز یکسان بسازد،
می‌توان در همان خانواده
مسئله کاهش $R$ \emph{یا}
کاهش زیان حاشیه را
\emph{با قید دقیق مراکز ثابت}
بررسی کرد:
\begin{equation}\label{cmp:eq:ch3-fixed-center-objectives}
\boxed{
\inf_{z\in\mathfrak F(\mathcal V^\star)}R(z)
\qquad\text{یا}\qquad
\inf_{z\in\mathfrak F(\mathcal V^\star)}
\mathcal L_{\mathrm{sig}}(z;\tau).
}
\end{equation}
این دو مسئله \emph{جایگزین‌های مستقل}
هستند و در اینجا هدف وزن‌دار
ترکیبی معرفی نمی‌کنیم.
به‌طور کلی، برای هر تابع هدف حقیقی
$f$ که روی مجموعه ناتهی
$\mathfrak F(\mathcal V^\star)$ تعریف شده باشد،
مسئله $\inf_{z\in\mathfrak F(\mathcal V^\star)}f(z)$
خوش‌تعریف است؛ اگر تابع از پایین کراندار باشد،
این اینفیمم متناهی است. بدون فرض تحقق مقدار بهینه
(مثلاً فشردگی مجموعه مجاز و پیوستگی تابع)،
نباید از وجود \emph{کمینه‌ساز} نتیجه گرفت.
اگر دو عضو این مجموعه مقدارهای متفاوتی از $f$
داشته باشند، عضو با مقدار کمتر نسبت به عضو دیگر
بهتر است؛ اما صرف وجود دو سیاست متمایز،
تفاوت مقادیر $f$ یا وجود جهت کاهشی موضعی
در هر نقطه را ثابت نمی‌کند.
وجود چند سیاست متمایز در این
مجموعه و امکان بهبود هر هدف،
به ساختار گراف و ظرفیت
خانواده سیاست وابسته است؛
از صرف تعریف مراکز میانگین
نمی‌توان وجود آن‌ها را نتیجه گرفت.
همچنین ثابت گذاشتن
$\mathcal V^\star$ صرفاً در
\emph{تابع زیان}، بدون کنترل
$\mathcal V(z)=\mathcal V^\star$،
به معنای ثابت ماندن واقعی
مراکز سیاست نیست.



\paragraph{نامعکوس‌پذیری یک‌به‌یک و تمایز پارامتر از سیاست مؤثر.}
از تعیین یکتای مراکز توسط
یک سیاست نمی‌توان نتیجه گرفت
که مراکز، سیاست سازنده را
به‌طور یکتا مشخص کنند.
در سطح پارامترهای خام،
برای هر سطر امتیازِ ورودی نرمال‌سازی با بیشینهٔ هموار،
افزودن یک ثابت مشترک به
تمام امتیازهای مجاز،
احتمال‌های انتقال را
تغییر نمی‌دهد؛ پس
ماتریس مراکز نیز عوض
نمی‌شود و نگاشت خام
$z\mapsto\mathcal V(z)$
یک‌به‌یک نیست.
این نمونه از عدم یکتایی،
\emph{صرفاً ناوردایی
پارامترسازی} است.
عدم یکتایی قوی‌تر،
یعنی دو سیاست با
احتمال‌های انتقال
واقعاً متفاوت اما
مراکز یکسان، به
درجه آزادی مؤثر و
ساختار نگاشت بستگی دارد؛
وجود آن برای هر گراف
از پنج فرض نتیجه نمی‌شود.
در چنین شرایطی
خانواده هم‌مرکز می‌تواند
بازنمایی‌های نمونه‌ای
و حاشیه‌های متفاوت
ایجاد کند، با آنکه
ماتریس مراکز یکسان
باقی می‌ماند.


\paragraph{شاهد سازنده برای عدم‌یکتایی مؤثر.}
عدم‌یکتایی مؤثر تنها یک امکان انتزاعی در فضای توزیع‌ها نیست:
در پیوست~\ref{app:ch3-inverse-policy-proof} یک مثال دقیق با
چهار نمونه، دو ویژگی موجود و $H=1$ ارائه شده است که در آن
دو سیاست \lr{shared} با احتمال‌های شروع متفاوت
$t=1/4$ و $t=3/4$، مراکز \emph{کاملاً یکسان}
اما بازنمایی‌های تک‌نمونه‌ای متفاوت می‌سازند
(رابطه~\eqref{eq:app-ch3-effective-nonunique-example}).
بنابراین وارون نگاشت سیاست \emph{مؤثر} به مراکز
در حالت کلی یک‌به‌یک نیست؛ این مثال نه عدم‌یکتایی
در هر گراف و هر افق را ثابت می‌کند و نه وجود سیاست بهتر
در هر فیبر هم‌مرکز را.



\subsection{بهبود مراکز و شرایط ضروری دستیابی به مراکز هدف}

\paragraph{روش دومرحله‌ای: بهبود مراکز و سپس بهبود سیاست با مراکز ثابت.}
نتیجه بنیادی روابط
\eqref{cmp:eq:ch3-fixed-center-I}
و \eqref{cmp:eq:ch3-fixed-center-objectives}
امکان جداسازی
\emph{کیفیت مراکز}
از \emph{کیفیت آرایش
نمونه‌ها در اطراف
مراکز} است.
در مرحله نخست،
یک سیاست $z^\star$
با هدف بهبود کیفیت
بازنمایی و هندسه مراکز
به‌دست می‌آوریم؛
مثلاً می‌توان از
بیشینه‌سازی $I$
یا یک معیار دیگر
که قبلاً معرفی شده
است بهره گرفت.
مراکز واقعاً
تولیدشده توسط آن
سیاست را به‌عنوان
هدف \emph{دست‌یافتنی}
ثابت می‌کنیم:
\[
\mathcal V^\star
=\mathcal F_{\beta,\lambda,H}(z^\star).
\]
در مرحله دوم،
تغییرات سیاست را
به مجموعه هم‌مرکز
محدود و زیان حاشیه
را بهینه می‌کنیم؛
در صورت نیاز،
پراکندگی درون‌طبقه‌ای
می‌تواند به‌عنوان
یک هدف مکمل
با ضریب نامنفی
در نظر گرفته شود:
\begin{equation}\label{cmp:eq:ch3-two-stage-fixed-center-optimization}
\boxed{
\inf_{z\in\mathfrak F(\mathcal V^\star)}
\Bigl[
\mathcal L_{\mathrm{sig}}(z;\tau)
+\gamma_R R(z)
\Bigr],
\qquad\gamma_R\ge0.
}
\end{equation}
با $\gamma_R=0$
فقط حاشیه بهینه می‌شود
و هیچ ضریب جدایی
در مرحله دوم لازم نیست؛
زیرا در صورت حفظ
\emph{دقیق} قید،
$I(\mathcal V(z))$
ثابت خواهد ماند.
انتخاب مراکز از خروجی
$z^\star$ سبب می‌شود
مجموعه مجاز مرحله دوم
دست‌کم یک عضو شناخته‌شده
داشته باشد، ولی
وجود یک سیاست
\emph{بهتر} در آن
تضمین نمی‌شود.
به‌علاوه، موفقیت
تقریبی حلگر با
حفظ دقیق مراکز
متفاوت است.
این چارچوب، مبنای
نظری مطالعه بهینه‌سازی
مقید سیاست است،
نه ادعای برتری همگانی
ترکیب معیارها. در این مرحله،
توان حلگر در یافتن سیاست بهتر،
کنترل واقعی قید مراکز، تغییر هندسه
بازنمایی‌ها و رابطه کاهش زیان با
دقت متوازن و تعمیم‌پذیری،
پرسش‌های محاسباتی و تجربی‌اند؛
کاهش زیان سیگموئیدی به‌تنهایی
بهبود \lr{BA} یا عملکرد آزمون را تضمین نمی‌کند.

\paragraph{کران‌های دوطرفهٔ ضروری برای دست‌یابی به مراکز.}
در این بحث داده و وزن‌های مثبت گراف ثابت‌اند، هر نمونه دست‌کم یک ویژگی دارد، $|D_b|\ge3$ برای هر ویژگی برقرار است و $H\ge1$، $0<\beta\le1$ و $0<\lambda\le1$ فرض می‌شوند؛ بنابراین همهٔ هسته‌های انتقالِ مجاز نرمال‌پذیرند. برای تمایز محدودیت ساختاریِ دست‌یابی از شکست یک حلگر عددی، \emph{کنترل‌گر انتخاب ویژگی} را آزادتر از خانوادهٔ واقعی سیاست در نظر می‌گیریم، بی‌آنکه هستهٔ انتقال را تغییر دهیم. این کنترل‌گر می‌تواند به نمونهٔ مبدأ $s$، شماره گام $t$، وضعیت فعلی و برچسب مورد بررسی $u$ وابسته باشد. هستهٔ ثابتِ بدون بازگشت به مبدأ $q_b^{(s)}(j\mid i)$ در رابطه~\eqref{cmp:eq:ch3-vertical-transition} و وزن‌های گشت سقف‌دار $\alpha_t^{(H)}$ در رابطه~\eqref{cmp:eq:ch3-capped-alpha} همان قرارداد مدل اصلی را دارند.

برای مبدأ $X_s$ و برچسب مقصد $u$، دو تابع
$G^-_{s,u,t}(i,a)$ و $G^+_{s,u,t}(i,a)$
را \emph{توابع ارزش کران‌ساز بلمن} می‌نامیم: آن‌ها به‌ترتیب کمترین و بیشترین \emph{سهم مورد انتظار برچسب $u$ از گام $t$ تا پایان گشت} را، پس از مشاهدهٔ نمونه $X_i$ با ویژگی ورودی $a$، در کنترل‌گر آزادشده تعیین می‌کنند. در این تعریف $i\ne s$، $y_i$ برچسب نمونهٔ فعلی، $A_{X_i}$ مجموعه ویژگی‌های مشاهده‌شدهٔ آن و $b\in A_{X_i}$ ویژگی انتخابیِ انتقال بعدی است. پاداش مشاهدهٔ گام $t$ برابر $\alpha_t^{(H)}\mathbf1\{y_i=u\}$ است. با شرط پایانی $G^\pm_{s,u,H+1}(i,a)=0$، \emph{بازگشت‌های کمینه و بیشینهٔ بلمن} برای $t=H,\ldots,1$ به‌صورت
\begin{equation}\label{cmp:eq:ch3-bellman-value-recursion}
\boxed{\begin{aligned}
G^-_{s,u,t}(i,a)
&=\alpha_t^{(H)}\mathbf1\{y_i=u\}
+\min_{b\in A_{X_i}}\sum_{j\ne i,s}
q_b^{(s)}(j\mid i)\,G^-_{s,u,t+1}(j,b),\\
G^+_{s,u,t}(i,a)
&=\alpha_t^{(H)}\mathbf1\{y_i=u\}
+\max_{b\in A_{X_i}}\sum_{j\ne i,s}
q_b^{(s)}(j\mid i)\,G^+_{s,u,t+1}(j,b).
\end{aligned}}
\end{equation}
محاسبه می‌شوند. در گام پایانی جملهٔ ادامه صفر است؛ انتخاب ویژگی تنها بر انتقال بعدی اثر می‌گذارد. اگرچه $a$ در نماد تابع ارزش باقی مانده است، در \emph{کنترل‌گر آزادشده} مجموعه اعمال مجاز به آن وابسته نیست؛ چنین استقلالی برای سیاست واقعی الزاماً برقرار نیست.

\emph{حرکت نخست} از مبدأ باید جداگانه اعمال شود. با تعریف
$B_b(s,j):=q_b^{(s)}(j\mid s)$
به‌عنوان توزیع مقصد نخست پس از انتخاب ویژگی $b$، \emph{کران‌های نمونه‌محور بلمن} را چنین به‌دست می‌آوریم:
\begin{equation}\label{cmp:eq:ch3-bellman-origin-bounds}
\boxed{\begin{aligned}
\ell_{s,u}
&=\min_{b\in A_{X_s}}\sum_{j\ne s}
B_b(s,j)\,G^-_{s,u,1}(j,b),\\
U_{s,u}
&=\max_{b\in A_{X_s}}\sum_{j\ne s}
B_b(s,j)\,G^+_{s,u,1}(j,b).
\end{aligned}}
\end{equation}
در اینجا $\ell_{s,u}$ و $U_{s,u}$ به‌ترتیب \emph{کران پایین و بالای سهم برچسب $u$ در بازنمایی نمونهٔ مبدأ $s$} هستند. برای هر طبقه $c$، با $C_c=\{s:y_s=c\}$ و $|C_c|>0$، میانگین این دو کمیت \emph{کران‌های مختصه‌ای مراکز} را می‌دهد:
\begin{equation}\label{cmp:eq:ch3-center-component-upper}
\boxed{\begin{gathered}
L_{c,u}:=\frac{1}{|C_c|}\sum_{s\in C_c}\ell_{s,u},
\qquad
\overline U_{c,u}:=\frac{1}{|C_c|}\sum_{s\in C_c}U_{s,u},\\
L_{c,u}\le\mathcal V_c(z;\beta,\lambda,H)(u)
\le\overline U_{c,u},
\qquad c,u\in\{1,\ldots,K\}.
\end{gathered}}
\end{equation}
که در آن $\mathcal V_c(z;\beta,\lambda,H)(u)$ مؤلفهٔ برچسب $u$ از مرکز طبقه $c$ تحت سیاست واقعی $z$ است. این نامساوی‌ها \emph{شرط لازم} هستند: اگر حتی یک مؤلفه از بازهٔ خود خارج باشد، آن هدف در داده، گراف، افق و خانوادهٔ سیاست ثابت دست‌یافتنی نیست؛ اما حضور همه مؤلفه‌ها در بازه‌ها، وجود یک سیاست واحد را تضمین نمی‌کند. زیرا کنترل‌گر آزادشده ممکن است برای هر $s$ و $u$ انتخاب متفاوتی داشته باشد، درحالی‌که سیاست واقعی باید قیود اشتراک پارامترهای خود را رعایت کند. شرط ناتهی‌بودن مقصدها و اثبات این شمول با استقرای معکوس در پیوست~\ref{app:ch3-inverse-policy-proof}، پیرامون بازگشت~\eqref{eq:app-ch3-feasibility-Bellman}، آمده‌اند.

مجموعهٔ واقعی مراکز $\mathfrak C_{\beta,\lambda,H}$ در رابطه~\eqref{cmp:eq:ch3-attainable-centers} به \emph{دادهٔ مرجع، وزن‌های گراف و خانوادهٔ سیاست} وابسته است. با آشکارکردن این وابستگی‌ها، همان مجموعه را چنین می‌نویسیم:
\begin{equation}\label{cmp:eq:ch3-attainable-centers-data-explicit}
\mathfrak R(D;\beta,\lambda,H,\mathcal Z)
=\{\mathcal V(z;D,\beta,\lambda,H):z\in\mathcal Z\}.
\end{equation}
تمام اعضای این مجموعه در بازه‌های ضروری رابطه~\eqref{cmp:eq:ch3-center-component-upper} قرار می‌گیرند. نه دست‌یافتنی‌بودن هدف بیزی در یک داده به همهٔ داده‌ها تعمیم دارد و نه برابرشدن مراکز بیزی و گشت، برابری پسین‌های نمونه‌ای یا بهبود دقت متوازن را تضمین می‌کند.

برای هدف $p=\mathcal V_c^\star$، اگر مختصهٔ $u$ بیرون بازهٔ $[L_{c,u},\overline U_{c,u}]$ باشد، برای هر مرکز مجاز $q$ از نامساوی داده‌پردازی و کران پینسکر با لگاریتم طبیعی نتیجه می‌شود
$\operatorname{JSD}(p,q)\ge \frac12\operatorname{dist}(p_u,[L_{c,u},\overline U_{c,u}])^2$.
صورت دقیق‌تر برنولی در پیوست اثبات شده است و \emph{کران‌های یک سطر را نمی‌توان بدون برهان روی مختصات جمع کرد.}
این قضیه در حساب دقیق برقرار است؛ تفسیر خروجی محاسبات شناور به‌عنوان گواه قطعی، مستلزم محصورسازی مستقل خطای هسته، ضرایب، بازگشت‌ها و هدف است. صرف شکست حلگر یا مشاهده شکاف در \lr{float64} چنین گواهی نیست.



\subsection{اصلاح سیاست بر پایه ژاکوبین و حفظ مراکز}

\paragraph{اصلاح تکراری از روی ژاکوبین مراکز.}
برای حذف مؤلفه احتمالیِ وابسته در هر سطر مرکز، بردار
$u(z)=\operatorname{vec}(\mathcal V(z)_{:,1:K-1})\in\mathbb R^r$
را با $r=K(K-1)$ تعریف می‌کنیم؛ در مسئله شرطیِ مستقل برای هر طبقه، $r=K-1$ کافی است. ژاکوبین این نگاشت در ناحیه‌ای که ماسک ویژگی‌ها ثابت است برابر است با
$J(z)=\partial u(z)/\partial z\in\mathbb R^{r\times p}$.
تغییرات کوچک مراکز را می‌توان با
$u(z+\delta)\simeq u(z)+J(z)\delta$
خطی‌سازی کرد. یک اصلاح میراشده از حل مسئله کمترین مربعات منظم‌شده به‌دست می‌آید:
\begin{equation}\label{cmp:eq:ch3-inverse-GN}
\boxed{\delta=
J^\top(JJ^\top+\mu I_r)^{-1}
\bigl(u^\star-u(z)\bigr),\qquad \mu>0.}
\end{equation}
این رابطه معادل جواب کمینه‌ساز
$\|J\delta-(u^\star-u(z))\|_2^2+\mu\|\delta\|_2^2$
است. می‌توان از سیاست حاصل از تقریب یک‌گامی $\mathsf Q$ در رابطه~\eqref{cmp:eq:ch3-feature-one-step-Q} به‌عنوان نقطه شروع استفاده کرد و سپس با محاسبه مراکز \emph{دقیق چندگامی}، اصلاح را به‌صورت تکراری انجام داد. برای کنترل خطای خطی‌سازی، تنها گامی پذیرفته می‌شود که پس از جست‌وجوی خط، مقدار واقعی $\mathcal L_{\mathrm{inv}}$ را کاهش دهد؛ کاهش، نسبت به مقدار پیش از همان گام سنجیده می‌شود. این قاعده کاهش خطای مراحل پذیرفته‌شده را تضمین می‌کند، اما تضمینی برای دستیابی به جواب سراسری، بازیابی یکتای سیاست یا نزدیک‌شدن به هر هدف دلخواه نیست. مشتق رابطه ماتریسی و نحوه محاسبه ژاکوبین در پیوست~\ref{app:ch3-inverse-policy-proof} آمده است.


\paragraph{درجه آزادی محلی و تفاوت قید خطی با قید دقیق.}
برای نگاشت مؤلفه‌های مستقل مراکز
$u(z)$، تغییر کوچک $\Delta z$
شرط خطی مرتبه اول
\begin{equation}\label{cmp:eq:ch3-fixed-center-tangent}
\boxed{
J(z)\Delta z=0
}
\end{equation}
را برای ثابت ماندن مراکز
\emph{تا مرتبه اول} برآورده می‌کند.
پس فضای صفر ژاکوبین،
جهت‌های بالقوه حرکت سیاست
بدون جابه‌جایی مرتبه اول مراکز را
مشخص می‌کند. اگر پارامترسازی را
پس از کنارگذاشتن ناوردایی‌های
نرمال‌سازی با بیشینهٔ هموار و مؤلفه‌های بی‌اثر،
با $p_{\mathrm{eff}}$ پارامتر مؤثر
در نظر بگیریم، بعد فضای صفر
برابر $p_{\mathrm{eff}}-\operatorname{rank}J_{\mathrm{eff}}$
خواهد بود؛ زیرا
$\operatorname{rank}J_{\mathrm{eff}}\le K(K-1)$،
تعداد پارامترهای مؤثر نسبت به
تعداد مؤلفه‌های مستقل مراکز
در ظرفیت این خانواده اهمیت دارد.
بااین‌حال زیادبودن تعداد
پارامترهای خام به‌تنهایی
وجود تغییرات \emph{مؤثر}
روی بازنمایی نمونه‌ها را
تضمین نمی‌کند.

شرط
\eqref{cmp:eq:ch3-fixed-center-tangent}
فقط تقریب خطی است:
حرکت متناهی در جهت
فضای صفر ممکن است به دلیل
جملات مرتبه دوم مراکز را
تغییر دهد. در نقاطی که
رتبه نگاشت در یک همسایگی
ثابت و دامنه پارامترها
باز باشد، قضیه رتبه ثابت
ساختار موضعی مجموعه
سیاست‌های هم‌مرکز را
به‌صورت یک خمینه با
بعد $p_{\mathrm{eff}}-\operatorname{rank}J_{\mathrm{eff}}$
توصیف می‌کند؛ در نقاط تکین
چنین نتیجه‌ای خودکار نیست.
برای حفظ \emph{دقیق} مراکز
باید پس از هر تغییر سیاست،
مراکز را دوباره محاسبه
و خطای قید را کنترل کرد
و در صورت نیاز از اصلاح
یا فرافکنی غیرخطی کمک گرفت.
اثبات و محدوده اعتبار
این تحلیل در پیوست
\ref{app:ch3-inverse-policy-proof}
آمده است.



\paragraph{جهت کاهش مرتبه اول با مراکز ثابت.}
برای تابع هدف مشتق‌پذیر
$f(z)=\mathcal L_{\mathrm{sig}}(z;\tau)
+\gamma_RR(z)$
و ژاکوبین $J(z)$
در مختصات مستقل مراکز،
فرافکن متعامد
روی $\ker J(z)$
را با وارون مور--پنروز
به‌صورت
\begin{equation}\label{cmp:eq:ch3-fixed-center-projection}
\boxed{
\Pi_J=I_p-J^\top(JJ^\top)^\dagger J,
\qquad
\Delta z=-\eta\,\Pi_J\nabla_z f(z),
\quad\eta>0
}
\end{equation}
تعریف می‌کنیم.
از ویژگی‌های فرافکن
نتیجه می‌شود
\begin{equation}\label{cmp:eq:ch3-fixed-center-projection-descent}
\boxed{
J\Delta z=0,
\qquad
\nabla_z f(z)^\top\Delta z
=-\eta\|\Pi_J\nabla_z f(z)\|_2^2\le0.
}
\end{equation}
اگر مؤلفه فرافکنده
گرادیان ناصفر باشد،
جهت فوق یک جهت
کاهشی مرتبه اول
با تغییر صفرِ
مرتبه اول مراکز
است. اگر رتبه ژاکوبین
در همسایگی ثابت باشد
و نقطه درون فضای
مجاز پارامترها قرار
داشته باشد، می‌توان
این جهت را مماس
یک مسیر موضعی
با \emph{مراکز
دقیقاً ثابت} دانست.
بااین‌حال گام خطی
$z+\Delta z$
عموماً قید غیرخطی
را به‌طور دقیق
حفظ نمی‌کند؛
پس اجرای عملی
نیازمند تصحیح قید،
محاسبه مجدد مراکز
و کنترل خطای
باقی‌مانده است.
صفر شدن گرادیان
فرافکنده نیز
صرفاً نبود جهت
کاهشی \emph{مرتبه اول}
در فضای مماس را
بیان می‌کند،
نه تضمین بهینگی
سراسری. اثبات در
پیوست
\ref{app:ch3-policy-center-projection}
آمده است.

\paragraph{پروتکل محاسباتیِ پیشنهادی برای گام هم‌مرکز.}
برای تبدیل جهت مماس به یک آزمون عددی قابل‌راستی‌آزمایی،
مختصات مستقل مراکز را با $u(z)$ و هدف را با
$u^\star=u(z^\star)$ مشخص می‌کنیم. باقی‌ماندهٔ واقعی قید
و تلرانس عددی مثبت را به‌ترتیب چنین تعریف می‌کنیم:
\begin{equation}\label{cmp:eq:ch3-fixed-center-residual}
r_C(z):=\|u(z)-u^\star\|_\infty,
\qquad \varepsilon_C>0.
\end{equation}
از سیاست آغازینِ دست‌یافتنی $z^\star$ شروع کرده و
در نقطه‌ای با $r_C(z)\le\varepsilon_C$،
یک جهت آزمایشی $v=-\Pi_J\nabla f(z)$
می‌سازیم؛ این جهت فقط $Jv=0$ را در مرتبه اول
برآورده می‌کند. برای هر طول گام آزمایشی $\eta$،
ابتدا $z_{\rm trial}=z+\eta v$ محاسبه می‌شود،
سپس با یک یا چند اصلاح میراشدهٔ ژاکوبینی
\[
\delta_{\rm corr}
=J(z_{\rm trial})^\top
\bigl(J(z_{\rm trial})J(z_{\rm trial})^\top+
\mu I\bigr)^{-1}
\bigl(u^\star-u(z_{\rm trial})\bigr)
\]
و ارزیابی دوبارهٔ \emph{امید دقیق}،
باقی‌مانده قید کنترل می‌شود.
این تصحیح غیرخطی به دلیل تغییر ژاکوبین
و رتبهٔ احتمالی آن، هیچ تضمین جهانی
برای رسیدن به فیبر نمی‌دهد.
تنها گامی پذیرفته می‌شود که هم
$r_C(z_{\rm new})\le\varepsilon_C$
و هم $f(z_{\rm new})<f(z)$
را پس از محاسبهٔ واقعی برآورده کند؛
در غیر این صورت طول گام کاهش می‌یابد.
در برابری رقیب‌ها که هدف حاشیه
مشتق‌پذیر نیست، باید زیرگرادیان مناسب
یا تقریب هموار مستقل از $\tau$ به کار برد.
معیار توقف می‌تواند کوچک‌شدن
$\|\Pi_J\nabla f\|$، نبود گام پذیرفته‌شده،
رسیدن به بودجه، یا عبور باقی‌مانده از
تلرانس باشد. بنابراین
$\varepsilon_C$ فقط یک \emph{تلرانس عددی}
است و حفظ \emph{ریاضیِ دقیق} مراکز
از $r_C\le\varepsilon_C$ نتیجه نمی‌شود.
این الگوریتم \emph{پیشنهادی} است؛
ایستایی مرتبه اول نیز تضمین کمینه سراسری،
بهبود دقت یا موفقیت بازسازی معکوس نیست.



\subsection{هزینهٔ محاسبات و مرز تضمین‌های نظری}

\paragraph{هزینه هر اصلاح.}
فرض کنید $B=e_v+e_h$ و $C_{\mathrm f}=O(nHB)$ هزینهٔ یک ارزیابی کامل مراکز با رابطهٔ ماتریسی باشد. برای سیاست مشترک با $r=K(K-1)$ خروجی مستقل، ساخت ژاکوبین با مشتق‌گیری معکوسِ هر خروجی حداکثر به مرتبه $O(rC_{\mathrm f})$ عملیات نیاز دارد؛ با نگهداری نمودار محاسبه و صرف‌نظر از ضرایب ثابت. پس از تشکیل $J$، ساخت $JJ^\top$ و حل دستگاه رابطه~\eqref{cmp:eq:ch3-inverse-GN} به‌ترتیب $O(r^2p)$ و $O(r^3)$ هزینه دارند. اگر جست‌وجوی خط $b_{\mathrm{ls}}$ ارزیابی دقیق انجام دهد، هزینه یک اصلاح در این پیاده‌سازی برابر است با:
\begin{equation}\label{cmp:eq:ch3-inverse-complexity}
\boxed{O\!\left((r+b_{\mathrm{ls}})nH(e_v+e_h)
+r^2p+r^3\right),\qquad
p=d(d+1).}
\end{equation}
در حالت شرطیِ $Z_c$، برای هر طبقه جداگانه $n_c=|C_c|$ و $r=K-1$ قرار می‌گیرند و هزینه روی طبقه‌ها جمع می‌شود. حافظه ژاکوبین $O(rp)$ است؛ حافظه میانجی مشتق‌گیری به شیوه ذخیره/بازمحاسبه وضعیت‌ها وابسته است و باید جدا از حافظه ساختاری گراف سنجیده شود. ازاین‌رو تعداد اصلاح‌ها، ارزیابی‌های تابع هدف، زمان واقعی و حافظه اوج، همگی در مقایسه محاسباتی روش‌ها اهمیت دارند. ارزیابی کارایی و محدودیت‌های این ابزار در فصل~\ref{chap:results} گزارش خواهد شد.


در نهایت، آنچه از
پنج فرض و ساختار
مارکوفی مستقیماً
نتیجه می‌شود،
وجود ماتریس گذار
معتبر و تعیین
یکتای بازنمایی
وزن‌دار برای
سیاست ثابت است،
نه ارگودیک بودن
الزامی زنجیره
بعد از حذف مبدأ.
از سوی دیگر،
روش دومرحله‌ای
می‌تواند با حفظ
هندسه مراکز،
آزادی باقیمانده
سیاست را برای
بهبود حاشیه
به کار گیرد؛
اما مقدار این آزادی
و امکان بهبود
باید برای خانواده
سیاست و گراف
مورد بحث بررسی
شوند.

\paragraph{توانمندی‌های روش محاسبه و مرز تضمین‌های نظری.}
نمایش بلوکی و بازگشت‌های
ماتریسی این فصل،
برای گشت سقف‌دار و
گراف ثابت، امید
بازنمایی‌ها و مراکز را
بدون خطای نمونه‌گیری
مونت‌کارلو محاسبه می‌کنند.
کران برنامه‌ریزی پویا
می‌تواند بعضی اهداف
\emph{ناممکن} را پیش از
بهینه‌سازی حذف کند و
ژاکوبین، حساسیت محلی
مراکز به سیاست را
در اختیار قرار می‌دهد.
در مقابل، هیچ‌یک از
این ابزارها به‌تنهایی
قابل‌دستیابی بودن هر
ماتریس مرکزی، یکتایی
سیاست بازیابی‌شده،
وجود جهت بهبود حاشیه
در خانواده هم‌مرکز،
همگرایی سراسری حلگر
یا تعمیم کیفیت تصمیم
به نمونه‌های جدید را
تضمین نمی‌کنند.
در سیاست‌های وابسته
به نمونه یا طبقه نیز
باید هزینه تولید امتیازها،
قابلیت تعریف سیاست برای
نمونه جدید و رعایت
محدودیت‌های اطلاعات
مجاز هنگام آزمون را
مستقلاً بررسی کرد.
این مرزبندی، ابزارهای
تحلیلی روش تحقیق را
از ادعاهای کارایی تجربی
جدا نگه می‌دارد.



\section{تحلیل نظری ویژگی‌های نویزی، زائد و وابسته}
\label{cmp:sec:ch3-noise-feature-analysis}
\label{cmp:sec:ch3-noise-and-outliers}% برچسب سازگاری با ارجاعات پیشین
در تحلیل رابطهٔ سیاست و مراکز،
گراف مرجع، فضای ویژگی‌ها و
خانوادهٔ سیاست ثابت بودند و
می‌پرسیدیم کدام مراکز از
این سازوکار قابل‌دستیابی‌اند.
اکنون فرض ثبات ساختار ورودی
را از زاویه‌ای دیگر بررسی
می‌کنیم: افزودن ستونی غیرآموزنده،
تغییر مقادیر یک ویژگی یا تغییر
الگوی مشاهدهٔ آن ممکن است
مجموعهٔ انتخاب‌های مجاز،
وزن یال‌های گراف و در نتیجه
مسیرهای گشت و بازنمایی‌ها
را تغییر دهد. حتی اگر
سیاست بدون تغییر بماند،
از ثبات مراکز یا عملکرد
طبقه‌بند نمی‌توان نتیجه گرفت.
بنابراین باید \emph{ماهیت تغییر
داده} را از \emph{واکنش سیاست
و کیفیت تصمیم} جدا کنیم.

نویز در یادگیری آماری به پدیده‌های متفاوتی
اطلاق می‌شود: \emph{نویز اندازه‌گیری}
مقادیر ویژگی‌ها را مخدوش می‌کند؛
\emph{نویز برچسب} به نادرستی یا آشفتگی
برچسب‌های آموزشی مربوط است؛
\emph{ویژگی زائد} ممکن است با وجود سایر
ویژگی‌ها اطلاعات افزوده نداشته باشد؛
و \emph{فقدان ویژگی} به مشاهده‌نشدن
مقدار آن اشاره می‌کند و لزوماً نویز نیست.
به‌علاوه، بی‌اطلاعی حاشیه‌ای یک ویژگی
از برچسب، بی‌اطلاعی آن در ترکیب
با سایر ویژگی‌ها را نتیجه نمی‌دهد.
ازاین‌رو، این وضعیت‌ها نباید
با یک مفهوم واحد جایگزین شوند.

ابتدا گونه‌های متفاوتِ بی‌اطلاعی، افزونگی،
تکرار، خطای شِمای داده و فقدان مقدار را
از نظر مفهوم آماری و ساختار گراف
تفکیک می‌کنیم. \emph{دامنهٔ تحلیل کمیِ
شاخص‌های مواجهه و نتایج موجود فصل چهارم}
محدودتر است: افزودن ستون‌های نویزیِ
دارای مرجع معلوم در طرح تولید داده.
چنین ستون‌هایی مجموعهٔ انتخاب‌ها،
نرمال‌سازی سیاست و مسیرهای بالقوهٔ
گشت نمونه--ویژگی را تغییر می‌دهند،
اما نه افت دقت را تضمین می‌کنند و نه
به‌معنای تشخیص خودکار نویز هستند.
سنجه‌های مواجهه با این ستون‌ها
نیز معیار مستقیم خطای پیش‌بینی نیستند.
سایر گونه‌ها در این فصل به‌عنوان
مسئله‌های متمایز و قابل‌آزمون معرفی
می‌شوند، نه نتایج تجربیِ به‌دست‌آمده.



\subsection{تمایز ویژگی‌های نویزی، زائد، تکراری و فقدان داده}
\label{cmp:subsec:ch3-feature-error-types}
برای توصیف خطاهای داده، باید دست‌کم سه پرسش را
از هم جدا کرد: آیا \emph{مقدارِ مشاهده‌شده}
صحت دارد، آیا \emph{ستونِ ویژگی} اطلاعات تازه‌ای
برای پیش‌بینی فراهم می‌کند، و آیا
\emph{نام ستون و ماسک مشاهده} به‌درستی
متغیر موردنظر را نمایش می‌دهند؟
یک ویژگی ممکن است در یک معنا زائد و در
معنایی دیگر اطلاع‌رسان باشد؛ همچنین
نبودن مقدار یک ویژگی لزوماً به معنای
نویزدار بودن آن نیست. طبقه‌بندی زیر
بر اساس این تمایزهاست و گونه‌ها
الزاماً مانعةالجمع نیستند.

\paragraph{ویژگی نامرتبط یا ستون نویزیِ فاقد اطلاعات افزوده.}
یک ستونِ افزوده‌شده ممکن است در قرارداد
مولد داده مستقل از برچسب $C$ و کل بلوک
ویژگی‌های اطلاع‌رسان باشد؛ نمونهٔ ساده،
ستونی از اعداد تصادفی است که مستقل از
متغیرهای مولد برچسب تولید می‌شود.
این حالت را از «بی‌اطلاعی حاشیه‌ای»
یک ویژگی از برچسب، که در
زیربخش~\ref{cmp:subsec:ch3-noise-selection}
صورت‌بندی می‌شود، متمایز می‌کنیم. اگر $A_0$ مجموعه‌ای از
ویژگی‌های دیگر و $X^{(A_0)}$ بردار مقادیر
آن‌ها باشد، \emph{اطلاعات افزودهٔ پیش‌بینی}
ستون $a$ در حضور آن ویژگی‌ها با
اطلاعات متقابل شرطی سنجیده می‌شود
\cite{CoverThomas2006}:
\begin{equation}\label{cmp:eq:ch3-feature-conditional-information}
I\bigl(C;X^{(a)}\mid X^{(A_0)}\bigr)=0,
\qquad A_0\subseteq A_D\setminus\{a\}.
\end{equation}
این رابطه به انتخاب مجموعهٔ شرطی $A_0$
وابسته است. بی‌اطلاعی \emph{حاشیه‌ای}
$I(C;X^{(a)})=0$ به‌تنهایی
بی‌اطلاعی \emph{شرطی} را نتیجه نمی‌دهد:
مثلاً در تابع برچسبی از نوع
\lr{XOR} ممکن است هر ورودی به‌تنهایی
مستقل از برچسب باشد، اما با دانستن
ورودی دیگر ضروری شود.
در مقابل، مستقل‌بودن ستون از
\emph{تمام بلوک اطلاع‌رسان و برچسب}
فرض قوی‌تری است.
بنابراین «ستون نویزی» در آزمایش‌های
دارای مولد معلوم، یک
\emph{قرارداد مبتنی بر طراحی داده} است
و از پایین‌بودن ارتباط نمونه‌ای با برچسب
یا کم‌انتخاب‌شدن توسط سیاست استنتاج نمی‌شود.
حتی ستون واقعاً نامرتبط می‌تواند
با افزودن گره و یال، مسیرهای گشت را تغییر دهد.

\paragraph{ویژگی زائد یا وابسته، بدون اطلاعات شرطی تازه.}
ویژگی $b$ ممکن است به‌تنهایی برای
تشخیص طبقه مفید باشد، اما پس از مشاهدهٔ
مجموعه‌ای دیگر از ویژگی‌ها، اطلاعات
جدیدی دربارهٔ $C$ ندهد؛ در این صورت
رابطه~\eqref{cmp:eq:ch3-feature-conditional-information}
برای $a=b$ و مجموعهٔ شرطیِ مناسب
برقرار است. این \emph{افزونگی اطلاعاتی}
با نویز مستقلِ فاقد اطلاعات یکسان نیست.
همچنین هم‌بستگی بالا یا وجود رابطهٔ
تقریبی میان ستون‌ها به‌تنهایی
اثبات‌کنندهٔ افزونگی کامل نیست:
ممکن است اختلاف کوچک دو ویژگی برای
برچسب تعیین‌کننده باشد.

یک صورتِ قوی‌تر، \emph{وابستگی قطعی} است.
برای مثال، اگر در تمام نمونه‌ها و
روی فضای مقادیرِ تعریف‌شده داشته باشیم
\begin{equation}\label{cmp:eq:ch3-feature-deterministic-redundancy}
\boxed{
X^{(3)}=2X^{(1)}-X^{(2)}
\quad\Longrightarrow\quad
I\bigl(C;X^{(3)}\mid X^{(1)},X^{(2)}\bigr)=0.
}
\end{equation}
در قرارداد متغیرهای \emph{گسسته}
(یا پس از یک گسسته‌سازی سازگار)،
همچنین
$H(X^{(3)}\mid X^{(1)},X^{(2)})=0$
است؛ اما برای متغیرهای پیوسته،
\emph{آنتروپی تفاضلی شرطی} را نمی‌توان
به‌جای این گزاره برابر صفر قرار داد.
بیان عمومی صحیح آن است که توزیع شرطی
$X^{(3)}$ با دانستن دو متغیر اول
به یک مقدار قطعی فرو می‌ریزد.
این رابطه نشان نمی‌دهد که
$X^{(3)}$ به‌تنهایی از برچسب مستقل است،
و همچنین ثابت نمی‌کند که گشت IMC
با افزودن آن بدون تغییر می‌ماند:
ستون تازه یک زیرگراف عمودی و
گزینهٔ انتخاب افقی تازه می‌سازد.

\paragraph{کپی کامل در برابر هم‌بستگی یا وابستگی تابعی.}
اگر دو ستون متمایز $a$ و $b$ در
\emph{تمام} نمونه‌ها مشاهده شوند و
$a(X_i)=b(X_i)$ برای همهٔ $i$،
ستون دوم یک \emph{کپی دقیق} از اولی است.
این حالت نمونهٔ خاص وابستگی قطعی است،
ولی با مثال ترکیب خطی چند ستون تفاوت
ساختاری دارد: در کپی کامل، دو مسیر
ویژگی با مقدار یکسان ایجاد می‌شود؛
در وابستگی تابعی، لزوماً مقادیر یا
وزن‌های کرنل دو ویژگی برابر نیستند.
وجود مقادیر برابر نیز برابری
متریک‌ها، کرنل‌ها و logits دو ستون را
به‌طور خودکار تضمین نمی‌کند.
بنابراین کپی حتی با وجود بی‌افزودگی
اطلاعات پیش‌بینی، ممکن است به‌علت
\emph{شمارش مضاعف یک مفهوم به‌صورت دو
گزینهٔ انتخاب}، سهم احتمال مسیرهای
گشت را تغییر دهد.

\paragraph{شکست شِما و ستون‌های مترادف با ماسک مکمل.}
حالت متفاوتی زمانی رخ می‌دهد که دو منبع
یک متغیر فیزیکی واحد، مثلاً «سن» در
یک تاریخ مرجع مشخص، را با دو نام ستون
\lr{age} و «سن» ثبت کرده باشند، ولی
فرایند تجمیع آن‌ها را دو ویژگی مستقل
تلقی کند. فرض کنیم $a$ به \lr{age}
و $b$ به «سن» اشاره کند و برای هر
نمونه $X_i$، ماسک مشاهده را با
$M_{i,a}=\mathbf1\{a\in A_{X_i}\}$
تعریف کنیم؛ در سناریوی مکمل داریم:
\begin{equation}\label{cmp:eq:ch3-feature-schema-complementary-mask}
\boxed{
M_{i,a}+M_{i,b}=1,\qquad
M_{i,a}M_{i,b}=0,\qquad i=1,\ldots,n.
}
\end{equation}
در این قرارداد، $M_{i,a}=1$ نشانهٔ
\emph{وجود} مقدار است؛ شاخص $m_{ij}$
در بحث آزمون جایگشتیِ بعدی، برعکس،
مقدار یک را برای \emph{فقدان} به کار می‌برد.
در این حالت هیچ نمونه‌ای هر دو ستون را
هم‌زمان ندارد، ولی اگر هر دو واقعاً
\emph{همان متغیر، با همان واحد و
زمان اندازه‌گیری} باشند، آن متغیر
در سطح مفهومی برای تمام نمونه‌ها
مشاهده شده است. خانه‌های خالیِ
ظاهریِ این دو ستون در جدول، به معنای فقدان
واقعیِ مقدار سنِ آن نمونه نیستند.
با بررسی معنای ستون‌ها و تبدیل
واحدها، صورت درست ادغام می‌تواند
یک ویژگی واحد $\widetilde a$ باشد:
\begin{equation}\label{cmp:eq:ch3-feature-schema-merge}
\widetilde a(X_i)=
\begin{cases}
a(X_i),&M_{i,a}=1,\\
b(X_i),&M_{i,b}=1.
\end{cases}
\end{equation}
این ادغام تنها پس از \emph{احراز
هم‌معنایی و سازگاری مقیاس} مجاز است؛
دو اندازه‌گیری سن در زمان‌های متفاوت،
یا دو سنجهٔ هم‌بسته اما متفاوت را
نمی‌توان صرفاً به‌دلیل شباهت نام
یکی دانست. اگر دو ستون گاهی
هم‌مشاهده شوند، سازگاری مقدارهای
مشترک و قاعدهٔ حل تعارض نیز
باید جداگانه تعیین شود.
این پدیده را می‌توان
\emph{خطای شناسایی و یکپارچه‌سازی شِما}
نامید، نه «کپی کامل مشاهده‌شده»
و نه «ویژگی نویزی مستقل».

\paragraph{مقدار واقعاً مفقود و سازوکار فقدان.}
در فقدان واقعی، برای یک متغیر
تعریف‌شده و مرتبط، مقدار لازم برای
نمونه واقعاً مشاهده نشده است؛
در حالت شکست شِما ممکن است همان
مقدار با نام دیگری ثبت شده باشد.
همچنین یک ویژگی ممکن است اساساً
برای یک نمونه تعریف‌پذیر نباشد،
که این \emph{نامربوط‌بودن ساختاری}
با گم‌شدن یک اندازه‌گیری موجود
متفاوت است. در گراف حاضر، هر دو
می‌توانند به حذف گرهٔ نمونه--ویژگی
منجر شوند؛ صرفِ ماسک نمی‌تواند
علتِ غیاب را تشخیص دهد.
جایگذاری یک مقدار تخمینی
(\lr{imputation}) نیز عملی مستقل
است و نباید بی‌ذکر قرارداد، با
مقدار مشاهده‌شده یا با حذف گره
یکسان گرفته شود.

برای توصیف آماری \emph{الگوی}
فقدان داده، سازوکارهای کلاسیک
برگرفته از صورت‌بندی داده‌های مفقود
\cite{Rubin1976MissingData}
قابل‌تفکیک‌اند: در \lr{MCAR}
وقوع فقدان از مقادیر کامل داده
مستقل است؛ در \lr{MAR} وابستگی
به متغیرهای مشاهده‌شده ممکن است،
اما پس از شرط‌کردن بر اطلاعات
مشاهده‌شدهٔ مناسب، به مقادیر
پنهانِ همان الگوی فقدان وابسته نیست؛
در \lr{MNAR} حتی پس از چنین
شرط‌گذاری‌ای، احتمال فقدان ممکن است
به خودِ مقدار مشاهده‌نشده وابسته بماند.
این دسته‌بندی دربارهٔ سازوکار
تولید ماسک است، نه دربارهٔ
اطلاع‌رسان‌بودن ویژگی یا کیفیت
سیاست؛ تمایز \lr{MAR} و
\lr{MNAR} نیز عموماً از دادهٔ
مشاهده‌شده به‌تنهایی قابل احراز
نیست. در نتیجه آموزش IMC با
ماسک‌های معلوم، به‌خودی‌خود
تشخیص سازوکار فقدان یا تصحیح
سوگیری آن را تضمین نمی‌کند.

\paragraph{نویز اندازه‌گیری، برچسب نادرست و عدم‌قطعیت ذاتی.}
در نویز اندازه‌گیری، یک ویژگی ممکن است
اطلاع‌رسان باشد، اما مقدار ثبت‌شده
نسخهٔ مخدوش مقدار موردنظر باشد؛
برای نمونه $X_{\mathrm{obs}}=X_{\mathrm{true}}
+\epsilon$ در یک مدل جمعی.
این با افزودن ستون مستقلِ تصادفی
متفاوت است. \emph{نویز برچسب} به
اشتباه در خود متغیر هدف، مانند
$\widetilde y_i\ne y_i$، اشاره
دارد و بر نظارت و تشکیل مراکز
طبقات اثر می‌گذارد، نه فقط بر
یال‌های ویژگی. گاهی نیز یک مقدار
ذاتاً دارای عدم‌قطعیت توصیف‌شده
با توزیع یا زوجِ میانگین و واریانس
است؛ این وضعیت نباید بدون قرارداد
متریک و کرنل مناسب به عدد قطعی
یا ستون نامرتبط فروکاسته شود.
مدل‌سازی اختصاصی این سه وضعیت
\emph{تعمیم‌های پژوهشی آینده}
هستند؛ از آزمایش افزودن ستون نویزی
نمی‌توان برای آن‌ها نتیجه گرفت.

\paragraph{پیامد گرافی و تفاوت سه سطح سیاست.}
طبق تعریف گراف در
بخش~\ref{cmp:subsec:ch3-graph}،
یال افقی فقط میان ویژگی‌هایی
قرار می‌گیرد که در \emph{همان
نمونه} موجود باشند و یال عمودی
فقط نمونه‌های دارای \emph{یک
نام ویژگی مشترک} را متصل می‌کند.
بنابراین در مثال ماسک مکمل
رابطه~\eqref{cmp:eq:ch3-feature-schema-complementary-mask}،
میان دو ستون مترادف هیچ یال افقی
مستقیمی از نوع
$(X_i,a)\leftrightarrow(X_i,b)$
وجود ندارد؛ بااین‌حال ممکن است
به کمک سایر ویژگی‌های مشترک
مسیرهای \emph{غیرمستقیم} بین
دو بخش گراف برقرار باشد.
اگر هیچ چنین واسطه‌ای نباشد،
شرط اتصال‌پذیری
رابطه~\eqref{cmp:eq:ch3-connectivity}
ممکن است نقض شود. همچنین باید
برای هر ستون \emph{واقعاً موجود
در گراف} فرض حداقل پوشش
$|D_a|\ge3$ را کنترل کرد؛
در غیر این صورت، اعتبار انتقال‌های
بدون مقصد نیازمند قرارداد تازه است.
شکست این فرض‌ها به‌معنای
\emph{نامعتبرشدن صورت‌بندی فعلی
در آن داده} است، نه اثبات
ناتوانی ذاتی هر نسخهٔ ممکن IMC.

حتی وقتی دو ستون معنای یکسان دارند،
باید سه سطح متمایز بماند:
\emph{(۱)} برابری امتیازهای خام
$z_{\cdot,a}$ و $z_{\cdot,b}$
در یک سطر مشخص؛
\emph{(۲)} برابری احتمال‌های
انتخاب \emph{مشروط به ماسک و
مجموعهٔ مجاز همان وضعیت}؛
و \emph{(۳)} برابری سهم واقعیِ
انتخاب آن‌ها در گشت‌ها و
بازنمایی نهایی. حتی برابری
logits، تنها وقتی هر دو گزینه
در \emph{همان} سطر و
مجموعهٔ مجاز حضور دارند
برابری احتمال Softmax را
نتیجه می‌دهد؛ در ماسک مکمل
هیچ‌گاه هر دو در یک نمونه
مجاز نیستند. برای سهم نهایی
نیز فراوانی حضور نمونه‌ها،
وزن‌های عمودی، وضعیت‌های
قابل‌دسترسی و مسیرهای گشت
مؤثرند. هویت‌های بدون ماسک
بخش~\ref{cmp:subsec:ch3-noise-logpolicy-sensitivity}
جای این تحلیل نمونه‌محور را
نمی‌گیرند. از یکسان‌بودن
معنای دو ویژگی، برابری logits،
احتمال شرطی یا استفادهٔ نهایی
توسط سیاست نتیجه نمی‌شود.

\paragraph{طرح‌های آزمون متمایز و مرز نتایج موجود.}
برای تفکیک علت‌ها در پژوهش‌های
بعدی، می‌توان با داده، تقسیم‌ها،
سیاست و بودجهٔ آموزش کنترل‌شده،
سناریوهای مستقل زیر را مقایسه کرد:
افزودن ستون نامرتبطِ مولد-مرجع؛
افزودن ستون وابستهٔ قطعی مانند
رابطه~\eqref{cmp:eq:ch3-feature-deterministic-redundancy}؛
کپی کامل با مشاهده در همهٔ ردیف‌ها؛
شکست شِما با ماسک کاملاً مکمل؛
همان شکست با میزان هم‌پوشانیِ
کنترل‌شده میان ستون‌ها؛ و
\emph{کنترل ادغام صحیح شِما}
طبق رابطه~\eqref{cmp:eq:ch3-feature-schema-merge}.
برای تحلیل معتبر، علاوه بر BA،
تعداد ویژگی‌ها و پارامترها،
پوشش، اتصال‌پذیری گراف،
مجموعه‌های مجاز، جرم انتخاب
و مسیرهای گشت باید گزارش شوند.
یکسان‌بودن دادهٔ بنیادی نیز
به معنای یکسان‌بودن ظرفیت
سیاست در نمایش‌های متفاوت نیست.

آزمایش حساسیت نویز
بخش~\ref{sec:ch4-noise-feature-sensitivity}
در وضعیت فعلی \emph{فقط}
افزودن ستون‌های نویزی مرجع‌دار
و شاخص‌های مواجههٔ سیاست را
در پیکربندی‌های معین بررسی کرده است.
این آزمایش، آزمون جامع
کپی، وابستگی قطعی، شکست شِما،
\lr{MCAR/MAR/MNAR}، نویز
اندازه‌گیری یا برچسب نیست.
سناریوهای فوق صرفاً
\emph{فرضیه‌ها و طرح‌های
آزمون‌پذیر آینده} هستند؛
نه مجوز اجرای آزمایش جدید و
نه شواهد تجربیِ حاصل‌شده برای IMC.

\subsection{مواجهه سیاست با ستون‌های نویزی}
\label{cmp:subsec:ch3-noise-selection}
فرض کنیم $A_N\subseteq A_D$ مجموعه
ستون‌های نویزی با مرجع معلوم
و $A_P$ مجموعه سایر ویژگی‌های
فضای مسئله باشد. با استفاده از
$d$، نماد پیشین شمار کل ویژگی‌ها،
داریم:
\begin{equation}\label{cmp:eq:ch3-noise-feature-sets}
\boxed{
A_P=A_D\setminus A_N,\qquad
A_D=A_P\mathbin{\dot\cup}A_N,\qquad
d=|A_D|=|A_P|+|A_N|.
}
\end{equation}
$A_P$ در اینجا \emph{ویژگی‌های غیرنویزی
طبق قرارداد طراحی مسئله} را مشخص می‌کند؛
لزومی ندارد تک‌تک اعضای آن به‌تنهایی
اطلاعاتی درباره برچسب داشته باشند.
نمادهای $|A_N|$ و $|A_P|$
برای شمار ستون‌های دو گروه کافی‌اند.

برای روشن‌شدن مرز تعریف، اگر $C$ برچسب تصادفی
نمونه و $X^{(a)}$ مقدار ستون $a$ باشد،
شرط بی‌اطلاعی \emph{حاشیه‌ای} را می‌توان به‌شکل
\begin{equation}\label{cmp:eq:ch3-noise-marginal-independence}
I(C;X^{(a)})=0
\quad\Longleftrightarrow\quad
X^{(a)}\ \perp\!\!\!\perp\ C
\end{equation}
نوشت. اما در داده واقعی، این شرط از نمونه
متناهی به‌طور قطعی احراز نمی‌شود و
به‌تنهایی برای اثبات نویزبودن ستون کافی نیست:
ویژگی‌ای ممکن است جداگانه بی‌اطلاع باشد ولی
در ترکیب با ویژگی دیگر اطلاع‌رسان شود؛
ویژگی تکراری هم ممکن است پس از شرط‌کردن بر سایر
ستون‌ها اطلاعات افزوده نداشته باشد، بی‌آنکه
ذاتاً نویز اندازه‌گیری باشد.
الگوی فقدان ویژگی نیز ممکن است به برچسب
وابسته باشد و باید جداگانه بررسی شود.
در نتیجه، $A_N$ در این تحلیل یک
\emph{مجموعه مرجع معلوم از طراحی مولد} است،
نه خروجی خودکار تشخیص نویز از سیاست آموزش‌دیده.
حتی استقلال حاشیه‌ای ستون از برچسب،
به‌معنای بی‌اثربودن آن بر مسیر گشت
نمونه--ویژگی نیست. در این بخش،
مجموعه $A_N$ را با اتکا به
قرارداد تولید داده مشخص می‌کنیم؛
بنابراین ادعای شناسایی نویز
در داده ناشناخته مطرح نمی‌شود.

برای یک سیاست \emph{ثابت عددی} از خانواده
رابطه~\eqref{cmp:eq:ch3-fixed-policy-matrix}،
سطر آغاز و سطرهای انتقال امتیاز را بدون اعمال
ماسک نمونه‌ها نرمال می‌کنیم:
\begin{equation}\label{cmp:eq:ch3-noise-policy-probabilities}
p_0(b)=\frac{\exp z_{\bot b}}{\sum_{u=1}^{d}\exp z_{\bot u}},
\qquad
P_{ab}=\frac{\exp z_{ab}}{\sum_{u=1}^{d}\exp z_{au}},
\qquad a,b\in\{1,\ldots,d\}.
\end{equation}
در این قرارداد، $p_0$ بردار سطری آغاز،
$P\in\mathbb R^{d\times d}$ یک ماتریس گذار سطری،
و $\mathbf 1_{A_N}$ بردار شاخص ستون‌های نویزی است.
این فرایند کمکی تنها \emph{انتخاب‌های متوالی ویژگی} را
با همان سیاست توصیف می‌کند؛ لازم نیست توزیع
نمونه جاری، ماسک ویژگی، توقف هندسی، سقف $H$ و
ممنوعیت بازگشت به مبدأ را بازسازی کند.
در سیاست شرطی طبقه، همین تعاریف برای هر فرض طبقه
$c_s$ به‌صورت $p_0^{(c_s)},P^{(c_s)}$ جداگانه اعمال می‌شوند.



\subsection{تحلیل حساسیت سیاست انتخاب ویژگی بدون ماسک}
\label{cmp:subsec:ch3-noise-logpolicy-sensitivity}
در این زیربخش، یک \emph{سطر ثابت} از سیاست کمکی
رابطه~\eqref{cmp:eq:ch3-noise-policy-probabilities}
را در نظر می‌گیریم: برای سطر آغاز، $p=p_0$ و
$z_b=z_{\bot b}$؛ برای سطر پس از ویژگی $a$،
$p_b=P_{ab}$ و $z_b=z_{ab}$ است.
برای \lr{class} همین تحلیل جداگانه در هر فرضیهٔ
طبقهٔ $c_s$ اعمال می‌شود.
در هر سطر، مجموعهٔ مجاز دقیقاً $A_D=A_P\mathbin{\dot\cup}A_N$
و \emph{ثابت} است و تمام logits آن متناهی‌اند؛
ازاین‌رو $p_b>0$ برای هر $b\in A_D$.

با مشتق‌گیری از
$\log p_b=z_b-\log\sum_{u\in A_D}e^{z_u}$،
هویت دیفرانسیلی عمومی زیر حاصل می‌شود:
\begin{equation}\label{cmp:eq:ch3-noise-logpolicy-differential}
\boxed{
\mathrm d\log p_b=\mathrm dz_b-
\sum_{u\in A_D}p_u\,\mathrm dz_u
=\mathrm dz_b-\sum_{u\in A_P}p_u\,\mathrm dz_u
-\sum_{u\in A_N}p_u\,\mathrm dz_u .
}
\end{equation}
برای تغییر \emph{متناهی} دلخواه از $z^{(t)}$
به $z^{(t+1)}$، با
$\Delta z_u^{(t)}=z_u^{(t+1)}-z_u^{(t)}$،
رابطهٔ \emph{دقیق} مستقل از تابع هدف و بهینه‌ساز برابر است با:
\begin{equation}\label{cmp:eq:ch3-noise-logpolicy-finite}
\boxed{
\Delta\log p_b^{(t)}
=\Delta z_b^{(t)}-
\log\!\left(\sum_{u\in A_D}
p_u^{(t)}e^{\Delta z_u^{(t)}}\right).
}
\end{equation}
جایگزینی عبارت لگاریتمی با
$\sum_u p_u^{(t)}\Delta z_u^{(t)}$
فقط تقریب مرتبهٔ اول برای آپدیت کوچک است.
برای دو ویژگی $a,b$ در \emph{همان سطر}،
هویت نسبت حتی در آپدیت بزرگ نیز دقیق باقی می‌ماند:
\begin{equation}\label{cmp:eq:ch3-noise-logpolicy-ratio}
\boxed{
\Delta\log\frac{p_a}{p_b}
=\Delta z_a-\Delta z_b,\qquad a,b\in A_D .
}
\end{equation}
بنابراین اگر صرفاً logits ستون‌های $A_N$ تغییر کنند
و logits دو ویژگی اصلی $a,b\in A_P$ ثابت بمانند،
\emph{نسبت} احتمال آن دو در این سطر ثابت است؛
بااین‌حال احتمال \emph{مطلق} هر دو ممکن است به علت
تغییر مخرج نرمال‌سازی جابه‌جا شود.
همچنین افزایش جرم سیاست روی نویز به‌تنهایی
کاهش نوسان یا گرادیان ویژگی‌های اصلی را اثبات نمی‌کند.

این سه رابطه هویت‌های \emph{سیاست بدون ماسک} هستند؛
مشتق موتور واقعیِ گشت ماسک‌دار، انتقال عمودی،
ممنوعیت بازگشت به مبدأ یا توقف هندسی را محاسبه نمی‌کنند.
در آموزش، گرادیان زیان مانند
رابطه~\eqref{cmp:eq:ch3-softplus-js-policy-gradient}
از تمام استفاده‌های logits مشترک در گشت و مراکز حاصل
می‌شود؛ بهینه‌ساز طبق
رابطه~\eqref{cmp:eq:ch3-adam-logit-update}
یا قرارداد جایگزین، $\Delta z$ را تعیین می‌کند؛
\emph{پس از آن} هویت‌های حاضر تغییر سیاست کمکی را بیان
می‌کنند. این زنجیره لزوماً به معنای اثر علّی تثبیت‌کنندهٔ
ستون‌های نویزی نیست. مقایسهٔ مدل‌های دارای تعداد متفاوت
ستون نویز نیز تغییر بُعد $A_D$ است، نه آپدیت متناهی
روی یک مجموعهٔ ثابت. اثبات و کنترل‌های جبری در
بخش~\ref{app:ch3-unmasked-logpolicy-proof}
از پیوست~\ref{app:ch3-proofs} آمده‌اند.

\subsection{شاخص‌های مواجهه سیاست با نویز و انتخاب ویژگی}

\paragraph{افق تحلیل انتخاب ویژگی.}
در این تحلیل، سه نوبت انتخاب ویژگی را
به‌عنوان \emph{فرض اولیه درباره عمق کوتاهِ مفید
ارتباطات گرافی} بررسی می‌کنیم: انتظار داریم
بخشی از اثر ویژگی‌های واسط در چند انتخاب نخست
آشکار شود. این انتخاب نه قضیه‌ای مبنی بر کافی‌بودن
سه گام برای همه داده‌هاست و نه به معنای برابر بودن
عمق مؤثر بازنمایی با سه است. عمق مؤثر و جرم
وزن گام‌ها طبق رابطه~\eqref{cmp:eq:ch3-effective-depth}
جداگانه تعریف شده‌اند؛ سقف $H$ گشت واقعی نیز
با افق $h$ این فرایند کمکی یکی نیست.
اگر داده برای آشکارشدن ارتباطات مفید به
انتخاب‌های بیشتری نیاز داشته باشد،
شاخص‌های سه‌نوبتی تنها رفتار اولیه سیاست را
نشان می‌دهند و برای توصیف تمام مسیر کافی نیستند.
ازاین‌رو افق عمومی $h\ge1$ را نیز مجاز می‌گیریم.

احتمال اینکه انتخاب ویژگی در نوبت $t$ام نویزی باشد،
با احتساب امکان انتخاب تکراری یک ویژگی، برابر است با
\begin{equation}\label{cmp:eq:ch3-noise-qt}
q_t=p_0P^{t-1}\mathbf 1_{A_N},\qquad t\ge1.
\end{equation}
برای روشن‌شدن معنای شاخص‌ها، نماد
$p_{\mathrm{hit},N}^{(3)}$ را برای احتمال حداقل
یک انتخاب نویزی در سه نوبت و
$\bar n_N^{(3)}$ را برای تعداد مورد انتظار
انتخاب‌های نویزی در همین سه نوبت برمی‌گزینیم.
بالانویس $(3)$ شمار نوبت‌های انتخاب و زیرنویس $N$
مجموعه ویژگی‌های نویزی را مشخص می‌کند.
شاخص سوم، $\mathrm{NFI}_N^{(3)}$،
نسبت احتمال برخورد به سیاست یکنواخت
هم‌بُعد است. تعریف دقیق این سه کمیت چنین است:
\begin{equation}\label{cmp:eq:ch3-noise-F3-E3}
\boxed{
\begin{aligned}
p_{\mathrm{hit},N}^{(3)}&=\Pr\{\text{دست‌کم یک انتخاب از }A_N\text{ در سه نوبت}\}
   =1-p_0[A_P]\,P_{A_P,A_P}^{\,2}\mathbf 1_{A_P},\\
\bar n_N^{(3)}&=\mathbb E[\text{تعداد انتخاب‌های متعلق به }N]
   =q_1+q_2+q_3,\\
\mathrm{NFI}_N^{(3)}&=\frac{p_{\mathrm{hit},N}^{(3)}}{p_{\mathrm{hit},N}^{(3),U}},\qquad |A_N|>0.
\end{aligned}}
\end{equation}
در اینجا $P_{A_P,A_P}$ زیرماتریس
انتقال میان ویژگی‌های $A_P$،
$p_0[A_P]$ زیر‌بردار آغاز
و $\mathbf 1_{A_P}$ بردار یک
با اندازه $|A_P|$ است.
بنابراین جمله کم‌شونده در تعریف
$p_{\mathrm{hit},N}^{(3)}$ دقیقاً
احتمال آن است که هر سه انتخاب در $A_P$ بمانند.
$p_{\mathrm{hit},N}^{(3)}$ \emph{احتمال برخورد دست‌کم یک‌باره} است،
درحالی‌که $\bar n_N^{(3)}$ \emph{تعداد مورد انتظار} برخوردها
با احتساب تکرار است. به همین دلیل
$0\le p_{\mathrm{hit},N}^{(3)}\le1$ و
$p_{\mathrm{hit},N}^{(3)}\le\bar n_N^{(3)}
\le3p_{\mathrm{hit},N}^{(3)}$.

\paragraph{تعمیم افق و نرمال‌سازی شمار انتخاب نویز.}
برای هر تعداد مثبت و ثابت $h$ از انتخاب‌های
فرایند کمکی بدون ماسک، شاخص‌های پیشین
به‌صورت زیر تعمیم می‌یابند:
\begin{equation}\label{cmp:eq:ch3-noise-general-horizon}
\boxed{
\begin{aligned}
p_{\mathrm{hit},N}^{(h)}
&=1-p_0[A_P]P_{A_P,A_P}^{\,h-1}\mathbf 1_{A_P},\\
\bar n_N^{(h)}
&=\sum_{t=1}^{h}q_t
=\sum_{t=1}^{h}p_0P^{t-1}\mathbf 1_{A_N}.
\end{aligned}}
\end{equation}
چون $\bar n_N^{(h)}$ حداکثر $h$ است،
برای مقایسه افق‌های متفاوت، سهم
\emph{مورد انتظار انتخاب نویز به‌ازای هر نوبت}
را تعریف می‌کنیم:
\begin{equation}\label{cmp:eq:ch3-noise-normalized-expectation}
\boxed{
s_N^{(h)}=\frac{\bar n_N^{(h)}}{h}
=\frac1h\sum_{t=1}^{h}q_t,\qquad
0\le s_N^{(h)}\le1.
}
\end{equation}
در سیاست یکنواخت بدون ماسک
(رابطه~\eqref{cmp:eq:ch3-uniform-policy})،
هر انتخاب ویژگی یکنواخت و
مستقل از انتخاب‌های پیشین است.
بنابراین برای هر افق $h\ge1$
مرجع تحلیلی دقیق زیر برقرار است:
\begin{equation}\label{cmp:eq:ch3-noise-uniform-reference}
\boxed{
\begin{aligned}
q_t^U&=\frac{|A_N|}{d},\qquad t\ge1,\\
p_{\mathrm{hit},N}^{(h),U}
&=1-\left(\frac{|A_P|}{d}\right)^h,\\
\bar n_N^{(h),U}&=\frac{h|A_N|}{d},
\qquad s_N^{(h),U}=\frac{|A_N|}{d}.
\end{aligned}}
\end{equation}
به‌ویژه برای دو و سه نوبت،
احتمال برخورد به‌ترتیب
$1-(|A_P|/d)^2$
و $1-(|A_P|/d)^3$ است.
شمار مورد انتظار انتخاب نویز
با $h$ خطی افزایش می‌یابد،
ولی سهم انتخاب نویز به‌ازای
هر نوبت برای تمامی افق‌ها
برابر $|A_N|/d$ می‌ماند.
این برابری‌های دقیق مربوط
به فرایند کمکی بدون ماسک‌اند،
نه گشت واقعی ماسک‌دار.
پس $\mathrm{NFI}_N^{(3)}$ میزان احتمال برخورد سیاست با نویز
\emph{نسبت به سیاست یکنواخت با همان تعداد ویژگی} است؛
این نسبت خود احتمال نیست و می‌تواند از یک بزرگ‌تر باشد.
برای $|A_N|=0$، $p_{\mathrm{hit},N}^{(3)}
=\bar n_N^{(3)}=0$ و $\mathrm{NFI}_N^{(3)}$ به علت مخرج
صفر \emph{تعریف‌نشده} است، نه صفر.
در سیاست‌های شرطی طبقه، شاخص‌ها ابتدا برای هر
$c_s$ محاسبه می‌شوند؛ میانگین برابر‌وزن طبقه‌ها
تنها یک خلاصه توصیفی است، مگر آنکه همان وزن‌دهی
در پروتکل نمونه‌گیری طبقه نیز برقرار باشد.

برای مقایسه سیاست‌ها در ابعاد متفاوت،
نسبت بی‌بعد دیگری نیز می‌توان تعریف کرد:
\begin{equation}\label{cmp:eq:ch3-noise-normalized-uniform-ratio}
\mathrm{NSR}_N^{(h)}
=\frac{s_N^{(h)}}{|A_N|/d}
=\frac{\bar n_N^{(h)}}{h|A_N|/d},
\qquad |A_N|>0.
\end{equation}
مقدار یک در این نسبت، مرجع یکنواخت
هم‌بُعد است؛ برای $|A_N|=0$ این نسبت
تعریف نمی‌شود. برخلاف
$\mathrm{NFI}_N^{(3)}$ که
\emph{احتمال حداقل یک برخورد} را مقایسه می‌کند،
$\mathrm{NSR}_N^{(h)}$ بر
\emph{میانگین دفعات انتخاب} استوار است.
برای نمونه، شرط $\bar n_N^{(3)}<1$
فقط معادل $s_N^{(3)}<1/3$ است
و به‌تنهایی نشان نمی‌دهد سیاست نسبت
به Uniform از نویز کمتر استفاده کرده است؛
این نتیجه به نسبت $|A_N|/d$ نیز وابسته است.

\paragraph{مقایسه موضعی در نزدیکی عمق مفید.}
اگر بر اساس ساختار گراف و جرم وزن‌های گشت
در رابطه~\eqref{cmp:eq:ch3-effective-depth}
عمق مفید را تقریباً در حدود دو یا سه انتخاب
در نظر بگیریم، ارزیابی اصلی این سنجه
باید بر $h=2$ و $h=3$ متمرکز باشد؛
نه اینکه صرفاً مقدار آن در یک افق بلند
گزارش شود. در این دو افق داریم:
\begin{equation}\label{cmp:eq:ch3-noise-local-depth-comparison}
\boxed{
\begin{aligned}
s_N^{(2)}&=\frac{q_1+q_2}{2},\\
s_N^{(3)}&=\frac{q_1+q_2+q_3}{3},\\
s_N^{(3)}-s_N^{(2)}
&=\frac{q_3-s_N^{(2)}}{3}.
\end{aligned}}
\end{equation}
پس علامت تغییر سهم نرمال‌شده از دو
به سه نوبت، تنها مشخص می‌کند احتمال
انتخاب نویز در نوبت سوم در مقایسه
با میانگین دو نوبت پیشین بیشتر است
یا کمتر؛ این تغییر به‌تنهایی
نشانه سودمندی یا آسیب نویز نیست.
در این مقایسه باید مرجع یکنواخت
$|A_N|/d$ و، در صورت مطالعه چند پیکربندی،
مجموعه ستون‌های نویزی $A_N$
و قرارداد سیاست نیز روشن بمانند.
افق‌های دیگر صرفاً برای تحلیل حساسیت
و زمانی که شواهد عمق مؤثر از ناحیه
دو یا سه فراتر می‌روند مناسب‌اند؛
این انتخاب یک \emph{قرارداد تحلیلی
وابسته به مسئله} است و نه قانون
ثابت عمق برای همه داده‌ها.
این مقایسه با آمار
\emph{انتخاب‌های بدون ماسک} تعریف
شده و نباید با توزیع عمق گشت واقعی
یا نتیجه تجربی اثبات‌شده اشتباه شود.

برای افق‌های بلند، احتمال دست‌کم یک برخورد
ممکن است به یک نزدیک شود و قدرت تمایز
$ p_{\mathrm{hit},N}^{(h)} $
کاهش یابد؛ در مقابل، سهم $s_N^{(h)}$
همچنان میانگین مصرف نویز در هر نوبت را
نشان می‌دهد. البته این تعمیم هنوز
\emph{فرایند کمکی بدون ماسک} است،
نه تخمین مواجهه در گشت واقعی؛
برای ارتباط با آن گشت باید طول‌های مؤثر،
ماسک و توزیع بازدید از گره‌ها نیز لحاظ شوند.
انتخاب افق مناسب و بررسی تغییر شاخص‌ها با
$h$، پرسش تجربی جداگانه است و از
تعریف‌های فوق نتیجه عددی به دست نمی‌آید.



\paragraph{شاخص استفاده از هر ستون و مسئله شناسایی نویز.}
مقادیر جمعی بالا به‌تنهایی نمی‌گویند
کدام ستون‌ها کمتر یا بیشتر انتخاب شده‌اند.
برای هر اندیس ویژگی $a\in\{1,\ldots,d\}$
می‌توان تعداد مورد انتظار انتخاب همان ویژگی
را در افق دلخواه $h$ تعریف کرد:
\begin{equation}\label{cmp:eq:ch3-noise-per-feature-selection}
\boxed{
\bar n_a^{(h)}
=\sum_{t=1}^{h}[p_0P^{t-1}]_a,\qquad
\bar n_N^{(h)}=\sum_{a\in A_N}\bar n_a^{(h)}.
}
\end{equation}
در سیاست یکنواخت بدون ماسک
$\bar n_a^{(h),U}=h/d$ است.
مقدار سه‌نوبتی این شاخص با قراردادن $h=3$
به‌دست می‌آید.
این شاخص‌ها را می‌توان از ماتریس امتیاز
سیاست تثبیت‌شده و بدون آموزش دوباره محاسبه کرد.
بااین‌حال، کم‌استفاده‌بودن یک ستون
\emph{معادل بی‌اطلاعی آماری آن نیست}؛
سیاست ممکن است یک ویژگی اطلاع‌رسان را
به‌دلیل هم‌پوشانی با ویژگی دیگر کمتر انتخاب کند،
یا از یک ستون فاقد اطلاعات مستقیم برچسب
به‌عنوان مسیر واسط در گراف استفاده کند.
در نتیجه، توانایی مدل در تحمل افزودن
$|A_N|$ ستون بی‌اطلاع و توانایی
\emph{شناسایی} همان ستون‌ها دو پرسش
مستقل‌اند. آزمون شناسایی به معیار مرجع
معلوم برای $A_N$ و ارزیابی جداگانه
رتبه‌بندی ستون‌ها نیاز دارد و
در این بخش نتیجه تجربی برای آن ادعا نمی‌شود.



\subsection{وابستگی ساختاری و اطلاعات تکراری در ویژگی‌ها}

تمایز مفهومیِ افزونگی، وابستگی قطعی و کپی کامل
در بخش~\ref{cmp:subsec:ch3-feature-error-types}
بیان شد. در اینجا تنها یکی از پیامدهای
\emph{شناسایی‌پذیری پارامترها} را با یک
مثال خطی بررسی می‌کنیم.

وابستگی ساختاری ویژگی‌ها
می‌تواند نگاشت پارامترها
به خروجی را غیرتزریقی کند.
برای نمونه، اگر
$x_1=t+r$،
$x_2=t-r$ و
$x_3=2t+r$ باشند،
رابطهٔ دقیق
$2x_3=3x_1+x_2$
برقرار است؛ بنابراین
در یک مدل امتیازدهی خطی،
تبدیل ضرایب
$(w_1,w_2,w_3)$
به
$(w_1-3a,w_2-a,w_3+2a)$
برای هر $a$،
امتیاز را تغییر نمی‌دهد.
این مثال، ناشناسایی در
یک نگاشت خطی را توضیح می‌دهد
و به‌تنهایی اثبات
هم‌ارزی سیاست‌های گشت
در مدل حاضر نیست.
در دادهٔ مصنوعی این آزمایش،
ستون‌های نویزی مستقل
از بلوک اطلاع‌رسان
تولید شده‌اند؛
بااین‌حال، در نمونه‌های متناهی
امکان هم‌بستگی و
جداسازی ظاهریِ تصادفی
وجود دارد.
ازاین‌رو باید
وابستگی ساختاری،
هم‌بستگی اتفاقی
و برابری صرف
معیار دقت را
از یکدیگر تفکیک کرد.



\subsection{تفکیک‌پذیری تصادفی ویژگی‌های نویزی در نمونه محدود}

\paragraph{تعریف آزمون‌های آماری مرجع.}
شاخص $\eta^2$ در این تحلیل،
\emph{اندازهٔ اثر شناخته‌شدهٔ
تحلیل واریانس یک‌طرفه}
است \cite{Richardson2011EtaSquared}؛
آزمون جایگشتی و اصلاح
\lr{single-step maxT}
از خانوادهٔ روش‌های
\lr{Westfall--Young}
هستند
\cite{WestfallYoung1993Resampling}.
این‌ها برخلاف معیارهای
هندسی و سیاست‌محور
مدل حاضر،
\emph{معیارهای جدید پیشنهادی
پایان‌نامه نیستند}.
تنها نماد
$p_{\mathrm{FWER7},j}$
در این گزارش به‌صورت
محلی برای مقدار
$p$ جایگشتی اصلاح‌شده
بر هفت ستون نویزی
انتخاب شده است؛
این نماد نام آزمون
جدید یا مقدار برآوردشدهٔ
نرخ خطای \lr{FWER} نیست.

برای ستون $j$،
$O_j$ مجموعهٔ ردیف‌های
دارای مقدار مشاهده‌شده،
$n_{jc}$ تعداد همان
ردیف‌ها در طبقه $c$،
$\bar x_{jc}$ میانگین طبقه‌ای
و $\bar x_j$ میانگین
تمام مشاهدات همان ستون است.
آمارهٔ آزمون با
اندازهٔ اثر تحلیل واریانس
برابر گرفته می‌شود:
\begin{equation}
\label{cmp:eq:ch4-e14-eta2}
T_j=\eta_j^2=
\frac{\sum_{c=1}^{K}
n_{jc}(\bar x_{jc}-\bar x_j)^2}
{\sum_{i\in O_j}
(x_{ij}-\bar x_j)^2}.
\end{equation}
آزمون \emph{مقادیر نویز}
فقط روی ردیف‌های
غیرمفقود محاسبه می‌شود؛
آزمون جداگانهٔ
\emph{ماسک فقدان}
همین آماره را روی
متغیر دودویی
$m_{ij}=\mathbf{1}
\{x_{ij}\text{ مفقود است}\}$
و تمام ردیف‌ها
به کار می‌برد.
این آماره تنها
تفکیک میانگین‌ها
را می‌سنجد
و آزمون کامل
تمام وابستگی‌های
چندمتغیره نیست.

تحت فرض صفرِ
تعویض‌پذیری برچسب‌ها
در برابر هفت ستون،
$B=2000$ جایگشت
تصادفی $\pi_b$
روی بردار برچسب‌ها
اعمال می‌شود.
در هر تکرار،
\emph{همان جایگشت}
برای هر هفت ستون
استفاده شده است؛
پس مقادیر،
الگوی فقدان
و وابستگی بین
ستون‌ها تغییر نمی‌کنند.
اگر
$A_N^{\max}$
مخزن هفت‌ستونی باشد،
\begin{equation}
\label{cmp:eq:ch4-e14-maxt}
T_j^{(b)}
=\eta^2(x_j,\pi_b(y)),
\qquad
M^{(b)}
=\max_{k\in A_N^{\max}}
T_k^{(b)}.
\end{equation}
مقدار $p$ تصحیح‌شدهٔ
تک‌مرحله‌ای برای ستون $j$
برابر است با
\begin{equation}
\label{cmp:eq:ch4-e14-p-fwer7}
p_{\mathrm{FWER7},j}
=\frac{
1+\sum_{b=1}^{2000}
\mathbf{1}
\left[M^{(b)}\ge
T_j^{\mathrm{obs}}\right]
}{2001}.
\end{equation}
اصلاح $+1$ در شمارنده
و مخرج،
برای جایگشت‌های
تصادفی نمونه‌برداری‌شده
در ادبیات آزمون جایگشتی
بررسی شده است
\cite{PhipsonSmyth2010Permutation}.
فرمول~\eqref{cmp:eq:ch4-e14-p-fwer7}
مقایسه‌های متعدد
\emph{داخل همان مجموعهٔ
هفت‌ستونی،
برای یک تقسیم و
یک نوع آماره}
را لحاظ می‌کند؛
بدون فرض‌های مناسب
نباید از آن
کنترل قوی خطای
خانوادگی برای
همهٔ حالات فرض صفر
یا تمام تحلیل‌های
اکتشافی استنتاج کرد.



\subsection{مقاومت در برابر افزودن نویز و محدودیت‌های استنباط}

\paragraph{انتظار مقاومت در برابر افزودن نویز.}
برای بررسی اثر افزایش شمار ستون‌های نویزی،
فرض کنیم خانواده‌ای از مسائل با مجموعه
ثابت ویژگی‌های مرجع $A_P$ و
مجموعه‌های نویزی با اندازه‌های متفاوت
$|A_N|$ در نظر گرفته شود و سیاست هر مسئله
\emph{جداگانه} آموزش ببیند.
از سیاستی که در برابر نویز
\emph{مقاومت در سطح انتخاب ویژگی}
نشان می‌دهد، انتظار داریم
سهم متوسط انتخاب نویز
$s_N^{(h)}$ با افزودن ستون‌های نویزی
افزایش شدیدی نداشته باشد؛
به‌خصوص در افق‌هایی که برای
بازنمایی مفیدند.
این یک \emph{فرضیه پژوهشی} است،
نه قضیه‌ای درباره همه سیاست‌ها.
برای مقایسه، مرجع یکنواخت
در این خانواده از مسائل برابر است با
\begin{equation}\label{cmp:eq:ch3-noise-resistance-uniform-growth}
\boxed{
s_N^{(h),U}(|A_N|)
=\frac{|A_N|}{|A_P|+|A_N|},
\qquad
\bar n_N^{(h),U}(|A_N|)
=\frac{h|A_N|}{|A_P|+|A_N|}.
}
\end{equation}
اگر $|A_P|>0$، هر دو مقدار
با افزایش شمار ستون‌های نویزی
صعودی‌اند؛ سهم یکنواخت
به یک و شمار مورد انتظار
به $h$ نزدیک می‌شوند.
با ثابت‌ماندن $A_P$،
اگر تنها \emph{یک ستون نویزی}
افزوده شود، افزایش سهم مرجع
دقیقاً برابر است با:
\begin{equation}\label{cmp:eq:ch3-noise-uniform-growth-derivative}
\boxed{
\begin{aligned}
\Delta s_N^{(h),U}
&=\frac{|A_N|+1}{|A_P|+|A_N|+1}
-\frac{|A_N|}{|A_P|+|A_N|}\\
&=\frac{|A_P|}
{(|A_P|+|A_N|)(|A_P|+|A_N|+1)}>0.
\end{aligned}}
\end{equation}
پس مرجع یکنواخت با
افزودن نویز افزایش می‌یابد،
اما مقدار افزایش برای
شمار نویزی بزرگ‌تر کاهش می‌یابد.
بنابراین \emph{کمترشدن شیب افزایش}
$s_N^{(h)}$ در مقایسه با مرجع
یکنواخت، یا باقی‌ماندن
$\mathrm{NSR}_N^{(h)}$ زیر یک،
می‌تواند شاهدی توصیفی برای
کاهش مواجهه نسبی سیاست
با نویز باشد. هیچ‌یک از این‌ها
به‌تنهایی اثبات نمی‌کند که
دقت تصمیم حفظ شده است؛
برای ادعای مقاومت پیش‌بینی
باید عملکرد و تغییرات آن
در یک پروتکل کنترل‌شده نیز
بررسی شود.


\paragraph{معنای حفظ جرم احتمال روی نویز.}
اگر $|A_N|>0$ و تمام امتیازهای ورودی نرمال‌سازی با بیشینهٔ هموار متناهی باشند،
احتمال انتخاب هر ویژگی در فرایند کمکی مثبت است؛
درنتیجه $\bar n_N^{(3)}>0$ حتی در سیاستی که
ویژگی‌های نویزی را بسیار کم‌اهمیت کرده باشد.
این مثبت‌بودن به‌تنهایی گواه سودمندی نویز
یا تصمیم آگاهانه سیاست برای حفظ آن نیست.
از سوی دیگر، شرط $\bar n_N^{(3)}<1$، اگر در یک سیاست
آموخته‌شده مشاهده شود، صرفاً می‌گوید
\emph{تعداد مورد انتظار} انتخاب‌های نویزی
در سه نوبت کمتر از یک است؛
نه اینکه هر مسیر حداکثر یک انتخاب نویزی دارد
و نه اینکه $\bar n_N^{(3)}<1$ برای همه سیاست‌ها تضمین شده است.
برای تشخیص سرکوب نسبی نویز باید مقدار آموخته‌شده
با مرجع هم‌بُعد $\bar n_N^{(3),U}$ و احتمال‌های
$q_1,q_2,q_3$ مقایسه شود؛ برای افق‌های
بزرگ‌تر نیز سنجه نرمال‌شده
$s_N^{(h)}$ یا نسبت
$\mathrm{NSR}_N^{(h)}$
مناسب‌تر از مقایسه مستقیم
شمار انتخاب‌های دو افق متفاوت است.
اینکه جرم غیرصفر نویز فقط از محدودیت نرمال‌سازی با بیشینهٔ هموار
باقی مانده یا بر اثر تغییر مسیرهای گراف
و کیفیت بازنمایی در آموزش حفظ شده است،
پرسشی تجربی است و از این روابط نتیجه نمی‌شود.

\paragraph{مرز استنباط از افزودن ویژگی نویزی.}
با تغییر $|A_N|$، هم اندازه ماتریس امتیاز
و هم توزیع احتمال حاصل از نرمال‌سازی با بیشینهٔ هموار تغییر می‌کند؛
تعداد امتیازهای آزاد در خانواده مشترک $d(d+1)$
و در خانواده شرطی طبقه $Kd(d+1)$ است.
علاوه بر این، ویژگی تازه می‌تواند توزیع انتقال عمودی
روی گراف را نیز تغییر دهد. ازاین‌رو حتی اگر
$p_{\mathrm{hit},N}^{(3)}$ یا
$\bar n_N^{(3)}$ در یک سیاست آموزش‌دیده کاهش یابد،
نمی‌توان نتیجه گرفت اثر واقعی نویز در
\emph{گشت ماسک‌دار} یا در دقت طبقه‌بندی کاهش یافته است.
همچنین در مقایسه سیاست‌هایی که برای ابعاد متفاوت
\emph{جداگانه آموزش دیده‌اند}،
هیچ تضمینی برای یکنواخت افزایش یافتن
$p_{\mathrm{hit},N}^{(3)}$
با تعداد ستون‌های نویزی وجود ندارد.
فرضیه‌های «بازتوزیع جرم احتمال در نرمال‌سازی با بیشینهٔ هموار»، «افزایش ظرفیت سیاست»
و «تغییر هندسه گراف» باید با کنترل‌های مجزا
از یکدیگر تفکیک شوند.
در سیاست مشترک، $p_{\mathrm{hit},N}^{(3)}$ و
$\bar n_N^{(3)}$ برای یک ماتریس
امتیاز ثابت، کمیت‌هایی در سطح \emph{سیاست} هستند
و به نمونه جاری وابسته نیستند؛ در سیاست شرطی طبقه
نیز برای هر فرض طبقه تعریف می‌شوند.
ازاین‌رو این شاخص‌های بدون ماسک به‌تنهایی
برای تعیین رفتار واقعی \emph{یک نمونه مشخص}
یا سنجش مقاومت پیش‌بینی کافی نیستند.
برای تحلیل نمونه‌محور،
مواجهه در گشت واقعی ماسک‌دار،
طول گشت و قاعده تصمیم
باید مستقل از این شاخص‌ها بررسی شوند.


\section{تحلیل پساآموزش بازنمایی‌ها و توسعهٔ مقاوم طبقه‌بند}
\label{cmp:sec:ch3-outliers}
این بخش به \emph{صورت‌بندی نظریِ توسعه‌های پساآموزش}
اختصاص دارد و روش‌های آن در آزمایش‌های این
پژوهش اجرا یا ارزیابی تجربی نشده‌اند.
هدف، تعریف ابزارهایی برای شناسایی نمونه‌های
نامعمول در فضای بازنمایی، بررسی هندسی تداخل
طبقه‌ها و مطالعهٔ امکان اصلاح مراکز طبقه‌بند است.
این ابزارها جزئی از الگوریتم آموزش یا استنتاج
اعتبارسنجی‌شدهٔ این پایان‌نامه نیستند؛
کاربرد عملی و اثر احتمالی آن‌ها بر دقت،
به آزمون مستقل در پژوهش‌های بعدی نیاز دارد.

اصطلاح \emph{نمونه پرت} به مشاهده‌ای
اشاره می‌کند که بر اساس یک معیار مشخص،
نسبت به ساختار مرجع داده نامعمول
یا دورافتاده است. این مفهوم، بدون
تعیین فضای سنجش، توزیع یا مجموعه
مرجع و آستانه تشخیص، تعریف عملیاتی
یکتایی ندارد. برای نمونه، یک مقدار
دور از میانگین در تنها یک ویژگی
لزومی ندارد مشاهده را در فضای
چندویژگی پرت کند: رابطه آن ویژگی
با سایر ویژگی‌ها و ساختار طبقه‌ها
می‌تواند چنین مقداری را توضیح دهد.
همچنین ممکن است هر مؤلفه به‌تنهایی
مقداری متعارف داشته باشد، اما
ترکیب مؤلفه‌ها نامعمول باشد.
درنتیجه، تشخیص دورافتادگی
بر پایه یک ویژگی منفرد، عموماً
جایگزین تحلیل مشترک ویژگی‌ها نیست.

از منظر آماری نیز باید میان
\emph{کم‌احتمال‌بودن خود مشاهده}
در توزیع داده، \emph{کم‌احتمال‌بودن
آن در طبقه واقعی} و
\emph{ابهام در پیش‌بینی برچسب}
تمایز گذاشت. این‌ها در صورت
داشتن یک مدل احتمالاتی مشخص
به‌ترتیب می‌توانند با کمیت‌های
$p_X(x)$، $p_{X\mid C}(x\mid c)$
و $P(C=c\mid X=x)$ سنجیده شوند؛
کوچکی احتمال یا چگالی در دو
مورد نخست با پایین‌بودن احتمال
پسین طبقه واقعی یکسان نیست.
حتی طبقه‌بند بهینه بیزی نیز
در ناحیه‌های هم‌پوشان طبقه‌ها
ممکن است با عدم قطعیت روبه‌رو شود،
بی‌آنکه مشاهده لزوماً پرت باشد.
همچنین ناتوانی یک مدل برازش‌شده
در توضیح یک نمونه، به‌تنهایی
پرت‌بودن ذاتی آن را اثبات نمی‌کند؛
ممکن است مدل یا فرض‌های آن
برای توزیع واقعی ناکافی باشند.

در چارچوب این پژوهش، تعریف
عملیاتی \emph{پرت‌بودن در بازنمایی}
را به کار می‌بریم: نمونه آموزشی
زمانی \emph{نامزد پرت} محسوب
می‌شود که توزیع برچسب همسایگی
آن، $\mathscr Y_i$، نسبت به
بازنمایی‌های طبقه واقعی خودش،
طبق معیار فاصله و آستانه تعیین‌شده،
دورافتاده باشد. این تعریف در
فضای \emph{بازنمایی برچسب همسایگی}
است، نه مستقیماً در فضای ویژگی
خام؛ بنابراین به سیاست، گراف،
مراکز طبقه و آستانه انتخابی
وابسته است. چنین نامزدی لزوماً
نشان‌دهنده اشتباه‌بودن برچسب،
ناکارآمدبودن نمونه یا خطای
پیش‌بینی آن نیست.

ابتدا فاصله از مرکز حسابی طبقه
و شعاع‌های چندکی KL را برای
غربال نامزدهای پرت تعریف می‌کنیم.
سپس با نگاشت هیلبرتی جنسن--شنون،
توپ‌های کمینه پوشاننده بازنمایی‌ها
و تأثیر نمونه‌های مرزی بر تداخل
توپ‌های طبقه‌ها را بررسی می‌کنیم.
در پایان، کاهش سهم نمونه‌های
دورافتاده \emph{فقط در برآورد مرکز}
را به‌عنوان مداخله‌ای نظری
صورت‌بندی می‌کنیم، بدون آنکه
نمونه از داده، گراف یا گشت
حذف شود.


\subsection{شعاع چندکی KL برای غربال نمونه‌های پرت}
\label{cmp:subsec:ch3-outlier-quantile-radius}
\paragraph{از مشاهدات گشت تا هزینه درست‌نمایی.}
در زیربخش~\ref{cmp:subsec:ch3-mean-center-justification}
دیدیم که مرکز حسابی توزیع‌های
برچسب همسایگی یک طبقه، بهترین
توزیع مشترک برای توضیح مشاهدات
آن طبقه در معنای درست‌نمایی
مورد انتظار است. اثبات این
نتیجه در پیوست~
\ref{app:ch3-mixture-likelihood}
بر مقایسه درست‌نمایی مشاهداتی
استوار بود که از توزیع‌های
متفاوت به دست می‌آیند.
اکنون همان استدلال را برای
\emph{یک نمونه مبدأ} بازخوانی
می‌کنیم تا معنای شعاع KL روشن شود.

فرض کنیم $i\in C_c$ یک نمونه
آموزشی از طبقه $c$ باشد و
$\mathscr Y_i\in\Delta_K$
توزیع برچسب مقصد مشاهده‌شده
در گشت تصادفیِ آغازشده از آن
را، تحت سیاست و قرارداد
توقف ثابت، نشان دهد.
مرکز همان طبقه مطابق
رابطه~\eqref{cmp:eq:ch3-class-center}
برابر است با
$\mathcal V_c=|C_c|^{-1}
\sum_{j\in C_c}\mathscr Y_j$.
برای توضیح برچسب‌های حاصل
از گشت این نمونه، می‌توان
از دو مدل احتمالاتی استفاده کرد:
\emph{توزیع اختصاصی} خود نمونه
یعنی $\mathscr Y_i$، یا
\emph{توزیع مشترک} طبقه
یعنی $\mathcal V_c$.

اگر $Z_{i,1},\ldots,Z_{i,T}$
برچسب‌های حاصل از $T$
\emph{تکرار مستقلِ آزمایش کامل گشت}
از مبدأ $i$ باشند، در قرارداد
مذکور هر $Z_{i,t}$ از
$\mathscr Y_i$ توزیع دارد.
درست‌نمایی این مشاهدات
با مدل اختصاصی و مدل مشترک
به‌ترتیب حاصل‌ضرب احتمال‌های
$\mathscr Y_i(Z_{i,t})$
و $\mathcal V_c(Z_{i,t})$
است. با همان استدلال قانون
اعداد بزرگ در رابطه~
\eqref{eq:app-ch3-mixture-log-likelihood-ratio}
از پیوست، داریم:
\begin{equation}\label{cmp:eq:ch3-outlier-kl-likelihood-decomposition}
\boxed{
\begin{aligned}
&\lim_{T\to\infty}\frac1T
\log\frac{
\prod_{t=1}^{T}\mathscr Y_i(Z_{i,t})
}{
\prod_{t=1}^{T}\mathcal V_c(Z_{i,t})
}\\
&\qquad
=\sum_{u=1}^{K}\mathscr Y_i(u)
\log\frac{\mathscr Y_i(u)}{\mathcal V_c(u)}
=D_{\mathrm{KL}}(\mathscr Y_i\|\mathcal V_c).
\end{aligned}}
\end{equation}
این حد برای تکرارهای مستقل
\emph{گشت کامل} است و به
استقلال گام‌های متوالی
\emph{درون یک گشت} نیازی ندارد.
ازآنجاکه مرکز از میانگین
همین بازنمایی‌ها تشکیل شده،
روی تکیه‌گاه $\mathscr Y_i$
مقدار مثبت دارد و واگرایی
فوق متناهی است.

بنابراین KL در اینجا فقط
یک سنجه اختلاف انتزاعی نیست:
\emph{افت متوسط لگاریتم درست‌نمایی
به‌ازای هر مشاهده برچسب مقصد}
هنگام جایگزین‌کردن توزیع
اختصاصی نمونه با مرکز مشترک
طبقه است. این نتیجه معادل
هویت آنتروپی متقاطع
$-\sum_u\mathscr Y_i(u)
\log\mathcal V_c(u)
=H(\mathscr Y_i)+
D_{\mathrm{KL}}(\mathscr Y_i\|\mathcal V_c)$
نیز هست. منظور، مدل‌کردن
\emph{برچسب‌های مشاهده‌شده در گشت}
است، نه درست‌نمایی برچسب
واقعی نمونه مبدأ یا مقادیر
ویژگی‌های خام.

برای هر عضو طبقه، همین
هزینه اضافی را به‌صورت
\begin{equation}\label{cmp:eq:ch3-outlier-forward-kl}
d_i^{(c)}
=D_{\mathrm{KL}}\bigl(
\mathscr Y_i\|\mathcal V_c\bigr),
\qquad i\in C_c
\end{equation}
تعریف می‌کنیم. هرچه این
مقدار بزرگ‌تر باشد، توضیح
الگوی مشاهده‌های گشت نمونه
با مرکز طبقه در مقایسه با
توزیع اختصاصی آن، هزینه
اضافی بیشتری دارد.
میانگین این هزینه‌ها همان
$R_c$ در رابطه~\eqref{cmp:eq:ch3-R-class}
است؛ ولی این میانگین
به‌تنهایی نشان نمی‌دهد کدام
نمونه‌ها هزینه غیرمعمول
بزرگی تحمیل می‌کنند.

بدین ترتیب، شعاع چندکی KL
آستانه‌ای برای \emph{هزینه
اضافی مدل مشترک} در پوشش
سهم معینی از نمونه‌های
یک طبقه خواهد بود؛ این
«شعاع» آستانه واگرایی
جهت‌دار است و الزاماً
شعاع توپ متریک نیست.
به‌کارگیری KL در این
مرحله از ارتباط آن با
درست‌نمایی و مرکز حسابی
ناشی می‌شود. در نگاه
هندسی بعدی، متریک
جذر جنسن--شنون را
برای توپ‌های کمینه پوشاننده
به کار می‌گیریم که مسئله
و هدف متفاوتی دارد.

فاصله‌های طبقه $c$ را به‌ترتیب صعودی
$d_{c,(1)}\le\cdots\le d_{c,(n_c)}$ مرتب می‌کنیم.
برای سطح پوشش $0<\alpha\le1$،
\emph{شعاع چندکی با مرکز ثابت} و
تابع پوشش تجربی را چنین تعریف می‌کنیم:
\begin{equation}\label{cmp:eq:ch3-outlier-quantile-coverage}
\boxed{
R_{c,\alpha}=d_{c,(\lceil\alpha n_c\rceil)},\qquad
\widehat G_c(r)=\frac1{n_c}
\sum_{i\in C_c}\mathbf 1\{d_i^{(c)}\le r\}.
}
\end{equation}
از تعریف آماره ترتیبی،
$\widehat G_c(R_{c,\alpha})\ge\alpha$؛
با وجود فاصله‌های مساوی، پوشش می‌تواند
بیش از $\alpha$ باشد.
منحنی $\alpha\mapsto R_{c,\alpha}$ غیرنزولی
و برای تعداد محدود نمونه پله‌ای است.
نمونه‌هایی با $d_i^{(c)}>R_{c,\alpha}$
صرفاً \emph{نامزد بررسی پرت‌بودن نسبت به مرکز فعلی}
هستند؛ بزرگی این فاصله، نه غلط‌بودن برچسب را
اثبات می‌کند و نه ضرورت حذف نمونه را.
این غربال در فضای \emph{بازنمایی برچسب همسایگی}
انجام می‌شود، نه مستقیماً در فضای مقادیر ویژگی خام.
از آنجا که $\mathscr Y_i$ و مرکز $\mathcal V_c$
هر دو به سیاست و گراف وابسته‌اند، یک نمونه ممکن است
با تغییر سیاست از مجموعه نامزدهای پرت خارج شود
یا به آن وارد شود. بنابراین پرت‌بودن در این قرارداد
ویژگی ذاتی و تغییرناپذیر خود نمونه نیست.

اگر خود مرکز نیز آزادانه انتخاب شود،
کمیت متفاوت زیر را می‌توان بررسی کرد:
\begin{equation}\label{cmp:eq:ch3-outlier-optimized-quantile}
C_{c,\alpha}=
\inf_{v\in\Delta_K}
\operatorname{Quantile}_{\alpha}
\bigl\{D_{\mathrm{KL}}(\mathscr Y_i\|v):i\in C_c\bigr\},
\qquad
C_{c,\alpha}\le R_{c,\alpha}.
\end{equation}
نابرابری به این دلیل برقرار است که مرکز
ثابت $\mathcal V_c$ نیز یکی از گزینه‌های مجاز است.
مرکز بهینه ممکن است با تغییر $\alpha$ تغییر کند؛
در $\alpha=1$، $C_{c,1}$ مسئله شعاع اطلاعاتی
کمینه--بیشینه است، نه میانگین پراکندگی $R_c$
و نه معیار جدایی مراکز $I$.
$C_{c,\alpha}$ را نباید با طبقه $C_c$ یا
شعاع \emph{ثابت} $R_{c,\alpha}$ یکسان دانست.

\paragraph{پرت‌بودن در برابر خطای تصمیم.}
پس از تثبیت یک سیاست و قاعده تصمیم، دو مجموعه
را برای اعضای طبقه $c$ تعریف می‌کنیم:
\begin{equation}\label{cmp:eq:ch3-outlier-error-sets}
\boxed{
\begin{aligned}
\mathcal O_{c,\alpha}
&=\{i\in C_c:\ d_i^{(c)}>R_{c,\alpha}\},\\
\mathcal E_c
&=\{i\in C_c:\ \widehat y_i\ne c\}.
\end{aligned}}
\end{equation}
چهار ناحیه
$\mathcal O_{c,\alpha}\cap\mathcal E_c$،
$\mathcal O_{c,\alpha}\setminus\mathcal E_c$،
$\mathcal E_c\setminus\mathcal O_{c,\alpha}$
و باقی‌مانده اعضای طبقه، به‌ترتیب
نمونه‌های دورافتاده و اشتباه، دورافتاده و درست،
غیرپرت طبق آستانه ولی اشتباه، و سایر نمونه‌ها
را متمایز می‌کنند. این تفکیک، مبنایی برای
بررسی مستقل رابطهٔ دورافتادگی و خطای تصمیم است.

وجود نمونه در
$\mathcal E_c\setminus\mathcal O_{c,\alpha}$
یادآور آن است که نزدیک‌تر بودن به مرکز واقعی
به‌تنهایی برتری نسبت به مراکز رقیب را تضمین نمی‌کند؛
باید حاشیه و امتیاز رقیبان نیز بررسی شود.
از سوی دیگر، اعضای
$\mathcal O_{c,\alpha}\setminus\mathcal E_c$
ممکن است نمونه‌های نامعمول ولی معتبر باشند.
نامشترک بودن این دو مجموعه، خودبه‌خود
اثبات امکان افزایش دقت با فشردگی بیشتر نیست.
نرخ خطا در ناحیه دنباله‌ای باید همراه با
شمار نمونه‌ها گزارش شود، زیرا اندازه آن
ممکن است کوچک و نتیجه ناپایدار باشد.

برچسب واقعی برای ساخت این مجموعه‌ها
صرفاً در مرحله \emph{ارزیابی} به کار می‌رود
و ورودی پیش‌بینی نمونه جدید نیست.
در سیاست شرطی طبقه، بازنمایی نمونه تحت
فرض هر طبقه ممکن است متفاوت باشد؛
بنابراین سنجش خطا و حاشیه باید طبق قرارداد
استنتاج رابطه~\eqref{cmp:eq:ch3-class-hypothesis-inference}
صورت گیرد. کران‌های هندسی مبتنی بر
یک بازنمایی مشترک، بدون اثبات جداگانه
به این حالت تعمیم داده نمی‌شوند.


\subsection{توپ‌های کمینهٔ پوشاننده، تداخل طبقه‌ها و ارتباط با SVDD}
\label{cmp:subsec:ch3-outlier-js-enclosing-balls}
در نگاه نخست، غربال چندکی KL
نمونه‌هایی را مشخص می‌کرد که
\emph{هزینه اضافی درست‌نمایی}
بیشتری در مقایسه با مرکز حسابی
ثابت طبقه خود تحمیل می‌کنند.
نگاه دوم پرسش متفاوتی دارد:
\emph{کدام نمونه‌های مرزی، هندسه
پوشش طبقه و تداخل آن با
طبقه‌های رقیب را تحت تأثیر
قرار می‌دهند؟}
برای پاسخ، کوچک‌ترین توپ
پوشاننده بازنمایی‌های هر طبقه
را در فضای هیلبرت پیدا می‌کنیم
و اثر نقاط مرزی آن را بر
تداخل توپ‌های طبقه‌ها می‌سنجیم.
در این نگاه، قرارگرفتن روی
مرز توپ یا اثرگذاری بر تداخل،
به‌تنهایی تعریف پرت‌بودن
آماری یا نادرستی نمونه نیست؛
همچنین هر نقطه مرزی لزوماً
در تعیین شعاع توپ مؤثر نیست.
پس این تحلیل، \emph{مکمل}
غربال نامزدهای پرت با KL است،
نه تعریف دوم و معادل آن.

\paragraph{صورت مسئله و مرز استنباط.}
سیاست، داده آموزش و بازنمایی
برچسب همسایگی نمونه‌ها را
تثبیت می‌کنیم. برای هر طبقه
$c$، مجموعه کامل نمونه‌های
آموزشی با برچسب همان طبقه را
$C_c$ و زیرمجموعه مورد استفاده
برای \emph{برآورد توپ پوشاننده}
را $T_c\subseteq C_c$ می‌نامیم.
هدف هندسی آن است که با حفظ
حداقل شمار از پیش تعیین‌شده
در هر طبقه، زیرمجموعه‌هایی
بیابیم که \emph{هزینه نقض شکاف
مطلوب میان توپ‌های کمینه
پوشاننده} را کاهش دهند.
بدین‌ترتیب کمیت مورد تحلیل
\emph{تداخل پوش‌های کروی طبقه‌ها}
است، نه هم‌پوشانی واقعی
توزیع‌های مولد، احتمال خطا،
یا احتمال هم‌برچسب‌بودن
نمونه‌ها. دو توپ می‌توانند
تقاطع داشته باشند، درحالی‌که
هیچ نمونه‌ای از طبقه اول
در نزدیکی نمونه‌های طبقه
دوم مشاهده نشده باشد.

برای آستانه شمارشی
$0<\alpha\le1$ و شکاف مطلوب
$\delta_{\mathrm{sep}}>0$،
پس از تعریف توپ و تابع هزینه
در ادامه این زیربخش، مسئله
موردنظر را به‌صورت کمینه‌سازی
همان هزینه روی خانواده
زیرمجموعه‌های مجاز صورت‌بندی
می‌کنیم. روش حریصانه بعدی
فقط یک راهکار پیشنهادی
برای بهبود موضعی این هدف است.
کنارگذاری نمونه صرفاً
از \emph{محاسبه توپ} انجام
می‌شود؛ هیچ نمونه‌ای از
داده، گراف، گشت یا آموزش
سیاست حذف نمی‌شود.
متریک
$d_{\mathrm{JS}}(P,Q)=\sqrt{\operatorname{JSD}(P,Q)}$
و نگاشت یک‌به‌یک
$\Phi:\Delta_K\to\mathcal H$
در روابط~\eqref{cmp:eq:ch3-js-Hilbert-map} و
\eqref{cmp:eq:ch3-js-distance} معرفی و
در پیوست~\ref{app:ch3-js-error-proof}
اثبات شده‌اند؛ بنابراین نیازی به
فرض اقلیدسی بودن فضای توزیع‌ها نیست.

\paragraph{شعاع و مرکز توپ کمینه.}
سیاست، داده آموزش و قرارداد بازنمایی را
ثابت می‌گیریم و برای هر نمونه آموزشی
$x_i=\Phi(\mathscr Y_i)\in\mathcal H$
می‌نویسیم. اگر خانواده سیاست برای
فرضیه‌های متفاوت طبقه، بازنمایی‌های
متفاوت تولید کند، مجموعه تصویرها باید
با یک قرارداد مشخص و سازگار تشکیل شود؛
جدایی توپ‌های این بازنمایی‌های آموزشی
به‌تنهایی کرانی برای قاعده استنتاج
شرطی نمونه جدید نیست. برای هر مجموعه ناتهی
$T_c\subseteq C_c$، مرکز توپ هیلبرتی
$\zeta_c(T_c)$ و شعاع آن
$\mathfrak r_c(T_c)$ را با مسئله
\begin{equation}\label{cmp:eq:ch3-outlier-js-meb-primal}
\boxed{
\begin{aligned}
\mathfrak r_c(T_c)^2
&=\min_{\zeta\in\mathcal H}
\max_{i\in T_c}\|x_i-\zeta\|_{\mathcal H}^2,\\
\zeta_c(T_c)
&\in\operatorname*{arg\,min}_{\zeta\in\mathcal H}
\max_{i\in T_c}\|x_i-\zeta\|_{\mathcal H}^2.
\end{aligned}}
\end{equation}
تعریف می‌کنیم و توپ متناظر را
$\mathcal B_c(T_c)=
\{z\in\mathcal H:
\|z-\zeta_c(T_c)\|_{\mathcal H}
\le\mathfrak r_c(T_c)\}$
می‌نامیم. برای مجموعه متناهی،
جواب وجود دارد و مرکز توپ یکتاست.
مرکز این توپ \emph{آزاد در فضای هیلبرت}
است و لزوماً به‌صورت
$\Phi(v)$ برای هیچ
$v\in\Delta_K$ به دست نمی‌آید؛
ازاین‌رو نباید آن را با
$\Phi(\mathcal V_c)$ یا مرکز
میانگین KL یکی گرفت.

مزیت این نگاشت آن است که شعاع
را می‌توان صرفاً با
$\operatorname{JSD}(\mathscr Y_i,\mathscr Y_j)$
محاسبه کرد. به‌طور دقیق، اگر
$\Delta(T_c)$ سیمپلکس وزن‌های
نامنفی با مجموع یک روی اعضای $T_c$ باشد،
\begin{equation}\label{cmp:eq:ch3-outlier-js-meb-dual}
\boxed{
\begin{aligned}
\mathfrak r_c(T_c)^2
&=\max_{w\in\Delta(T_c)}
\frac12\sum_{i,j\in T_c}w_iw_j
\operatorname{JSD}(\mathscr Y_i,\mathscr Y_j),\\
\zeta_c(T_c)
&=\sum_{i\in T_c}w_i^*x_i
\quad\text{برای یک وزن بهینه }w^*.
\end{aligned}}
\end{equation}
این یک بیشینه‌سازی تابع مقعر
روی سیمپلکس است و به ساخت
مختصات هیلبرتی بی‌نهایت‌بعدی
نیازی ندارد. اثبات دوگانگی و
روش محاسبه فاصله میان مراکز
دو توپ در
پیوست~\ref{app:ch3-js-enclosing-balls}
آمده است.

\paragraph{ارتباط با توصیف دامنه داده با بردارهای پشتیبان (\lr{SVDD}).}
دوگان فوق از نظر ساختار
دقیقاً با صورت \emph{مرز سخت}
روش \lr{Support Vector Data Description}
یا \lr{SVDD} کرنلی متناظر است.
در نگاشت هیلبرتی انتخاب‌شده،
می‌توان کرنل مثبت‌نیمه‌معین
زیر را به کار برد:
\begin{equation}\label{cmp:eq:ch3-outlier-js-svdd-kernel}
\kappa_{\mathrm{JS}}(P,Q)
=\langle\Phi(P),\Phi(Q)\rangle_{\mathcal H}
=\frac{\log 2-\operatorname{JSD}(P,Q)}{2}.
\end{equation}
از ثابت‌بودن
$\kappa_{\mathrm{JS}}(P,P)=\log2/2$
و $\sum_iw_i=1$ نتیجه می‌شود که
تابع هدف دوگان استاندارد
توپ پوشاننده در \lr{SVDD}،
\[
\sum_iw_i\kappa_{\mathrm{JS}}(P_i,P_i)
-\sum_{i,j}w_iw_j\kappa_{\mathrm{JS}}(P_i,P_j),
\]
با قرار دادن $P_i=\mathscr Y_i$
دقیقاً به تابع هدف
رابطه~\eqref{cmp:eq:ch3-outlier-js-meb-dual}
تبدیل می‌شود.
بنابراین این همان مسئله
\emph{توپ کمینه پوشاننده کرنلی}
است، نه \lr{SVM} دودویی
با ابرصفحه جداکننده طبقه‌ها.
این هم‌ارزی امکان بهره‌گیری
از شیوه‌های محاسباتی
مبتنی بر مجموعه فعال و
حل‌گرهای \lr{SVDD} را
در مطالعات بعدی مطرح می‌کند.

روش \lr{SVDD} در پژوهش
\lr{Tax} و \lr{Duin}~
\cite{TaxDuin2004SVDD}
با الهام از بردارهای پشتیبان
در \lr{SVM} مطرح شده است.
هم‌ارزی حاضر، مربوط به
\emph{توصیف یک طبقه با توپ}
است، نه ابرصفحه جداکننده
در \lr{SVM} دودویی.

برای مقایسه با صورت مرز سخت،
صورت مرز نرم \lr{SVDD}
با کرنل~\eqref{cmp:eq:ch3-outlier-js-svdd-kernel}
را نیز می‌توان نوشت.
برای مجموعه ناتهی $T_c$
و ضریب جریمه
$C_{\mathrm{SVDD}}>0$،
مسئله اولیه عبارت است از:
\begin{equation}\label{cmp:eq:ch3-outlier-js-svdd-soft-primal}
\boxed{
\begin{aligned}
\min_{\zeta,t,\xi}\quad&
t+C_{\mathrm{SVDD}}\sum_{i\in T_c}\xi_i\\
\text{\lr{s.t.}}\quad&
\|x_i-\zeta\|_{\mathcal H}^2
\le t+\xi_i\quad(i\in T_c),\\
&t\ge0,\qquad \xi_i\ge0.
\end{aligned}}
\end{equation}
با ضرب‌کننده‌های $w_i$
برای قیود پوشش،
شرط بهینگی نسبت به
$\xi_i$ کران
$0\le w_i\le C_{\mathrm{SVDD}}$
و شرط بهینگی نسبت به
$\zeta$ رابطه
$\zeta=\sum_iw_ix_i$
را در حالتی که
$\sum_iw_i=1$ می‌دهد.
برای $C_{\mathrm{SVDD}}\ge
1/|T_c|$، می‌توان
یک جواب بهینه دوگان با
$\sum_iw_i=1$ برگزید؛
در این صورت دوگان
به شکل زیر است:
\begin{equation}\label{cmp:eq:ch3-outlier-js-svdd-soft-dual}
\boxed{
\begin{aligned}
\max_{w}\quad&
\frac12\sum_{i,j\in T_c}
w_iw_j\operatorname{JSD}
(\mathscr Y_i,\mathscr Y_j),\\
\text{\lr{s.t.}}\quad&
\sum_{i\in T_c}w_i=1,\qquad
0\le w_i\le C_{\mathrm{SVDD}}.
\end{aligned}}
\end{equation}
اگر $C_{\mathrm{SVDD}}\ge1$،
کران بالای وزن‌ها زائد
و دوگان همان صورت
مرز سخت است.
شرایط مکملی همچنین
نشان می‌دهند برای
$0<w_i<C_{\mathrm{SVDD}}$
نمونه روی مرز توپ
و دارای $\xi_i=0$ است؛
اما $w_i=C_{\mathrm{SVDD}}$
به‌تنهایی تخطی مثبت
را تضمین نمی‌کند.

این روابط صرفاً ارتباط
با روش‌های حل کرنلی را
روشن می‌کنند. \lr{SVDD}
نرم مستقیماً تداخل
توپ‌های \emph{طبقه‌های متفاوت}
را کمینه نمی‌کند و
قید حداقل پوشش شمارشی
زیرمجموعه‌های منتخب را نیز،
که در ادامه تعریف می‌شود،
خودبه‌خود تضمین نمی‌کند. در ادامه، به‌جای این
مدل مرز نرم، غربال حریصانه‌ای
را بر اساس کاهش مستقیم هزینهٔ
تداخل توپ‌ها صورت‌بندی می‌کنیم.

\paragraph{نقاط مرزی و نقش آن‌ها.}
مجموعه نقاط قرارگرفته روی مرز
توپ کمینه را تعریف می‌کنیم:
\begin{equation}\label{cmp:eq:ch3-outlier-js-meb-boundary}
\mathcal S_c(T_c)=
\left\{i\in T_c:
\|x_i-\zeta_c(T_c)\|_{\mathcal H}
=\mathfrak r_c(T_c)\right\}.
\end{equation}
وزن‌های مثبت یک جواب دوگان
تنها به نقاط این مجموعه تعلق می‌گیرند.
برخلاف تشبیه دایره دوبعدی که
با حداکثر سه نقطه پشتیبان
تعیین می‌شود، تعداد نقاط
پشتیبان یک توپ در فضای هیلبرت
لزوماً سه نیست. همچنین هر نقطه
مرزی الزاماً پرت یا زیان‌آور
برای جدایی طبقه‌ها نیست.
اگر یک نمونه از مجموعه پوشش
کنار گذاشته شود، شعاع توپ کمینه
افزایش نمی‌یابد؛ ولی ممکن است
اصلاً کاهش نیابد یا مرکز توپ
به جهتی جابه‌جا شود که تداخل
با طبقه دیگر را افزایش دهد.
از دید تحلیلی، سه مفهوم
را تفکیک می‌کنیم:
\emph{نمونه مرزی} عضوی از
$\mathcal S_c(T_c)$ است؛
\emph{نمونه مؤثر بر توپ}
عضوی است که حذف آن، مرکز یا
شعاع توپ کمینه را تغییر می‌دهد؛
و \emph{نمونه مؤثر بر هزینه
تداخل} عضوی است که حذف آن،
$\Omega_{\delta_{\mathrm{sep}}}$
را تغییر می‌دهد. تعریف سوم
فقط با سنجش دوباره توپ‌ها
و شکاف‌های طبقه‌ها مشخص
می‌شود؛ مثبت بودن سود
$\Delta_{c,i}$ شرط قوی‌تر
برای پذیرش حذف در الگوریتم
حریصانه است.
همان‌طور که در پیوست~
\ref{app:ch3-js-enclosing-balls}
اثبات شده، اگر یک جواب
بهینه دوگان برای نمونه
$i$ وزن صفر داشته باشد،
حذف آن توپ را عوض نمی‌کند.
بنابراین مرزی‌بودن به‌تنهایی
حتی مؤثربودن بر توپ را
تضمین نمی‌کند.

\paragraph{تعریف تداخل توپ‌های طبقه‌ها.}
با قرار دادن
$T=(T_1,\ldots,T_K)$،
برای هر دو طبقه متفاوت $c,r$
شکاف هندسی توپ‌ها را چنین می‌سنجیم:
\begin{equation}\label{cmp:eq:ch3-outlier-js-ball-gap}
g_{cr}(T)=
\|\zeta_c(T_c)-\zeta_r(T_r)\|_{\mathcal H}
-\mathfrak r_c(T_c)-\mathfrak r_r(T_r).
\end{equation}
شرط $g_{cr}>0$ به معنای
جدا بودن دو توپ بسته،
$g_{cr}=0$ به معنای تماس
و $g_{cr}<0$ به معنای
هم‌پوشانی آن‌هاست.
برای مطالبه یک فاصله مثبت
$\delta_{\mathrm{sep}}>0$،
هزینه تداخل کل را تعریف می‌کنیم:
\begin{equation}\label{cmp:eq:ch3-outlier-js-total-overlap}
\boxed{
\Omega_{\delta_{\mathrm{sep}}}(T)=
\sum_{1\le c<r\le K}
\bigl[\delta_{\mathrm{sep}}-g_{cr}(T)\bigr]_+,
\qquad [u]_+=\max\{u,0\}.
}
\end{equation}
در نتیجه $\Omega_{\delta_{\mathrm{sep}}}=0$
شرط کافی برای جدایی همه توپ‌ها
با شکاف دست‌کم
$\delta_{\mathrm{sep}}$ است.
اگر $\delta_{\mathrm{sep}}=0$،
صفر شدن هزینه، تماس مرزی را
نیز مجاز می‌شمارد.
این تابع، میزان \emph{نقض
شکاف مطلوب میان توپ‌های
پوشاننده} را می‌سنجد و
به‌تنهایی معیاری از اشتراک
واقعی چگالی یا پسین طبقه‌ها
نیست.

\paragraph{مسئله انتخاب زیرمجموعه با حداقل پوشش شمارشی.}
با فرض ثابت بودن
$\alpha$ و $\delta_{\mathrm{sep}}$،
خانواده زیرمجموعه‌های مجاز را
$\mathcal T_\alpha$ تعریف
می‌کنیم و هدف هندسی را می‌نویسیم:
\begin{equation}\label{cmp:eq:ch3-outlier-js-subset-problem}
\boxed{
\begin{aligned}
\mathcal T_\alpha
&=\left\{(T_1,\ldots,T_K):
 T_c\subseteq C_c,\;
 |T_c|\ge\lceil\alpha n_c\rceil
 \ \forall c\right\},\\
\min_{T\in\mathcal T_\alpha}
&\quad\Omega_{\delta_{\mathrm{sep}}}(T).
\end{aligned}}
\end{equation}
در اینجا $n_c=|C_c|$ و
$0<\alpha\le1$ تضمین می‌کند
که هر $T_c$ ناتهی بماند.
خانواده مجاز متناهی و ناتهی
است؛ ازاین‌رو کمینه سراسری
وجود دارد، اما پیدا کردن
آن با روش حریصانه تضمین
نمی‌شود. همچنین قید پوشش
صرفاً \emph{شمار اعضای
باقی‌مانده} را تضمین
می‌کند، نه حفظ زیرگروه‌ها،
تنوع آماری یا نمایندگی
توزیع واقعی طبقه.
اگر چند جواب هزینه یکسان
داشته باشند، ترجیح نگه‌داشتن
نمونه‌های بیشتر می‌تواند
معیار ثانویه باشد؛ چنین
ترجیحی بخشی از تابع هزینه
فعلی نیست.

\paragraph{کنارگذاری تدریجی نمونه‌های مرزی متداخل.}
هدف \emph{کمینه کردن شعاع همه طبقه‌ها}
نیست؛ زیرا یک نقطه دورافتاده
ممکن است به‌درستی تنوع طبقه خودش
را نمایندگی کند و به توپ رقیب
نزدیک نباشد. ابتدا فقط طبقه‌هایی
را بررسی می‌کنیم که با دست‌کم
یک توپ رقیب شکاف کمتر از
$\delta_{\mathrm{sep}}$ دارند.
نامزدها از نقاط مرزی
$\mathcal S_c(T_c)$ همین طبقه‌ها
انتخاب می‌شوند. برای اولویت‌بندی
نمونه‌های نزدیک یا واردشده به
توپ رقیب می‌توان کمیت
\begin{equation}\label{cmp:eq:ch3-outlier-js-rival-intrusion}
a_{i,c}(T)=
\min_{\substack{r\ne c\\g_{cr}(T)<\delta_{\mathrm{sep}}}}
\bigl[
\|x_i-\zeta_r(T_r)\|_{\mathcal H}
-\mathfrak r_r(T_r)
\bigr]
\end{equation}
را محاسبه کرد.
مقدار منفی یعنی نمونه درون
دست‌کم یکی از توپ‌های رقیب است.
بااین‌حال، مثبت‌بودن $a_{i,c}$
دلیل کافی برای بی‌تأثیربودن
نمونه بر تداخل نیست؛ بنابراین
آن را \emph{اولویت نرم} می‌دانیم،
نه شرط اجباری نامزدی.

برای هر نامزد $(c,i)$،
فقط $i$ را \emph{به‌طور آزمایشی
از مجموعه برآورد توپ} خارج می‌کنیم،
$T_c^{(-i)}=T_c\setminus\{i\}$،
توپ همین طبقه را از نو می‌یابیم
و تغییر هزینه سراسری را می‌سنجیم:
\begin{equation}\label{cmp:eq:ch3-outlier-js-greedy-gain}
\boxed{
\Delta_{c,i}(T)=
\Omega_{\delta_{\mathrm{sep}}}(T)
-\Omega_{\delta_{\mathrm{sep}}}
(T_1,\ldots,T_c^{(-i)},\ldots,T_K).
}
\end{equation}
در هر دور، از میان نامزدهای
مجاز، تنها یک نمونه با
\emph{بیشترین کاهش مثبت هزینه}
کنار گذاشته می‌شود و پس از
بازمحاسبه توپ‌ها، نامزدهای
دور بعد دوباره تعیین می‌شوند.
هنگام تساوی سود، می‌توان
به نمونه دارای $a_{i,c}$
کمتر اولویت داد. در هر گام، نامزد
$i\in T_c$ تنها زمانی
\emph{مجاز} است که
$|T_c|-1\ge
\lceil\alpha n_c\rceil$؛
بنابراین قید شمارشی
رابطه~\eqref{cmp:eq:ch3-outlier-js-subset-problem}
در تمام گام‌ها برقرار
می‌ماند. اگر هزینه سراسری
به صفر برسد یا در \emph{هیچ
یک از طبقه‌ها} نامزد مجاز
با $\Delta_{c,i}>0$ باقی
نماند، الگوریتم متوقف می‌شود.
رسیدن یک طبقه به حداقل
پوشش، به‌تنهایی شرط توقف
سراسری نیست.
این یک روش \emph{حریصانه پیشنهادی}
است: بهینه‌بودن شمار کنارگذاری‌ها،
یافتن جدایی در صورت امکان، یا
بهبود دقت طبقه‌بندی را تضمین
نمی‌کند. به‌ویژه جدایی توپ‌ها
به‌خودی‌خود تضمین درستی تصمیم
نزدیک‌ترین \emph{مرکز توزیعی}
نیست؛ این مراکز ممکن است با
مراکز توپ‌ها متفاوت باشند.

شکل~\ref{cmp:fig:ch3-js-overlap-boundary-screening}
سازوکار کنارگذاری یک نمونه مرزی
و بازبهینه‌سازی توپ را نشان می‌دهد.
توپ آبی طبقه سوم از دو توپ دیگر
جداست و به همراه توپ نارنجی طبقه نخست
در هر دو قاب \emph{کاملاً ثابت می‌ماند}.
نقطه مشخص‌شده، یک نمونه \emph{سبز}
از طبقه دوم است که روی مرز توپ سبز
و درون توپ نارنجی قرار دارد.
پس از کنارگذاری همین نمونه از
\emph{محاسبه توپ سبز}، مرکز و شعاع
فقط توپ طبقه دوم دوباره تعیین می‌شوند.
نقطه در داده و گراف باقی می‌ماند،
اما بیرون توپ جدید قرار می‌گیرد.

\begin{figure}[htbp]
\centering
\resizebox{0.98\linewidth}{!}{%
% تصویر کاملاً شماتیک: دو برش دوبعدی از مفهوم توپ پوشاننده در هندسه هیلبرت.
% این مختصات ادعای تصویر دوبعدی دقیق نگاشت JS یا نتیجه یک آزمایش را ندارند.
% با وجود این، دایره‌های ترسیم‌شده واقعاً کوچک‌ترین دایره پوشاننده
% نقاط نمایشی خود هستند (کنترل هندسی مستقل در بازبینی انجام شده است).
% ۱) توپ مستقلِ بزرگ کلاس ۳ دست‌نخورده باقی می‌ماند.
% ۲) در پنل نخست، نقطه قرمز عضو کلاس ۱ است، روی پوش آن قرار دارد
%    و در توپ کلاس ۲ هم قرار گرفته است.
% ۳) در پنل دوم، همان نمونه هنوز رسم شده و در مجموعه داده/گراف وجود دارد
%    ولی از مجموعه محاسبه توپ کلاس ۱ کنار گذاشته شده است.
% ۴) مراکز و شعاع‌ها از روی نقاط حاضر در هر محاسبه بازبهینه شده‌اند.
\begingroup
\definecolor{jsballblue}{HTML}{286C9E}
\definecolor{jsballorange}{HTML}{CE7B28}
\definecolor{jsballgreen}{HTML}{25816B}
\definecolor{jsballred}{HTML}{BD3443}
\definecolor{jsballink}{HTML}{334453}
\definecolor{jsballpanel}{HTML}{F7F9FC}

\begin{tikzpicture}[x=1cm,y=1cm,line cap=round,line join=round]
% سه توپ: کلاس ۳ (آبی)، کلاس ۱ (نارنجی)، کلاس ۲ (سبز).
% نمونه‌های کلاس ۱ پیش از کنارگذاری:
% (-.5,0),(-.05,.46),(0,-.45),(.4,.15),(1.4,.08)
% MEB_1: مرکز (.45,.04)، شعاع .95084.
% پس از کنارگذاری (1.4,.08):
% MEB_1: مرکز (-.03906,.00937)، شعاع .46103.
% MEB_2: مرکز (2.16,0)، شعاع 1؛ MEB_3: مرکز (-3.1,0)، شعاع 1.22.
% شکاف کلاس‌های ۱ و ۲: پیش -0.24037 و پس +0.73805.
% اعداد صرفاً کنترل هندسیِ چیدمان شماتیک‌اند، نه نتیجه تجربی.
\foreach \panel/\shift in {1/0,2/9.1}{%
  \begin{scope}[xshift=\shift cm]
    \fill[jsballpanel,rounded corners=6pt] (0,-1.92) rectangle (8.55,1.94);
    \draw[jsballink!25,line width=.45pt,rounded corners=6pt]
        (0,-1.92) rectangle (8.55,1.94);
    % انتقال دستگاه مختصات نمونه‌ها به کادر
    \begin{scope}[shift={(4.56,0)}]
      % کلاس بزرگ جدا: در هر دو قاب بدون تغییر
      \fill[jsballblue!11] (-3.10,0) circle (1.22);
      \draw[jsballblue,line width=1.1pt] (-3.10,0) circle (1.22);
      % کلاس رقیب: در هر دو قاب بدون تغییر
      \fill[jsballgreen!12] (2.16,0) circle (1.00);
      \draw[jsballgreen,line width=1.1pt] (2.16,0) circle (1.00);
      \ifnum\panel=1
        \fill[jsballorange!14] (.45,.04) circle (.95084);
        \draw[jsballorange,line width=1.2pt] (.45,.04) circle (.95084);
        % مرکز پیشین (مربع نارنجی)
        \filldraw[fill=jsballorange,draw=white,line width=.75pt]
          (.45,.04) rectangle ++(.13,.13);
        % نمونه مرزی کنارگذاری‌شونده، هم داخل توپ رقیب
        \filldraw[fill=jsballred,draw=white,line width=.7pt]
          (1.4,.08) circle (.092);
        \draw[jsballred,line width=.8pt]
          (1.4,.08) circle (.145);
      \else
        \fill[jsballorange!14] (-.03906,.00937) circle (.46103);
        \draw[jsballorange,line width=1.2pt]
          (-.03906,.00937) circle (.46103);
        % مرکز تازه پس از بازبهینه‌سازی
        \filldraw[fill=jsballorange,draw=white,line width=.75pt]
          (-.03906,.00937) rectangle ++(.13,.13);
        % نمونه همچنان وجود دارد؛ فقط در محاسبه توپ منظور نشده است.
        \draw[jsballred,line width=1.1pt] (1.4,.08) circle (.11);
        \draw[jsballred,line width=1.2pt]
          (1.34,.02)--(1.46,.14)
          (1.34,.14)--(1.46,.02);
      \fi

      % نقاط کلاس ۳ (دو نمونه روی قطرِ تعیین‌کننده پوش)
      \foreach \x/\y in {-4.32/0,-1.88/0,-3.20/.91,-2.70/-.70,-3.65/.40,-2.92/.10}{%
        \filldraw[fill=jsballblue,draw=white,line width=.65pt]
          (\x,\y) circle (.073);
      }
      % نقاط کلاس ۲ (دو نمونه روی قطرِ تعیین‌کننده پوش)
      \foreach \x/\y in {1.16/0,3.16/0,2.15/.78,2.45/-.73,1.85/-.60}{%
        \filldraw[fill=jsballgreen,draw=white,line width=.65pt]
          (\x,\y) circle (.073);
      }
      % نقاط باقی‌مانده کلاس ۱، شامل نقاط روی و داخل پوش
      \foreach \x/\y in {-.50/0,-.05/.46,0/-.45,.40/.15}{%
        \filldraw[fill=jsballorange,draw=white,line width=.65pt]
          (\x,\y) circle (.073);
      }
      % مراکز توپ‌های تغییرنکرده
      \filldraw[fill=jsballblue,draw=white,line width=.7pt]
          (-3.10,0) rectangle ++(.13,.13);
      \filldraw[fill=jsballgreen,draw=white,line width=.7pt]
          (2.16,0) rectangle ++(.13,.13);
      % مشخصه کلاس‌ها خارج از سطح توپ‌ها
      \node[font=\footnotesize,text=jsballblue]
          at (-3.10,1.51) {$\mathcal B_3$};
      \node[font=\footnotesize,text=jsballgreen]
          at (2.16,1.36) {$\mathcal B_2$};
      \ifnum\panel=1
        \node[font=\footnotesize,text=jsballorange]
          at (.23,1.33) {$\mathcal B_1$};
      \else
        \node[font=\footnotesize,text=jsballorange]
          at (-.06,.79) {$\mathcal B_1^{(-i)}$};
      \fi
    \end{scope}
    \ifnum\panel=1
      \node[font=\small\bfseries,text=jsballink]
         at (4.28,1.73) {پیش از کنارگذاری نمونه مرزی};
      \node[font=\footnotesize,text=jsballred]
         at (4.25,-1.62) {هم‌پوشانی توپ‌های کلاس‌های ۱ و ۲};
    \else
      \node[font=\small\bfseries,text=jsballink]
         at (4.28,1.73) {پس از بازبهینه‌سازی توپ};
      \node[font=\footnotesize,text=jsballgreen]
         at (4.25,-1.62) {جدایی دو توپ با حفظ نمونه در گراف};
    \fi
  \end{scope}
}
% فلش بین دو قاب
\draw[->,line width=1.4pt,jsballink!70] (8.65,0)--(9.0,0);
% راهنمای نمادها (به‌همراه متن فارسی کوتاه)
\begin{scope}[shift={(3.15,-2.32)}]
  \filldraw[fill=jsballorange,draw=white,line width=.4pt] (0,0) circle (.085);
  \node[anchor=west,font=\scriptsize,text=jsballink] at (.17,0) {نمونه فعال};
  \filldraw[fill=jsballblue,draw=white,line width=.4pt]
      (2.05,-.075) rectangle +( .15,.15);
  \node[anchor=west,font=\scriptsize,text=jsballink] at (2.35,0) {مرکز توپ};
  \draw[jsballred,line width=1pt] (4.65,0) circle (.085);
  \draw[jsballred,line width=.8pt] (4.60,-.05)--(4.70,.05)
                                   (4.60,.05)--(4.70,-.05);
  \node[anchor=west,font=\scriptsize,text=jsballink] at (4.88,0)
      {نمونه کنارگذاشته از محاسبه توپ};
\end{scope}
\end{tikzpicture}
\endgroup
%
}
\caption{نمایش شماتیک دوبعدی غربال نمونه مرزی با توپ‌های کمینه پوشاننده. هر دایره سه نمونه روی محیط و حداکثر سه نمونه داخلی دارد؛ مربع‌ها مراکز توپ‌اند. نقطه سبز مشخص‌شده در قاب اول بر محیط توپ طبقه دوم و درون توپ نارنجی طبقه نخست قرار دارد. در قاب دوم همان نقطه در موقعیت قبلی با علامت خط‌چین/ضربدر باقی مانده، ولی فقط از \emph{محاسبه توپ سبز} کنار گذاشته شده است؛ مرکز و شعاع سبز تغییر می‌کنند و سه نمونه دیگر مرز توپ جدید را تعیین می‌کنند. توپ‌های آبی و نارنجی و تمام نقاطشان در هر دو قاب ثابت‌اند. این آرایش صرفاً یک مثال هندسی دوبعدی است و داده آزمایشی یا نگاشت دقیق هیلبرتی جنسن--شنون را نمایش نمی‌دهد.}
\label{cmp:fig:ch3-js-overlap-boundary-screening}
\end{figure}

در تمام این مراحل، نمونه‌های
کنارگذاشته‌شده همچنان در
\emph{داده آموزش، گراف و گشت}
حضور دارند؛ فقط در محاسبه
توپ پوشاننده همان طبقه دخیل نیستند.
بنابراین \emph{خود غربال هندسی}
بدون تغییر قاعده تصمیم یا مراکز
طبقه‌بند، هیچ پیش‌بینی و دقتی
را تغییر نمی‌دهد.
اگر بعدها از این غربال برای
تغییر مرکز طبقه‌بند استفاده شود،
آن مداخله باید مطابق قرارداد
زیربخش~\ref{cmp:subsec:ch3-outlier-robust-centers}
و با ارزیابی مستقل انجام شود.
بنابراین غربال هندسی باید
از اصلاح مراکز و قاعدهٔ تصمیم
به‌روشنی تفکیک شود.


\subsection{برآورد مقاوم مرکز بدون حذف نمونه از گراف}
\label{cmp:subsec:ch3-outlier-robust-centers}
در ادامه، امکان کاهش سهم نمونه‌های
دورافتاده را تنها در \emph{محاسبه مرکز طبقه}
کاهش دهیم. برای وزن‌های نامنفی با مجموع مثبت تعریف می‌کنیم:
\begin{equation}\label{cmp:eq:ch3-outlier-weighted-center}
\widetilde{\mathcal V}_c(w)=
\frac{\sum_{i\in C_c}w_i\mathscr Y_i}
{\sum_{i\in C_c}w_i},
\qquad
0\le w_i\le1,\quad \sum_{i\in C_c}w_i>0.
\end{equation}
وزن‌های یکسان، مرکز اصلی را بازتولید می‌کنند؛
وزن صفر تنها سهم نمونه را از \emph{میانگین مرکز}
حذف می‌کند و به معنای پاک‌کردن نمونه از
داده، گراف، گشت یا یادگیری سیاست آن نیست.
این تمایز ساختاری ضروری است: یک نمونه دورافتاده
ممکن است با سیاست اختصاصی خود به‌درستی
بازنمایی یا طبقه‌بندی شود و حضورش همچنان
برای انتقال‌های نمونه--ویژگی اهمیت داشته باشد.

برای وزن‌های \emph{ثابت}، مرکز وزن‌دار در جهت
$D_{\mathrm{KL}}(\mathscr Y_i\|v)$ نیز
میانگین بهینه است؛ به‌طور دقیق اگر
$\bar w_i=w_i/\sum_{j\in C_c}w_j$،
برای هر $v$ با واگرایی‌های متناهی داریم:
\begin{equation}\label{cmp:eq:ch3-outlier-weighted-kl-decomposition}
\sum_{i\in C_c}\bar w_i
D_{\mathrm{KL}}(\mathscr Y_i\|v)
=
\sum_{i\in C_c}\bar w_i
D_{\mathrm{KL}}(\mathscr Y_i\|
\widetilde{\mathcal V}_c)
+
D_{\mathrm{KL}}(\widetilde{\mathcal V}_c\|v).
\end{equation}
در مقابل، اگر وزن‌ها بر اساس فاصله از
\emph{همان مرکز در حال تغییر} تعیین شوند،
این همانی برای هر گام با وزن ثابت معتبر است،
اما به‌تنهایی همگرایی کل فرایند بازوزن‌دهی
یا افزایش دقت را تضمین نمی‌کند.

\paragraph{حفظ تکیه‌گاه در برآورد مقاوم.}
صفرکردن وزن برخی نمونه‌ها ممکن است مؤلفه‌ای
را که در بازنمایی آن نمونه‌ها مثبت است،
در $\widetilde{\mathcal V}_c$ صفر کند؛
در این حالت، اگرچه نمونه در گراف باقی مانده،
ممکن است
$D_{\mathrm{KL}}(\mathscr Y_i\|
\widetilde{\mathcal V}_c)=+\infty$
شود. راه پیشنهادی برای جلوگیری از این وضعیت،
آمیختن مرکز وزن‌دار با مرکز حسابی اصلی است:
\begin{equation}\label{cmp:eq:ch3-outlier-support-preserving-center}
\widetilde{\mathcal V}_{c,\varepsilon}
=(1-\varepsilon)\widetilde{\mathcal V}_c
+\varepsilon\mathcal V_c,\qquad
0<\varepsilon\le1.
\end{equation}
برای هر $i\in C_c$، مثبت بودن یک مؤلفه
$\mathscr Y_i$ مثبت بودن مؤلفه متناظر
$\mathcal V_c$ و در نتیجه
$\widetilde{\mathcal V}_{c,\varepsilon}$
را تضمین می‌کند؛ پس KL پیشرو
برای تمام اعضای آموزشی طبقه متناهی می‌ماند.
این مرکز آمیخته الزاماً کمینه‌کننده
\emph{همان} زیان KL وزن‌دار قبلی نیست
و انتخاب $\varepsilon$ نیاز به اعتبارسنجی دارد.
این پیشنهاد جایگزین بی‌صدای
مرکز حسابی، تعریف فعلی $R_c$ یا قرارداد
محاسباتی KL دقیق نمی‌شود.

\paragraph{پروتکل سنجش و پرسش‌های باز.}
در بررسی تجربی آینده می‌توان
شعاع‌های پوشش مانند $90$، $95$ و $99$ درصد،
وزن‌دهی نرم در برابر حذف صرفاً از میانگین مرکز
و مراکز ثابت در برابر مراکز بهینه‌شده را مقایسه کرد.
آستانه‌ها باید صرفاً با آموزش و اعتبارسنجی
گزینش شوند و مجموعه آزمون برای انتخاب آن‌ها
استفاده نشود. برای پرهیز از ارزیابی دوری،
محاسبه فاصله‌های برون‌چین یا طرح‌های مستقل
می‌تواند لازم باشد.
تعریف نامزد پرت در بازنمایی و غربال
هندسی نقاط مرزی، هیچ‌کدام به‌تنهایی
نادرستی برچسب یا خطای پیش‌بینی
را اثبات نمی‌کنند؛ از این تعریف‌ها
هیچ رابطه علّی، مزیت قطعی برآورد
مقاوم یا بهبود تضمین‌شده دقت
طبقه‌بندی نتیجه نمی‌شود.

\section{جمع‌بندی فصل و گذار به ارزیابی تجربی}
در این فصل، مسیر محاسباتی از سیاست انتخاب ویژگی
و گشت روی گراف تا بازنمایی برچسب همسایگی،
نماینده‌های احتمالی طبقات و قاعدهٔ پیش‌بینی
صورت‌بندی شد. انتخاب مرکز میانگین از دید
درست‌نمایی مدل مشترک توجیه شد، بی‌آنکه
بهینگی آن برای هر سنجهٔ تصمیم ادعا شود.
سپس تراکم و جدایی بازنمایی‌ها، حاشیه و خطای
تصمیم و دو خانوادهٔ زیان احتمالاتی و حاشیه‌ای
به‌صورت متمایز بررسی شدند.

روابط کرانی این فصل حدود قابل‌اثبات هر معیار
را مشخص می‌کنند و به‌تنهایی تضمینی برای
بهبود دقت در تمام گام‌های یادگیری نیستند.
تحلیل مسئلهٔ معکوس، مواجهه با ویژگی‌های نویزی
و صورت‌بندی‌های پساآموزش نیز از روش اصلی
پیش‌بینی تفکیک شدند.
فصل~\ref{chap:results} رفتار واقعی این معیارها
و سیاست‌های آموخته‌شده را روی مجموعه‌های
آموزش، اعتبارسنجی و آزمون بررسی می‌کند.

\endgroup

